\documentclass[11pt]{article}
\usepackage[T1]{fontenc}
\usepackage[utf8]{inputenc}
\usepackage{lmodern}
\usepackage[margin=1in]{geometry}
\usepackage{amsmath,amssymb,amsthm,mathtools,mathrsfs,bm}
\usepackage{booktabs,array,longtable,enumitem,microtype,aliascnt}
\usepackage[colorlinks=true,linkcolor=blue,citecolor=blue,urlcolor=blue]{hyperref}
\usepackage[nameinlink,noabbrev]{cleveref}
\hypersetup{pdftitle={Arithmetic Loading and Holonomy Duality: Cumulative Research Notes V13},pdfauthor={Arithmetic--Geometry Research Programme},pdfsubject={Curve-selection firewall and native cotangent Frobenius normalization}}
\setlength{\emergencystretch}{3em}
\allowdisplaybreaks
\newtheorem{theorem}{Theorem}[section]
\newaliascnt{proposition}{theorem}
\newtheorem{proposition}[proposition]{Proposition}
\aliascntresetthe{proposition}
\newaliascnt{lemma}{theorem}
\newtheorem{lemma}[lemma]{Lemma}
\aliascntresetthe{lemma}
\newaliascnt{corollary}{theorem}
\newtheorem{corollary}[corollary]{Corollary}
\aliascntresetthe{corollary}
\theoremstyle{definition}
\newaliascnt{definition}{theorem}
\newtheorem{definition}[definition]{Definition}
\aliascntresetthe{definition}
\newaliascnt{construction}{theorem}
\newtheorem{construction}[construction]{Construction}
\aliascntresetthe{construction}
\theoremstyle{remark}
\newaliascnt{remark}{theorem}
\newtheorem{remark}[remark]{Remark}
\aliascntresetthe{remark}
\newaliascnt{scope}{theorem}
\newtheorem{scope}[scope]{Scope}
\aliascntresetthe{scope}
\newcommand{\C}{\mathbb C}
\newcommand{\R}{\mathbb R}
\newcommand{\F}{\mathbb F}
\newcommand{\Q}{\mathbb Q}
\newcommand{\Z}{\mathbb Z}
\newcommand{\HH}{\mathcal H}
\newcommand{\Tr}{\operatorname{Tr}}
\newcommand{\rank}{\operatorname{rank}}
\newcommand{\Span}{\operatorname{span}}
\newcommand{\spec}{\operatorname{spec}}
\newcommand{\End}{\operatorname{End}}
\newcommand{\dist}{\operatorname{dist}}
\newcommand{\HS}{\mathrm{HS}}
\newcommand{\op}{\mathrm{op}}
\newcommand{\id}{\mathrm{id}}
\newcommand{\dd}{\mathrm d}
\newcommand{\norm}[1]{\lVert#1\rVert}
\newcommand{\ip}[2]{\langle#1,#2\rangle}
\newcommand{\ketbra}[1]{|#1\rangle\langle#1|}
\newcommand{\AR}{\mathcal E_{\mathrm{tet}}}
\newcommand{\Mtet}{\mathscr M_{\mathrm{tet}}}
\newcommand{\Jq}{\mathcal J}
\newcommand{\one}{\mathbf1}
\newcolumntype{L}[1]{>{\raggedright\arraybackslash}p{#1}}
\title{\textbf{Arithmetic Loading and Holonomy Duality}\\[0.35em]
\large Cumulative Research Notes, V13\\[0.3em]
\normalsize Pure Theta-Source Obstruction and Signed Descent,\\
with Marked Crystalline Frobenius,\\
Minimal Two-Source Completion and Horizontal Transport}
\author{Arithmetic--Geometry Research Programme}
\date{20 September 2026}
\begin{document}
\maketitle
\begin{abstract}
We initiate a new append-only proof-note lineage for the arithmetic-duality programme. The normalized arithmetic loading, the pointed genus-two theta packet, the integral arithmetic--$K_4$ lattice, and the generic five-orbit Gram/Krylov recognition theorems are imported rather than rederived. The first unresolved comparison is attacked at its actual source: five singular theta displacements of a single squared theta section in the four-dimensional level-two theta representation. In its Heisenberg-unitary Hermitian metric, the complete normalized Gram of these five vectors always has spectrum $(2,2,1,0,0)$. It therefore cannot be the Gram of an ordinary pointed $S_4$ source, even after arbitrary changes of the five ray phases. For the retained real theta lifts, the squared distance to the invariant-Gram subspace is exactly $4-2\kappa^2\ge2$, with a sharp constant. Thus the proposed ordinary five-parameter comparison fails on this precise pure-source branch before any period calculation or transfer estimate. The ten unreduced imaginary overlaps nevertheless admit an exact polynomial pure-state recognition and reconstruction theorem. We also classify all invariant mixed-state displacement Grams, prove a minimum reduced-state rank of two, and construct a canonical rank-two quaternionic completion with explicit information loss and word-completion cost. These are finite source statements, not a completed arithmetic--Wilson or arithmetic--matter duality.
\end{abstract}
\paragraph{V2 continuation.}
Sections~\ref{sec:v2-context}--\ref{sec:v2-status} classify the coherent signed
projective tetrahedral lifts, prove their fixed-pure-source and quartic-frame
obstructions, identify an exact completely positive averaging alternative,
and give four theta-constant coordinates with certified acquisition bounds
for the actual arithmetic source. The V1 mathematical body below is unchanged.
\paragraph{V3 continuation.}
Sections~\ref{sec:v3-context}--\ref{sec:v3-status} acquire the actual
curve's complete displacement magnitudes from algebraic branch weights,
prove two exact zeros, and classify the thirty-two-state pure-source phase
fibre with certified separation and a five-sign selection rule. Both earlier
mathematical bodies remain unchanged.
\paragraph{V4 continuation.}
Sections~\ref{sec:v4-context}--\ref{sec:v4-status} retain complex Thomae
phases to select the actual sixteen-element Heisenberg source orbit,
reduce marked acquisition to four signs, prove a complete projective
source test, and reconstruct the arithmetic real involution. All three
earlier mathematical bodies remain unchanged.
\paragraph{V5 continuation.}
Sections~\ref{sec:v5-context}--\ref{sec:v5-status} acquire actual local
Frobenius data, prove their separation from a fixed finite theta-frame
normalizer, construct the native level-raising theta maps and the ordinary
split tower, and identify a positive local response and a central twist
obstruction. All four earlier mathematical bodies remain unchanged.
\paragraph{V6 continuation.}
Sections~\ref{sec:v6-context}--\ref{sec:v6-status} identify the actual
source span with quadratic Hodge data by second jets, construct its native
Cartier transfer, prove the Frobenius jet-order and parity boundaries,
and certify the geometric endomorphism ring from the two acquired primes.
The cited abelian-surface classification then identifies the full symplectic
algebraic monodromy and its small-rank tensor obstruction.
All five earlier mathematical bodies remain unchanged.
\paragraph{V7 continuation.}
Sections~\ref{sec:v7-context}--\ref{sec:v7-status} reconstruct the
actual conformal second-jet metric algebraically, prove its non-product
character, and construct the explicit Gauss--Manin connection on an
independently declared family through the arithmetic curve. The original
quadratic source has a certified minimal rank-ten differential completion;
its local monodromy exposes a separate positive-metric transport obstruction.
All six earlier mathematical bodies remain unchanged.
\clearpage
\paragraph{V9 continuation.}
Sections~\ref{sec:v9-context}--\ref{sec:v9-status} construct the Chern
connections of the already acquired Hodge and theta-source bundles. Exact
kernel rotation and certified theta osculation give full restricted joint
holonomy $U(2)\times U(3)$. A strict determinant-curvature test separates
this from a determinant-one physical gauge group, and a projective-curvature
bound distinguishes the actual theta connection from the quadratic Hodge
connection. All eight earlier mathematical bodies remain unchanged.
\paragraph{V10 continuation.}
Sections~\ref{sec:v10-context}--\ref{sec:v10-status} identify the native
cotangent determinant line, construct its root-cover and root-free
arithmetic incidences, and test the actual Hodge quotient metric. A unique
explicit determinant-compatibility normalization on the existing cotangent
line gives a full $S(U(2)\times U(3))$ Chern holonomy completion. This does
not assert that the original metrics or a physical renewal process select
that normalization. All nine earlier mathematical bodies remain unchanged.

\tableofcontents
\clearpage

\paragraph{V8 continuation.}
Sections~\ref{sec:v8-context}--\ref{sec:v8-status} identify the hypergeometric
period system and its arithmetic monodromy, acquire a complete marked Betti
comparison and the actual Hodge metric, determine the actual theta-frame
signs by validated evaluation, and construct a positive rank-ten completion
under an explicit minimum-change rule. All seven earlier mathematical
bodies remain unchanged.

\paragraph{V11 continuation.}
Sections~\ref{sec:v11-context}--\ref{sec:v11-status} acquire the actual
crystalline Frobenius in the marked integral differential basis at three
and five. The original quadratic source has prime-cyclic dimension nine
and eight, while the two already rational branch sources minimally
complete the rank-ten system at depth four. A convergent horizontal
Frobenius germ is fixed by the acquired matrices and the actual connection.
All ten earlier mathematical bodies remain unchanged.
\paragraph{V12 continuation.}
Sections~\ref{sec:v12-source-moduli}--\ref{sec:v12-status} prove that the
independently marked pure theta source is locally moduli-complete, construct
a positive projective source metric with an explicit observability margin,
and reduce the remaining theta-fibre Frobenius ambiguity to one native line
scalar. All eleven earlier mathematical bodies remain unchanged.
\paragraph{V13 continuation.}
Sections~\ref{sec:v13-context}--\ref{sec:v13-status} prove that the bare
pointed tetrahedral conditions admit an infinite family of inequivalent
arithmetic curves, isolate a source-readable local curve selector, and use
the already occurring eighth-power cotangent tensor to fix the residual
theta-line Frobenius scalar under a normalized finite-flat rule. All twelve
earlier mathematical bodies remain unchanged.

% ===================== BEGIN IMMUTABLE V1 MATHEMATICAL BODY =====================
\section{Context, inherited results, and the precise first attack}
\label{sec:context}

\subsection{Purpose and version convention}
The intended paper concerns arithmetic loading and arithmetic--geometric or arithmetic--gauge duality. The spacetime, spectral, action--multiplicity, and Einstein--Standard-Model variational constructions are treated as citable companion results \cite{Spacetime,Spectral,Duality,Variational}. They are not reconstruction targets of this lineage.

The initial programme suggested a progression from normalized chronology to an arithmetic source, a tetrahedral source comparison, transfer equivalence, and ultimately Frobenius--Wilson interface equivalence. An audit of the supplied monographs and the recovered earlier arithmetic--spacetime notes \cite{GT77,Arithmetic,AD9} shows that substantial portions of the first stages have already been completed. This version therefore does not restart Peano loading, finite Tannaka reconstruction, the $K_4$ lattice calculation, or the general Gram/Krylov criterion. It attacks the concrete theta-vector source which those works leave unevaluated.

This is the first version of \path{Arithmetic_Loading_Duality_Research_Notes_V1.tex}. A later version is to retain the complete V1 mathematical body and its historical status ledger, append a separately dated continuation before the bibliography, and increment the filename. A correction to a result must be stated explicitly in that continuation rather than silently overwriting the earlier argument.

\subsection{Dependency audit}
\label{subsec:audit}
The following are inherited statements, not new theorems of this note. Original theorem labels are retained so that the dependencies can be located without relying on changing pagination.

\begin{longtable}{L{0.25\textwidth}L{0.67\textwidth}}
\toprule
Imported result & Source and role\\
\midrule
\endfirsthead
\toprule Imported result & Source and role\\\midrule
\endhead
Normalized arithmetic loading & \cite{GT77}, labels \path{thm:ar-Peano-multiplication}, \path{thm:ar-loading-deformations}, \path{thm:ar-Peano-stability}, and \path{thm:ar-record-cutoff}. Uniqueness, deformation classification, finite stability, forward cutoff descent, and absorbing-corner distinctions already exist.\\
Arithmetic source discipline & \cite{Arithmetic}, Sections on canonical arithmetic loading, expansion--pullback, and reader restoration. These do not furnish the theta-vector metric or a Frobenius--Wilson identification. The completed Dirichlet theta kernel is not the genus-two theta-section source studied here.\\
Pointed genus-two packet & \cite{AD9}, inherited V8 labels \path{thm:v8-pointed-theta-packet} and \path{thm:v8-five-singular}. The curve, rational odd characteristic, temporal class, residual $S_4$ Galois action, and pairings of the five singular classes are established there.\\
Integral arithmetic--$K_4$ comparison & \cite{AD9}, \path{thm:v9-integral-arithmetic-k4} and \path{cor:v9-tetrahedron-gram}. The discriminant-twisted integral cycle lattice and its four-source Gram $4I_4-\one\one^{\mathsf T}$ are already matched. This is not the theta-vector Gram.\\
Operator-source comparison & \cite{AD9}, \path{thm:v9-operator-source-comparison}. The five displacement \emph{operators}, with normalized Hilbert--Schmidt inner product, have Gram $I_5$ and admit the marked operator comparison on the frame cover. This is not the Gram of five operators applied to one vector.\\
Five-orbit recognition & \cite{AD9}, \path{thm:v9-five-parameter-gram}, \path{thm:v9-gram-unitary-test}, \path{thm:v9-finite-transfer-reconstruction}; also \cite{GT77}, \path{thm:five-parameter-tetrahedral}. These are conditional recognition theorems once the actual source and its ordinary unitary $S_4$ covariance are available.\\
Signature comparison & \cite{AD9}, \path{thm:v9-one-number-signature}. The relative central signature is reduced to a common-source expectation. Its actual value is not supplied by the separate Pauli packets.\\
Physical companion theories & \cite{Spacetime,Spectral,Duality,Variational}. The geometric carrier, differential data, represented action--multiplicity duality, and analytic continuum closure have separate hypotheses. None identifies the present theta source with a physical source.\\
\bottomrule
\end{longtable}

Recomputing the abstract tetrahedral lattice Gram would therefore duplicate \cite{AD9}. The unevaluated vector source is instead
\begin{equation}
 C_{\mathrm{tet}}:\quad y^2=x(x^4-x-1),\qquad
 J=\operatorname{Jac}(C_{\mathrm{tet}}),\qquad
 \HH_\Theta=H^0(J,\mathcal L_\kappa^{\otimes2}),
 \qquad \dim_\C\HH_\Theta=4,
 \label{eq:theta-input}
\end{equation}
with the ray
\begin{equation}
 [\eta_\kappa]=[\vartheta_\kappa^{\otimes2}],\qquad
 \xi_a=R_a\eta_\kappa,\qquad
 a\in\Sigma=\{t,v_1,v_2,v_3,v_4\}.
 \label{eq:actual-source}
\end{equation}
The ray is the object labeled \path{eq:v9-theta-source-ray} in \cite{AD9}. The five displaced vectors are its source family \path{eq:v9-theta-source-family}. They are not replaced here by a counting-space source, a lattice frame, or the normalized trace state.

\subsection{What is assumed, and what is proved here}
For the retained odd $\mu_2$ refinement in \cite{AD9}, the five nonzero singular classes satisfy
\begin{equation}
 q(a)=0,\qquad \omega(a,b)=1\quad(a\ne b).
 \label{eq:singular-input}
\end{equation}
Their lifts therefore obey $R_a^2=I$ and $R_aR_b=-R_bR_a$ in the central-character-minus-one representation. We work in a positive Hermitian metric in which that finite theta group is unitary. Such a metric exists by finite averaging, and is unique up to scale on its irreducible four-dimensional representation: if two invariant metrics differ by a positive operator, that operator commutes with the representation and hence is scalar. Fixing $\norm{\eta_\kappa}=1$ removes this scale.

This is the \emph{Heisenberg-unitary source branch}. If a proposed polarized metric is not invariant under these theta displacements, it is a different marked metric and must be compared independently; the metric is not allowed to be changed to force the desired Gram. The new conclusions are universal for every nonzero vector in the branch, so no unknown period or theta-constant value is needed for the obstruction below.

The finite Clifford algebra used in the proofs is standard background; binary symplectic/Clifford methods and spinorial Bloch geometry have substantial prior literature \cite{DehaeneDeMoor,Hasebe}. We give all identities needed here explicitly. The programme-specific advance is their application to \eqref{eq:actual-source}, the exact obstruction to the proposed five-parameter test, and the resulting concrete source-reconstruction and enlargement alternatives. No priority claim is made for the background Clifford identities.

\section{The pure five-displacement Gram}
\label{sec:pure-gram}

\subsection{A fixed convention and its quaternionic structure}
Relabel $t$ by $0$ and $v_i$ by $i$, and write $\gamma_a=R_a$.
Thus on $\HH\cong\C^4$,
\begin{equation}
 \gamma_a^*=\gamma_a,\qquad
 \gamma_a\gamma_b+\gamma_b\gamma_a=2\delta_{ab}I,
 \qquad 0\le a,b\le4.
 \label{eq:clifford}
\end{equation}
The source Gram and the Hermitian bivectors are
\begin{equation}
 G(\rho)_{ab}=\Tr(\rho\gamma_a\gamma_b),\qquad
 B_{ab}=i\gamma_a\gamma_b\quad(a<b),
 \label{eq:gram-rho}
\end{equation}
where $\rho\succeq0$, $\Tr\rho=1$. For a pure source $\rho=\ketbra{\eta}$,
$G(\rho)_{ab}=\ip{\gamma_a\eta}{\gamma_b\eta}$. Inner products are conjugate-linear in the first variable. Put
\begin{equation}
 n_a=\Tr(\rho\gamma_a),\qquad b_{ab}=\Tr(\rho B_{ab})\in\R.
 \label{eq:n-b}
\end{equation}
Then $G_{aa}=1$ and $G_{ab}=-ib_{ab}$ for $a<b$.

\begin{lemma}[Explicit quaternionic and trace identities]
\label{lem:quaternionic}
There is an antiunitary $\Jq$ on $\HH$ such that
\begin{equation}
 \Jq^2=-I,\qquad \Jq\gamma_a=\gamma_a\Jq,
 \qquad \Jq B_{ab}\Jq^{-1}=-B_{ab}.
 \label{eq:J-properties}
\end{equation}
It is unique up to a scalar phase. The sixteen Hermitian matrices
\begin{equation}
 I,\quad \gamma_0,\ldots,\gamma_4,\quad B_{ab}\ (a<b)
 \label{eq:operator-basis}
\end{equation}
are trace-orthogonal and each has squared Hilbert--Schmidt norm four. Hence
\begin{equation}
 \rho=\frac14\left(I+\sum_a n_a\gamma_a
                     +\sum_{a<b}b_{ab}B_{ab}\right).
 \label{eq:rho-expansion}
\end{equation}
Finally, for $\Phi(A)=\sum_{a=0}^4\gamma_aA\gamma_a$,
\begin{equation}
 \Phi(I)=5I,\qquad \Phi(\gamma_b)=-3\gamma_b,
 \qquad \Phi(B_{bc})=B_{bc}.
 \label{eq:Phi}
\end{equation}
\end{lemma}
\begin{proof}
Let $X,Y,Z$ be the standard Pauli matrices. A convention realizing
\eqref{eq:clifford} is
\begin{equation}
 \gamma_0=Z\otimes Z,\quad
 \gamma_1=X\otimes I_2,\quad
 \gamma_2=Y\otimes I_2,\quad
 \gamma_3=Z\otimes X,\quad
 \gamma_4=Z\otimes Y.
 \label{eq:pauli-realization}
\end{equation}
Let $K_0$ denote entrywise complex conjugation and set
$\Jq=C K_0$, where $C=\gamma_2\gamma_4$. The matrix $C$ is real,
$C^2=-I$, and $C\overline{\gamma_a}=\gamma_a C$ for all $a$.
Thus $\Jq$ is antiunitary, has square $-I$, and commutes with every
$\gamma_a$. Antiunitarity changes the scalar $i$ to $-i$, giving the last
identity in \eqref{eq:J-properties}.

These computations cover every realization in dimension four. Indeed, the
sixteen products of the four generators $\gamma_1,\ldots,\gamma_4$ are
trace-orthogonal: every nonempty such monomial anticommutes with one of the
four generators, so has trace zero, and products of distinct monomials
reduce, up to a phase, to nonempty monomials. They therefore span $M_4(\C)$.
Two four-generator realizations define a unital $*$-isomorphism of this
matrix algebra and are unitarily conjugate. The fifth generator is, up to
sign, the product of the other four; changing that sign does not affect
\eqref{eq:J-properties}. If $\Jq'$ has the same properties, the linear unitary
$\Jq'\Jq^{-1}$ commutes with $M_4(\C)$ and is scalar.

The same monomial argument, or the explicit Pauli realization, proves trace
orthogonality of \eqref{eq:operator-basis}. Expansion in that basis gives
\eqref{eq:rho-expansion}. For the last three formulas, conjugation fixes the
identity; one generator commutes with itself while the other four
anticommute with it; and conjugation negates $B_{bc}$ for the two indices
$b,c$ and fixes it for the other three. This gives respectively the
coefficients $5,-3,1$.
\end{proof}

\begin{lemma}[Pure-state quaternionic identities]
\label{lem:pure-identities}
For $\rho=\ketbra{\eta}$ with $\norm\eta=1$, set
$\rho^\sharp=\Jq\rho\Jq^{-1}=\ketbra{\Jq\eta}$. Then
\begin{align}
 \ip{\eta}{\Jq\eta}&=0,\label{eq:orthogonal-twin}\\
 P:=\rho+\rho^\sharp&=\frac12(I+n\cdot\gamma),
 &P^2&=P,\quad \rank P=2,\label{eq:ranktwo-P}\\
 \sum_a n_a^2&=1,
 &\sum_{a<b}b_{ab}^2&=2,\label{eq:pure-sums}\\
 D_b:=\frac12\sum_{a<b}b_{ab}B_{ab}&=\rho-\rho^\sharp.
 \label{eq:Db-pure}
\end{align}
\end{lemma}
\begin{proof}
For any antiunitary with square $-I$, its defining inner-product identity
applied to $\eta,\Jq\eta$ gives
$\ip{\Jq\eta}{\Jq^2\eta}=\ip{\Jq\eta}{\eta}$; the left side is the negative
of the right side. This proves \eqref{eq:orthogonal-twin}. Hence $P$ is a
rank-two orthogonal projection. By \cref{lem:quaternionic}, the vector terms
in \eqref{eq:rho-expansion} are fixed by $\Jq$, while the bivector terms change
sign. Adding and subtracting the two state expansions gives
\eqref{eq:ranktwo-P} and \eqref{eq:Db-pure}. Since
$(n\cdot\gamma)^2=\norm n^2 I$ and $(2P-I)^2=I$, one has $\norm n^2=1$.
Moreover, $(\rho-\rho^\sharp)^2=P$, whose trace is two. Trace orthogonality
in \eqref{eq:Db-pure} gives $\Tr D_b^2=\sum b_{ab}^2=2$.
\end{proof}

\begin{theorem}[Universal pure theta-source Gram spectrum]
\label{thm:pure-spectrum}
For every unit vector $\eta\in\C^4$ and every five-generator packet
\eqref{eq:clifford}, the source synthesis
\begin{equation}
 V_\eta:\C^5\longrightarrow\C^4,\qquad V_\eta e_a=\gamma_a\eta
 \label{eq:source-synthesis}
\end{equation}
has rank three, and
\begin{equation}
 \boxed{\spec G(\ketbra\eta)=\{2,2,1,0,0\}.}
 \label{eq:universal-spectrum}
\end{equation}
In particular,
\begin{equation}
 \Tr G=5,\qquad \Tr G^2=9,\qquad
 \norm G_{\HS}=3,
 \qquad G(G-I)(G-2I)=0.
 \label{eq:Gram-polynomial}
\end{equation}
For an unnormalized nonzero source, the Gram and its eigenvalues are multiplied
by $\norm\eta^2$.
\end{theorem}
\begin{proof}
The frame operator is $V_\eta V_\eta^*=\Phi(\rho)$. By
\eqref{eq:rho-expansion} and \eqref{eq:Phi},
\begin{equation}
 \Phi(\rho)=\rho+I-n\cdot\gamma
           =\rho+2(I-P).
 \label{eq:frame-operator}
\end{equation}
On $\C\eta$, $\C\Jq\eta$, and $P^\perp\HH$, this operator is respectively
$1$, $0$, and $2I_2$. The nonzero spectra of $V_\eta V_\eta^*$ and
$V_\eta^*V_\eta$ agree, including multiplicity. Since the latter acts on
$\C^5$, its spectrum is exactly \eqref{eq:universal-spectrum}. All remaining
formulas follow.
\end{proof}

\begin{corollary}[Evaluation of the inherited theta candidate without period data]
\label{cor:actual-theta}
For the actual source \eqref{eq:actual-source}, in the Heisenberg-unitary
metric and after normalization, \eqref{eq:universal-spectrum} holds
independently of the theta constants of $C_{\mathrm{tet}}$. The candidate is
therefore a rank-three vector-source system, not the rank-five
Hilbert--Schmidt operator-source system of \cite{AD9}.
\end{corollary}
\begin{proof}
The inherited odd quadratic refinement gives \eqref{eq:singular-input}.
Unitarity and $R_a^2=I$ make every $R_a$ Hermitian, and the commutator pairing
gives \eqref{eq:clifford}. The source ray is nonzero. Apply
\cref{thm:pure-spectrum}.
\end{proof}

\section{An exact replacement for the unevaluated five-parameter test}
\label{sec:recognition}
The correct unsymmetrized source data are the ten real numbers $b_{ab}$,
not five invariant parameters. They are not arbitrary: purity imposes
explicit polynomial equations.

\begin{theorem}[Polynomial recognition and reconstruction of a pure source]
\label{thm:pure-recognition}
Fix the operator frame \eqref{eq:clifford}. Given ten real numbers $b_{ab}$,
form
\begin{equation}
 D_b=\frac12\sum_{a<b}b_{ab}B_{ab},\qquad
 G_b{}_{aa}=1,\quad G_b{}_{ab}=-ib_{ab}\quad(a<b),
 \label{eq:candidate-G}
\end{equation}
with the lower triangle filled by Hermitian symmetry. The following are equivalent:
\begin{enumerate}[label=\textup{(\roman*)},leftmargin=2.7em]
\item $G_b$ is the displacement Gram of a normalized pure state;
\item $D_b^3=D_b$ and $\Tr D_b^2=2$.
\end{enumerate}
When these conditions hold, the pure density matrix is unique and equals
\begin{equation}
 \boxed{\rho_b=\frac12(D_b^2+D_b).}
 \label{eq:state-reconstruction}
\end{equation}
Thus the marked full displacement Gram determines the pure source ray.
No positive-semidefinite assumption on the initially proposed $G_b$ is
needed in clause \textup{(ii)}: positivity follows from the conclusion.
\end{theorem}
\begin{proof}
Necessity and \eqref{eq:state-reconstruction} follow from
$D_b=\rho-\rho^\sharp$, where the two rank-one projections are orthogonal.
Conversely, $D_b$ is Hermitian and $\Jq D_b\Jq^{-1}=-D_b$. The polynomial
identity forces its eigenvalues into $\{-1,0,1\}$. Its squared trace is two;
its trace is zero because every $B_{ab}$ is traceless. Therefore the nonzero
eigenvalues are one copy of $1$ and one copy of $-1$. Formula
\eqref{eq:state-reconstruction} is the rank-one projection onto the positive
eigenspace, and its quaternionic conjugate is $(D_b^2-D_b)/2$.

It remains to verify that its overlaps are the proposed ones. The matrix
$D_b^2$ is fixed by $\Jq$, whereas every $B_{ab}$ is negated. The real
Hilbert--Schmidt inner product of a fixed Hermitian matrix and a negated
Hermitian matrix is zero: conjugation by an antiunitary preserves that real
inner product. Hence $\Tr(D_b^2B_{ab})=0$. Trace orthogonality gives
$\Tr(D_b B_{ab})=2b_{ab}$, so
$\Tr(\rho_bB_{ab})=b_{ab}$. The resulting Gram is $G_b$.
For uniqueness, any realizing pure state satisfies the reconstruction formula.
\end{proof}

\begin{corollary}[Uniform marked-state comparison bound]
\label{cor:state-stability}
Let $G,G'$ be two actual normalized pure displacement Grams in the same
operator frame, and let $\rho,\rho'$ be their reconstructed states. Then
\begin{equation}
 \boxed{\norm{\rho-\rho'}_{\HS}
       \le\frac{3}{2\sqrt2}\norm{G-G'}_{\HS}.}
 \label{eq:state-stability}
\end{equation}
After an independently established operator-frame identification, this is a
quantitative common-source comparison. It does not establish that identification.
\end{corollary}
\begin{proof}
The diagonal Gram entries agree, so trace orthogonality gives
$\norm{D_b-D_{b'}}_{\HS}=\norm{G-G'}_{\HS}/\sqrt2$.
The recognized matrices $D_b,D_{b'}$ have operator norm one. Therefore
\[
 \norm{D_b^2-D_{b'}^2}_{\HS}
 \le(\norm{D_b}_{\op}+\norm{D_{b'}}_{\op})
       \norm{D_b-D_{b'}}_{\HS}
 \le2\norm{D_b-D_{b'}}_{\HS}.
\]
Apply \eqref{eq:state-reconstruction} and the triangle inequality.
\end{proof}

\begin{remark}[What has and has not been numerically evaluated]
The ten individual overlaps of the squared theta section have not been
numerically computed in this version. What has been evaluated exactly is
their universal spectral and purity constraint, which suffices to decide the
ordinary tetrahedral comparison negatively. The reconstruction theorem also
specifies the correct finite data to acquire for a subsequent \emph{projective
or signed} common-frame comparison. No theta period matrix or period
normalization has been invented.
\end{remark}

\section{The ordinary pointed tetrahedral branch is obstructed}
\label{sec:obstruction}

\subsection{A phase-independent obstruction}
Let $P_\pi$ be the permutation representation of $S_4$ on $\C^5$ which fixes
$e_0$ and permutes $e_1,\ldots,e_4$. Let $\Mtet$ be the real vector space of
Hermitian matrices commuting with every $P_\pi$, and let
\begin{equation}
 \AR(A)=\frac1{24}\sum_{\pi\in S_4}P_\pi^*AP_\pi.
 \label{eq:Reynolds}
\end{equation}
This is the Hilbert--Schmidt orthogonal projection onto $\Mtet$.
The inherited five-orbit theorem gives
\begin{equation}
 \C^5\cong\one\oplus\one\oplus W,\qquad \dim W=3,
 \qquad
 \End_{S_4}(\C^5)\cong M_2(\C)\oplus\C I_W.
 \label{eq:tet-decomp}
\end{equation}
In particular, every matrix in $\Mtet$ has an eigenvalue of multiplicity at
least three. This is a consequence of an inherited representation theorem,
not a new derivation of the five-parameter normal form.

\begin{theorem}[No ordinary pointed $S_4$ linearization of the pure theta source]
\label{thm:no-S4}
For the source in \cref{thm:pure-spectrum}, no choices of five phases
$z_a\in\C$, $|z_a|=1$, make the Gram of $z_a\gamma_a\eta$ belong to
$\Mtet$. Consequently there is no ordinary unitary $S_4$ action on its source
span which fixes its nonzero temporal vector and permutes its four residual
vectors exactly, after any fixed rephasing.
\end{theorem}
\begin{proof}
Rephasing conjugates the Gram by a diagonal unitary and preserves its spectrum.
By \eqref{eq:universal-spectrum}, all eigenvalue multiplicities are at most
two. This contradicts \eqref{eq:tet-decomp}. The Gram of an ordinary
unitary-covariant pointed source would belong to $\Mtet$.
\end{proof}

\begin{scope}[Precise correction to the proposed route]
The five-parameter theorem remains valid. Its ordinary unitary-covariance
premise fails for the specific pure displacement source
\eqref{eq:actual-source} in the Heisenberg-unitary metric. The inherited
$S_4$ action on arithmetic torsion labels and the integral $K_4$ lattice
remain untouched. A Galois action on algebraic labels, a projective action on
rays, and an ordinary complex-unitary action on oriented source vectors are
different assertions. This result rules out only the last assertion for
this source; it does not rule out a signed/projective comparison, a different
marked source, or arithmetic--gauge duality in another category.
\end{scope}

\subsection{A sharp real-lift gap}

\begin{theorem}[Exact distance to the pointed invariant-Gram space]
\label{thm:sharp-gap}
For a normalized pure displacement source with retained real lifts, define
\begin{equation}
 Q=\frac12\sum_{i=1}^4 B_{0i},\qquad
 \kappa=\Tr(\rho Q),\qquad
 \beta=\frac14\sum_{i=1}^4b_{0i}=\frac\kappa2.
 \label{eq:kappa}
\end{equation}
Then $Q^*=Q$, $Q^2=I$, and
\begin{equation}
 \AR(G)=
 \begin{pmatrix}
 1&-i\beta&-i\beta&-i\beta&-i\beta\\
 i\beta&1&0&0&0\\
 i\beta&0&1&0&0\\
 i\beta&0&0&1&0\\
 i\beta&0&0&0&1
 \end{pmatrix}.
 \label{eq:averaged-Gram}
\end{equation}
Moreover,
\begin{equation}
 \boxed{
 \dist_{\HS}(G,\Mtet)^2
 =\norm{G-\AR(G)}_{\HS}^2
 =4-2\kappa^2\ge2.}
 \label{eq:sharp-gap}
\end{equation}
The bound is attained. In the normalization $\norm G_{\HS}=3$, the relative
distance is at least $\sqrt2/3$. All choices of signs of the real theta lifts
obey the same bound.
\end{theorem}
\begin{proof}
The four $B_{0i}$ are Hermitian involutions and anticommute pairwise. Their
half-sum therefore squares to the identity, implying $|\kappa|\le1$.
Every off-diagonal Gram entry is purely imaginary. Permutation averaging
pairs each ordered spatial pair with its reversal, so the averaged spatial
off-diagonal entry is zero; the temporal entries average to $-i\beta$.
This gives \eqref{eq:averaged-Gram}. Orthogonal projection and
\eqref{eq:Gram-polynomial} now give
\[
 \norm{G-\AR(G)}_{\HS}^2
 =9-\bigl(5+8\beta^2\bigr)=4-2\kappa^2.
\]
A unit eigenvector of $Q$ has $|\kappa|=1$, so the bound is sharp. Changing
any lift sign preserves all Clifford identities and the same argument.
\end{proof}

\begin{proposition}[Phase-independent uniform gap]
\label{prop:phase-gap}
For every diagonal unitary $D$ and every $M\in\Mtet$,
\begin{equation}
 \boxed{\norm{D^*GD-M}_{\HS}^2\ge\frac23.}
 \label{eq:phase-gap}
\end{equation}
Thus arbitrary phase optimization cannot produce even an asymptotic ordinary
pointed source identification within this normalized pure-source class.
\end{proposition}
\begin{proof}
For Hermitian matrices $A,M$ with increasingly ordered eigenvalues $\lambda_j,
\mu_j$, one has
\begin{equation}
 \norm{A-M}_{\HS}^2\ge\sum_j(\lambda_j-\mu_j)^2.
 \label{eq:HW}
\end{equation}
For completeness, write $\Tr(AM)=\sum_{jk}\lambda_j\mu_k
|\ip{u_j}{v_k}|^2$ in eigenbases. The overlap-squared matrix is doubly
stochastic. A finite doubly stochastic matrix is a convex combination of
permutation matrices; the rearrangement inequality bounds the displayed
sum by $\sum_j\lambda_j\mu_j$. Expanding the squared norm proves
\eqref{eq:HW}.

The spectrum of $D^*GD$ is $(0,0,1,2,2)$. By
\eqref{eq:tet-decomp}, at least three consecutive eigenvalues of $M$ are
one common number $\mu$. The possible consecutive triples of the former
spectrum are $(0,0,1)$, $(0,1,2)$, and $(1,2,2)$. Their minimum squared
distances to a constant triple are respectively $2/3$, $2$, and $2/3$.
Discarding the other nonnegative terms in \eqref{eq:HW} proves the result.
\end{proof}

\begin{corollary}[Noise cannot conceal the obstruction below the explicit margin]
\label{cor:noise}
If $\norm{\widehat G-G}_{\HS}\le\varepsilon$, then
\[
 \dist_{\HS}(\widehat G,\Mtet)\ge\sqrt2-\varepsilon
\]
in the retained real-lift convention. Allowing arbitrary phases gives
\[
 \inf_{\substack{D\ \mathrm{diagonal\ unitary}\\M\in\Mtet}}\norm{D^*\widehat GD-M}_{\HS}
 \ge\sqrt{2/3}-\varepsilon,
\]
where $D$ ranges over diagonal unitaries. The right sides are useful when
positive.
\end{corollary}
\begin{proof}
Apply the triangle inequality to \cref{thm:sharp-gap,prop:phase-gap};
diagonal unitary conjugation preserves the error norm.
\end{proof}

\subsection{The operator signs cannot all be discarded}
\begin{proposition}[Central-volume constraint on a permutation lift]
\label{prop:volume}
Suppose a unitary $U$ normalizing the five Clifford rays satisfies
\[
 U\gamma_aU^*=\epsilon_a\gamma_{\pi(a)},\qquad
 \epsilon_a\in\{\pm1\},\qquad \pi(0)=0.
\]
Then
\begin{equation}
 \boxed{\prod_{a=0}^4\epsilon_a=\operatorname{sgn}(\pi).}
 \label{eq:sign-product}
\end{equation}
In particular an odd spatial permutation cannot be implemented with all five
lift signs positive on the irreducible four-dimensional module.
\end{proposition}
\begin{proof}
The product $\Omega=\gamma_0\gamma_1\gamma_2\gamma_3\gamma_4$ commutes with
all generators and is therefore scalar. It is a nonzero Hermitian
involution, so $\Omega=\pm I$. Conjugating this ordered product by $U$ gives
$(\prod_a\epsilon_a)\operatorname{sgn}(\pi)\Omega$. Equality to $\Omega$
proves \eqref{eq:sign-product}.
\end{proof}

This is an operator-level restriction, separate from
\cref{thm:no-S4}. A scalar source Gram can accidentally have ordinary
permutation symmetry without those source isometries intertwining the full
Clifford action. The next section constructs exactly such a distinction.

\section{Minimal mixed-source alternatives and their information cost}
\label{sec:mixed}

\subsection{All ordinary invariant displacement Grams}

\begin{theorem}[Invariant mixed-state Gram classification]
\label{thm:mixed-classification}
For a density matrix $\rho$ on the same irreducible four-dimensional module,
$G(\rho)\in\Mtet$ if and only if
\begin{equation}
 \Tr(\rho B_{ij})=0\quad(1\le i<j\le4),\qquad
 \Tr(\rho B_{0i})=\beta\quad(1\le i\le4)
 \label{eq:mixed-conditions}
\end{equation}
for one real $\beta$. The complete set of possible Grams is
\eqref{eq:averaged-Gram} with
\begin{equation}
 |\beta|\le\frac12,
 \qquad
 \spec G=\{1,1,1,1+2|\beta|,1-2|\beta|\}.
 \label{eq:mixed-spectrum}
\end{equation}
Every such value occurs; one explicit realizing state is
\begin{equation}
 \rho_\beta=\frac14(I+2\beta Q).
 \label{eq:rhobeta}
\end{equation}
Every state realizing an invariant Gram obeys
\begin{equation}
 \boxed{\norm\rho_{\op}\le\frac12,\qquad
        \Tr\rho^2\le\frac12,\qquad \rank\rho\ge2.}
 \label{eq:purity-cost}
\end{equation}
Its von Neumann entropy is at least $\log2$.
\end{theorem}
\begin{proof}
The off-diagonal entries are purely imaginary, while ordinary $S_4$
invariance makes all spatial off-diagonal entries equal to one real number.
That number must be zero. Temporal entries must be equal. This proves the
first equivalence. The invariant Gram is therefore exactly
\eqref{eq:averaged-Gram}. It has the three-dimensional spatial-augmentation
block $I_3$ and the remaining block
$\begin{psmallmatrix}1&-2i\beta\\2i\beta&1\end{psmallmatrix}$.
Positivity gives \eqref{eq:mixed-spectrum}. Conversely, \eqref{eq:rhobeta}
is positive with trace one when $|\beta|\le1/2$, since $Q$ has eigenvalues
$\pm1$, each twice. Trace orthogonality verifies all moments in
\eqref{eq:mixed-conditions}.

For the stronger purity assertion, $\gamma_0$ is, up to sign, the product
of the four spatial generators. Its two eigenspaces have dimension two.
The eight matrices
\[
 I,\ \gamma_0,\ B_{ij}\quad(1\le i<j\le4)
\]
span the entire block-diagonal algebra for this decomposition: they commute
with $\gamma_0$, are independent, and that algebra has dimension eight.
The six vanishing moments in \eqref{eq:mixed-conditions} consequently imply
that, in this decomposition,
\begin{equation}
 \rho=\begin{pmatrix}aI_2&C\\C^*&bI_2\end{pmatrix},
 \qquad a,b\ge0,\qquad a+b=\frac12.
 \label{eq:chiral-block}
\end{equation}
Positivity gives $\norm C_{\op}^2\le ab$. When either diagonal coefficient
is zero, positivity forces $C=0$; otherwise the inequality is the elementary
Schur-complement condition. Singular-value decomposing $C$ reduces
\eqref{eq:chiral-block} to two matrices
$\begin{psmallmatrix}a&s_j\\s_j&b\end{psmallmatrix}$.
Their largest eigenvalues are at most
\[
 \frac{a+b+\sqrt{(a-b)^2+4ab}}2=a+b=\frac12.
\]
This proves the operator-norm bound, from which
$\Tr\rho^2\le\norm\rho_{\op}\Tr\rho\le1/2$ and $\rank\rho\ge2$ follow.
For eigenvalues $\lambda_j\le1/2$, the entropy satisfies
$-\sum_j\lambda_j\log\lambda_j\ge\sum_j\lambda_j\log2=\log2$.
\end{proof}

\subsection{A canonical rank-two completion of the actual pure source}

\begin{construction}[Quaternionic completion]
\label{con:completion}
For the actual normalized pure theta state $\rho$, define
\begin{equation}
 \rho_+=\frac12(\rho+\rho^\sharp)
        =\frac14(I+n\cdot\gamma).
 \label{eq:completion}
\end{equation}
The conjugate state $\rho^\sharp$ is independent of the phase ambiguity in
$\Jq$ and of the phase of $\eta$. This is a new state associated to the old
one, not an assertion that it physically occurred.
\end{construction}

\begin{theorem}[Exact completion, minimal rank, and lost pure-state fibre]
\label{thm:completion}
The state \eqref{eq:completion} has the following properties.
\begin{enumerate}[label=\textup{(\roman*)},leftmargin=2.7em]
\item It has rank two and nonzero eigenvalues $1/2,1/2$, and
\begin{equation}
 \boxed{G(\rho_+)=I_5.}
 \label{eq:completed-G}
\end{equation}
Its rank is the minimum compatible with any ordinary invariant displacement
Gram.
\item It is the unique Hilbert--Schmidt closest $\Jq$-invariant density
matrix to $\rho$. Its distances from $\rho$ are
\begin{equation}
 \norm{\rho-\rho_+}_{\HS}=\frac1{\sqrt2},\qquad
 \frac12\norm{\rho-\rho_+}_1=\frac12.
 \label{eq:completion-distance}
\end{equation}
\item Writing $P=2\rho_+$, the pure source rays with the same completion are
exactly $\mathbb P(P\HH)\cong\mathbb{CP}^1$.
\item A purification uses one two-dimensional auxiliary factor:
\begin{equation}
 \Psi_\eta=\frac1{\sqrt2}
 (\eta\otimes e_0+\Jq\eta\otimes e_1).
 \label{eq:purification}
\end{equation}
The five vectors $(\gamma_a\otimes I_2)\Psi_\eta$ are orthonormal. An auxiliary
factor of dimension less than two cannot purify any state with an ordinary
invariant displacement Gram.
\end{enumerate}
\end{theorem}
\begin{proof}
The two pure summands in \eqref{eq:completion} are orthogonal by
\cref{lem:pure-identities}, so the rank and spectrum follow. Their bivector
moments have opposite signs, proving \eqref{eq:completed-G}. Minimality is
\cref{thm:mixed-classification}.

On the real Hilbert space of Hermitian matrices, conjugation by $\Jq$ is an
orthogonal involution. Averaging with that involution is its orthogonal
projection onto the fixed subspace. The projected matrix \eqref{eq:completion}
is positive with trace one, so it is also the unique closest point in the
subset of fixed density matrices. The difference is
$(\rho-\rho^\sharp)/2$, with eigenvalues $1/2,-1/2,0,0$, proving
\eqref{eq:completion-distance}.

The plane $P\HH$ is $\Jq$-invariant. For every unit vector $\zeta$ in it,
$\zeta,\Jq\zeta$ form an orthonormal basis, so their averaged density is
$P/2$. Conversely, a pure state whose quaternionic completion is $P/2$ has
its range contained in $P\HH$. This proves the projective-line fibre.

Taking the partial trace of \eqref{eq:purification} gives $\rho_+$, and its
displacement Gram is therefore \eqref{eq:completed-G}. Any purification of a
density matrix has Schmidt rank equal to the rank of that density, which
is at least two by \eqref{eq:purity-cost}. Equivalently, the reduced density
of a vector in $\C^4\otimes\C^m$ has rank at most $m$.
\end{proof}

\begin{scope}[This is not the old operator-source calculation]
The identity Gram of the full normalized trace state was already established
in \cite{AD9}. The new statement is that the particular theta pure state has
a canonical \emph{rank-two} completion producing that Gram, with a sharp
minimum-rank cost, a fixed nonzero distance from the original state, and a
$\mathbb{CP}^1$ ambiguity. It is not legitimate to substitute this completion
for the original source without recording the change.
\end{scope}

\begin{proposition}[Quaternionic averaging is not an automatic physical channel]
\label{prop:not-CP}
The state transformation in \eqref{eq:completion}, extended complex-linearly
to matrices as
\begin{equation}
 \mathcal T(A)=\frac12(A+C A^{\mathsf T}C^*),
 \qquad C=\gamma_2\gamma_4
 \label{eq:positive-notCP}
\end{equation}
in the convention \eqref{eq:pauli-realization}, is positive and trace
preserving but is not completely positive. Hence this mathematical
completion is not a universally available operation on an unknown source.
\end{proposition}
\begin{proof}
Transpose and unitary conjugation preserve positivity and trace, proving the
first assertion. Let $\Omega=\sum_{j=1}^4e_j\otimes e_j$ and let $F$ be the
swap on $\C^4\otimes\C^4$. Applying $\id\otimes\mathcal T$ to
$\ketbra\Omega$ gives
\[
 \frac12\left(\ketbra\Omega+(I\otimes C)F(I\otimes C^*)\right).
\]
The second term inside parentheses has a negative eigenspace of dimension
six, the conjugate of the antisymmetric tensor space. Its intersection
with $\Omega^\perp$ has dimension at least five. On every nonzero vector
in that intersection the displayed quadratic form is strictly negative.
Thus the map fails positivity on this bipartite input and is not completely
positive. Preparation of the particular state \eqref{eq:completion} may of
course use additional independently specified information or resources;
that is a different occurrence statement.
\end{proof}

\subsection{The identity source exposes the information erased by \texorpdfstring{$I_5$}{the identity Gram}}

\begin{theorem}[Complete invariant-Gram fibre and an identity-inserted test]
\label{thm:identity-read}
The density matrices whose displacement Gram is exactly $I_5$ are
\begin{equation}
 \boxed{\rho_n=\frac14(I+n\cdot\gamma),\qquad n\in\R^5,
              \norm n\le1.}
 \label{eq:Bloch-fibre}
\end{equation}
Their eigenvalues are $(1+\norm n)/4$ and $(1-\norm n)/4$, each twice.
Appending the identity to the five displacement sources in the
Hilbert--Schmidt realization $S_A=A\sqrt{\rho_n}$ gives the complete
six-source Gram
\begin{equation}
 \boxed{\mathcal G_1(\rho_n)=
 \begin{pmatrix}1&n^{\mathsf T}\\n&I_5\end{pmatrix}.}
 \label{eq:identity-Gram}
\end{equation}
Thus this inserted panel recovers all five previously invisible coordinates.
For the completion of a pure theta source, $\norm n=1$ and
$\rank\mathcal G_1=5$.
\end{theorem}
\begin{proof}
The condition $G=I_5$ is equivalent to $b_{ab}=0$ for every pair.
The operator expansion \eqref{eq:rho-expansion} then gives
\eqref{eq:Bloch-fibre}. Since $(n\cdot\gamma)^2=\norm n^2I$ and
$n\cdot\gamma$ is traceless, its eigenvalues are $\pm\norm n$, each twice
when $n\ne0$; the zero case is immediate. Positivity is exactly
$\norm n\le1$. The inner product with the identity source is
$\Tr(\rho_n\gamma_a)=n_a$, giving \eqref{eq:identity-Gram}. Its eigenvalues
are $1+\norm n$, $1-\norm n$, and four copies of $1$. Apply
\eqref{eq:pure-sums} for the final assertion.
\end{proof}

\begin{corollary}[A first held-out ordinary covariance test]
\label{cor:heldout}
For the six-source panel \eqref{eq:identity-Gram}, let $S_4$ fix the identity
and temporal labels and permute the four residual labels. It is invariant
if and only if
\begin{equation}
 n_1=n_2=n_3=n_4.
 \label{eq:next-test}
\end{equation}
With $\bar n=(n_1+\cdots+n_4)/4$, its squared Hilbert--Schmidt distance to
that invariant-matrix space is
\begin{equation}
 2\sum_{i=1}^4(n_i-\bar n)^2.
 \label{eq:identity-defect}
\end{equation}
The fact that the five-displacement marginal always equals $I_5$ does not
supply \eqref{eq:next-test}.
\end{corollary}
\begin{proof}
Only the identity-to-residual row in \eqref{eq:identity-Gram} can fail
covariance. Its invariant projection replaces the four entries by their
mean, and the transpose row contributes the same squared norm.
\end{proof}

\begin{theorem}[Exact first-word completion cost]
\label{thm:word-cost}
Use the source Hilbert space with inner product
$\ip{A}{B}_\rho=\Tr(\rho A^*B)$ modulo its null space, and let $d_{\le r}$
be the rank of the Gram of all words of length at most $r$ in the five
generators. Then
\begin{equation}
 \begin{array}{c|cc|c}
 \text{state}&d_{\le1}&d_{\le2}&\text{new directions at length two}\\\hline
 \rho=\ketbra\eta&3&4&1\\
 \rho_+=\tfrac12(\rho+\rho^\sharp)&5&8&3
 \end{array}.
 \label{eq:word-ranks}
\end{equation}
In both rows, words of length two already span the complete represented
matrix algebra before passing to the source quotient.
\end{theorem}
\begin{proof}
For a pure source, the five-generator span has rank three by
\cref{thm:pure-spectrum}. Since $(n\cdot\gamma)\eta=\eta$, the identity
source adds no new direction. For the rank-two completion,
\cref{thm:identity-read} gives $d_{\le1}=5$.

The words $I,\gamma_a,\gamma_a\gamma_b$ span $M_4(\C)$ by
\cref{lem:quaternionic}. The quotient map is represented isometrically by
$A\mapsto A\sqrt\rho$. If $\rho$ has rank $r$, this image is the full space
of maps from its $r$-dimensional support to $\C^4$, of dimension $4r$.
For the two states in question, $r=1$ and $r=2$. Subtracting the first-word
ranks gives the last column.
\end{proof}

\section{Consequences for the arithmetic-duality programme}
\label{sec:consequences}

\subsection{What this iteration actually decides}
The ordinary pointed-$S_4$ comparison proposed for the pure theta source is
not an unevaluated positive branch after this calculation. In the
Heisenberg-unitary source category it is excluded by
\cref{thm:no-S4}, with uniform quantitative margins
\cref{thm:sharp-gap,prop:phase-gap}. The conclusion does not require a
choice of theta periods and does not become weaker at larger arithmetic
cutoffs: each normalized instance belongs to the same obstructed
four-dimensional source class.

This does not contradict the already matched integral lattice Gram, the
operator-source Gram, or the marked operator equivalence on the quartic
frame cover. Those are different sources and different statements. In
particular, a rank-three first-word span is not a physical generation
carrier: \cref{thm:word-cost} shows that it is not even closed under the
complete matrix action.

\subsection{The next noncircular routes}
Two mathematically distinct continuations remain available.

\paragraph{Retain the actual pure theta source.}
Use the inherited marked operator equivalence on the frame cover and retain
its signs/cocycle. Compare the complete ten real bivector overlaps of the
arithmetic and renewal \emph{pure} sources, without first imposing an ordinary
$S_4$ Gram. \Cref{thm:pure-recognition,cor:state-stability} give an exact
recognition algorithm and a stable state-comparison bound in that common
frame. The actual theta coefficients, renewal source state, and arithmetic
transfer still have to be independently supplied or computed. A semilinear
Galois descent problem cannot be replaced by an ordinary unitary action
without a separate theorem.

\paragraph{Demand an ordinary pointed source.}
Then the reduced state cannot be pure. \Cref{thm:mixed-classification}
classifies the possible displacement Grams and proves the minimum mixed-state
cost. \Cref{thm:completion} gives a canonical mathematical rank-two source
attached to the original theta ray, but \cref{prop:not-CP} prevents treating
it as a free universal physical channel. Its occurrence, the
identity-inserted moments, and the three missing length-two source directions
must be retained. The scalar identity Gram alone does not determine the
state, much less its dynamics or its tensor/holonomy structure.

Neither route supplies a Frobenius--Wilson comparison yet. The finite
Tannaka and interface identities of \cite{GT77} should be used only after a
common marked operator/state packet has passed the relevant tests. Repeating
those generic reconstruction theorems before the source problem is resolved
would not advance this programme.

\subsection{A concrete next-iteration target}
A next note should keep \eqref{eq:actual-source} fixed and work on the signed
common-frame comparison. Its minimal tasks are to determine a consistent
level-two theta coordinate frame and its arithmetic descent, express
$\vartheta_\kappa^2$ in that frame with a stated invariant metric, and compute
the ten $b_{ab}$ rather than fitting five orbit parameters. The polynomial
certificate \eqref{eq:state-reconstruction} then checks the resulting
source. A separately reconstructed renewal pure state can be compared under
the already inherited operator intertwiner. If the intended renewal source
is mixed or lives on another carrier, that change is to be stated before
comparing moments.

The relative signature expectation remains a separate common-source row.
The present Gram obstruction does not evaluate it, and its vanishing would
not undo the pure-source symmetry obstruction.

\section{Verification and V1 proof-status ledger}
\label{sec:status}

\subsection{Exact computational checks}
The proofs above are symbolic and hold for all admissible sources. A
separate exact-rational complex-matrix audit was also executed using
SymPy, with no floating-point tolerance. It verifies the Clifford relations,
trace orthogonality of the sixteen operator coordinates, the antiunitary
identities, all sixteen conjugation-superoperator identities, the pure
reconstruction polynomial, the first two word ranks, and the invariant
mixed-state family.

Three sample rays were checked in the explicit Pauli convention:
$e_1$, $(1,1+i,2-i,-1+2i)^{\mathsf T}$, and $(I+Q)e_1$, normalized through
their exact rank-one density matrices. Their squared ordinary symmetry
defects are respectively
\begin{equation}
 4,\qquad \frac{668}{169},\qquad 2,
\end{equation}
with $\kappa=0,-2/13,1$. In every case the Gram characteristic polynomial is
$z^2(z-1)(z-2)^2$. The unnormalized Choi matrix of
\eqref{eq:positive-notCP} has eigenvalues $3/2$ once, $-1/2$ five times, and
$1/2$ ten times. These computations check conventions and explicit witnesses;
they are not substitutes for the universal proofs or a formal proof-assistant
verification.

\subsection{Proved in this version}
\begin{enumerate}[label=\textup{(P\arabic*)},leftmargin=3em]
\item The actual pure five-displacement theta candidate has universal Gram
spectrum $(2,2,1,0,0)$ in the Heisenberg-unitary metric.
\item Its full ten imaginary overlaps satisfy a terminating polynomial
pure-source criterion and reconstruct its density matrix uniquely, with
an explicit comparison bound.
\item No fixed phase choices make this source an ordinary pointed $S_4$
source. The real-lift squared obstruction has sharp floor two, and arbitrary
rephasing retains a squared floor of at least $2/3$.
\item An operator lift of an odd residual permutation necessarily carries a
nontrivial product of ray signs.
\item All ordinary invariant mixed-state displacement Grams are classified;
every realizing state has largest eigenvalue at most $1/2$, purity at most
$1/2$, and rank at least two.
\item Quaternionic completion is canonical as a mathematical state
construction, attains the minimum rank, and has an explicit
$\mathbb{CP}^1$ pure-state fibre. It is not a universally completely positive
operation on an unknown input.
\item Identity-inserted moments recover the five coordinates concealed by an
identity displacement Gram; full length-two completion costs one additional
direction on the pure branch and three on the rank-two branch.
\end{enumerate}

\subsection{Inherited and not redone}
Normalized arithmetic loading and its cutoff calculus; the pointed
arithmetic torsion packet; the integral arithmetic--$K_4$ lattice and its
static Gram; the refined operator-source comparison; the generic five-orbit
Gram/Krylov recognition theorems; the separate signature criterion; finite
Tannaka/interface machinery; and the four companion physical constructions.

\subsection{Still unresolved}
The ten numerical overlaps of the specified squared theta section in a fully
marked complex realization; its signed/projective arithmetic descent; an
independently reconstructed renewal state on the common frame; the actual
central signature expectation; a naturally selected arithmetic transfer and
its complete word moments; and any source-complete Frobenius--Wilson,
cofinal, or physical matter comparison. In particular, no critical-line,
prime-correlation, Standard-Model-parameter, or Einstein-continuum statement
is inferred here.

\subsection{Version integrity}
The logical endpoint of V1 is a decided obstruction for one specific pure
comparison category, together with constructive state-level alternatives.
It is not a claim that arithmetic--geometry duality is impossible. Future
work is to append to this record, not repeat the loading or lattice proofs
and not restore the excluded ordinary pure-source branch without explicitly
changing and justifying its hypotheses.

% ====================== END IMMUTABLE V1 MATHEMATICAL BODY ======================
% APPEND NEW ITERATIONS HERE, preserving the complete V1 body above.

% ===================== BEGIN IMMUTABLE V2 MATHEMATICAL BODY =====================
\clearpage
\section{V2 continuation: signed descent before symmetry reduction}
\label{sec:v2-context}
\noindent\textbf{20 September 2026.}

V1 excluded ordinary pointed-$S_4$ covariance of the actual pure theta
source, but explicitly left a signed/projective comparison open. The present
continuation tests that alternative instead of repeating the arithmetic
loading, lattice, or five-parameter recognition theorems. It distinguishes
three statements: an action on the five Clifford \emph{rays}, a coherent
projective action on their four-dimensional complex module, and preservation
of the particular squared theta-section ray. The first two do not imply the
third.

The main new finite result is a complete classification of coherent signed
lifts in this fixed-fibre unitary category. Up to fixed changes of the real
lift signs there are two classes. Both admit explicit projective unitary
implementers, and neither has a fixed pure source ray. More strongly, the
$S_3$ stabilizer of a vertex on the quartic frame cover has no fixed pure ray
in either class. Thus passing to that cover solves the inherited
\emph{operator} frame problem, but does not solve this stronger fixed-fibre
\emph{source} problem. The resulting obstruction has short two-generator
witnesses and exact quantitative versions.

A constructive arithmetic branch is also advanced. The actual squared
section is expressed in a marked level-two theta basis using four
second-order theta constants. Explicit Gaussian-tail, period-error, and
normalization estimates turn this expression into a certified finite
acquisition procedure for all ten overlaps in V1. This does not assume a
five-parameter symmetry or substitute a generic spinor for the arithmetic
section.

\paragraph{Dependency and scope audit.}
The source \eqref{eq:actual-source}, the quartic cover, and the marked
operator comparison are inherited from \cite{AD9}, in particular
\path{prop:v9-arithmetic-source-candidate},
\path{thm:v9-arithmetic-frame-cover}, and
\path{thm:v9-operator-source-comparison}. V1 supplies
\cref{lem:quaternionic,thm:pure-recognition,thm:completion}. The Schur double
covers of symmetric groups and theta addition formulas are classical
\cite{ReichsteinShukla,DLMFTheta}; they are not claimed as discoveries here.
The new programme-level content is the signed-lift/source classification,
its frame-cover obstruction, the precise completely positive repair, and
the certified acquisition of the previously unevaluated theta source.
Neither a semilinear Galois action nor a family over varying complex
embeddings is silently identified with a complex-linear action on one
Hermitian fibre.

\subsection{The exact comparison category}
Put $\Sigma=\{0,1,2,3,4\}$ and let $S_4$ fix $0$ and permute the other
labels. Write $p(g)\in\F_2$ for permutation parity, and retain the
Hermitian Clifford generators \eqref{eq:clifford}.

\begin{definition}[Coherent signed projective normalizer]
\label{def:v2-normalizer}
A coherent signed projective normalizer is a family $U_g\in U(4)$ with
$U_e=I$ such that
\begin{align}
 U_gU_h&=c(g,h)U_{gh},\\
 U_g\gamma_0U_g^*&=\epsilon_0(g)\gamma_0,\\
 U_g\gamma_iU_g^*&=\epsilon_i(g)\gamma_{g(i)},\qquad 1\le i\le4,
 \label{eq:v2-normalizer}
\end{align}
where $|c(g,h)|=1$ and $\epsilon_a(g)\in\{\pm1\}$. A pure source is
fixed projectively when $U_g\eta\in\C\eta$ for every $g$, equivalently
$U_g\rho U_g^*=\rho$ for $\rho=\ketbra\eta$.
\end{definition}

Hermiticity forces the ray phases in the conjugation equations to be real
signs. Scalar phases in $U_g$ remain allowed. A fixed reorientation
$\gamma_i\mapsto d_i\gamma_i$, $d_i\in\{\pm1\}$, is an allowed change
of lift convention; it is not a change of source state. The induced signed
matrix $r(g)$ on $\R^5$ is an ordinary representation because scalar
multipliers disappear under conjugation. By V1's central-volume constraint,
$\det r(g)=1$.

\begin{scope}[What this category excludes]
The definition requires the product defect of two implementers to be scalar.
It does not include a chosen section whose product defect is a noncentral
Heisenberg operator, an antiunitary action, a nonunitary semilinear Galois
map, or transport between different complex fibres. Those are distinct
possible continuations, not instances ruled out by the theorem below.
\end{scope}

\subsection{There are exactly two coherent signed actions}

\begin{theorem}[Classification of signed pointed tetrahedral lifts]
\label{thm:v2-signed-classification}
Every action induced by \cref{def:v2-normalizer} is, after fixed real
reorientations of the four spatial generators, exactly one of
\begin{align}
 r_0(g)e_0&=(-1)^{p(g)}e_0,&
 r_0(g)e_i&=e_{g(i)},\\
 r_1(g)e_0&=(-1)^{p(g)}e_0,&
 r_1(g)e_i&=(-1)^{p(g)}e_{g(i)}.
 \label{eq:v2-two-actions}
\end{align}
The two classes are inequivalent under such reorientations. Conversely,
both are implemented projectively by unitaries. Thus the remaining
coherent real-sign ambiguity is a two-class structural choice, not a
continuous source-fitting freedom.
\end{theorem}
\begin{proof}
Write $\epsilon_i(g)=(-1)^{s(g,i)}$. The multiplication law for the signed
matrices is exactly
\begin{equation}
 s(gh,i)=s(h,i)+s(g,h(i))\quad\text{in }\F_2.
 \label{eq:v2-sign-cocycle}
\end{equation}
Fix the spatial label $4$ and its stabilizer $H\cong S_3$. The map
$h\mapsto s(h,4)$ is a homomorphism $H\to\F_2$. Since the transpositions
generate $S_3$, are conjugate, and its three-cycles are their products, this
homomorphism is either zero or parity. Let $k\in\{0,1\}$ denote that
choice and set $t(g,i)=s(g,i)+kp(g)$. It satisfies the same cocycle law
and vanishes at $(h,4)$ for every $h\in H$.

Choose $h_i\in S_4$ with $h_i(4)=i$, and put $d_i=t(h_i,4)\in\F_2$.
The element $h_{g(i)}^{-1}gh_i$ fixes $4$. Applying
\eqref{eq:v2-sign-cocycle} to $gh_i$ in these two ways yields
\[
 t(g,i)+d_i=d_{g(i)}.
\]
Consequently
\begin{equation}
 s(g,i)=kp(g)+d_i+d_{g(i)}.
 \label{eq:v2-sign-coboundary}
\end{equation}
Reorienting $\gamma_i$ by $(-1)^{d_i}$ removes the last two terms.
Moreover, summing \eqref{eq:v2-sign-coboundary} over $i=1,\ldots,4$
gives zero. Hence the product of the four spatial lift signs is always
$+1$. The determinant-one condition therefore forces
$\epsilon_0(g)=(-1)^{p(g)}$. This gives \eqref{eq:v2-two-actions}.

The sign attached to a fixed spatial label under its stabilizer is unchanged
by reorientation, so it distinguishes $k=0$ from $k=1$. The explicit
implementers in \cref{thm:v2-double-cover} below prove the converse.
\end{proof}

\begin{remark}[An exact finite count]
With the labelled real frame fixed, there are sixteen such signed
representations: eight coboundaries in \eqref{eq:v2-sign-coboundary} for
each of the two classes. Adding the same bit to all four $d_i$ does not
change the cocycle, and these are the only duplicate choices. This count
will be checked independently by enumerating the signed Coxeter relations.
\end{remark}

\section{Projective lift obstruction and the cost of an invariant source}
\label{sec:v2-lifts}
Define
\begin{equation}
 H=\frac12(\gamma_1+\gamma_2+\gamma_3+\gamma_4),\qquad
 Q=i\gamma_0H,\qquad H^2=Q^2=I,
 \label{eq:v2-HQ}
\end{equation}
so $Q$ agrees with the operator used in V1's sharp-gap calculation.
Both $H$ and $Q$ are traceless Hermitian involutions. For $1\le j\le3$ put
\begin{equation}
 d_j=\frac{\gamma_j-\gamma_{j+1}}{\sqrt2},\qquad
 u_j^{(0)}=\gamma_0d_j,\qquad u_j^{(1)}=d_j.
 \label{eq:v2-adjacent-lifts}
\end{equation}
These are lifts of the adjacent transpositions $(j,j+1)$.

\begin{theorem}[Explicit double-cover implementation and its nonsplitting witness]
\label{thm:v2-double-cover}
The matrices \eqref{eq:v2-adjacent-lifts} implement respectively $r_0$
and $r_1$. They satisfy
\begin{align}
 (u_j^{(k)})^2&=(-1)^{1-k}I,\\
 u_j^{(k)}u_{j+1}^{(k)}u_j^{(k)}
 &=u_{j+1}^{(k)}u_j^{(k)}u_{j+1}^{(k)},\\
 u_1^{(k)}u_3^{(k)}&=-u_3^{(k)}u_1^{(k)}.
 \label{eq:v2-Schur-relations}
\end{align}
Each generated group has order $48$ and maps onto $S_4$ with kernel
$\{\pm I\}$. Choosing one lift for each permutation gives the projective
unitaries of \cref{def:v2-normalizer}. Their scalar multiplier cannot be
removed by rephasing the implementers, even by arbitrary $U(1)$ phases.
\end{theorem}
\begin{proof}
The $d_j$ are Hermitian involutions, anticommute with $\gamma_0$, and obey
\[
 d_jd_{j+1}+d_{j+1}d_j=-I,\qquad d_1d_3=-d_3d_1.
\]
Conjugation by a unit Clifford vector $\gamma(v)$ sends a coefficient vector
$x$ to $2\ip{v}{x}v-x$. Thus conjugation by $d_j$ negates $\gamma_0$,
sends $\gamma_j$ to $-\gamma_{j+1}$ and vice versa, and negates the other
two spatial generators. This is $r_1$ on an adjacent transposition.
Additional conjugation by $\gamma_0$ removes all spatial minus signs and
leaves the temporal minus sign, giving $r_0$.

The squares and distant anticommutation follow immediately. For adjacent
indices, $d_jd_{j+1}d_j=-d_{j+1}-d_j$, which is symmetric in the two
indices and proves the braid relation for $k=1$. For $k=0$ the same product
is $\gamma_0(d_j+d_{j+1})$, again symmetric.

Every generated matrix normalizes the five generators and induces one of the
faithful signed permutations $r_k(g)$. A matrix in the kernel commutes with
all generators, hence is scalar by their irreducibility. Every generator in
\eqref{eq:v2-adjacent-lifts} commutes with the quaternionic antiunitary
$\Jq$ of V1, so a scalar in this kernel is real. Unitarity makes it $\pm I$.
The squares for $k=0$, and the distant commutator for $k=1$, show that $-I$
occurs. There are therefore exactly two lifts per permutation and $48$
elements. The displayed relations are the familiar two Schur double-cover
presentations \cite{ReichsteinShukla}, but the kernel calculation here is
independent of a classification of those groups.

Finally, $(12)$ and $(34)$ commute in $S_4$, whereas their lifts anticommute.
Scalar rephasings leave this commutator unchanged. A genuine linear
representation of $S_4$ would map commuting elements to commuting matrices,
which is impossible here.
\end{proof}

\begin{theorem}[No fixed pure source, including on the quartic vertex stabilizer]
\label{thm:v2-no-pure}
For every coherent signed projective normalizer in
\cref{def:v2-normalizer}, there is no nonzero pure source ray fixed by all
of $S_4$. There is also no fixed pure source ray for any vertex stabilizer
$S_3\subset S_4$.

For the full group, two disjoint transpositions already give the
quantitative witness
\begin{equation}
 \norm{\rho-U_{(12)}\rho U_{(12)}^*}_{\HS}^2+
 \norm{\rho-U_{(34)}\rho U_{(34)}^*}_{\HS}^2\ge2
 \label{eq:v2-two-swap-gap}
\end{equation}
for every normalized pure state $\rho$.
\end{theorem}
\begin{proof}
By \cref{thm:v2-signed-classification}, it suffices to use the two normal
forms. Anticommuting invertible matrices cannot have a common eigenvector:
if $A\eta=\lambda\eta$ and $B\eta=\mu\eta$, their anticommutation gives
$\lambda\mu\eta=-\lambda\mu\eta$, a contradiction. The two distant
lifts in \eqref{eq:v2-Schur-relations} therefore exclude a fixed ray.

For the quantitative assertion, multiply the two lifts by scalar phases to
obtain anticommuting Hermitian involutions $A,B$. For every real $x,y$,
$(xA+yB)^2=(x^2+y^2)I$. Its expectation implies
$\langle A\rangle_\rho^2+\langle B\rangle_\rho^2\le1$.
For a pure state,
$\norm{\rho-A\rho A}_{\HS}^2=2(1-\langle A\rangle_\rho^2)$.
Adding the two identities proves \eqref{eq:v2-two-swap-gap}.

A vertex stabilizer can be relabelled to contain $(12)$ and $(23)$.
Their two lifts have commutator
$\pm[d_1,d_2]$, and the anticommutator relation above gives
\begin{equation}
 [u_1^{(k)},u_2^{(k)}]^2=-3I.
 \label{eq:v2-S3-commutator}
\end{equation}
This commutator is invertible. A common eigenvector would lie in its
kernel, so no such vector exists. Fixed reorientations of the Clifford
frame and scalar changes of lifts preserve the conclusion.
\end{proof}

\begin{corollary}[Minimum pure-ray orbit size]
\label{cor:v2-orbit-size}
Every pure ray has at least six distinct images under either projective
$S_4$ action. The lower bound is attained. Thus no one-ray, two-ray,
three-ray, or four-ray orbit completion gives a fixed-fibre realization of
this coherent signed source action.
\end{corollary}
\begin{proof}
Let $K$ be a ray stabilizer. If $|K|>4$, its order is $6,8,12$, or $24$.
A subgroup of order $6$ has a normal Sylow $3$ subgroup; the normalizer in
$S_4$ of that three-cycle group is its vertex stabilizer $S_3$ of order six.
Thus this case is excluded by \cref{thm:v2-no-pure}. An order-eight subgroup
is a Sylow $2$ subgroup, conjugate to the square group generated by
$(1234),(13)$; it contains the disjoint transpositions $(13),(24)$, again
excluded.

An index-two subgroup of $S_4$ is $A_4$: every homomorphism to $\{\pm1\}$
is fixed by its value on the conjugate transpositions. The two commuting
double transpositions $(12)(34)$ and $(13)(24)$ in $A_4$ have lifts,
up to scalar,
\[
 A=d_{12}d_{34},\qquad B=d_{13}d_{24},\qquad
 d_{ij}=\frac{\gamma_i-\gamma_j}{\sqrt2}.
\]
They satisfy $AB=-BA$. One can verify this by expansion. Equivalently, put
$w_1=(1,1,-1,-1)/2$, $w_2=(1,-1,1,-1)/2$,
$w_3=(1,-1,-1,1)/2$ in the spatial coefficient space. These are
orthonormal; the two lifts are, up to signs,
$\gamma(w_2)\gamma(w_3)$ and $\gamma(w_1)\gamma(w_3)$, which anticommute.
Both order $12$ and order $24$ are therefore excluded. Consequently
$|K|\le4$, and orbit--stabilizer gives at least $24/4=6$ images.

A lift of a four-cycle is unitary and has an eigenvector. Its eigenray is
fixed by the cyclic order-four subgroup in the projective action. Its
stabilizer has order at least four and, by the preceding argument, at most
four. Its orbit therefore has exactly six rays.
\end{proof}

\subsection{The entire invariant-state fibre and an implementable repair}

\begin{theorem}[Invariant-state classification and completely positive averaging]
\label{thm:v2-twirl}
Let $C_0=H$ and $C_1=Q$ as in \eqref{eq:v2-HQ}. For the two normal forms,
\begin{equation}
 \{U_g:g\in S_4\}'=\Span_\C\{I,C_k\}.
 \label{eq:v2-commutant}
\end{equation}
Their invariant density matrices are exactly
\begin{equation}
 \boxed{\sigma_{k,q}=\frac14(I+qC_k),\qquad -1\le q\le1.}
 \label{eq:v2-invariant-states}
\end{equation}
They have eigenvalues $(1+q)/4,(1-q)/4$, each twice, and minimum rank two.

The group average is the completely positive, trace-preserving map
\begin{equation}
 \boxed{\mathcal T_k(A)=\frac1{24}\sum_g U_g A U_g^*
 =\frac{\Tr A}{4}I+\frac{\Tr(AC_k)}4 C_k.}
 \label{eq:v2-twirl}
\end{equation}
It is the Hilbert--Schmidt orthogonal projection onto
\eqref{eq:v2-commutant}. In particular it gives the closest invariant
density matrix to any given density matrix.
\end{theorem}
\begin{proof}
Expand an arbitrary matrix in the sixteen-coordinate basis
\eqref{eq:operator-basis}. The coefficients of $\gamma_a$ transform by
$r_k$, and those of $B_{ab}$ transform by its exterior square.
For $k=0$, an invariant vector coefficient has zero temporal component and
four equal spatial components, hence is a multiple of $H$. A temporal
bivector $B_{0i}$ is negated by a transposition fixing $i$, so its invariant
coefficient vanishes. A spatial bivector $B_{ij}$ is negated by the
transposition $(ij)$, so those coefficients also vanish.
For $k=1$, every vector coefficient vanishes: the temporal one changes
sign, and each spatial one is negated by some transposition fixing its label.
Temporal bivectors are permuted without a sign and so leave precisely the
line spanned by $Q$. Spatial bivectors again have no invariant coefficient.
This proves \eqref{eq:v2-commutant}, including for non-Hermitian matrices
by complex linearity.

The two $C_k$ are traceless Hermitian involutions with two-dimensional
positive and negative eigenspaces. Trace one, Hermiticity, and positivity
therefore give exactly \eqref{eq:v2-invariant-states} and its stated
spectrum. A finite average of unitary conjugations is completely positive
and trace preserving. Because conjugation is an ordinary group action, its
average is the orthogonal projection onto its fixed space. Trace
orthogonality of $I,C_k$ gives \eqref{eq:v2-twirl}. The projected image of a
density matrix is a density matrix, which proves the nearest-point claim.
\end{proof}

\begin{corollary}[Exact purity cost and orbit defect]
\label{cor:v2-twirl-gap}
For a pure state $\rho$ put $q_k=\Tr(\rho C_k)$. Then
\begin{align}
 \norm{\rho-\mathcal T_k(\rho)}_{\HS}^2
 &=\frac{3-q_k^2}{4}\ge\frac12,\label{eq:v2-state-gap}\\
 \frac1{24}\sum_g\norm{\rho-U_g\rho U_g^*}_{\HS}^2
 &=\frac{3-q_k^2}{2}\ge1.\label{eq:v2-orbit-gap}
\end{align}
The average has rank two exactly when $|q_k|=1$, and rank four otherwise.
The minimum-rank repair is therefore not generic.
\end{corollary}
\begin{proof}
Orthogonal projection gives
$\norm{\rho-\mathcal T_k(\rho)}_{\HS}^2
=\Tr\rho^2-\Tr\mathcal T_k(\rho)^2$.
The two terms are $1$ and $(1+q_k^2)/4$. Expanding each term of the orbit
average gives twice that difference. The rank assertion follows from the
spectrum in \cref{thm:v2-twirl}.
\end{proof}

\begin{proposition}[Relation to, and distinction from, V1's quaternionic completion]
\label{prop:v2-two-completions}
For the pure-state completion $\rho_+$ of \cref{thm:completion},
\begin{equation}
 \mathcal T_0(\rho)=\rho_+
 \quad\Longleftrightarrow\quad
 n=\pm(0,1,1,1,1)/2.
 \label{eq:v2-completion-equality}
\end{equation}
The equality $\mathcal T_1(\rho)=\rho_+$ never holds for a pure state.
The map $\mathcal T_k$ is an available random-unitary construction only
when its unitaries are themselves admitted operations. Its complete
positivity does not prove their occurrence in the arithmetic or physical
process.
\end{proposition}
\begin{proof}
V1 gives $\rho_+=(I+n\cdot\gamma)/4$ with $\norm n=1$.
The first twirl retains only the component of $n$ in the unit spatial
average direction. Equality holds exactly when $n$ lies in that direction,
which gives the two signs in \eqref{eq:v2-completion-equality}. The second
twirl has only a bivector component, while $\rho_+$ has a nonzero vector
component. The two sectors of the operator basis are orthogonal, so equality
is impossible. The operational qualification is separate from the matrix
identity.
\end{proof}

\begin{proposition}[Complete signed-Gram covariance still excludes a pure source]
\label{prop:v2-Gram-gap}
Let $\mathcal E_k(G)=24^{-1}\sum_g r_k(g)G r_k(g)^{\mathsf T}$.
For every pure displacement Gram of V1,
\begin{align}
 \mathcal E_0(G)&=I_5,&
 \norm{G-\mathcal E_0(G)}_{\HS}^2&=4,\\
 \norm{G-\mathcal E_1(G)}_{\HS}^2&=4-2q_1^2\ge2.
 \label{eq:v2-signed-Gram-gaps}
\end{align}
Thus neither coherent signed normal form gives an invariant pure
five-displacement Gram. A fixed real reorientation conjugates these
statements and does not remove the lower bounds.
\end{proposition}
\begin{proof}
The covariance of $G$ under operator conjugation gives
$\mathcal E_k(G(\rho))=G(\mathcal T_k(\rho))$.
For $k=0$ the twirled state has no bivector coordinate, so its Gram is $I_5$.
V1's identity $\sum b_{ab}^2=2$ then gives squared distance four.
For $k=1$, $r_1(g)=(-1)^{p(g)}P_g$, where $P_g$ is the ordinary pointed
permutation matrix of V1. The scalar sign cancels in matrix conjugation,
so this is exactly the invariant projection already evaluated in
\cref{thm:sharp-gap}, with $q_1=\Tr(\rho Q)$.
\end{proof}

\section{The arithmetic descent consequence: keep fibres, not a fictitious fixed ray}
\label{sec:v2-descent}

\subsection{A precise boundary for the earlier frame-cover argument}
The quartic root cover in \cite{AD9} selects one residual vertex and reduces
its permutation stabilizer to $S_3$. This is an inherited algebraic frame
statement. \Cref{thm:v2-no-pure} shows that it cannot additionally make the
pure theta ray invariant under the corresponding coherent complex-linear
unitary normalizer. \Cref{cor:v2-orbit-size} gives the stronger finite orbit
bound of six. This is an obstruction to that \emph{fixed-fibre}
interpretation of the cover, not an obstruction to the quartic cover or to
the rationally defined section itself.

In particular, the statement in
\path{prop:v9-arithmetic-source-candidate} that Galois permutes the
arithmetic rays must be used at its algebraic/semilinear level. Finite
averaging of Hermitian forms is legitimate once a finite projective
\emph{complex-linear} action on the same vector space has actually been
constructed. The arithmetic permutation of torsion labels alone does not
supply that hypothesis or identify an invariant pure source for it.

\begin{proposition}[Semilinear descent has a different multiplication law]
\label{prop:v2-semilinear}
Let $K/k$ be a finite Galois extension with group $\Gamma$, let $V$ be a
$K$-vector space with a chosen $K$-basis, and let $D_g$ be semilinear maps
satisfying $D_g(av)=g(a)D_g(v)$. Write $D_g(v)=A_g\,g(v)$ in that basis.
Then exact descent $D_gD_h=D_{gh}$ is equivalent to
\begin{equation}
 A_g\,g(A_h)=A_{gh},
 \label{eq:v2-semilinear-law}
\end{equation}
and invariance of a line $K\eta$ is the condition
\begin{equation}
 A_g\,g(\eta)=\lambda_g\eta\qquad(\lambda_g\in K^\times).
 \label{eq:v2-semilinear-ray}
\end{equation}
Neither equation implies $A_gA_h=c(g,h)A_{gh}$ or $A_g\eta\in K\eta$.
For a complex embedding $\iota:K\hookrightarrow\C$, $D_g$ canonically
induces a complex-linear map
\begin{equation}
 V\otimes_{K,\iota\circ g}\C\longrightarrow
 V\otimes_{K,\iota}\C,
 \label{eq:v2-embedding-map}
\end{equation}
not an endomorphism of the one chosen fibre unless additional
identifications are supplied.
\end{proposition}
\begin{proof}
Composition gives
$D_gD_h(v)=A_g g(A_h)gh(v)$, proving
\eqref{eq:v2-semilinear-law}; the line condition is its definition.
For a concrete failure of ordinary line invariance take $K=\Q(i)$,
$g$ complex conjugation, $A_g=\operatorname{diag}(1,-1)$, and
$\eta=(1,i)^{\mathsf T}$. Then $A_g g(A_g)=I$ and
$A_g g(\eta)=\eta$, while $A_g\eta=(1,-i)^{\mathsf T}$ is not proportional
to $\eta$. For a failure of the ordinary multiplication inference, choose
any $B\in\operatorname{GL}_2(K)$ and put
$A_g=B g(B)^{-1}$. This satisfies the semilinear law. For example
$B=\operatorname{diag}(1,2+i)$ gives
$A_g=\operatorname{diag}(1,(2+i)/(2-i))$, whose square is not scalar;
so these $A_g$ do not give a projective ordinary action of the order-two
group.

For \eqref{eq:v2-embedding-map}, send $v\otimes z$ to $D_g(v)\otimes z$.
The tensor relation for $av$ on the source has scalar $\iota(g(a))$,
exactly the scalar obtained from $D_g(av)=g(a)D_g(v)$ on the target. The
map is therefore well defined and complex linear between the displayed
fibres. Identifying those fibres, their metrics, and their markings is
additional data.
\end{proof}

\subsection{A nonempty fibrewise alternative without changing the pure source}

\begin{construction}[Projective source-orbit transport]
\label{con:v2-orbit}
For either complex projective lift choose $U_g$ once and retain a normalized
pure source $\eta$. Define $\eta_g=U_g\eta$ and
$\rho_g=\ketbra{\eta_g}$. Let
$V_g e_a=\gamma_a\eta_g$ and $G_g=V_g^*V_g$.
The source remains pure at each frame; it is not required to be the same ray
at every frame.
\end{construction}

\begin{theorem}[Exact covariance of the source-orbit family]
\label{thm:v2-orbit}
For the preceding construction,
\begin{align}
 U_hV_g&=c(h,g)V_{hg}r_k(h),\\
 G_{hg}&=r_k(h)G_g r_k(h)^{\mathsf T},\\
 \frac1{24}\sum_g\rho_g&=\mathcal T_k(\rho).
 \label{eq:v2-orbit-transport}
\end{align}
Every $G_g$ has V1's spectrum $(2,2,1,0,0)$, and the projective group
multipliers obey the cocycle identity. Thus exact fibrewise covariance is
compatible with the absence of an invariant pure source in a fixed frame.
\end{theorem}
\begin{proof}
For each label,
$U_h\gamma_a\eta_g
=\sum_b r_k(h)_{ba}\gamma_b U_hU_g\eta
=c(h,g)V_{hg}r_k(h)e_a$.
Taking Grams and using orthogonality of $r_k(h)$ gives the second formula.
The third is the definition of the twirl. V1 applies to each normalized
vector. Associativity of $U_gU_hU_\ell$ gives
$c(g,h)c(gh,\ell)=c(h,\ell)c(g,h\ell)$.
\end{proof}

\begin{scope}[Constructed transport is not identified arithmetic transport]
This constructs the correct mathematical comparison \emph{type} for a
moving pure source. It does not prove that the actual Galois descent of
$\vartheta_\kappa^2$ is this orbit transport, nor that its arrows physically
occur on a renewal carrier. Matching the arithmetic descent, its metric
transport, and its signs/cocycle to the independently supplied renewal
packet remains necessary. The mean state in
\eqref{eq:v2-orbit-transport} is a changed ensemble and is not substituted
for any individual arithmetic fibre.
\end{scope}

\section{A certified acquisition formula for the actual theta-section source}
\label{sec:v2-theta}

This section does not choose a new arithmetic source. It evaluates the
already inherited ray \eqref{eq:actual-source} in standard analytic
coordinates. The addition identity used in the proof is classical
\cite[\S21.6(ii)]{DLMFTheta}; the contribution here is to retain the
level-two coordinate marking, source normalization, and finite error budget
required for V1's ten-overlap test.

\subsection{Four theta constants determine the squared section}
Let $\tau=\tau^{\mathsf T}\in M_2(\C)$ with $Y=\operatorname{Im}\tau\succ0$
be the period matrix in an independently chosen symplectic homology marking
of $J(\C)$. Encode the marked characteristic $\kappa$ by
$\alpha,\beta\in\{0,1\}^2$, with $\alpha\cdot\beta=1\pmod2$ on the odd
branch. Define
\begin{align}
 \theta\!\begin{bmatrix}\alpha/2\\\beta/2\end{bmatrix}(z\mid\tau)
 &=\sum_{m\in\Z^2}
 e^{\pi i(m+\alpha/2)^{\mathsf T}\tau(m+\alpha/2)
       +2\pi i(m+\alpha/2)^{\mathsf T}(z+\beta/2)},\\
 \Theta_\sigma(z\mid\tau)
 &=\theta\!\begin{bmatrix}\sigma/2\\0\end{bmatrix}(2z\mid2\tau),
 &t_\sigma(\tau)&=\Theta_\sigma(0\mid\tau),
 \qquad\sigma\in\{0,1\}^2.
 \label{eq:v2-theta-definitions}
\end{align}
The four $\Theta_\sigma$ are the marked level-two theta basis. The standard
Heisenberg shift and sign operators permute this basis and act diagonally
by signs, respectively. Consequently its invariant Hermitian metric is
scalar identity; we fix that common scale to one, as in V1.

\begin{proposition}[Marked squared-section coordinates]
\label{prop:v2-theta-coordinates}
In this basis, the source ray $[\vartheta_\kappa^{\otimes2}]$ has coefficient
vector
\begin{equation}
 \boxed{r_\sigma=(-1)^{\sigma\cdot\beta}
                    t_{\sigma+\alpha}(\tau),}
 \label{eq:v2-rtheta}
\end{equation}
where the subscript is reduced modulo two. In particular, $r\ne0$ and
$\norm r=\norm t$. If the independently aligned Clifford matrices in this
same theta basis are $\Gamma_a$, its actual marked density and overlaps are
\begin{equation}
 \boxed{\rho_\Theta=\frac{rr^*}{r^*r},\qquad
 b_{ab}=\frac{r^*(i\Gamma_a\Gamma_b)r}{r^*r},\qquad
 n_a=\frac{r^*\Gamma_a r}{r^*r}.}
 \label{eq:v2-actual-moments}
\end{equation}
The matrices $\Gamma_a$ must come from the retained theta-lift marking;
substituting a preferred Clifford frame without transporting $r$ is not
permitted.
\end{proposition}
\begin{proof}
The Fourier frequencies of $\Theta_\sigma$ have parity $\sigma$, so these
four functions are independent. The level-two dimension is four by the
inherited principal-polarization statement, hence they form a basis. The
Heisenberg sign operators have distinct joint characters on these basis
vectors, making them orthogonal in an invariant metric; shifts permute them
transitively, making their norms equal.

Square the first series in \eqref{eq:v2-theta-definitions}. Absolute
convergence on compact $z$-sets allows regrouping. For the two summation
indices $m,n$, put $u=m+n+\alpha$ and $v=m-n$. Then
$u+v=\alpha\pmod2$, and
\[
 (m+\alpha/2)^{\mathsf T}\tau(m+\alpha/2)
 +(n+\alpha/2)^{\mathsf T}\tau(n+\alpha/2)
 =\tfrac12(u^{\mathsf T}\tau u+v^{\mathsf T}\tau v).
\]
Write $u=2k+\sigma$ and $v=2\ell+\sigma+\alpha$, reducing the latter
characteristic by an integer reindexing if necessary. The linear part of
the exponent is $2\pi i u^{\mathsf T}z+\pi i u^{\mathsf T}\beta$.
The sum therefore factors as
\begin{equation}
 \theta\!\begin{bmatrix}\alpha/2\\\beta/2\end{bmatrix}(z\mid\tau)^2
 =\sum_{\sigma\in\{0,1\}^2}
 (-1)^{\sigma\cdot\beta}
 \Theta_\sigma(z\mid\tau)t_{\sigma+\alpha}(\tau).
 \label{eq:v2-addition}
\end{equation}
This proves \eqref{eq:v2-rtheta}. The original theta section is nonzero, so
its square is nonzero and the coefficient vector cannot vanish. Signs and
permutation preserve its Euclidean norm. Finally, the invariant metric has
been fixed to the Euclidean one in this basis, and the ordinary density and
expectation formulas give \eqref{eq:v2-actual-moments}.
\end{proof}

\begin{remark}[Oddness does not annihilate the source]
At $z=0$ the left side of \eqref{eq:v2-addition} vanishes for an odd
characteristic. The right side vanishes by pairwise cancellation of
$\sigma$ and $\sigma+\alpha$, since $\alpha\cdot\beta$ is odd. This is
one quadratic relation among the theta constants; it does not say that the
coefficient vector $r$ or the squared theta section vanishes.
\end{remark}

\subsection{Explicit truncation and period-error bounds}
For $c>0$ and an integer $N\ge1$, put
\begin{equation}
 B(c)=2+\sqrt{\pi/c},\qquad R_N=N+\tfrac12,\qquad
 T_N(c)=\left(2+\frac1{cR_N}\right)e^{-cR_N^2}.
 \label{eq:v2-tail-definitions}
\end{equation}

\begin{theorem}[Certified four-coordinate theta acquisition]
\label{thm:v2-theta-tail}
Assume $\lambda_{\min}(\operatorname{Im}\tau)\ge\lambda>0$ and set
$c=2\pi\lambda$. Define the finite sums
\begin{equation}
 t^{[N]}_\sigma=
 \sum_{\substack{m\in\Z^2\\ |m_1|,|m_2|\le N}}
 e^{2\pi i(m+\sigma/2)^{\mathsf T}\tau(m+\sigma/2)}.
 \label{eq:v2-theta-finite}
\end{equation}
Then
\begin{equation}
 |t_\sigma-t^{[N]}_\sigma|\le2B(c)T_N(c),\qquad
 \boxed{\norm{t-t^{[N]}}_2\le E_N:=4B(c)T_N(c).}
 \label{eq:v2-theta-tail}
\end{equation}
If instead $\tau$ is replaced by $\widehat\tau$ with
$\norm{\tau-\widehat\tau}_{\op}\le\delta$, and every point of the
straight segment between them has imaginary part at least $\lambda I$,
then an additional valid vector error bound is
\begin{equation}
 \boxed{E_{\mathrm{per}}=
       \frac{4}{e\lambda}B(\pi\lambda)^2\,\delta.}
 \label{eq:v2-period-error}
\end{equation}
Certified errors in the finite exponential sums are added to these bounds.
No floating-point equality test is needed to assert them.
\end{theorem}
\begin{proof}
Every summand is bounded in modulus by
$e^{-c\norm{m+\sigma/2}^2}$. For $s\in\{0,1/2\}$, split the one-dimensional
Gaussian sum at its two nearest half-lines. Monotone integral comparison
gives
\[
 \sum_{m\in\Z}e^{-c(m+s)^2}\le2+\int_{-\infty}^{\infty}e^{-cx^2}\,\dd x=B(c).
\]
Both tails outside $[-N,N]$ are bounded by
$\sum_{j=N+1}^\infty e^{-c(j-1/2)^2}$. With $R_N=N+1/2$,
\[
 \sum_{j=N+1}^\infty e^{-c(j-1/2)^2}
 \le e^{-cR_N^2}+\int_{R_N}^{\infty}e^{-cx^2}\,\dd x
 \le\left(1+\frac1{2cR_N}\right)e^{-cR_N^2}.
\]
Their sum is $T_N(c)$. The complement of the two-dimensional square is
contained in the union of the two coordinate-tail events. The other
coordinate has full sum at most $B(c)$, so the scalar tail is at most
$2B(c)T_N(c)$. Four scalar coordinates give the vector bound
\eqref{eq:v2-theta-tail}.

For the period perturbation, differentiate along the straight segment. For
$x=m+\sigma/2$ the derivative has modulus at most
$2\pi\delta\norm x^2e^{-c\norm x^2}$. The inequality
$y e^{-cy}\le 2(ce)^{-1}e^{-cy/2}$ gives a scalar derivative sum at most
\[
 \frac{4\pi\delta}{ce}B(c/2)^2
 =\frac{2\delta}{e\lambda}B(\pi\lambda)^2.
\]
Integrate over the unit parameter interval, and multiply by two for the norm
of the four-coordinate vector. Absolute convergence of the differentiated
series follows from the same Gaussian bound, justifying the argument.
\end{proof}

\begin{theorem}[Certified normalized source and overlap errors]
\label{thm:v2-normalized-theta}
Let $\widehat r$ be obtained from a finite, possibly interval-evaluated theta
vector by \eqref{eq:v2-rtheta}. Suppose a valid error budget gives
$\norm{r-\widehat r}_2\le E$, and let $m=\norm{\widehat r}_2>E$.
Use the same exact aligned Clifford frame for both calculations. Then
\begin{align}
 \norm{\rho_\Theta-\widehat\rho}_{\HS}
 &\le\frac{2\sqrt2 E}{m},&
 \norm{\rho_\Theta-\widehat\rho}_1&\le\frac{4E}{m},\\
 |b_{ab}-\widehat b_{ab}|&\le\frac{4E}{m},&
 |n_a-\widehat n_a|&\le\frac{4E}{m},\\
 &\boxed{\norm{G(\rho_\Theta)-G(\widehat\rho)}_{\HS}\le\frac{8E}{m}.}
 \label{eq:v2-normalized-errors}
\end{align}
Here $\widehat\rho=\widehat r\widehat r^*/m^2$. If the period is known
with arbitrarily refinable enclosures and the marked input remains fixed,
this gives an eventually terminating certificate for every prescribed
positive overlap accuracy. It is not a terminating decision procedure for
exact equality of arbitrary transcendental overlaps.
\end{theorem}
\begin{proof}
Let $a=\norm r>0$. The normalization estimate
\[
 \norm{r/a-\widehat r/m}
 \le\norm{r/a-r/m}+\norm{(r-\widehat r)/m}
 =\frac{|m-a|}{m}+\frac E m\le\frac{2E}{m}
\]
uses only the reverse triangle inequality. For unit vectors $x,y$, the
associated projectors satisfy
$\norm{\ketbra x-\ketbra y}_{\HS}\le\sqrt2\norm{x-y}$ and
$\norm{\ketbra x-\ketbra y}_1\le2\norm{x-y}$, obtained directly from
$\Tr(\ketbra x\ketbra y)=|\ip{x}{y}|^2$. This proves the density bounds.
The moment bounds follow because every $\Gamma_a$ and $i\Gamma_a\Gamma_b$
has operator norm one.

To obtain the sharper complete Gram bound, expand the traceless Hermitian
matrix $\Delta\rho$ in the orthogonal vector/bivector basis of V1. Then
$\sum_{a<b}(\Delta b_{ab})^2\le4\norm{\Delta\rho}_{\HS}^2$, while
$\norm{\Delta G}_{\HS}^2=2\sum_{a<b}(\Delta b_{ab})^2$.
This gives the last inequality in \eqref{eq:v2-normalized-errors}.
Finally, $r\ne0$ by \cref{prop:v2-theta-coordinates}; refining the period
enclosure and increasing $N$ makes $E\to0$ and $m\to\norm r>0$.
Consequently $m>E$ and every fixed positive accuracy eventually pass.
\end{proof}

\begin{scope}[The actual remaining arithmetic inputs]
For $C_{\mathrm{tet}}:y^2=x(x^4-x-1)$, this version has not evaluated a
certified period matrix, its symplectic marking, the coordinates of
$\kappa_\infty$ in that marking, or the five aligned matrices $\Gamma_a$.
No arbitrary period matrix or example spinor is reported as that curve's
source. The proven advance is that once those independently fixed inputs
are supplied, four convergent sums with the displayed error bounds compute
the actual source and all ten overlaps. Numerical equality of source
moments is not inferred from a finite precision calculation.
\end{scope}

\section{V2 verification, status ledger, and next arithmetic target}
\label{sec:v2-status}

\subsection{Exact checks and their limits}
The accompanying V2 check script verifies the two signed generator actions,
all square/braid/distant relations, the invertible $S_3$ commutator, the
normal Klein-four anticommutation witness, and both two-dimensional
commutants with exact symbolic matrices. It also enumerates every signed
adjacent-generator assignment satisfying the determinant and Coxeter
conditions and obtains exactly the sixteen actions classified above.
The finite projective conjugation averages are tested on the complete
sixteen-element operator basis, not only on one source.

The theta-addition formula and analytic tail bounds are proved above. A
separately labelled numerical smoke test of the evaluator uses an explicitly
chosen Siegel matrix \emph{unrelated to the curve} to check implementation
conventions. Such a test is not an interval proof and is not an evaluation
of $\vartheta_{\kappa_\infty}^2$ for $C_{\mathrm{tet}}$. Any floating-point
evaluator must include certified roundoff and period enclosures before its
reported error becomes a numerical certificate.

\subsection{Proved in V2}
\begin{enumerate}[label=\textup{(P2.\arabic*)},leftmargin=3.6em]
\item Coherent signed projective normalizers of the five Hermitian Clifford
rays, with one temporal and four permuted residual labels, have exactly two
real-sign classes, with sixteen labelled actions before sign quotienting.
\item Both classes have explicit order-$48$ double-cover implementations;
distant swaps give a scalar-phase-invariant obstruction to an ordinary
linear lift.
\item No pure source ray is fixed by either $S_4$ action or by its quartic
vertex stabilizer $S_3$. Every pure-ray orbit has at least six elements,
and this minimum is attained.
\item Invariant density matrices form an explicit interval with rank-two
endpoints. Exact group averaging is completely positive, has a measured
purity cost, and differs from V1's quaternionic completion except on the
specified exceptional branch.
\item Signed source-orbit families have exact fibrewise Gram and cocycle
transport without an invariant pure ray. This is a mathematical
construction, not an identification with the actual arithmetic descent.
\item The inherited squared theta section has four explicit second-order
theta-constant coordinates; Gaussian tails, period perturbations, and
normalization give a finite acquisition theorem for its actual ten overlaps.
\end{enumerate}

\subsection{Inherited and not redone}
All of V1, including its historical ledger, is retained without alteration.
The chronology/Peano arithmetic loading, its cutoff calculus, the
arithmetic torsion and quartic frame packet, the integral $K_4$ lattice,
operator-source equivalence, generic source/Krylov and Tannakian/interface
criteria, and the companion physical theories remain inherited. Classical
Schur-cover presentations and theta addition are cited as background tools.

\subsection{Open and now more precisely localized}
The unresolved first arithmetic computation is the \emph{marked} period and
theta-structure realization of the specified curve, followed by the four
constants in \eqref{eq:v2-rtheta}. The next comparison must retain the
actual semilinear or fibrewise transport of the theta ray. The two coherent
unitary signed classes have been completely tested and cannot make that ray
fixed in one fibre, even after merely choosing a quartic-cover vertex.
A noncentral extension defect, semilinear action, or moving-fibre model must
be declared and compared on its own terms.

An independently acquired renewal source in the same transported operator
frame, the central Store/chronology signature expectation, a native
arithmetic transfer, Frobenius--Wilson compatibility, and cofinal estimates
remain open. No new matter multiplicity, Standard-Model parameter, prime
estimate, or continuum field equation follows from the orbit size, invariant
rank, or theta tail bound.

\paragraph{Next noncircular target.}
Construct a certified symplectic period marking of the actual genus-two
curve, identify the characteristic and five torsion lifts in its level-two
basis, and evaluate the four theta constants with the preceding acquisition
theorem. Use the resulting ten overlaps and their error intervals with V1's
pure-state reconstruction, then compare the actual arithmetic descent maps
with the inherited renewal frame-cover maps. Do not reopen an ordinary or
coherent signed fixed-pure-source symmetry test which this version has
already decided negatively.

\paragraph{Version integrity.}
The V1 mathematical body is preserved byte for byte. This V2 continuation is
appended before the bibliography, with new references added afterward. The
accompanying verification records the SHA-256 digest of the retained V1 body.
Future iterations should preserve both bodies and append their work before
the bibliography rather than rewriting either historical ledger.
% ====================== END IMMUTABLE V2 MATHEMATICAL BODY ======================
% APPEND V3 AND LATER ITERATIONS HERE.


% ===================== BEGIN IMMUTABLE V3 MATHEMATICAL BODY =====================
\clearpage
\section{V3 continuation: algebraic acquisition before period integration}
\label{sec:v3-context}
\noindent\textbf{20 September 2026.}

V2 gives a convergent period/theta-constant procedure for the actual source,
but it does not require that every new comparison begin with a numerical
period matrix. This continuation goes upstream of that calculation. The
classical first Thomae formula expresses the absolute fourth powers of the
ten even theta constants in terms of the six branch points. Combining it
with the already proved squared-section coordinates and a finite two-copy
Pauli calculation eliminates the common period factor. It determines all
fifteen squared Hermitian displacement expectations, including every
absolute entry of the five-source Gram, without computing a period.

For the specified curve, these are actual arithmetic numbers, not values at
an artificial Siegel matrix. Two expectations vanish exactly. The other
thirteen are enclosed by rational intervals from certified disks for the
quartic roots. V1's pure-source recognition theorem then reduces the full
oriented source to exactly thirty-two algebraic density matrices. An exact
finite sign action and fifteen interval exclusion witnesses prove this count.
Five independently acquired binary signs suffice to select the marked
source. The earlier general claim that finite precision cannot decide
arbitrary transcendental equality is unchanged: this particular source has
now been reduced to a finite algebraic solution set.

\paragraph{What is imported, and what is new.}
The curve, odd characteristic, temporal class, and displacement labeling are
inherited from \cite{AD9}. The source is still
$[\vartheta_{\kappa_\infty}^{\otimes2}]$ in the Heisenberg-unitary metric.
The coordinate formula is \cref{prop:v2-theta-coordinates}; the cubic
pure-state recognition and reconstruction are \cref{thm:pure-recognition}.
None is rederived. The first Thomae formula is a classical external input,
used in the precise absolute-value form below \cite[\S2.6, Eq.~(18)]{BernatskaThomae}.
It is not claimed as a new theorem. New here are its normalized source
consequence, the exact acquisition and phase ambiguity of this particular
arithmetic packet, and the finite certificates. No physical source,
Frobenius map, connection, or transfer is supplied by a theta identity.

\section{Branch weights and the complete phase-free theta source}
\label{sec:v3-thomae-source}

\subsection{Intrinsic branch weights and the classical input}
Let $C$ be a smooth genus-two hyperelliptic curve over a subfield of $\C$,
with its six distinct labeled branch points $\mathcal B$. For homogeneous
representatives $v_i\in\C^2\setminus\{0\}$ write
$[ij]=\det(v_i,v_j)$. Let $\mathcal P$ be the set of ten unordered
partitions $S\mid S^c$ of $\mathcal B$ into two triples. Define
\begin{equation}
 w_S=\left|\prod_{\{i,j\}\subset S}[ij]
                \prod_{\{i,j\}\subset S^c}[ij]\right|,
 \qquad Z_{\mathcal B}=\sum_{S\mid S^c\in\mathcal P}w_S,
 \qquad p_S=\frac{w_S}{Z_{\mathcal B}}.
 \label{eq:v3-branch-weights}
\end{equation}
The products use either fixed ordering of each pair; their absolute values
remove its signs. Every $w_S$ is strictly positive.

\begin{lemma}[Coordinate independence and algebraicity]
\label{lem:v3-weight-intrinsic}
The labeled probabilities $p_S$ are independent of the homogeneous
representatives and invariant under a common fractional-linear coordinate
change. If the branch points are algebraic in a fixed complex embedding,
then every $p_S$ is a positive real algebraic number. Relabeling the branch
points permutes the $p_S$.
\end{lemma}
\begin{proof}
Each vertex occurs in exactly two brackets of a partition product. Replacing
$v_i$ by $c_i v_i$ multiplies every $w_S$ by the same factor
$\prod_i|c_i|^2$. A matrix $A\in\operatorname{GL}_2(\C)$ multiplies each
of its six brackets by $\det A$, again giving the same factor for every
partition. Both factors cancel in \eqref{eq:v3-branch-weights}.
For algebraic $z$, $|z|$ is the nonnegative square root of the real
algebraic number $z\overline z$. Closure of algebraic numbers under sums,
products, division by nonzero numbers, and square roots proves the claim.
\end{proof}

\paragraph{Classical Thomae input, with the ambiguity retained.}
In a marked analytic theta realization, let $m_S$ be the even
characteristic corresponding to the partition $S\mid S^c$. First Thomae
implies that there is one positive number $c_\tau$, independent of $S$,
such that
\begin{equation}
 \big|\theta[m_S](0\mid\tau)\big|^4=c_\tau w_S.
 \label{eq:v3-thomae-absolute}
\end{equation}
Indeed, the usual formula is an eighth-root-of-unity times a common period
factor and the fourth root of the two Vandermonde products. Taking absolute
fourth powers gives \eqref{eq:v3-thomae-absolute}. This loses phase
information, deliberately. With one branch at infinity it is exactly the
finite $2+3$ partition formula in \cite[\S2.6]{BernatskaThomae}; the
homogeneous version follows by the common transformation factors in
\cref{lem:v3-weight-intrinsic}. The characteristic--partition
correspondence is part of this standard input, not an identification chosen
from the desired source values.

\subsection{A finite two-copy identity}
Use $E=\F_2^2$ and the ordered basis $(e_\sigma)_{\sigma\in E}$ of
$\C^4$. For $m=(\alpha,\beta)\in E\oplus E$ set
\begin{equation}
 (M_mx)_\sigma=(-1)^{\sigma\cdot\beta}x_{\sigma+\alpha},
 \quad q(m)=\alpha\cdot\beta\pmod2,
 \quad \omega(m,n)=\alpha\cdot\beta'+\alpha'\cdot\beta,
 \quad P_m=i^{q(m)}M_m.
 \label{eq:v3-pauli}
\end{equation}
Thus $M_m$ is real orthogonal, $M_m^{\mathsf T}=(-1)^{q(m)}M_m$,
and $P_m$ is a Hermitian involution. An overall sign in a retained
Hermitian theta lift does not change any squared expectation below.

For the ten even $m$, put
\[
 E_m=\frac12\sum_{\sigma,\nu}(M_m)_{\sigma\nu}
                     e_\sigma\otimes e_\nu.
\]

\begin{lemma}[Two-copy diagonalization]
\label{lem:v3-bell}
The $E_m$ for $q(m)=0$ form an orthonormal basis of
$\operatorname{Sym}^2(\C^4)$. If $x\ne0$ and
$Q_m(x)=x^{\mathsf T}M_mx$, then
\begin{align}
 \sum_{q(m)=0}|Q_m(x)|^2&=4\norm{x}^4,\
 \left(\frac{x^*P_hx}{x^*x}\right)^2
 &=\sum_{q(m)=0}(-1)^{q(h)+\omega(h,m)}
               \frac{|Q_m(x)|^2}{4\norm{x}^4}.
 \label{eq:v3-bell-identities}
\end{align}
\end{lemma}
\begin{proof}
The matrices $M_m$ are trace-orthogonal with squared Hilbert--Schmidt norm
four. The ten even matrices are symmetric, so their normalized
vectorizations give ten orthonormal vectors in the ten-dimensional
symmetric square. The coefficient of $x\otimes x$ along $E_m$ is
$Q_m(x)/2$. Parseval gives the first identity.

Pauli commutation and $P_h^{\mathsf T}=(-1)^{q(h)}P_h$ give
\[
 P_hM_mP_h^{\mathsf T}
   =(-1)^{q(h)+\omega(h,m)}M_m.
\]
Consequently $P_h\otimes P_h$ is diagonal in the displayed symmetric
basis with these signs. Its expectation in the normalized vector
$x\otimes x/\norm{x}^2$ is also the square of the single-copy
expectation. This proves the second identity.
\end{proof}

\begin{scope}[Two copies are a mathematical realization]
The projectors onto the $E_m$ define valid mathematical probabilities. Their
use in this proof does not assert that an independently prepared two-copy
physical experiment occurs in the renewal system. Such an occurrence would
require its own source-complete preparation and comparison.
\end{scope}

\subsection{A period-free arithmetic source formula}
For distinct branch points $i,j$, let $P_{ij}$ denote either Hermitian
unitary lift of the nonzero two-torsion class $[W_i-W_j]$. Let
$\rho_\Theta$ be the normalized squared-theta-section state in the
Heisenberg-unitary metric. Define
\begin{equation}
 \mu_{ij}=\big(\Tr(\rho_\Theta P_{ij})\big)^2.
 \label{eq:v3-mu-definition}
\end{equation}
These expectations are real, so the right side is also their absolute
square.

\begin{theorem}[Branch-weight reconstruction of all displacement magnitudes]
\label{thm:v3-thomae-source}
For every smooth genus-two curve and every squared theta section as above,
\begin{equation}
 \boxed{\displaystyle
 \mu_{ij}=1-2\sum_{\substack{S\mid S^c\in\mathcal P\\
                      i,j\text{ in the same triple}}}p_S.}
 \label{eq:v3-magnitude-formula}
\end{equation}
The ten $p_S$ are the two-copy probabilities in \cref{lem:v3-bell} for the
actual section source. In particular, all fifteen squared displacement
expectations are determined by the branch points, without periods or theta
phase choices. They satisfy
\begin{equation}
 0\le\mu_{ij}<1,\qquad
 \sum_{i<j}\mu_{ij}=3,\qquad
 \sum_{j\ne i}\mu_{ij}=1\quad\text{for each fixed }i.
 \label{eq:v3-sums}
\end{equation}
\end{theorem}
\begin{proof}
Write $t$ for V2's four second-order theta constants. V2's addition formula
at zero gives
$Q_m(t)=\theta[m](0\mid\tau)^2$ for every even characteristic $m$.
The actual squared section has coefficients $r=M_\kappa t$ by
\cref{prop:v2-theta-coordinates}. Real orthogonality and Pauli commutation
imply $Q_m(r)=\pm Q_m(t)$ and $\norm r=\norm t$. Thus
\eqref{eq:v3-thomae-absolute} and Parseval give
\[
 \frac{|Q_{m_S}(r)|^2}{4\norm r^4}
 =\frac{c_\tau w_S}{\sum_Tc_\tau w_T}=p_S.
\]
This is the cancellation of the common period factor; the theta constants
themselves are not declared algebraic.

It remains to identify the eigenvalue sign intrinsically. In characteristic
coordinates the six branch points correspond to the six odd
characteristics $\kappa_i$, and
$h_{ij}=\kappa_i+\kappa_j$, $m_S=\sum_{k\in S}\kappa_k$.
For reference the six odd characteristics are
\[
 (01\mid01),\ (01\mid11),\ (10\mid10),\ (10\mid11),\
 (11\mid01),\ (11\mid10).
\]
Their sum is zero and every sum of three distinct entries is even, as is
verified directly by $q(\alpha,\beta)=\alpha\cdot\beta$.
For $h=h_{ij}$ the quadratic identity therefore gives
\[
 q(h)+\omega(h,\kappa_k)
 =\begin{cases}0,&k=i\text{ or }k=j,\\1,&k\ne i,j.\end{cases}
\]
Summing over the three indices in $S$ shows that
$q(h)+\omega(h,m_S)=1$ exactly when $i,j$ are in the same triple.
Substitution in \eqref{eq:v3-bell-identities} proves
\eqref{eq:v3-magnitude-formula}.

Nonnegativity follows from its squared-expectation realization. Exactly
four partitions place a fixed pair in one triple, and their weights are
positive; hence $\mu_{ij}<1$. Each partition contains six same-triple
pairs, so summing its signs over all fifteen pairs gives $15-2\cdot6=3$.
For a fixed vertex each partition places exactly two of the five other
vertices in its triple, giving $5-2\cdot2=1$. Averaging these identities
with the probabilities proves \eqref{eq:v3-sums}.
\end{proof}

\begin{corollary}[Exact inverse transform and no displacement stabilizer]
\label{cor:v3-inverse}
Set $L_{ij,S}=1$ for separated pairs and $L_{ij,S}=-1$ for same-triple
pairs. Then
\begin{equation}
 \mu=Lp,\qquad L^{\mathsf T}L=16I_{10}-\one\one^{\mathsf T},
 \qquad \boxed{p=\frac{\one+L^{\mathsf T}\mu}{16}.}
 \label{eq:v3-inverse}
\end{equation}
Thus the branch-weight packet and the complete squared-expectation packet
are equivalent data. No nonzero two-torsion displacement fixes the actual
theta source ray; its Heisenberg translation orbit has sixteen rays.
\end{corollary}
\begin{proof}
A sign column has squared norm fifteen. Two different $3+3$ partitions
disagree on eight pair signs, so their inner product is $15-16=-1$.
This proves the matrix identity, and $\sum_Sp_S=1$ gives the inverse.
A Hermitian involution fixes a pure ray only when its expectation has
absolute value one. This is excluded by \cref{thm:v3-thomae-source}.
The sixteen projective Heisenberg displacements therefore act freely on
that ray.
\end{proof}

\begin{corollary}[The five-source magnitudes in the inherited labeling]
\label{cor:v3-five-source}
Choose the branch at infinity as the distinguished odd characteristic.
Let $a,b$ run through the five finite branch points and let
$\Gamma_a$ be the inherited singular lift for $[W_a-W_\infty]$.
Then
\begin{equation}
 \boxed{\displaystyle
 |\Tr(\rho_\Theta\Gamma_a)|^2=\mu_{\infty a},\qquad
 |G_\Theta(a,b)|^2=\mu_{ab}\quad(a\ne b).}
 \label{eq:v3-five-source}
\end{equation}
Equivalently, $n_a^2=\mu_{\infty a}$ and $b_{ab}^2=\mu_{ab}$.
\end{corollary}
\begin{proof}
For $a\ne b$, the Hermitian bivector $i\Gamma_a\Gamma_b$ is, up to a
real sign, the Hermitian lift of
$[W_a-W_\infty]+[W_b-W_\infty]=[W_a-W_b]$.
Use \eqref{eq:gram-rho} and \cref{thm:v3-thomae-source}.
\end{proof}

\section{The actual curve: exact zeros and certified algebraic overlaps}
\label{sec:v3-actual-curve}

\subsection{A branch labeling fixed independently of the target}
Retain $g(x)=x^4-x-1$. Denote its negative and positive real roots by
$u,v$, respectively, and its upper and lower nonreal roots by $z,\bar z$.
The branch set is
\begin{equation}
 \mathcal B=(\infty,0,u,v,z,\bar z).
 \label{eq:v3-branch-label}
\end{equation}
The original temporal class is $[W_0-W_\infty]$; the four quartic roots
remain the residual tetrahedral labels. This labeling uses root location,
not a target overlap or fitted physical source.

Every partition is represented uniquely as $\{\infty,a,b\}\mid
(\mathcal B\setminus\{\infty,a,b\})$. Its weight is
\begin{equation}
 w_{ab}=|a-b|\prod_{\{c,d\}\subset\{0,u,v,z,\bar z\}\setminus\{a,b\}}
                   |c-d|.
 \label{eq:v3-explicit-weights}
\end{equation}
This is an explicit real algebraic definition of all numbers used below.

\begin{proposition}[Two exact zeros from the real branch geometry]
\label{prop:v3-zero}
For the actual source of $C_{\rm tet}$,
\begin{equation}
 \boxed{\mu_{\infty,0}=0,\qquad \mu_{u,v}=0.}
 \label{eq:v3-zero}
\end{equation}
Hence $n_t=0$ and the two source vectors $\Gamma_u\eta$ and
$\Gamma_v\eta$ are orthogonal. These statements hold more generally for a
branch configuration $(\infty,0,u,v,z,\bar z)$ with $u<0<v$ and
$\operatorname{Im}z>0$.
\end{proposition}
\begin{proof}
Put $U=-u>0$, $V=v>0$, $L=U+V$, $y=\operatorname{Im}z$,
$R_0=|z|$, $R_u=|z-u|$, $R_v=|z-v|$, and $R=R_0R_uR_v$.
The ten weights in \eqref{eq:v3-explicit-weights} are
\begin{align*}
 A=w_{0u}=2yUR_v^2,&\qquad B=w_{0v}=2yVR_u^2,\\
 C=w_{0z}=w_{0\bar z}=LR,&\qquad E=w_{uv}=2yLR_0^2,\\
 F=w_{uz}=w_{u\bar z}=VR,&\qquad
 G=w_{vz}=w_{v\bar z}=UR,\\
 H=w_{z\bar z}=2yUVL.&
\end{align*}
Expanding $|z+U|^2$ and $|z-V|^2$ gives
$A+B=E+H$, and $C=F+G$. Therefore
$Z_{\mathcal B}=2(A+B+2C)$.
The same-triple weights for $(\infty,0)$ sum to $A+B+2C$.
Those for $(u,v)$ sum to $E+H+2C$, the same number.
Both sums are $Z_{\mathcal B}/2$; apply
\eqref{eq:v3-magnitude-formula}.
\end{proof}

\begin{scope}[A zero temporal moment is not the central signature scalar]
The identity $n_t=0$ concerns the Hermitian displacement
$[W_0-W_\infty]$. It is not an evaluation of
$\langle\Omega,-\mathfrak j_{\rm St}\mathfrak j_{\rm chr}\Omega\rangle$
from \cite{AD9}. The latter involves a different common-source mixed
operator and remains uncomputed.
\end{scope}

\subsection{A rational certificate, rather than an uncertified root solver}
All four roots are isolated in disks of radius $r=10^{-28}$ centered at
\begin{align}
 c_u&=-0.724491959000515611588372282187,\notag\\
 c_v&=\phantom{-}1.220744084605759475361685349109,\notag\\
 c_z&=-0.248126062802621931886656533461
       +1.033982060975967756716863165809\,i,\notag\\
 c_{\bar z}&=\overline{c_z}.
 \label{eq:v3-root-centers}
\end{align}
Every displayed decimal is an exact rational, not a floating-point premise.
For a center $c$ set
\[
 M_c=|\operatorname{Re}c|+|\operatorname{Im}c|,\qquad
 d_c=\max\{|\operatorname{Re}(4c^3-1)|,
                  |\operatorname{Im}(4c^3-1)|\}.
\]
Direct rational arithmetic gives, at all four centers,
\begin{equation}
 \boxed{\displaystyle
 \frac{|\operatorname{Re}g(c)|+|\operatorname{Im}g(c)|
           +6M_c^2r^2+4M_cr^3+r^4}{d_cr}<\frac1{100}.}
 \label{eq:v3-rouche}
\end{equation}
The centers are pairwise more than $2r$ apart. The rational inequalities
\eqref{eq:v3-rouche}, including their unreduced fractions, are checked and
recorded by the accompanying certificate.

\begin{lemma}[Certified roots and arithmetic interval operations]
\label{lem:v3-interval}
Each disk in \eqref{eq:v3-root-centers} contains exactly one root of $g$.
The first two roots are real, of the indicated signs; the other two are
complex conjugates in the indicated half-planes. The following procedure
provides rigorous enclosures for \eqref{eq:v3-explicit-weights} and
\eqref{eq:v3-magnitude-formula} using only integer arithmetic.

For rational $x\ge0$ and $Q=10^{36}$ define
\[
 k=\left\lfloor\sqrt{\lfloor Q^2x\rfloor}\right\rfloor.
\]
Then $k/Q\le\sqrt x\le(k+1)/Q$, with the upper endpoint reduced to the
lower one when its square is $x$. Enclose distances between disk centers
by this rule, enlarge by the sum of their radii, and propagate interval
products, positive divisions, and sums outwards. Rounding every endpoint
outwards to the grid $Q^{-1}\Z$ preserves enclosure at each step.
\end{lemma}
\begin{proof}
On $|w|=r$,
\[
 g(c+w)=(4c^3-1)w+\bigl(g(c)+6c^2w^2+4cw^3+w^4\bigr).
\]
The first term has modulus at least $d_cr$; the second is bounded by the
numerator of \eqref{eq:v3-rouche}. Rouch\'e's theorem therefore gives
exactly one root in the disk. Four disjoint disks account for all four
roots. Complex conjugation preserves each real-centered disk, so its
unique root is real; it exchanges the remaining disks.

The integer-square-root inequality follows by squaring its endpoints.
Triangle and reverse triangle inequalities bound the change of distance
under perturbation inside the two disks. A product of real intervals is
enclosed by the minimum and maximum of its four endpoint products;
reciprocal endpoints reverse order on a positive interval. These rules,
and outward rounding, prove the propagation claim by induction over the
finite arithmetic expression.
\end{proof}

\begin{theorem}[Certified actual-curve source magnitudes]
\label{thm:v3-numbers}
The fifteen $\mu_{ij}$ of the actual theta source obey the enclosures in
\cref{tab:v3-moments}. The zeros are exact by \cref{prop:v3-zero}; every
other entry lies strictly between $1/50$ and $17/50$.
\end{theorem}
\begin{proof}
Apply \cref{lem:v3-interval} to the ten explicit weights and sum the four
same-triple terms for each pair. The rational interval certificate evaluates
these formulas with outward $10^{-36}$ rounding. The wider table endpoints
are rational outward roundings to $10^{-12}$. The exact-zero substitution
is licensed only by \cref{prop:v3-zero}, not by a small numerical value.
\end{proof}

\begin{table}[htbp]
\centering\small
\caption{Certified squared Hermitian displacement expectations. All
endpoints are rational; the pair labels are branch points, not Clifford
matrix indices. In the five-source Gram, the ten finite--finite pairs are
squared off-diagonal magnitudes; the five pairs with $\infty$ are $n_a^2$.}
\label{tab:v3-moments}
\begin{tabular}{lll}
\toprule
Pair & Lower bound & Upper bound\\\midrule
$(\infty,0)$ & $0$ & $0$\\
$(\infty,u)$ & $0.113137916901$ & $0.113137916902$\\
$(\infty,v)$ & $0.337737936950$ & $0.337737936951$\\
$(\infty,z),(\infty,\bar z)$ & $0.274562073073$ & $0.274562073074$\\
$(0,u)$ & $0.303217611526$ & $0.303217611527$\\
$(0,v)$ & $0.267258881088$ & $0.267258881089$\\
$(0,z),(0,\bar z)$ & $0.214761753692$ & $0.214761753693$\\
$(u,v)$ & $0$ & $0$\\
$(u,z),(u,\bar z)$ & $0.291822235786$ & $0.291822235787$\\
$(v,z),(v,\bar z)$ & $0.197501590980$ & $0.197501590981$\\
$(z,\bar z)$ & $0.021352346466$ & $0.021352346467$\\
\bottomrule
\end{tabular}
\end{table}

\begin{corollary}[Actual archimedean magnitude symmetry]
\label{cor:v3-symmetry}
The permutation group of the phase-free five-source packet that fixes the
temporal label $0$ is exactly
$\{1,(z\ \bar z)\}\cong C_2$.
\end{corollary}
\begin{proof}
The temporal row has three distinct squared magnitudes: those for $u$, $v$,
and the common value for $z,\bar z$. Their disjoint certified intervals
force $u$ and $v$ to be fixed. Complex conjugation of the branch configuration
exchanges $z,\bar z$ and leaves all absolute weights unchanged, so this
swap is indeed allowed. These are the only possibilities.
\end{proof}

This is a symmetry statement about one fixed complex embedding and its
positive source metric. It does not reduce the inherited arithmetic Galois
group $S_4$; its other elements act between marked arithmetic/complex
realizations rather than preserving these particular numbers in one fibre.

\section{Exact phase retrieval: thirty-two algebraic states and five signs}
\label{sec:v3-phase-retrieval}

\subsection{The finite solution problem}
From now on, Clifford indices $0,1,2,3,4$ refer respectively to the finite
branch labels $0,u,v,z,\bar z$. Retain exactly the matrices of
\cref{lem:quaternionic}, with $B_{ab}=i\gamma_a\gamma_b$.
Write
\begin{equation}
 a_{ab}=\sqrt{\mu_{\beta_a\beta_b}},\quad
 (\beta_0,\beta_1,\beta_2,\beta_3,\beta_4)=(0,u,v,z,\bar z).
 \label{eq:v3-amplitudes}
\end{equation}
These are explicitly defined nonnegative real algebraic numbers.
Exactly $a_{12}$ vanishes. Every possible marked pure-source Gram has
$b_{ab}=\pm a_{ab}$ on the other nine pairs. There are therefore at most
$2^9$ sign assignments before pure-state consistency.

V1 gives the exact test, which we use rather than reprove:
\begin{equation}
 D_b=\frac12\sum_{a<b}b_{ab}B_{ab},\qquad
 D_b^3=D_b,\quad\Tr(D_b^2)=2,\qquad
 \rho_b=\frac{D_b^2+D_b}{2}.
 \label{eq:v3-cubic-use}
\end{equation}
The trace condition holds for all the candidate signs, since
$\sum b_{ab}^2=2$ by \eqref{eq:v3-sums} with the distinguished infinity
vertex removed. The cubic condition remains essential.

\begin{lemma}[A free thirty-two-element sign action]
\label{lem:v3-sign-action}
Conjugation by projective Heisenberg operators acts on the bivector
expectations by
\begin{equation}
 b_{ab}\longmapsto\delta_a\delta_b b_{ab},
 \qquad \delta_a\in\{\pm1\},\quad\prod_{a=0}^4\delta_a=1.
 \label{eq:v3-delta}
\end{equation}
All sixteen such patterns occur. Adjoining quaternionic conjugation
$\rho\mapsto\Jq\rho\Jq^{-1}$ gives
\begin{equation}
 b_{ab}\longmapsto\lambda\delta_a\delta_b b_{ab},
 \qquad\lambda\in\{\pm1\},
 \label{eq:v3-sign-action}
\end{equation}
and a free action of order thirty-two on the present sign assignments. Each
orbit contains exactly one representative with
\begin{equation}
 b_{01},b_{02},b_{03},b_{04},b_{13}>0.
 \label{eq:v3-sign-normalization}
\end{equation}
Thus only sixteen normalized assignments remain, indexed by the signs of
$(b_{14},b_{23},b_{24},b_{34})$.
\end{lemma}
\begin{proof}
Conjugation by $\gamma_a$ fixes that generator and reverses the other four.
These five sign patterns generate the even-sign subgroup of rank four;
their product is identity and any four are independent. They therefore
produce all patterns in \eqref{eq:v3-delta}. V1 gives
$\Jq B_{ab}\Jq^{-1}=-B_{ab}$, proving \eqref{eq:v3-sign-action}.

For freeness, a stabilizing pattern must satisfy
$\lambda\delta_a\delta_b=1$ on the graph $K_5$ minus its single edge
$12$. A triangle forces $\lambda=1$. Connectedness then forces all
$\delta_a$ equal, and their product being one in odd dimension forces
them all to be $1$.

For existence and uniqueness of the normalized representative, let
$s_{ab}=\operatorname{sgn}b_{ab}$. Set
\[
 \lambda=s_{01}s_{03}s_{13},\qquad
 \delta_0=\prod_{j=1}^4s_{0j},\qquad
 \delta_j=\lambda\delta_0s_{0j}\quad(1\le j\le4).
\]
These are the unique signs obeying the product condition and making the
five entries in \eqref{eq:v3-sign-normalization} positive. Four other
nonzero pair signs remain; the missing edge $12$ is zero.
All these transformations preserve pure-state realizability, because they
are induced by unitary or antiunitary conjugation of density matrices.
\end{proof}

\subsection{Finite exclusions and an exact constructive solution}
For a normalized assignment use
$(\varepsilon_{14},\varepsilon_{23},\varepsilon_{24},\varepsilon_{34})$
for its four remaining signs. In \cref{tab:v3-witnesses}, each row is a
rational enclosure for one real or imaginary matrix entry of $D_b^3-D_b$.
Matrix entries are numbered $1$ through $4$ in the explicit Pauli basis of
V1. The rational arithmetic rules are precisely those of
\cref{lem:v3-interval}, now also applied to $\sqrt{\mu_{ij}}$ and to
complex matrix products. The intervals exclude zero.

\begin{table}[htbp]
\centering\small
\caption{Fifteen exact interval exclusion witnesses. Decimal bounds are
rational outward enclosures, not floating-point equality tests.}
\label{tab:v3-witnesses}
\begin{tabular}{ccc}
\toprule
$(\varepsilon_{14},\varepsilon_{23},\varepsilon_{24},\varepsilon_{34})$
& Entry of $D_b^3-D_b$ & Enclosure\\\midrule
$(-,-,-,-)$ & $\operatorname{Re}(2,2)$ & $[0.158,0.159]$\\
$(-,-,-,+)$ & $\operatorname{Im}(2,3)$ & $[-0.096,-0.095]$\\
$(-,-,+,-)$ & $\operatorname{Im}(1,3)$ & $[-0.037,-0.036]$\\
$(-,+,-,-)$ & $\operatorname{Im}(1,3)$ & $[-0.161,-0.160]$\\
$(-,+,-,+)$ & $\operatorname{Im}(1,3)$ & $[-0.125,-0.124]$\\
$(-,+,+,-)$ & $\operatorname{Re}(1,1)$ & $[-0.159,-0.158]$\\
$(-,+,+,+)$ & $\operatorname{Re}(1,4)$ & $[0.095,0.096]$\\
$(+,-,-,-)$ & $\operatorname{Re}(2,2)$ & $[0.137,0.138]$\\
$(+,-,-,+)$ & $\operatorname{Re}(1,1)$ & $[0.137,0.138]$\\
$(+,-,+,-)$ & $\operatorname{Re}(1,4)$ & $[0.150,0.151]$\\
$(+,-,+,+)$ & $\operatorname{Re}(1,4)$ & $[0.151,0.152]$\\
$(+,+,-,-)$ & $\operatorname{Im}(2,3)$ & $[0.151,0.152]$\\
$(+,+,-,+)$ & $\operatorname{Im}(2,3)$ & $[0.150,0.151]$\\
$(+,+,+,-)$ & $\operatorname{Im}(2,3)$ & $[0.247,0.248]$\\
$(+,+,+,+)$ & $\operatorname{Re}(1,4)$ & $[0.247,0.248]$\\
\bottomrule
\end{tabular}
\end{table}

\begin{theorem}[Complete pure-state solution set of the actual magnitude data]
\label{thm:v3-thirtytwo}
The squared bivector data acquired for $C_{\rm tet}$ have exactly
thirty-two pure-state realizations in the fixed Clifford frame. They form
one orbit under \eqref{eq:v3-sign-action}. Its unique representative
satisfying \eqref{eq:v3-sign-normalization} has signs
\begin{equation}
 \boxed{(b_{01},b_{02},b_{03},b_{04},b_{12},b_{13},b_{14},b_{23},b_{24},b_{34})
       =(+,+,+,+,0,+,-,-,+,+)\odot(a_{ab})_{a<b}.}
 \label{eq:v3-canonical-signs}
\end{equation}
In particular, putting these explicit algebraic coefficients into
\eqref{eq:v3-cubic-use} gives an exact rank-one density matrix
$\rho_\star$. Every realization is
\begin{equation}
 \boxed{\rho_{\delta,\lambda}
    =\frac14\left(I+\sum_a\delta_a n_a^\star\gamma_a
           +\lambda\sum_{a<b}\delta_a\delta_b b_{ab}^\star B_{ab}\right),
       \quad \prod_a\delta_a=1,\quad\lambda=\pm1.}
 \label{eq:v3-states}
\end{equation}
The actual marked arithmetic source is one of them. The theorem does not
identify which marking corresponds to each of the other thirty-one states.
\end{theorem}
\begin{proof}
By \cref{thm:v3-thomae-source}, an actual pure theta source with these
magnitudes exists. The inherited irreducible theta module can be put in
V1's Clifford frame after the real lift signs and central-volume convention
are fixed. This is an operator marking, not a fit of source coefficients.
Its sign vector lies among the $2^9$ assignments. By
\cref{lem:v3-sign-action} some normalized representative must be pure.

\Cref{tab:v3-witnesses} excludes fifteen of the sixteen normalized
representatives by the necessary exact cubic equation. The only remaining
one is $(-,-,+,+)$, giving \eqref{eq:v3-canonical-signs}. Existence forces
this remaining representative to satisfy the cubic equation \emph{exactly}.
The trace equation was already established from the acquired magnitudes.
V1's theorem now proves that $\rho_\star=(D_\star^2+D_\star)/2$ is
positive, trace one and rank one, and that it reconstructs the designated
$b^\star$. This step does not infer equality from an interval containing
zero.

The sign action preserves pure-state realizability and is free, so it
produces thirty-two distinct solutions. Any further solution would have a
normalized representative among the fifteen excluded rows, which is
impossible. Formula \eqref{eq:v3-states} follows from the trace-orthogonal
operator expansion in V1: unitary reorientation changes $n_a$ by $\delta_a$
and $b_{ab}$ by $\delta_a\delta_b$, while quaternionic conjugation fixes
$n$ and reverses $b$.
\end{proof}

\begin{remark}[What has become algebraic]
The entries of $D_\star$ and $\rho_\star$, and those of all the other
realizations, are algebraic in the fixed real/complex embedding. They are
given by root-isolated branch distances, positive square roots, and the
polynomial formula \eqref{eq:v3-cubic-use}. This is an exact specification,
not a decimal reconstruction. It neither asserts that a period matrix is
algebraic nor selects the five residual marking signs.
\end{remark}

\subsection{A sharp finite ambiguity and a coarse acquisition threshold}

\begin{theorem}[Separation and five-sign selection]
\label{thm:v3-five-bits}
The minimum Hilbert--Schmidt distance between distinct states in
\eqref{eq:v3-states} is
\begin{equation}
 d_{\min}=\sqrt{2(\mu_{\infty u}+\mu_{0u})},\qquad
 \boxed{0.912530030659<d_{\min}<0.912530030660.}
 \label{eq:v3-separation}
\end{equation}
The five signs of $b_{01},b_{02},b_{03},b_{04},b_{13}$ identify exactly one
state. Every one of these five expectations has absolute value greater
than $0.46$, so an independent estimate of each with absolute error less
than $9/20$ determines its sign. Five is the minimum possible worst-case
number of binary answers for identifying an arbitrary member of this
thirty-two-element solution set. Alternatively, a density estimate known
to lie within Hilbert--Schmidt distance $9/20$ of the actual state selects
at most one member.
\end{theorem}
\begin{proof}
The transitive unitary/antiunitary sign action preserves the real trace
pairing. It therefore suffices to compare $\rho_\star$ to its thirty-one
nonidentity images. Orthogonality of the sixteen matrices in V1 gives
\begin{equation}
 F_{\delta,\lambda}
 :=\Tr(\rho_\star\rho_{\delta,\lambda})
 =\frac14\left(1+\sum_a\delta_a(n_a^\star)^2
       +\lambda\sum_{a<b}\delta_a\delta_b(b_{ab}^\star)^2\right).
 \label{eq:v3-fidelity}
\end{equation}
The squared entries are precisely the certified $\mu$ values. Direct
rational interval comparison over the thirty-one sign choices shows that
the unique maximum occurs at
$\delta=(-1,-1,1,1,1)$ and $\lambda=-1$ and belongs to
$[0.583644471572,0.583644471573]$. To simplify its value, use
$\sum n_a^2=1$, $\sum b_{ab}^2=2$, and
\[
 \sum_{b\ne a}b_{ab}^2=1-n_a^2,
\]
the last identity following from the fixed-vertex sum in
\eqref{eq:v3-sums}. Substitution into \eqref{eq:v3-fidelity} gives
$F_{\max}=1-n_0^2-n_1^2-b_{01}^2
=1-\mu_{\infty u}-\mu_{0u}$, since $n_0=0$.
For two pure states, $\norm{\rho-\sigma}_{\HS}^2=2(1-\Tr\rho\sigma)$.
This proves the exact formula and the certified enclosure in
\eqref{eq:v3-separation}.

The normalization calculation in \cref{lem:v3-sign-action} shows that the
five selected signs give a bijection from the sign action to $\{\pm1\}^5$.
The corresponding five magnitudes are bounded below by $0.46$ using
\cref{tab:v3-moments}. An error less than $0.45$ therefore cannot reverse
their signs. At most $2^r$ leaves can occur in a binary decision tree of
worst-case depth $r$, so distinguishing thirty-two possibilities needs
$r\ge5$. This is a bound for binary answers, not a claim that every
real-valued scalar observation conveys at most one bit.
Finally $d_{\min}>0.9$ makes the open balls of radius $0.45$ about these
states disjoint.
\end{proof}

\begin{scope}[No hidden source replacement]
$\rho_\star$ is a declared canonical representative of the finite
magnitude fibre. It is not asserted to be the actual independently marked
$\vartheta_{\kappa_\infty}^2$ state. Quaternionic conjugation is used here
to classify possible pure density matrices, not as an assumed physical
channel. Selection of the actual sheet requires the five marked signs or
equivalent independently aligned theta/source data. Matching magnitudes
alone is insufficient for a source-fixing unitary comparison.
\end{scope}

\section{V3 consequences, verification, and remaining gates}
\label{sec:v3-status}

\subsection{The revised arithmetic target}
The next acquisition no longer requires high-precision reconstruction of
all ten complex overlaps. Their magnitudes are exact algebraic quantities
already acquired here. One may use V2's theta series in any independently
fixed period/Heisenberg marking merely to determine the five signs in
\cref{thm:v3-five-bits}. Its certified approximation needs only cross the
stated coarse sign thresholds, after the displacement frame has been aligned.
An algebraic theta-structure descent that determines those same signs would
be an alternative. Neither route may choose its phase sheet by inspecting
the desired renewal source.

For an independently acquired renewal source, the first comparison is now
an explicit labeled matrix of algebraic magnitudes. A discrepancy cannot be
removed by real reorientation or phase changes of the source rays. If the
magnitudes agree and both sources are in the stated pure Clifford category,
the finite solution set and the five signs decide their full marked density
comparison. Transfer equality still requires the actual arithmetic and
renewal operators and their mixed source moments. No transfer is manufactured
from a successful static calculation.

The residual Galois $S_4$, the fixed-fibre magnitude stabilizer $C_2$, and
the thirty-two-member phase-retrieval fibre are different groups/sets. Their
orders are not dimensions of matter carriers. Absolute values in
\eqref{eq:v3-branch-weights} depend on the chosen archimedean embedding;
an arbitrary arithmetic Galois map is not assumed to commute with this
Hermitian normalization. This is consistent with V2's distinction between
semilinear descent and unitary transport.

\subsection{Verification and its trust boundary}
The accompanying file
\path{Arithmetic_Loading_Duality_Research_Notes_V3_checks.py} uses SymPy
for exact finite Pauli and characteristic identities. Its root, moment,
sign-exclusion, and separation certificates use only rational arithmetic,
integer square roots, and outward rounding. No floating-point or
transcendental period/theta evaluation is used in these certificates. The
machine-readable output
\path{Arithmetic_Loading_Duality_Research_Notes_V3_certificate.json}
retains the root disks, rational Rouch\'e inequalities, moment enclosures,
and unrounded rational exclusion witnesses.

The analytic input remains the explicitly cited classical Thomae theorem;
it is not verified by a numerical test. The proof of exact existence of the
surviving sign representative is the existence/exclusion argument of
\cref{thm:v3-thirtytwo}, not an assertion that a small residual equals zero.
The certificate is a reproducible finite computer-assisted calculation, not
a proof-assistant formalization. The script also compares the retained V1
and V2 mathematical bodies byte for byte and records their SHA-256 digests.

\subsection{Proved in V3}
\begin{enumerate}[label=\textup{(P3.\arabic*)},leftmargin=3.6em]
\item The ten normalized branch Vandermonde magnitudes are the actual
squared-theta-section's symmetric two-copy probabilities. A finite signed
incidence transform recovers every squared Hermitian displacement
expectation, with an exact inverse transform.
\item For $C_{\rm tet}$ all fifteen squared expectations are now acquired as
explicit real algebraic numbers with rational enclosures. The temporal
first moment and the $u$--$v$ source overlap vanish exactly.
\item The pointed archimedean magnitude packet has exactly the residual
permutation symmetry exchanging the conjugate quartic roots.
\item Its full pure-state magnitude fibre has exactly thirty-two algebraic
states, explicitly reconstructed by V1's polynomial formula. Fifteen
nonzero interval witnesses prove exhaustion of the possible sign orbits.
\item Five marked signs select one state, with quantitative acquisition and
state-separation margins. This replaces an unevaluated continuous overlap
problem by a finite phase-selection problem for the stated actual curve.
\end{enumerate}

\subsection{Inherited and not rederived}
V1 and V2, including both historical ledgers, remain unchanged. The
chronology/Peano loading, arithmetic cutoff calculus, pointed genus-two
torsion packet, integral arithmetic--$K_4$ lattice, operator-source
comparison, fixed-source symmetry obstructions, squared-section addition
formula, generic Krylov/Tannaka/interface theorems, and all companion
physical reconstructions are imported at their original scope. First
Thomae is a standard analytic/algebraic input, not a new claim of this
programme.

\subsection{Unresolved}
The actual five signs in an independently aligned arithmetic theta frame
have not been selected here. Neither the actual arithmetic descent maps nor
their common Hermitian transport have been identified with the renewal
frame-cover maps. The central Store/chronology signature scalar, a native
arithmetic transfer, Frobenius--Wilson compatibility, open-interface
comparison, and regulator-uniform estimates remain unresolved. No physical
generation count, parameter value, prime-distribution estimate, or Einstein
field equation follows from this static theta-source calculation.

\paragraph{Version integrity.}
This continuation is appended after the complete V2 mathematical body.
Both earlier bodies and their historical status ledgers are preserved
byte for byte. Later iterations must append their new proofs and explicit
corrections, if needed, before the bibliography rather than silently
rewriting any earlier mathematical body.
% ====================== END IMMUTABLE V3 MATHEMATICAL BODY ======================
% APPEND V4 AND LATER ITERATIONS HERE.


% ===================== BEGIN IMMUTABLE V4 MATHEMATICAL BODY =====================
\section{V4 continuation: complex arithmetic phase and real descent}
\label{sec:v4-context}

V3 acquired absolute branch weights and classified the resulting pure-source
magnitude fibre. The fibre has thirty-two elements, in two Heisenberg orbits
exchanged by quaternionic conjugation. It did not use the complex arguments
of the branch differences. This continuation retains that discarded
arithmetic information before performing any period integration.

The principal result is that the complex first Thomae identity selects
\emph{one} of those two orbits for the actual embedded curve
$y^2=x(x^4-x-1)$. Thus the arithmetic source is now completely acquired up to
the sixteen-element Heisenberg marking action. Its five displaced source
\emph{rays}, with their ordered labels and positive Hermitian geometry, have
one unitary equivalence class. Four independently marked signs, rather than
five, specify its representative in an external fixed operator frame. A
phase-invariant triangle overlap distinguishes this class from the excluded
quaternionic class. We also reconstruct the actual archimedean conjugation
on the source, modulo the same marking action, as an explicit antiunitary
involution. No claim of a fixed-source $S_4$ unitary symmetry is reinstated.

\paragraph{Dependencies and nonduplication.}
The V3 absolute Thomae formula, fifteen squared expectations, two exact
zeros, thirty-two-state classification, and separation theorem are used,
not rederived. The new classical input is the complex phase retained in
first Thomae: the theta constant has an eighth-root-of-unity multiplier, so
its eighth power is a common period factor times the square of a branch
Vandermonde product. We use exactly the first Thomae theorem in
\cite[Section 2.6, equation (13)]{BernatskaThomae}; it was checked in the
original paper, including the condition $\epsilon^8=1$. The finite
bilinear identity used with it is the specialization of the already proved
V2 addition formula. Neither the companion double-centralizer theorems nor
the generic Gram/Krylov recognition machinery supplies the complex phase
calculated here. The pointed arithmetic source itself remains the one
specified in \cite[\texttt{eq:v9-theta-source-ray}]{AD9}.

\paragraph{Conventions that remain fixed.}
The branch labels are those of \cref{eq:v3-branch-label},
\[
 (\mathfrak b_0,\ldots,\mathfrak b_5)=(\infty,0,u,v,z,\bar z),
 \qquad u<0<v,\quad\operatorname{Im}z>0.
\]
Clifford indices $0,1,2,3,4$ still denote $t,u,v,z,\bar z$ respectively;
branch index $0$ and Clifford index $0$ therefore have different meanings.
We retain the explicit Pauli realization \eqref{eq:pauli-realization}, its
volume $\gamma_0\gamma_1\gamma_2\gamma_3\gamma_4=-I$, and V3's algebraic
representative $\rho_\star$. The Hermitian metric is Heisenberg-unitary as
in V1. All transposes below refer to the displayed Schr\"odinger basis;
bilinear forms must be transported by congruence under a basis change.

\section{Complex Thomae ratios as marking-invariant source data}
\label{sec:v4-complex-thomae}

\subsection{Bilinear covariants and the branch-characteristic dictionary}

Write a binary characteristic as $m=(a_1,a_2,b_1,b_2)=(a,b)$, put
$q(m)=a\cdot b\pmod2$, and retain the real matrices and Hermitian
Pauli displacements of V3:
\begin{equation}
 (M_m r)_\sigma=(-1)^{\sigma\cdot b}r_{\sigma+a},
 \qquad P_m=i^{q(m)}M_m,
 \qquad \sigma\in\F_2^2.
 \label{eq:v4-real-pauli}
\end{equation}
For $q(m)=0$, $M_m$ is real, symmetric and orthogonal. Define the
bilinear covariant
\begin{equation}
 Q_m(r)=r^{\mathsf T}M_mr.
 \label{eq:v4-bilinear}
\end{equation}
This is not the Hermitian expectation $r^*P_mr$. Its projective ratios will
supply information absent from their absolute values.

The binary labels of the five Clifford rays, in their fixed order, are
\begin{equation}
 h_0=0011,\quad h_1=1000,\quad h_2=1010,\quad
 h_3=0110,\quad h_4=0111.
 \label{eq:v4-h-labels}
\end{equation}
Here a four-digit word is ordered as $(a_1,a_2,b_1,b_2)$. The characteristic
of the marked odd theta divisor is
\[
 k_\infty=1101.
\]
Indeed this is the unique odd $k$ for which all $k+h_a$ are odd. Accordingly
the six odd characteristics attached to the six branch points are
\begin{equation}
 (k_0,k_1,k_2,k_3,k_4,k_5)
 =(1101,1110,0101,0111,1011,1010).
 \label{eq:v4-odd-dictionary}
\end{equation}
They give the following ten even characteristics $m_S=\sum_{j\in S}k_j$:
\begin{center}
\begin{tabular}{cc@{\qquad}cc}
\toprule
$S$ & $m_S$ & $S$ & $m_S$\\\midrule
$012$&$0110$&$024$&$0011$\\
$013$&$0100$&$025$&$0010$\\
$014$&$1000$&$034$&$0001$\\
$015$&$1001$&$035$&$0000$\\
$023$&$1111$&$045$&$1100$\\\bottomrule
\end{tabular}
\end{center}
Each $S$ denotes its unordered partition $S\mid S^c$. The dictionary is
fixed from the operator labels and odd characteristic before evaluating the
source; it is not obtained by fitting theta constants to desired overlaps.
It follows by binary addition and the branch-characteristic correspondence
already used in V3. In particular,
\begin{equation}
 S_0=012,\qquad T_0=014,\qquad
 M_{m_{S_0}}=\gamma_3,\qquad M_{m_{T_0}}=\gamma_1.
 \label{eq:v4-test-covariants}
\end{equation}
An analytic symplectic marking compatible with these binary labels can be
used to express the squared section in the V2 second-order theta basis.
Changing its higher theta marking while keeping the Pauli labels fixed
amounts projectively to the Heisenberg action characterized below.

\begin{lemma}[Heisenberg invariance and quaternionic conjugation]
\label{lem:v4-phase-transform}
For even $m,n$ with $Q_n(r)\ne0$, define
\begin{equation}
 \mathcal R_{mn}([r])=\frac{Q_m(r)^4}{Q_n(r)^4}.
 \label{eq:v4-source-ratio}
\end{equation}
This is independent of the scalar representative of $[r]$, invariant under
$[r]\mapsto[P_h r]$ for every $h\in\F_2^4$, and satisfies
\begin{equation}
 \mathcal R_{mn}([\Jq r])
 =\overline{\mathcal R_{mn}([r])}.
 \label{eq:v4-ratio-conjugation}
\end{equation}
For a normalized pure density matrix $\rho=rr^*$, put
\begin{equation}
 H_{mn}(\rho)=\Tr(M_m\rho M_n^*\rho^{\mathsf T}),
 \qquad \mathcal C_{mn}(\rho)=H_{mn}(\rho)^4.
 \label{eq:v4-coherence}
\end{equation}
Then
\begin{equation}
 H_{mn}(\rho)=Q_m(r)\overline{Q_n(r)},\qquad
 \mathcal C_{mn}(\rho)=|Q_n(r)|^8\mathcal R_{mn}([r]).
 \label{eq:v4-coherence-ratio}
\end{equation}
The same invariance and conjugation laws hold for $\mathcal C_{mn}$.
\end{lemma}
\begin{proof}
A scalar multiplication of $r$ multiplies both $Q_m$ and $Q_n$ by its
square. Direct multiplication of the signed permutation matrices
\eqref{eq:v4-real-pauli} gives
$P_h^{\mathsf T}M_mP_h=\epsilon(h,m)M_m$, with
$\epsilon(h,m)\in\{1,-1\}$. Fourth powers remove these signs. By V1,
$\Jq=C K_0$, where $C=\gamma_2\gamma_4$ is a real Pauli matrix up to
sign. The same congruence identity gives
$Q_m(\Jq r)=\epsilon_m\overline{Q_m(r)}$, proving
\eqref{eq:v4-ratio-conjugation}.
Finally,
\[
 \Tr(M_mrr^*M_n^*\bar r r^{\mathsf T})
 =(r^{\mathsf T}M_mr)(r^*M_n^*\bar r),
\]
which is \eqref{eq:v4-coherence-ratio}. The transformations of
$\mathcal C_{mn}$ follow either from this expression or from the same
congruence calculation.
\end{proof}

\begin{scope}[Bilinear marking is not an unmentioned Hermitian observable]
The transpose in \eqref{eq:v4-coherence} is part of a fixed theta marking.
Under $r\mapsto Ur$, the matrix of $Q_m$ becomes
$U^{-\mathsf T}M_mU^{-1}$, not $UM_mU^*$. The invariant is therefore
portable under a transported bilinear frame, not under an arbitrary
unrecorded choice of transpose. Its independence of Heisenberg marking is
the precise invariance proved above. It is not asserted to be an
unacquired physical measurement in the renewal system.
\end{scope}

\subsection{Removing the classical root-of-unity ambiguity}
For each partition define the complex product
\begin{equation}
 \Delta_S=\prod_{\{i,j\}\subset S,\ i<j}[ij]
           \prod_{\{i,j\}\subset S^c,\ i<j}[ij].
 \label{eq:v4-complex-delta}
\end{equation}
The homogeneous brackets are those of V3. Thus $w_S=|\Delta_S|$. Reversing
individual bracket conventions changes a sign, which will disappear on
squaring. A projective change of the six points multiplies every
$\Delta_S$ by the same nonzero factor, as in V3's projective-invariance
proof.

\begin{theorem}[Period-free complex source ratio]
\label{thm:v4-complex-thomae}
For the actual squared theta section in a compatible marked Schr\"odinger
realization, all ten $Q_{m_S}$ are nonzero and
\begin{equation}
 \boxed{\mathcal R_{m_Sm_T}(\rho_\Theta)
       =\left(\frac{\Delta_S}{\Delta_T}\right)^2.}
 \label{eq:v4-thomae-phase}
\end{equation}
In the normalization $\Tr\rho_\Theta=1$ one also has
\begin{equation}
 \boxed{\mathcal C_{m_Sm_T}(\rho_\Theta)
       =(4p_T)^4\left(\frac{\Delta_S}{\Delta_T}\right)^2.}
 \label{eq:v4-thomae-coherence}
\end{equation}
The right sides require only the embedded branch points. No period matrix,
choice of an eighth root, or target-dependent sign is used in these ratios.
\end{theorem}
\begin{proof}
Let $t=(t_\sigma)$ be the second-order theta-null vector of V2. Its addition
identity at $z=0$ gives, for even $m$, $Q_m(t)=\theta[m](0;\tau)^2$, with
at most the harmless real sign from the fixed binary characteristic
convention. V2 also gives $r=M_{k_\infty}t$ for the squared odd section,
up to one scalar. Congruence by this signed permutation changes each
$Q_m$ by a sign and the same common scalar square.

The classical first Thomae formula is
\[
 \theta[m_S](0;\tau)
 =\epsilon_S c(\tau)\Delta_S^{1/4},\qquad \epsilon_S^8=1,
\]
where the period factor $c(\tau)$ is the same for every partition
\cite[Section 2.6]{BernatskaThomae}. Taking eighth powers removes both the
partition-dependent root of unity and the fourth-root choice. Dividing
by the corresponding expression for $T$ proves
\eqref{eq:v4-thomae-phase}. This is a use of the classical theorem, not a
new proof of it. Nonvanishing follows from distinctness of the six branch
points. By V3, $|Q_{m_T}(r)|^2=4p_T$ after normalizing $r$; now use
\eqref{eq:v4-coherence-ratio} to obtain
\eqref{eq:v4-thomae-coherence}.
\end{proof}

\section{Algebraic selection of the actual holomorphic phase orbit}
\label{sec:v4-orbit-selection}

\subsection{A separating ratio of two explicit partitions}
For \eqref{eq:v4-test-covariants}, cancellation of common factors in
\eqref{eq:v4-complex-delta} gives
\begin{equation}
 R_{\mathrm{ar}}
 :=\left(\frac{\Delta_{012}}{\Delta_{014}}\right)^2
 =\left[
 \frac{u(v-z)(z-\bar z)}{z(u-v)(u-\bar z)}
 \right]^2.
 \label{eq:v4-explicit-ratio}
\end{equation}
The chosen embedding, especially $\operatorname{Im}z>0$, is retained.
Conjugating that embedding conjugates this ratio.

\begin{lemma}[Strict complex-phase certificate]
\label{lem:v4-phase-certificate}
The actual branch ratio and V3's algebraic density representative satisfy
\begin{align}
 1.092297864426026&<\operatorname{Re}R_{\mathrm{ar}}
                  <1.092297864426027,\notag\\
 0.716248184244455&<\operatorname{Im}R_{\mathrm{ar}}
                  <0.716248184244456,
 \label{eq:v4-ratio-interval}\\
 0.074288820985750&<\operatorname{Re}\mathcal C_{m_{S_0}m_{T_0}}(\rho_\star)
                  <0.074288820985751,\notag\\
 0.048713116516679&<\operatorname{Im}\mathcal C_{m_{S_0}m_{T_0}}(\rho_\star)
                  <0.048713116516680.
 \label{eq:v4-coherence-interval}
\end{align}
All displayed endpoints are terminating rational decimals. In particular,
\begin{equation}
 \operatorname{Im}R_{\mathrm{ar}}>7/10,
 \qquad \operatorname{Im}\mathcal C_{m_{S_0}m_{T_0}}(\rho_\star)>6/125.
 \label{eq:v4-phase-margins}
\end{equation}
\end{lemma}
\begin{proof}
Use the four rational root disks already certified in V3. Propagate their
rectangular complex enclosures through \eqref{eq:v4-explicit-ratio}, using
$w^{-1}=\bar w/|w|^2$ with an independently positive lower bound on
$|w|^2$. This gives \eqref{eq:v4-ratio-interval} by outward-rounded
rational arithmetic.

For an independently checkable polynomial certificate of the second
calculation, write V3's canonical bivector entries as
\[
 (b_{01},b_{02},b_{03},b_{04},b_{12},b_{13},b_{14},b_{23},b_{24},b_{34})
 =(a,b,c,c,0,d,-d,-e,e,f),
\]
where the six letters denote their positive algebraic magnitudes, not new
free fitted parameters. Substitution of
$\rho_\star=(D_\star^2+D_\star)/2$ into
$H=\Tr(\gamma_3\rho_\star\gamma_1\rho_\star^{\mathsf T})$ gives the exact
polynomials
\begin{align}
 \operatorname{Re}H
 &=\tfrac12(abe f-2ace^2+ac+b^2df-2bcde+df),\notag\\
 -\operatorname{Im}H
 &=\tfrac14(a^2d+2abe+2acf+3b^2d+6c^2d
                    +2d^3+2de^2+df^2).
 \label{eq:v4-H-polynomials}
\end{align}
They follow by multiplication of the four-by-four matrices in V1 and are
also symbolically verified in the accompanying script. V3's rational
intervals for the six positive square roots give
\[
 0.078<\operatorname{Re}H<0.080,
 \qquad 0.539<-\operatorname{Im}H<0.542.
\]
If $H=x-iy$ with these positive $x,y$, then
$\operatorname{Im}H^4=4xy(y^2-x^2)>0$. Full outward-rounded substitution
gives \eqref{eq:v4-coherence-interval}. No square root of an uncertified
sign or eigenvector phase is required.

The companion machine-readable certificate records the exact rational
endpoints before decimal coarsening. An optional comparison of its intervals
with the independently computed expression in
\eqref{eq:v4-thomae-coherence} is a consistency check only. It is not used
as a proof that two almost equal algebraic numbers are equal.
\end{proof}

\begin{remark}[The branch-phase sign has a direct geometric explanation]
Write $z=\alpha+i\beta$, with $\beta>0$, and $r=|z|^2$. The quartic gives
$\alpha=-(u+v)/2$ and $uvr=-1$. The rough root bounds
$-3/4<u<-7/10$ and $6/5<v<5/4$ imply $\alpha<0$ and $-uv<1$.
Up to the common real sign convention,
\[
 \Delta_{012}=-u|v-z|^2(z-\bar z),\qquad
 \Delta_{014}=(v-u)\{\alpha(uv+3r)+i\beta(uv-r)\}.
\]
The first is purely imaginary and nonzero. Both the real and imaginary
parts of the second are negative, since $uv+3r>0$ and $uv-r<0$.
Consequently $\operatorname{Im}(\Delta_{012}/\Delta_{014})^2>0$.
This argument proves the sign without the fine interval calculation;
\eqref{eq:v4-ratio-interval} supplies the quantitative margin.
\end{remark}

\subsection{The phase fibre shrinks from thirty-two states to sixteen}

\begin{theorem}[The actual arithmetic source occupies one Heisenberg orbit]
\label{thm:v4-holomorphic-orbit}
In the independently fixed embedded branch labeling and compatible Pauli
convention above, the actual normalized squared-theta-section state belongs
to precisely the following half of V3's magnitude fibre:
\begin{equation}
 \boxed{\mathcal F_{\mathrm{ar}}
       =\{\rho_{\delta,+}:\delta_a\in\{\pm1\},\ \prod_{a=0}^4\delta_a=1\}
       =\{P_h\rho_\star P_h^*:h\in\F_2^4\}.}
 \label{eq:v4-selected-orbit}
\end{equation}
It has sixteen distinct elements. Every member obeys all the complex
ratios \eqref{eq:v4-thomae-phase}. None of the sixteen states
$\rho_{\delta,-}$ obeys even the single $S_0,T_0$ phase comparison.
In particular, the quaternionic alternative has been excluded by arithmetic
data alone, without computing a period matrix or using a renewal target.
\end{theorem}
\begin{proof}
V3 places the actual source among the thirty-two states
$\rho_{\delta,\lambda}$. By \cref{lem:v4-phase-transform}, the quantity
$\mathcal C_{m_{S_0}m_{T_0}}$ is constant on each of the two Heisenberg
orbits, and the two values are complex conjugates. By
\cref{lem:v4-phase-certificate}, its imaginary part is positive on the
orbit of $\rho_\star$ and negative on the orbit of $\Jq\rho_\star\Jq^{-1}$.
On the actual theta source, \cref{thm:v4-complex-thomae} makes this imaginary
part $(4p_{T_0})^4\operatorname{Im}R_{\mathrm{ar}}>0$. The negative orbit
is therefore impossible. Existence of the theta source now forces it into
the positive orbit.

Every member of that orbit has the same full set of ratios as the actual
source, by Heisenberg invariance. Thus all of those equalities are exact
consequences of existence, the finite V3 classification, and the strict
phase separation, rather than numerical equality tests. Distinctness and
the orbit size are inherited from V3's free sign action.
\end{proof}

\begin{corollary}[Four independent marked signs and a forced fifth sign]
\label{cor:v4-four-signs}
Let $s_i=\operatorname{sgn}b_{0i}$ for $1\le i\le4$. On
$\mathcal F_{\mathrm{ar}}$ these four signs determine the state uniquely,
with
\begin{equation}
 \delta_0=\prod_{i=1}^4s_i,\qquad \delta_i=s_i\delta_0.
 \label{eq:v4-four-to-delta}
\end{equation}
In particular,
\begin{equation}
 \boxed{\operatorname{sgn}b_{13}=s_1s_3,
 \qquad b_{01}b_{03}b_{13}>0.}
 \label{eq:v4-forced-fifth}
\end{equation}
All sixteen four-sign patterns occur. Four is consequently the minimum
worst-case number of binary answers for selection inside this phase-refined
fibre. Each of the four magnitudes is greater than $0.46$, so the V3
absolute error allowance $9/20$ still suffices for sign acquisition.
\end{corollary}
\begin{proof}
For $\lambda=+1$, V3 gives $b_{0i}=\delta_0\delta_i b_{0i}^\star$, with
$b_{0i}^\star>0$. Hence $s_i=\delta_0\delta_i$ and
$\prod_i s_i=\delta_0$ because $\prod_a\delta_a=1$.
This proves \eqref{eq:v4-four-to-delta}. Since $b_{13}^\star>0$,
its sign is $\delta_1\delta_3=s_1s_3$, proving
\eqref{eq:v4-forced-fifth}. The inverse formulas realize every sign word.
The binary counting bound and magnitude margins then follow exactly as
in V3, on the smaller, now arithmetically selected fibre.
\end{proof}

\begin{lemma}[The residual sixteenfold ambiguity is the marking action]
\label{lem:v4-marking}
A complex-linear unitary $U$ which preserves every labelled Pauli ray
$[P_m]$ is, up to one scalar phase, a unique $P_h$. Equivalently, for the
five Clifford generators its signs are exactly the sixteen even sign
patterns $\delta$ in \eqref{eq:v4-selected-orbit}. This action is free on
the present arithmetic source.
\end{lemma}
\begin{proof}
Hermiticity and $P_m^2=I$ force
$UP_mU^*=\epsilon(m)P_m$ with $\epsilon(m)\in\{\pm1\}$.
Multiplying the Pauli relations shows
$\epsilon(m+n)=\epsilon(m)\epsilon(n)$. Nondegeneracy of the symplectic
pairing gives a unique $h$ with $\epsilon(m)=(-1)^{\omega(h,m)}$.
Conjugation by $P_h$ has exactly those signs. Thus $P_h^*U$ commutes with
the full matrix algebra and is scalar. The volume identity forces
$\prod\delta_a=1$; the sixteen characters realize every such sign pattern.
Freeness on the source follows from V3, or because a stabilizing pattern
would obey $\delta_a\delta_b=1$ on every nonzero bivector edge. Those
edges form the connected graph $K_5\setminus\{12\}$, forcing all
$\delta_a$ equal, and their product forces $\delta_a=1$.
\end{proof}

\begin{scope}[Gauge fixing versus acquisition of an external frame]
The curve, embedded branch labels, and complex covariants now determine one
Heisenberg orbit, not the actual four signs of an independently fixed
physical frame. A source-adapted convention $b_{0i}>0$ chooses
$\rho_\star$ algebraically and without reference to the renewal target.
It is a valid gauge convention, but not a measurement of the relative
arithmetic--renewal marking. Any data invariant under the free marking
action have the same values on all sixteen states and cannot distinguish
them in a fixed external frame. The original five-sign statement in V3
remains correct for the larger magnitude-only fibre; the new complex
arithmetic input removes one of those binary choices.
\end{scope}

\section{Projective source recognition and an invariant triangle phase}
\label{sec:v4-projective-recognition}

\subsection{One arithmetic ray-source class}
For normalized source vectors $\xi_a=\gamma_a\eta$, define the cyclic
three-source overlap
\begin{equation}
 \mathcal B_{013}(G)=G_{01}G_{13}G_{30},\qquad
 \chi(G)=i\mathcal B_{013}(G)=b_{01}b_{13}b_{03}.
 \label{eq:v4-triangle}
\end{equation}
The last equality uses the V1 Hermitian convention. Both the cyclic
product and $\chi$ are invariant under arbitrary separate phase changes
of the three source vectors and under a common complex-linear unitary.
For the actual source, \cref{thm:v4-holomorphic-orbit} gives
\begin{equation}
 \boxed{\chi(G_\Theta)
   =\sqrt{\mu_{0u}\mu_{0z}\mu_{uz}}
   \in(0.137852571639642,\,0.137852571639643).}
 \label{eq:v4-actual-triangle}
\end{equation}
Here the subscripts on $\mu$ are branch names, not Clifford indices.

\begin{theorem}[Complete projective recognition on the actual arithmetic fibre]
\label{thm:v4-projective-test}
Let $\rho$ be an independently supplied normalized pure state in the same
five-generator Clifford category, after an independently specified
operator-label identification. Suppose its ten squared bivector
expectations are exactly V3's arithmetic values, including $b_{12}=0$.
Then the following are equivalent:
\begin{enumerate}[label=\textup{(\roman*)},leftmargin=2.8em]
\item its five ordered displacement source rays are related to the actual
      arithmetic source rays by one complex-linear unitary;
\item $\chi(G(\rho))>0$;
\item $\rho\in\mathcal F_{\mathrm{ar}}$.
\end{enumerate}
In the positive case the unitary may be chosen from the Heisenberg marking
action, with the corresponding five real column signs. A full
source-\emph{vector} comparison in an external marked frame additionally
requires the four signs of \cref{cor:v4-four-signs}.
\end{theorem}
\begin{proof}
By V3, $\rho=\rho_{\delta,\lambda}$ for one of the thirty-two classified
states. The triangle product transforms as
\[
 \chi(G(\rho_{\delta,\lambda}))
 =\lambda^3\delta_0^2\delta_1^2\delta_3^2\chi(G(\rho_\star))
 =\lambda\chi(G(\rho_\star)).
\]
Its positive value on $\rho_\star$ gives (ii)$\Leftrightarrow$(iii).
For (iii)$\Rightarrow$(i), if $\eta'=P_h\eta$, then
\[
 \gamma_a\eta'=\delta_aP_h\gamma_a\eta,
\]
so the five rays are intertwined by $P_h$. The actual source is in the
same orbit by \cref{thm:v4-holomorphic-orbit}. Conversely, (i) preserves
the cyclic overlap, including its imaginary sign. The arithmetic sign is
positive by \eqref{eq:v4-actual-triangle}, proving (ii). The marked claim
is \cref{cor:v4-four-signs}.
\end{proof}

This theorem closes a static \emph{projective} source-recognition gate.
It is not the earlier ordinary pointed-$S_4$ covariance statement. In
particular, the fixed-fibre ordinary and coherent projective $S_4$
obstructions of V1--V2 are unchanged.

\subsection{Quantitative separation under arbitrary ray rephasing}

\begin{proposition}[A phase-independent separation from the quaternionic class]
\label{prop:v4-triangle-margin}
Let $G_+$ and $G_-$ be any two normalized Grams in the selected and excluded
classes respectively. For every diagonal unitary $D$ on the five source
labels,
\begin{equation}
 \boxed{\norm{G_+-D^*G_-D}_{\HS}
 \ge 2\sqrt{\frac23}\sqrt{\mu_{0u}\mu_{0z}\mu_{uz}}
 >\frac9{40}.}
 \label{eq:v4-rephasing-margin}
\end{equation}
The explicit expression on the right belongs to
$(0.225112306831725,\,0.225112306831726)$.
\end{proposition}
\begin{proof}
Every entry of a normalized Gram has modulus at most one. For two such
Hermitian matrices $A,B$, telescope the three factors in
$\mathcal B_{013}$ to get
\[
 |\mathcal B_{013}(A)-\mathcal B_{013}(B)|
 \le |A_{01}-B_{01}|+|A_{13}-B_{13}|+|A_{30}-B_{30}|
 \le\sqrt{\frac32}\norm{A-B}_{\HS}.
\]
The last factor uses Hermiticity: the Hilbert--Schmidt square includes both
orientations of each of the three different edges. The two triangle
products have opposite imaginary signs and the same modulus, and diagonal
rephasing changes neither. Their difference is therefore
$2\sqrt{\mu_{0u}\mu_{0z}\mu_{uz}}$. Rearranging proves the first
inequality; outward-rounded substitution of V3's intervals proves the
strict rational margin.
\end{proof}

\begin{proposition}[A fixed-frame error bound for the complex coherence]
\label{prop:v4-coherence-stability}
For density matrices $\rho,\sigma$ in the same bilinear theta frame,
\begin{equation}
 |\mathcal C_{m_Sm_T}(\rho)-\mathcal C_{m_Sm_T}(\sigma)|
 \le 8\norm{\rho-\sigma}_{\HS}.
 \label{eq:v4-coherence-stability}
\end{equation}
Consequently an error less than $3/500$ about the actual source cannot
reverse the sign of the imaginary coherence for $S_0,T_0$.
\end{proposition}
\begin{proof}
The symmetric matrices $M_{m_S},M_{m_T}$ are unitary and every density
matrix has Hilbert--Schmidt norm at most one. Subtracting the two quadratic
expressions \eqref{eq:v4-coherence}, and using Hilbert--Schmidt
Cauchy--Schwarz, gives
$|H(\rho)-H(\sigma)|\le2\norm{\rho-\sigma}_{\HS}$ and $|H(\rho)|\le1$.
The elementary factorization of $x^4-y^4$ now gives
\eqref{eq:v4-coherence-stability}. On the actual orbit the imaginary
part is greater than $6/125$ by \cref{lem:v4-phase-certificate}, whereas
$8(3/500)=6/125$.
\end{proof}

\section{The arithmetic real involution on the acquired source}
\label{sec:v4-real-descent}

Complex conjugation is the Galois operation that can be interpreted
canonically as an antilinear operation at the fixed archimedean realization.
For this curve it fixes $\infty,0,u,v$ and exchanges $z,\bar z$.
It must not be identified with the quaternionic operator $\Jq$: the latter
fixes all Clifford labels, squares to $-I$, and takes the source to the
excluded phase orbit without exchanging the branch labels.

\subsection{An explicit involution in the canonical source gauge}
Set
\begin{equation}
 V=\frac{\gamma_3-\gamma_4}{\sqrt2},\qquad
 \mathcal K_{\mathrm{ar}}=\gamma_0V\Jq.
 \label{eq:v4-Kar}
\end{equation}
This is an antiunitary; $V^2=I$ and $\gamma_0V=-V\gamma_0$.

\begin{theorem}[Source-fixing real descent and its uniqueness]
\label{thm:v4-real-descent}
The operator \eqref{eq:v4-Kar} obeys
\begin{equation}
 \boxed{\mathcal K_{\mathrm{ar}}^2=I,
 \qquad\mathcal K_{\mathrm{ar}}\rho_\star
                     \mathcal K_{\mathrm{ar}}^{-1}=\rho_\star,}
 \label{eq:v4-real-fix}
\end{equation}
and its action on the five generators is
\begin{equation}
 \boxed{
 (\gamma_0,\gamma_1,\gamma_2,\gamma_3,\gamma_4)
 \longmapsto(-\gamma_0,\gamma_1,\gamma_2,\gamma_4,\gamma_3).}
 \label{eq:v4-real-rays}
\end{equation}
It is the unique projective antiunitary with this branch-label permutation
which fixes $\rho_\star$, even when the signs on the five rays are not
prescribed in advance. On the state $P_h\rho_\star P_h^*$ the corresponding
unique projective operation is $P_h\mathcal K_{\mathrm{ar}}P_h^*$.
\end{theorem}
\begin{proof}
The quaternionic structure commutes with each $\gamma_a$ and with the
real-coefficient vector $V$. Consequently
\[
 (\gamma_0V\Jq)^2=(\gamma_0V)^2\Jq^2=(-I)(-I)=I.
\]
Conjugation by $V$ negates $\gamma_0,\gamma_1,\gamma_2$ and sends
$\gamma_3\mapsto-\gamma_4$, $\gamma_4\mapsto-\gamma_3$.
Conjugating once more by $\gamma_0$ gives
\eqref{eq:v4-real-rays}. The antiunitarity negates the factor $i$ in
$B_{ab}=i\gamma_a\gamma_b$.

V3's exact conjugate-root equalities and canonical sign pattern give the
bivector array $(a,b,c,c,0,d,-d,-e,e,f)$ used in
\eqref{eq:v4-H-polynomials}. Applying \eqref{eq:v4-real-rays} to this
array, including the sign of $i$, leaves its density expectations unchanged:
$01,02$ are fixed, $03$ and $04$ are exchanged, $13$ goes to $-14$,
$23$ to $-24$, $34$ is fixed, and $12$ is zero. The transformed pure state
therefore has the same full bivector packet. V1's polynomial reconstruction
makes it the same density matrix, proving \eqref{eq:v4-real-fix}.

If $\mathcal K'$ were another antiunitary with the same label permutation
and invariant source, the linear unitary
$\mathcal K'\mathcal K_{\mathrm{ar}}^{-1}$ would preserve each labelled
Pauli ray and stabilize $\rho_\star$. By \cref{lem:v4-marking}, it is a
scalar. This proves projective uniqueness. Conjugating by $P_h$ proves
the last statement; its antilinear square remains one.
\end{proof}

In the concrete basis of V1 the same operation is $A K_0$, where
\begin{equation}
 A=-\frac1{\sqrt2}
 \begin{pmatrix}
 0&0&1+i&0\\
 0&0&0&1-i\\
 1+i&0&0&0\\
 0&1-i&0&0
 \end{pmatrix},\qquad A\overline A=I.
 \label{eq:v4-real-matrix}
\end{equation}
The source ray thus admits a representative fixed by
$\mathcal K_{\mathrm{ar}}$: if
$\mathcal K_{\mathrm{ar}}\eta=e^{i\theta}\eta$, multiplying $\eta$ by
$e^{i\theta/2}$ gives a fixed vector. This choice of scalar phase is not
a new state or an additional source direction.

\subsection{Why this is the arithmetic conjugation, rather than an engineered symmetry}

\begin{lemma}[Heisenberg-unitary metric is compatible with real descent]
\label{lem:v4-real-metric}
Let a geometrically defined antilinear involution $K$ on the theta module
normalize its finite Heisenberg group. Then its Heisenberg-invariant
positive Hermitian metric is $K$-invariant; in particular $K$ is antiunitary.
\end{lemma}
\begin{proof}
Write $h_K(x,y)=\overline{h(Kx,Ky)}$. Normalization of the finite group
makes $h_K$ Heisenberg invariant. Irreducibility makes it $c h$ for one
positive real $c$. Applying the involution twice gives $c^2=1$, so $c=1$.
This is exactly the antiunitary identity for $K$.
\end{proof}

\begin{corollary}[Identification of the actual archimedean descent row]
\label{cor:v4-actual-real-descent}
The curve and symmetric theta divisor defined by
$\kappa_\infty=\mathcal O(\infty)$ are defined over $\R$. Their natural
conjugation on $H^0(J,\mathcal L_\kappa^{\otimes2})$, after the same
Heisenberg marking and normalization used throughout these notes, is
projectively $P_h\mathcal K_{\mathrm{ar}}P_h^*$ on the actual source
$P_h\rho_\star P_h^*$. Thus one nontrivial arithmetic descent operation
has been reconstructed, modulo precisely the same four marking signs.
\end{corollary}
\begin{proof}
The real structure of the curve induces real structures on its Jacobian,
on the divisor selected by the real Weierstrass point at infinity, and on
the associated line bundle. On global sections this gives an antilinear
involution. It fixes the unique section line of $\mathcal L_\kappa$ and
hence its squared source ray. It transports two-torsion displacements by
conjugating their branch labels, and preserves the finite torsion lifts
in the theta group: conjugation preserves the defining multiplication and
orders, while Hermitian involutive lifts over a fixed label differ by a
real sign. Thus the finite Heisenberg group is normalized.
\Cref{lem:v4-real-metric} makes the operation antiunitary in the retained
metric. Its label permutation is $(z\ \bar z)$ and its source lies in
\eqref{eq:v4-selected-orbit}. The uniqueness statement in
\cref{thm:v4-real-descent} identifies it projectively as claimed.
\end{proof}

\begin{scope}[What one real descent row does not establish]
The conclusion is an actual archimedean conjugation statement, not a
unitary representation of the full residual $S_4$ on one fixed source.
It is not a native time-evolution operator, a Frobenius operator at a finite
prime, or a solution of the central Store/chronology signature comparison.
The scalar phase of an antilinear involution is immaterial for the density
and ray descent proved here; coherent tensor and finite-prime comparisons
must retain any further scalar cocycles they require. In particular the
rational temporal torsion label is fixed although its retained Hermitian
lift changes sign in \eqref{eq:v4-real-rays}. Label invariance and invariance
of a chosen lifted operator are different assertions.
\end{scope}

\section{V4 audit, precise progress, and next gates}
\label{sec:v4-status}

\subsection{Verification}
The script \path{Arithmetic_Loading_Duality_Research_Notes_V4_checks.py}
imports V3's rational interval implementation and independently calls its
root and branch-weight certificates. It checks the binary characteristic
dictionary, all even bilinear congruences, the sixteen Heisenberg sign
actions, the four-sign inverse, the explicit real structure, and its
formal bivector covariance with exact symbolic matrices. It computes
\eqref{eq:v4-ratio-interval} and \eqref{eq:v4-coherence-interval} with
outward-rounded rational intervals, and certifies the triangle and ray
separation bounds. The complete numerical output is retained in
\path{Arithmetic_Loading_Duality_Research_Notes_V4_certificate.json}.

No floating-point arithmetic, period quadrature or numerical theta series
enters the certificate. The classical complex Thomae formula remains a
cited input. Exact phase-orbit membership is the existence and strict
exclusion proof in \cref{thm:v4-holomorphic-orbit}; it is not a numerical
near-zero argument. The script is a reproducible computer-assisted check,
not a proof-assistant formalization. The three earlier immutable mathematical
bodies are compared byte for byte and their SHA-256 digests are retained.

\subsection{Proved in V4}
\begin{enumerate}[label=\textup{(P4.\arabic*)},leftmargin=3.6em]
\item Complex eighth-power Thomae ratios are exact, period-free,
      Heisenberg-marking-invariant source quantities. Unlike their absolute
      values, they distinguish quaternionic conjugates at this embedded curve.
\item A strict algebraic phase certificate selects the positive Heisenberg
      orbit of $\rho_\star$ among V3's thirty-two states. The actual source is
      acquired up to sixteen lift-marking choices; four marked signs select
      a representative and force the former fifth sign.
\item The complete ordered displacement ray source has one unitary
      equivalence class. An invariant triangle phase gives a complete
      projective source test within the certified pure magnitude fibre and
      a positive separation from the excluded class under arbitrary ray phases.
\item The actual real conjugation on the theta source is the explicitly
      reconstructed source-fixing antiunitary involution, modulo the same
      four signs. This supplies one genuine arithmetic descent row without
      creating an ordinary fixed-source tetrahedral symmetry.
\end{enumerate}

\subsection{Inherited}
All V1--V3 mathematical bodies, their hypotheses and historical status
ledgers remain unchanged. The chronology loading, finite source and
operator reconstruction results, integral $K_4$ bridge, ordinary and
signed-projective fixed-source obstructions, pure-state recognition,
absolute branch-weight acquisition, and generic transfer/Tannakian/interface
criteria are inherited. First Thomae and its eighth-root multiplier are
standard arithmetic-geometric input; the new result is their use to close
the concrete source-phase gate and reconstruct its real descent.

\subsection{Unresolved and the revised next attack}
The four relative lift signs in an \emph{independently prescribed renewal
frame} have not been measured. They are no longer needed to identify the
arithmetic source rays up to unitary equivalence; they remain needed for a
common marked source-vector and operator comparison. An independently
supplied renewal source must first pass the explicit arithmetic magnitude
and positive triangle tests, rather than have its frame or state selected
to match $\rho_\star$ by definition.

The next genuinely stronger arithmetic task is transport beyond complex
conjugation: identify the labelled semilinear descent and scalar cocycle on
the arithmetic frame cover and its independently supplied renewal analogue.
Alternatively, acquire a native arithmetic transfer and its inserted source
moments. Static source acquisition and one involution do not define such
a transfer. The central Store/chronology signature scalar, finite-prime
Frobenius--Wilson identification, open scattering/interface comparison and
regulator-uniform estimates remain open. No Standard-Model multiplicity,
physical coupling, prime-distribution estimate or Einstein equation is
inferred from the present phase and descent theorems.

\paragraph{Version integrity.}
This continuation is appended after the complete V3 mathematical body.
Every previous body and historical ledger is preserved byte for byte.
Subsequent iterations are to append their results and any explicit
corrections before the bibliography and advance the version filename.
% ====================== END IMMUTABLE V4 MATHEMATICAL BODY ======================
% ===================== BEGIN IMMUTABLE V5 MATHEMATICAL BODY =====================
\clearpage
\section{V5 continuation: actual Frobenius and the correct arithmetic carrier}
\label{sec:v5-context}

\subsection{Why this continuation goes upstream}
V4 fixes the actual complex theta-source orbit and its real involution. It
would be circular to call one of its finite frame permutations an arithmetic
transfer merely because both are represented by four-dimensional matrices.
We therefore go back to the independently specified rational curve
\[
 C:\quad y^2=f(x),\qquad f(x)=x^5-x^2-x=xg(x),\qquad g(x)=x^4-x-1,
\]
and acquire its actual finite-prime Frobenius. This produces two constructive
outputs and a separation theorem. First, exact point counts determine the
local weight-one Frobenius polynomials. Second, the native pullback on theta
sections is a level-raising map, with a canonical splitting at one of the
examined primes but not the other. The local weight-one operator cannot be
represented by any fixed level-two finite theta-frame normalizer. A positive
tracial local response can nevertheless be reconstructed from its spectrum,
with its remaining source-coupling ambiguity explicitly classified.

The inherited results are the normalized chronology loading, the pointed
arithmetic two-torsion packet, the integral $K_4$ bridge, the finite
Heisenberg operator comparison, V1--V4 source recognition and phase
selection, and the generic word/Krylov and interface recognition calculus
\cite{GT77,Arithmetic,AD9}. None is reproved. In particular, recovering
characteristic polynomials from complete power traces is standard Newton
calculus and was already part of the earlier Frobenius programme; it is not
claimed as a new general theorem here. The new computations concern this
actual Jacobian, the mismatch with its finite theta-frame normalizer, the
native theta-level maps, the local source fibre, and a concrete central
arithmetic datum forgotten by the projective theta source.

\subsection{Three distinct Frobenius realizations}
For an odd prime $p\ne283$, the square-free reduction of $f$ gives a smooth
projective genus-two curve $C_p$ and its principally polarized Jacobian
$J_p=\operatorname{Jac}(C_p)$. Let $\pi_p:J_p\to J_p$ be the $p$-power
Frobenius endomorphism over $\F_p$. We distinguish
\[
 \begin{array}{lll}
 \text{two-torsion action}&\pi_p|_{J_p[2]}&\text{a finite label action},\\
 \text{weight-one action}&\pi_p|_{V_\ell(J_p)},\quad \ell\ne p
     &\text{a rank-four characteristic-zero action},\\
 \text{theta pullback}&\pi_p^*:H^0(L_p^m)\to H^0(\pi_p^*L_p^m)
     &\text{a characteristic-$p$ section map}.
 \end{array}
\]
Here $V_\ell(J_p)=T_\ell(J_p)\otimes_{\Z_\ell}\Q_\ell$ and $L_p$ is the
principal theta line selected by the reduction of $\kappa_\infty$.
The convention for the weight-one action is the endomorphism acting on the
Tate module; we do not substitute the inverse Galois convention without
changing formulas. A complex matrix introduced below is a spectral model
of its algebraic eigenvalues, not an identification of the Tate module with
V4's Hermitian theta space.

\section{Exact local Frobenius data of the actual curve}
\label{sec:v5-counts}

\subsection{Classical input and a finite count certificate}
The standard Weil and Jacobian trace formulas state that the characteristic
polynomial of $\pi_p$ is integral of degree four, its eigenvalues
$\alpha_1,\ldots,\alpha_4$ have complex absolute value $\sqrt p$, the
polarization pairs them as $\alpha$ and $p/\alpha$, and
\[
 \#C_p(\F_{p^r})=p^r+1-\sum_{i=1}^4\alpha_i^r.
\]
These arithmetic-geometric facts are imported from the classical theory;
see \cite{MilneAV}, Chapters II, Section~1 and III, Section~11. The following
coefficient extraction and finite counts are explicit applications.

\begin{proposition}[Two exact counts determine the local polynomial]
\label{prop:v5-polynomial-counts}
Put $N_r=\#C_p(\F_{p^r})$. Then
\begin{equation}
 \boxed{
 P_p(T)=T^4-a_pT^3+b_pT^2-pa_pT+p^2,\quad
 a_p=p+1-N_1,\quad b_p=\frac{a_p^2+N_2-p^2-1}{2}.}
 \label{eq:v5-count-polynomial}
\end{equation}
For $p=3$ and $p=5$, respectively,
\begin{align}
 &(N_1,N_2)=(2,14),&
 P_3(T)&=T^4-2T^3+4T^2-6T+9,\label{eq:v5-P3}\\
 &(N_1,N_2)=(10,30),&
 P_5(T)&=T^4+4T^3+10T^2+20T+25.\label{eq:v5-P5}
\end{align}
\end{proposition}
\begin{proof}
The trace formula gives $\sum\alpha_i=a_p$ and
$\sum\alpha_i^2=p^2+1-N_2$. Thus
$2\sum_{i<j}\alpha_i\alpha_j=a_p^2-p^2-1+N_2$.
Polarization reciprocity gives the remaining coefficients, proving
\eqref{eq:v5-count-polynomial}.

For a finite field $k$ of odd cardinality $q$, let $\chi_k$ be its quadratic
character, including $\chi_k(0)=0$. The odd-degree model has one point at
infinity, so direct counting gives
\begin{equation}
 N(k)=q+1+\sum_{x\in k}\chi_k(f(x)).
 \label{eq:v5-character-count}
\end{equation}
For both primes choose $k=\F_p[w]/(w^2-2)$. Two is a nonsquare in each
case. Since
\[
 \chi_{p^2}(u)=\chi_p(\operatorname{Norm}_{\F_{p^2}/\F_p}u),
 \qquad \operatorname{Norm}(a+bw)=a^2-2b^2,
\]
all character evaluations use arithmetic in $\F_p$. The character identity
follows for $u\ne0$ by raising $u^{p+1}$ to $(p-1)/2$, and also holds at
zero. In row $a$, column $b$, the exact arrays
$\chi_{p^2}(f(a+bw))$ are
\begin{equation}
 \begin{pmatrix}0&1&1\\1&-1&-1\\1&1&1\end{pmatrix}
 \quad(p=3),\qquad
 \begin{pmatrix}
 0&1&-1&-1&1\\
 1&-1&-1&-1&-1\\
 1&1&1&1&1\\
 1&-1&-1&-1&-1\\
 1&1&1&1&1
 \end{pmatrix}
 \quad(p=5).
 \label{eq:v5-character-arrays}
\end{equation}
Their sums are $4$ in both cases. The prime-field sums are $-2$ and $4$.
Equation~\eqref{eq:v5-character-count} proves the counts and hence the two
polynomials. The accompanying script independently enumerates the fibres
of the squaring map, using a separate Horner evaluation of $f$, and obtains
identical counts. This second audit is not needed to replace the displayed
finite calculation.
\end{proof}

\begin{remark}[Additional finite checks]
The same exact procedure gives the following audit table. Only $p=3,5$ are
needed for the subsequent theorems.
\begin{center}
\begin{tabular}{rrrrr}
\toprule $p$&$N_1$&$N_2$&$a_p$&$b_p$\\\midrule
3&2&14&2&4\\5&10&30&-4&10\\7&5&61&3&10\\
11&15&157&-3&22\\13&19&169&-5&12\\17&20&306&-2&10\\\bottomrule
\end{tabular}
\end{center}
No point count at a bad prime or characteristic two is asserted.
\end{remark}

\subsection{The torsion permutation is only a reduction}
\begin{proposition}[The same residual four-cycle, different weight-one data]
\label{prop:v5-fourcycles}
Both $g\bmod3$ and $g\bmod5$ are irreducible quartics. Thus both local
Frobenius elements permute the four residual singular two-torsion labels
as four-cycles. Their weight-one polynomials are nevertheless
\eqref{eq:v5-P3} and \eqref{eq:v5-P5}.
\end{proposition}
\begin{proof}
The lists of $g(a)$ modulo $p$, for $a=0,\ldots,p-1$, are
$(2,2,1)$ and $(4,4,3,2,1)$, so there is no linear factor. Modulo $g$ the
remainders of $x^{p^2}-x$ and their multiplicative inverses are
\[
 \begin{array}{c|cc}
 p&(x^{p^2}-x)\bmod g&\text{inverse modulo }g\\\hline
 3&x^3-x^2&x^3+x\\
 5&2x^3+2x^2+x+1&x^2+2x.
 \end{array}
\]
Multiplication verifies that each product is one modulo $g$. A reducible
quartic without linear factors is a product of two irreducible quadratics;
such a factor would divide $x^{p^2}-x$. The displayed inverse excludes that
case. The factorization--Frobenius correspondence therefore gives a
four-cycle.

For completeness, the inherited five singular classes satisfy
$e_2(v_i,v_j)=-1$ for $i\ne j$ and $t=\sum_{i=1}^4v_i$. The $4\times4$
polarization matrix with zero diagonal and ones off diagonal is invertible
over $\F_2$ (its square is the identity), so the $v_i$ form a basis.
Frobenius is the ordinary permutation matrix on this basis. Its
characteristic polynomial modulo two is $T^4+1$, in agreement with both
$P_3\bmod2$ and $P_5\bmod2$. This congruence checks the connection between
the torsion and Tate-module representations, not an equality of the
characteristic-zero operators.
\end{proof}

\section{A quantitative obstruction to a fixed theta-frame Frobenius}
\label{sec:v5-normalizer}

\subsection{A scalar-invariant finite-projective-order test}
\begin{lemma}[Trace obstruction to finite projective order]
\label{lem:v5-projective-integrality}
For an invertible complex matrix $F$, define
\[
 \mathfrak i(F)=\Tr(F)\Tr(F^{-1}).
\]
It is invariant under similarity and nonzero scalar multiplication. If
$F^m=cI$ for some $m\ge1$ and $c\ne0$, then $\mathfrak i(F)$ is an
algebraic integer.
\end{lemma}
\begin{proof}
The invariances follow directly from the traces. The polynomial $T^m-c$
has distinct roots in characteristic zero, so $F$ is diagonalizable.
Writing its eigenvalues as $\alpha_i$, every $\alpha_i/\alpha_j$ is an
$m$th root of unity. Therefore
$\mathfrak i(F)=\sum_{i,j}\alpha_i/\alpha_j$ is a sum of algebraic
integers.
\end{proof}

\begin{theorem}[The actual local Frobenius has infinite projective order]
\label{thm:v5-infinite-projective}
For the weight-one Frobenius of $J_p$,
\begin{equation}
 \boxed{\mathfrak i(\pi_p)=a_p^2/p.}
 \label{eq:v5-projective-trace}
\end{equation}
In particular $\pi_3$ and $\pi_5$ have infinite projective order. No
nonzero scalar multiple of either can be similar to a matrix of finite
projective order.
\end{theorem}
\begin{proof}
The reciprocal eigenvalue pairing gives
$\Tr(\pi_p^{-1})=a_p/p$. At the two examined primes the product is $4/3$
and $16/5$. A rational algebraic integer is an integer, so
\cref{lem:v5-projective-integrality} excludes finite projective order.
The proof applies to any characteristic-zero realization of the algebraic
Frobenius polynomial; it makes no choice of a theta-space metric.
\end{proof}

\subsection{The finite Heisenberg normalizer}
Let $\{P_v:v\in\F_2^4\}$ be the sixteen Hermitian Pauli representatives in
V1's four-dimensional Heisenberg representation, with $P_0=I$. The
projective theta-frame normalizer consists of unitaries $U$ for which
\begin{equation}
 UP_vU^*=\epsilon_U(v)P_{g_Uv},\qquad \epsilon_U(v)\in\{1,-1\}.
 \label{eq:v5-Pauli-normalizer}
\end{equation}
The label map $g_U$ is symplectic because conjugation preserves products
and commutators. This is a normalizer of a specified finite theta frame,
not the full unitary group of a four-dimensional space.

\begin{lemma}[Finite normalizer and an exact trace character]
\label{lem:v5-normalizer-trace}
The normalizer modulo scalar unitaries is finite. If
$K_U=\ker(g_U-I)$, then
\begin{equation}
 \boxed{|\Tr U|^2\in\{0,2^{\dim_{\F_2}K_U}\}.}
 \label{eq:v5-trace-character}
\end{equation}
When $g_U$ is a four-cycle on the residual $v_1,\ldots,v_4$, the two
possible values are $0$ and $2$.
\end{lemma}
\begin{proof}
By \cref{lem:v4-marking}, the kernel of the action on the sixteen labels
consists projectively of the sixteen Pauli displacements. The image is a
subgroup of the finite group $\operatorname{Sp}_4(\F_2)$; hence the
projective normalizer is finite. More explicitly, if $g_U$ has order $m$,
then $U^m$ is a scalar Pauli displacement and $U^{2m}$ is scalar.

The sixteen operators are an orthogonal basis of $\End(\C^4)$. On this
basis $\operatorname{Ad}_U$ is a signed permutation, giving
\[
 \Tr(\operatorname{Ad}_U)=\sum_{v\in K_U}\epsilon_U(v).
\]
For $v,w\in K_U$, conjugating the identity
$P_vP_w=c(v,w)P_{v+w}$ shows that
$\epsilon_U(v+w)=\epsilon_U(v)\epsilon_U(w)$: the fixed multiplication
coefficient $c(v,w)$ cancels. Thus the signs on $K_U$ form a character.
Its sum is zero unless it is trivial, when its sum is $|K_U|$.
Finally $\operatorname{Ad}_U\cong U\otimes\overline U$, so its trace is
$|\Tr U|^2$. For a four-cycle in the residual basis the fixed subspace is
the one-dimensional constant line, proving the last assertion.
\end{proof}

\begin{theorem}[No fixed-frame lift of the actual weight-one Frobenius]
\label{thm:v5-no-fixed-theta}
Neither $\pi_3$ nor $\pi_5$, even after arbitrary scalar normalization and
change of characteristic-zero basis, is an element of the fixed
level-two theta-frame normalizer.
For a unitary spectral representative $U_3$ of $\pi_3/\sqrt3$, every
normalizer unitary $V$ inducing a residual four-cycle satisfies
\begin{equation}
 \boxed{
 \inf_{R\in U(4),\ |z|=1}\norm{RU_3R^*-zV}_{\HS}
 \ge \frac1{\sqrt2}-\frac1{\sqrt3}>0.1297.}
 \label{eq:v5-normalizer-distance}
\end{equation}
\end{theorem}
\begin{proof}
The first conclusion follows from finite projective order in
\cref{lem:v5-normalizer-trace} and the obstruction in
\cref{thm:v5-infinite-projective}. For the quantitative conclusion,
$|\Tr U_3|=2/\sqrt3$, while $|\Tr V|$ is zero or $\sqrt2$.
For $4\times4$ matrices,
$|\Tr(A-B)|\le2\norm{A-B}_{\HS}$. The reverse triangle inequality for
the two traces, uniform over both $R$ and $z$, gives
\[
 2\norm{RU_3R^*-zV}_{\HS}
 \ge\bigl||\Tr U_3|-|\Tr V|\bigr|
 \ge\sqrt2-2/\sqrt3.
\]
This proves the stated bound. It is a certified lower bound, not a claim of
optimal distance.
\end{proof}

\begin{scope}[What this obstruction does not exclude]
The result excludes a fixed complex-linear finite-theta-frame model of the
weight-one Frobenius. It does not exclude semilinear descent between
conjugate fibres, an operator outside the finite frame normalizer, a
higher-level theta construction, or an independently established common
arithmetic--gauge carrier. In particular, the finite $S_4$ torsion action
and the $\ell$-adic Tate-module action remain legitimate, different
outputs of the same arithmetic object.
\end{scope}

\section{The native Frobenius raises theta level}
\label{sec:v5-level}

\subsection{The actual section map and its source}
For $m\ge1$ set $H_{m,p}=H^0(J_p,L_p^{\otimes m})$, as an $\F_p$ vector
space. All statements in this section are in characteristic $p$. The
principal theta line is ample with Euler characteristic one; the
Riemann--Roch and ample-line cohomology theorems for abelian varieties give
\begin{equation}
 \dim H_{m,p}=m^2.
 \label{eq:v5-theta-dimension}
\end{equation}
These standard inputs hold also in positive characteristic
\cite{MilneAV}, Chapter I, Section~11.

\begin{theorem}[Source-preserving level-raising Frobenius]
\label{thm:v5-level-raising}
There is a canonical line-bundle identification
$\pi_p^*L_p\cong L_p^{\otimes p}$ and an injective $\F_p$-linear map
\begin{equation}
 \boxed{j_{m,p}:H_{m,p}\longrightarrow H_{mp,p},\qquad s\longmapsto s^p.}
 \label{eq:v5-level-map}
\end{equation}
These maps preserve products, and their iterates give an exact directed
system. The squared theta source is carried as
\begin{equation}
 \boxed{\vartheta_\kappa^2\longmapsto\vartheta_\kappa^{2p}.}
 \label{eq:v5-theta-source-pullback}
\end{equation}
Thus the native level-two pullbacks have dimensions $4\to36$ at $p=3$ and
$4\to100$ at $p=5$, rather than $4\to4$.
\end{theorem}
\begin{proof}
In local trivializations, the transition functions of $\pi_p^*L_p$ are
the $p$th powers of those of $L_p$. These are exactly the transition
functions of $L_p^p$. A section is sent to its $p$th power under this
identification. Additivity follows from characteristic $p$, and scalar
linearity holds over $\F_p$. The variety is reduced, so $s^p=0$ implies
$s=0$, proving injectivity. The identity $(st)^p=s^pt^p$ proves product
compatibility, and iteration gives $s\mapsto s^{p^r}$ independently of
parenthesization. Formula~\eqref{eq:v5-theta-source-pullback} is the case
$s=\vartheta_\kappa^2$; its source line is nonzero. Dimensions follow from
\eqref{eq:v5-theta-dimension}.
\end{proof}

\begin{remark}[Scalar extension and marked displacement covariance]
After extending $\F_p$, one uses the base-changed relative Frobenius map;
the naive $p$th-power operation on all new scalars is semilinear, not
complex linear. For a geometric point $x$ the group identity
$\pi_p\circ t_x=t_{\pi_p(x)}\circ\pi_p$ supplies the commutative
pullback square for translated line bundles. When actual theta-group
lifts are retained, their natural pullbacks obey this square. This is a
cross-level covariance statement and does not identify the codomain with
$H_{m,p}$. Indeed $L_p^{mp}$ and $L_p^m$ have different nonzero numerical
classes and different self-intersection numbers for $p>1$; ampleness
prevents a line-bundle identification that would erase the level change.
\end{remark}

\subsection{The actual ordinary and nonordinary primes}
We next determine whether the native injection has an intrinsic geometric
retraction. This is more informative than choosing an arbitrary
four-dimensional compression of its codomain.

\begin{proposition}[Exact Cartier--Manin matrices]
\label{prop:v5-Cartier}
In the ordered differential basis $(dx/y,x\,dx/y)$, use the Cartier--Manin
convention in which the $ij$ entry is $c_{pi-j}^{1/p}$, where
$f(x)^{(p-1)/2}=\sum c_kx^k$. At the two primes the matrices are
\begin{equation}
 \boxed{
 M_3=\begin{pmatrix}2&2\\1&0\end{pmatrix},\qquad
 M_5=\begin{pmatrix}1&2\\0&0\end{pmatrix}.}
 \label{eq:v5-Cartier-matrices}
\end{equation}
The corresponding Hasse--Witt matrices on $H^1(C_p,\mathcal O)$ are their
transposes over $\F_p$. Hence $J_3$ is ordinary and $J_5$ is not ordinary.
\end{proposition}
\begin{proof}
We recall the coefficient computation to fix the otherwise ambiguous
transpose convention; see \cite{AchterHowe}, Sections~2.4 and 3.1.
For $\omega_j=x^{j-1}dx/y$,
\[
 \omega_j=y^{-p}\sum_kc_kx^{k+j-1}\,dx.
\]
The Cartier operator annihilates monomials unless $k+j-1\equiv p-1$
modulo $p$, and sends
$h^px^{p-1}dx$ to $h\,dx$. The term contributing to $\omega_i$ therefore
has $k=pi-j$ and coefficient $c_{pi-j}^{1/p}$. Over $\F_p$ that root is
the coefficient itself. At $p=3$ the required power is $f$. At $p=5$ it is
\[
 f^2=x^{10}-2x^7-2x^6+x^4+2x^3+x^2.
\]
Reading the four entries gives \eqref{eq:v5-Cartier-matrices}. Cartier and
Frobenius on $H^1(\mathcal O)$ are Serre-dual, giving the indicated
transpose. The Abel--Jacobi identification of $H^1(\mathcal O)$ commutes
with Frobenius. The determinants are $1$ modulo three and zero modulo
five. Invertibility of this Frobenius is the standard ordinary criterion
for an abelian variety, proving the assertion. We need only these
invertibility conclusions, not a classification of all Newton slopes.
\end{proof}

\begin{lemma}[Classical unique Frobenius splitting criterion]
\label{lem:v5-splitting-input}
For a smooth proper geometrically connected variety $X/\F_p$ with trivial
canonical bundle, the map $\mathcal O_X\to F_*\mathcal O_X$ admits an
$\mathcal O_X$-linear splitting if and only if Frobenius on
$H^{\dim X}(X,\mathcal O_X)$ is bijective. When a splitting exists it is
unique after imposing $\sigma(1)=1$.
\end{lemma}
\begin{proof}
This is the standard Frobenius-splitting criterion, stated with uniqueness
in \cite{AchingerZdanowicz}, Proposition~2.3.1. Its short duality argument
is useful here. Finite duality for Frobenius identifies the space of
$\mathcal O_X$-linear maps $F_*\mathcal O_X\to\mathcal O_X$ with
$H^0(X,\omega_X^{1-p})$, which is one-dimensional. Evaluation at one is
nonzero exactly when the dual top-cohomology Frobenius is nonzero, by
Serre duality. The latter cohomology is one-dimensional, so nonzero means
bijective. In that case there is exactly one map with evaluation one,
and its composite with the unit $\mathcal O_X\to F_*\mathcal O_X$ is
the identity. If evaluation vanishes, no splitting is possible.
\end{proof}

\begin{theorem}[A canonical split theta tower at three, but not at five]
\label{thm:v5-split-tower}
There is a unique normalized Frobenius splitting
$\sigma_3:(\pi_3)_*\mathcal O_{J_3}\to\mathcal O_{J_3}$. For every
$m\ge1$ it induces a canonical map
\begin{equation}
 \boxed{r_{m,3}:H_{3m,3}\longrightarrow H_{m,3},\qquad
 r_{m,3}j_{m,3}=I.}
 \label{eq:v5-retraction}
\end{equation}
Consequently $E_{m,3}=j_{m,3}r_{m,3}$ is an idempotent of rank $m^2$ and
\begin{equation}
 \boxed{H_{3m,3}=j_{m,3}H_{m,3}\oplus\ker r_{m,3},\qquad
 \dim\ker r_{m,3}=8m^2.}
 \label{eq:v5-splitting-dimension}
\end{equation}
In particular the level-six space is $36$ dimensional with a canonical
four-dimensional split image and a $32$-dimensional kernel. At $p=5$ no
$\mathcal O_{J_5}$-linear Frobenius splitting exists.
\end{theorem}
\begin{proof}
For an abelian surface,
$H^2(J_p,\mathcal O)=\Lambda^2H^1(J_p,\mathcal O)$, compatibly with
Frobenius. The top operation is thus the determinant of the Hasse--Witt
matrix. By \cref{prop:v5-Cartier} it is nonzero at three and zero at five.
Apply \cref{lem:v5-splitting-input}.

The projection formula identifies
\[
 (\pi_3)_*(\pi_3^*L_3^m)
 \cong L_3^m\otimes(\pi_3)_*\mathcal O_{J_3}.
\]
Tensor $\sigma_3$ with $L_3^m$ and take global sections. This gives
$r_{m,3}$; tensoring the split unit identity gives
$r_{m,3}j_{m,3}=I$. The direct sum and dimensions follow by elementary
linear algebra and \eqref{eq:v5-theta-dimension}. At five the vanishing
top Frobenius excludes the sheaf splitting by the same criterion.
\end{proof}

\begin{corollary}[Exact iterated compatibility and the identity roundtrip]
\label{cor:v5-iterate-split}
The maps in \cref{thm:v5-split-tower} iterate to compatible injections and
retractions between $H_{m,3}$ and $H_{3^rm,3}$ for every $r\ge1$.
Their lower-level roundtrip is always the identity. It is therefore not
the weight-one Frobenius of \eqref{eq:v5-P3}.
\end{corollary}
\begin{proof}
Compose the unit maps and their splittings. The composite for $F^r$ is
again normalized and $\mathcal O$-linear. The same duality argument uses
$\omega^{1-3^r}$ and shows uniqueness, so staged and direct splittings
agree. Projection formula transports this identity to every $L_3^m$.
Telescoping $r_{m,3}j_{m,3}=I$ proves the roundtrip statement. Its
projective order is one, whereas \cref{thm:v5-infinite-projective}
shows that the weight-one Frobenius has infinite projective order.
\end{proof}

\begin{scope}[A geometric retraction is not yet a physical transfer]
The construction proves canonical maps of section spaces; no numerical
$36\times4$ theta matrix is acquired here. The idempotent is not called
an orthogonal or completely positive projection: no positive complex
metric has been transported to these characteristic-$3$ spaces. At five,
abstract vector-space left inverses still exist for each injective map;
what fails is the geometrically natural sheaf splitting specified above.
The Cartier/duality trace used by the splitting must not be replaced by
the field-extension trace of the purely inseparable Frobenius extension,
which is zero. None of these maps identifies the level-two theta source
with a Tate-module source or a renewal boundary source.
\end{scope}

\section{An explicit positive local arithmetic response and its source fibre}
\label{sec:v5-response}

\subsection{The local spectral class is constructive}
At $p=3$, factor the actual polynomial over $\Q(\sqrt3)$ as
\begin{equation}
 P_3(T)=
 (T^2-(1+\sqrt3)T+3)(T^2-(1-\sqrt3)T+3).
 \label{eq:v5-P3-factor}
\end{equation}
The two real middle coefficients have absolute value less than $2\sqrt3$.
The four normalized roots are consequently distinct and are
$e^{i\theta_+},e^{-i\theta_+},e^{i\theta_-},e^{-i\theta_-}$, with
\begin{equation}
 \cos\theta_\pm=\frac{1\pm\sqrt3}{2\sqrt3},\qquad 0<\theta_\pm<\pi.
 \label{eq:v5-angles}
\end{equation}
Thus they specify one conjugacy class in the compact symplectic group
$\operatorname{USp}(4)$ and one ordinary unitary conjugacy class in $U(4)$.
This is the spectral type of this local arithmetic operator, not a proposed
Standard-Model gauge group. Also
$P_5(T)=(T^2+5)(T^2+4T+5)$, giving another explicit local class.

Let $\tau=2/\sqrt3$, $\beta=4/3$, and set
\begin{equation}
 \ell_3(z)=z^4-\tau z^3+\beta z^2-\tau z+1,
 \qquad L_3(z)=1-\tau z+\beta z^2-\tau z^3+z^4.
 \label{eq:v5-normalized-polynomial}
\end{equation}
The two expressions are the same reciprocal polynomial; different
letters only distinguish the characteristic and determinant uses.
Write $\lambda_i$ for its four unit-circle roots. On
$\mathscr H_3=\C[z]/(\ell_3)$ define the normalized tracial spectral form
\begin{equation}
 \ip{q}{r}_{\rm tr}=\frac14\sum_{i=1}^4
       \overline{q(\lambda_i)}r(\lambda_i).
 \label{eq:v5-trace-form}
\end{equation}
The choice is intrinsic to the finite spectral algebra: it gives equal
weight to each simple spectral point. It is not asserted to be the
relative source state supplied by $\vartheta_\kappa^2$.

\begin{theorem}[Exact arithmetic companion and positive source Gram]
\label{thm:v5-tracial-realization}
In the basis $(1,z,z^2,z^3)$, multiplication by $z$ has matrix
\begin{equation}
 C_3=\begin{pmatrix}
 0&0&0&-1\\1&0&0&\tau\\0&1&0&-\beta\\0&0&1&\tau
 \end{pmatrix},
 \quad
 G_3=\begin{pmatrix}
 1&\tau/4&-1/3&\tau/12\\
 \tau/4&1&\tau/4&-1/3\\
 -1/3&\tau/4&1&\tau/4\\
 \tau/12&-1/3&\tau/4&1
 \end{pmatrix}
 \label{eq:v5-companion-Gram}
\end{equation}
is its source/Krylov Gram, and
\begin{equation}
 \boxed{C_3^*G_3C_3=G_3,\qquad
 \det G_3=\frac{13}{36},\qquad
 G_3\succeq\frac{39}{256}I.}
 \label{eq:v5-positive-Gram}
\end{equation}
The source $e_0=(1,0,0,0)^{\mathsf T}$ is cyclic. This constructs one
source-minimal positive unitary realization of the actual local
Frobenius spectrum without using a renewal comparison target.
\end{theorem}
\begin{proof}
Multiplication modulo $\ell_3$ gives the last column of $C_3$ and its
three shift columns. Distinct roots imply that evaluation at the four
$\lambda_i$ is an isomorphism; hence \eqref{eq:v5-trace-form} is positive
definite. Multiplication by $z$ is unitary in this form because
$|\lambda_i|=1$.

The normalized power traces for $j=0,1,2,3$ are
\[
 \frac14\sum_i\lambda_i^j=1,\quad \tau/4,\quad-1/3,\quad\tau/12.
\]
They follow by the first three Newton identities for \eqref{eq:v5-P3};
complex-conjugate pairing gives the negative moments. These entries give
$G_3$. Direct determinant evaluation, or the Vandermonde calculation in
\cref{thm:v5-weight-fibre} below, yields $13/36$. Its leading principal
minors are $1,11/12,2/3,13/36$, providing a second exact positivity check.
For the least eigenvalue $\gamma$, the other three positive eigenvalues
have sum at most $\Tr G_3=4$ and hence product at most $(4/3)^3$.
Therefore
$\gamma\ge(13/36)(3/4)^3=39/256$. Finally
$e_0,C_3e_0,C_3^2e_0,C_3^3e_0$ are the four displayed basis vectors, so
the source is cyclic.
\end{proof}

\begin{proposition}[A native tracial Carath\'eodory response]
\label{prop:v5-response}
For $|z|<1$ define
\begin{equation}
 \boxed{
 H_3(z)=1-\frac z2\frac{L_3'(z)}{L_3(z)}
 =\frac{(1-z^2)(1-z/\sqrt3+z^2)}{L_3(z)}.}
 \label{eq:v5-response}
\end{equation}
Then
\begin{align}
 H_3(z)&=\frac14\sum_{i=1}^4\frac{1+z\lambda_i}{1-z\lambda_i}
 =2\ip{e_0}{(I-zC_3)^{-1}e_0}_{G_3}-1,\label{eq:v5-response-realization}\\
 \operatorname{Re}H_3(z)&=\frac14\sum_{i=1}^4
          \frac{1-|z|^2}{|1-z\lambda_i|^2}>0.\label{eq:v5-response-positive}
\end{align}
In particular $(H_3-1)/(H_3+1)$ is a holomorphic scalar Schur function
on the unit disk.
\end{proposition}
\begin{proof}
The determinant is $L_3(z)=\prod_i(1-z\lambda_i)$, because the
normalized polynomial is reciprocal. Logarithmic differentiation gives
\eqref{eq:v5-response-realization}; expansion of the numerator gives
\eqref{eq:v5-response}. The resolvent formula is the spectral form
\eqref{eq:v5-trace-form}. Multiplying numerator and denominator of each
summand by $1-\overline z\overline\lambda_i$ proves
\eqref{eq:v5-response-positive}. The Cayley transform of a number in
the open right half-plane has modulus less than one, proving the last
assertion. No analytic continuation beyond this local disk or global
number-field zero statement is used.
\end{proof}

\subsection{What a Frobenius polynomial does not determine about a source}
\begin{theorem}[Exact source-weight fibre and its conditioning]
\label{thm:v5-weight-fibre}
Let $U$ have the four distinct normalized eigenvalues of $P_3$, and let
$\eta$ be a unit source. Up to simultaneous ordinary unitary equivalence
of the pair $(U,\eta)$, the possibilities are exactly the weights
\begin{equation}
 \boxed{w_i=\norm{\Pi_i\eta}^2\ge0,\qquad \sum_{i=1}^4w_i=1,}
 \label{eq:v5-source-simplex}
\end{equation}
where the spectral labels are retained. The source is cyclic precisely
when all weights are positive. For its four-column Krylov Gram,
\begin{equation}
 K_{jk}(w)=\sum_iw_i\overline{\lambda_i}^{\,j}\lambda_i^k,
 \qquad 0\le j,k\le3,
\end{equation}
one has
\begin{equation}
 \boxed{\det K(w)=\frac{832}{9}\prod_iw_i.}
 \label{eq:v5-weight-determinant}
\end{equation}
The tracial weights $w_i=1/4$ uniquely maximize this determinant and give
\cref{thm:v5-tracial-realization}. If every $w_i\ge\delta>0$, then
\begin{equation}
 \lambda_{\min}K(w)\ge39\delta^4.
 \label{eq:v5-source-weight-floor}
\end{equation}
Thus even an exact local polynomial leaves a three-real-dimensional
source-moment fibre before any additional theta marking is compared.
\end{theorem}
\begin{proof}
Diagonalize $U$. Its simple spectral spaces are one-dimensional, so a
unitary commuting with $U$ can adjust each nonzero component phase of
$\eta$ independently. Two pairs have the same weights exactly when such
a unitary maps one source to the other. Every weight vector is attained
by a source with components $\sqrt{w_i}$. Vandermonde invertibility shows
that the source is cyclic exactly when all four components are nonzero.

Let $V_{ij}=\lambda_i^j$, $0\le j\le3$. Then
$K=V^*\operatorname{diag}(w_i)V$, so
$\det K=|\det V|^2\prod_iw_i$. The discriminant of $P_3$ is
$67392$. One exact way to evaluate it is to use
\eqref{eq:v5-P3-factor}: the two quadratic discriminants have product
$52$, and their resultant is $36$, giving $52\cdot36^2=67392$.
Normalizing each root by $\sqrt3$ divides the squared Vandermonde by
$3^6$, hence $|\det V|^2=67392/3^6=832/9$.
Arithmetic--geometric mean on the weights proves the unique maximum.
Finally $\Tr K=4$, so the product of three other eigenvalues is at most
$(4/3)^3$. Combining this with the determinant gives
$\lambda_{\min}K\ge(832/9)(3/4)^3\delta^4=39\delta^4$.
\end{proof}

\begin{remark}[The relative acquisition still required]
The classification concerns the ordinary unitary pair $(U,\eta)$, not a
larger packet retaining Clifford rays, tensor intertwiners, or a prescribed
real structure. Such extra markings can leave additional relative data.
Even after the four V4 lift signs are settled, identifying V4's theta
source with a weight-one carrier needs an independent comparison map and
its source weights. The exact tracial form is an arithmetic benchmark,
not evidence that the actual theta source has equal spectral weights.
The lower bound in \eqref{eq:v5-positive-Gram} is an observability bound
for this finite cyclic source. It is not a dissipative gap, an arithmetic
cancellation estimate, or a regulator-uniform physical coercivity bound.
\end{remark}

\section{A central arithmetic datum invisible to the projective theta source}
\label{sec:v5-twist}

\subsection{The same complex source, different signed Frobenius}
The previous sections separate realization types for one fixed rational
model. A quadratic-twist test now determines what is lost when the
arithmetic input is compressed to the projective complex theta source and
its real involution. This is an obstruction for that forgetful interface,
not a failure of the full rational curve to determine its arithmetic.

\begin{theorem}[Positive quadratic twists survive all complex source tests]
\label{thm:v5-twist}
For a positive nonsquare rational $d$, let
\[
 C^{(d)}:\ y^2=d f(x),\qquad J^{(d)}=\operatorname{Jac}(C^{(d)}).
\]
The labelled complex theta-source packet of V4 and its real involution are
isomorphic to those of $C$. Its projective level-two theta descent also
forgets the quadratic twisting cocycle. In contrast the Tate modules obey
\begin{equation}
 \boxed{V_\ell(J^{(d)})\cong V_\ell(J)\otimes\chi_d.}
 \label{eq:v5-twist-representation}
\end{equation}
For every odd good prime $p$ away from $d$, in the same Frobenius convention,
\begin{equation}
 \boxed{P_p^{(d)}(T)=P_p(\chi_d(p)T).}
 \label{eq:v5-twist-polynomial}
\end{equation}
For $d=2$ and $p=3$ this is the explicit separation
\begin{equation}
 \boxed{
 P_3^{(2)}(T)=T^4+2T^3+4T^2+6T+9,\qquad
 (N_1^{(2)},N_2^{(2)})=(6,14).}
 \label{eq:v5-twist-example}
\end{equation}
\end{theorem}
\begin{proof}
The map $(x,y)\mapsto(x,y/\sqrt d)$ is an isomorphism over
$\Q(\sqrt d)$ and over $\R$ because $d>0$. It preserves $\infty$, zero,
all four labelled branch points, the principal polarization, the odd
characteristic, and its squared theta-section line. It therefore induces
the labelled Heisenberg representation isomorphism. The invariant positive
Hermitian form is unique up to scalar by irreducibility, so normalized
source Grams and all phase invariants agree. The isomorphism is real;
hence it also intertwines the real descent of V4.

For $\sigma\in G_\Q$, the ratio
$\sigma(\sqrt d)/\sqrt d=\chi_d(\sigma)$ shows that the two choices of
this isomorphism differ by the hyperelliptic involution precisely when
$\chi_d(\sigma)=-1$. On the Jacobian that involution is $[-1]$:
indeed, $P+\iota(P)\sim2\infty$ implies
$[P-\infty]+[\iota(P)-\infty]=0$, and these classes generate the
Jacobian. Thus the cocycle acts by $-I$ on $T_\ell J$, proving
\eqref{eq:v5-twist-representation}.

On $J[2]$, $[-1]$ is the identity. It is also projectively the identity
on the level-two theta sections. In the second-order theta basis used
in V2, each function is even under $z\mapsto-z$ (replace its integer
summation variable by its negative minus the binary characteristic).
Thus the symmetric lift of inversion is scalar on this four-dimensional
space; changing its linearization changes only that scalar. Consequently
the projective level-two descent cannot see this quadratic cocycle.
The Tate-module action, however, is multiplied by $\chi_d(\sigma)$.
For a good prime the character value is the Legendre symbol $(d/p)$;
taking characteristic polynomials in dimension four gives
\eqref{eq:v5-twist-polynomial}. At three, $(2/3)=-1$, and
\eqref{eq:v5-P3} gives the displayed polynomial. Its first two power
traces, or direct application of \eqref{eq:v5-character-count} to $2f$,
give $(6,14)$.
\end{proof}

\begin{corollary}[A necessary central datum for signed local duality]
\label{cor:v5-central-datum}
No rule depending only on the labelled projective complex level-two
source, its full V4 invariant data, and its real involution can recover
the signed first weight-one Frobenius trace for every rational model in
the positive quadratic-twist family. The missing distinction is a
rank-one quadratic descent character, not an additional choice among
V4's four Heisenberg frame signs.
\end{corollary}
\begin{proof}
The inputs agree for $d=1$ and $d=2$ by
\cref{thm:v5-twist}, whereas the first trace at three is respectively
$2$ and $-2$. A function of the common inputs cannot return both values.
The sign action on the Tate module is central and projectively invisible;
V4's frame reorientations act within one fixed complex source orbit and
do not restore this rational central cocycle.
\end{proof}

\begin{scope}[What is and is not forgotten]
The curve $y^2=f(x)$, its Jacobian over $\Q$, and its complete rational
descent are fixed throughout the principal computation. They do distinguish
its twists. The counterexample applies only after the stated projective
source compression. It does not assert that all arithmetic descent data
are already included in that compressed interface, nor does it deny that
the rational model determines the local polynomials. Since $\pi_p$ and
$-\pi_p$ have the same projective class, the twist test specifically
concerns signed traces and central lifts, not projective spectral ratios.
\end{scope}

\section{V5 verification, proof status, and the next arithmetic gate}
\label{sec:v5-status}

\subsection{Exact verification and reproducibility}
The accompanying \path{Arithmetic_Loading_Duality_Research_Notes_V5_checks.py}
uses integer arithmetic, finite-field pairs, and exact symbolic matrices.
It acquires each listed point count twice: by the norm-character table and
by independent direct squaring fibres with Horner polynomial evaluation.
It verifies the two quartic irreducibility witnesses, both Cartier matrices,
the twist at three, the positive companion Gram, its principal minors and
unitarity, the rational response, and the discriminant and source-weight
constant. It also generates all $11520$ projective two-qubit Clifford
normalizer actions as signed permutations of the sixteen Pauli matrices;
all adjoint traces belong to $\{0,1,2,4,8,16\}$, and the $96$ actions
inducing a residual four-cycle have trace zero or two. This exhaustive
check supplements, rather than substitutes for, the character proof in
\cref{lem:v5-normalizer-trace}.

The machine-readable audit is retained in the accompanying V5 certificate
JSON file. No numerical periods, fitted matrix entries, or floating-point
sign tests are used. Weil theory, abelian Riemann--Roch, Cartier--Serre duality, and
Frobenius splitting are cited classical inputs, not conclusions of the
finite script. The script is not a proof-assistant formalization. It also
compares all four earlier mathematical bodies byte for byte and records
their SHA-256 digests.

\subsection{Proved in this continuation}
\begin{enumerate}[label=\textup{(P5.\arabic*)},leftmargin=3.6em]
\item The actual curve has the exact local polynomials
      \eqref{eq:v5-P3} and \eqref{eq:v5-P5}. Their identical residual
      four-cycle type does not determine their weight-one data.
\item Both local Frobenius operators have infinite projective order,
      excluding every scalar-normalized fixed finite theta-frame lift.
      A positive Hilbert--Schmidt separation is explicit at three.
\item Native theta Frobenius is level raising. At three, the classical
      splitting criterion yields a unique compatible split theta tower;
      at five the required sheaf splitting is obstructed. The canonical
      lower-level roundtrip at three is identity, not the weight-one action.
\item The actual local spectrum at three has an explicit positive,
      source-minimal tracial companion, a uniform finite Gram bound, and
      a rational positive-real response. All ordinary unitary source
      couplings of that simple spectrum form the explicitly conditioned
      weight simplex \eqref{eq:v5-source-simplex}.
\item Positive quadratic twists preserve the complete projective complex
      theta source and real involution while changing signed local
      Frobenius traces. This identifies a central arithmetic descent
      datum absent from the compressed theta interface.
\end{enumerate}

\subsection{Inherited inputs and what remains unresolved}
The V1--V4 bodies and their historical ledgers are unchanged. The present
results do not repair the ordinary or signed fixed-pure-$S_4$ symmetry
obstructions by redefining the source. They do not rederive the already
established arithmetic loading, tetrahedral torsion geometry, integral
cycle lattice, operator-Gram comparison, or generic monoidal/interface
recognition theorems. The companion SM--spacetime and variational results
remain cited downstream infrastructure \cite{Duality,Variational}.

The revised positive target is now a cross-realization comparison, not a
search for another fixed four-dimensional theta permutation. One must
retain the rational central descent, declare whether the target is the
Tate-module conjugacy class or the level-raising theta tower, and acquire
an independently supplied mixed source map. On the weight-one route,
the actual theta-to-spectral source weights must be acquired rather than
set to the tracial values. On the section route, a useful complex or
$p$-adic realization must preserve the level maps and the relevant source;
the characteristic-three splitting alone does not provide that comparison.

Neither one local polynomial nor the positive scalar response determines
relative Frobenius matrices at different primes, tensor intertwiners,
Wilson loops, or open Redheffer interfaces. These require a common
arithmetic realization and its independently defined physical counterpart.
The four relative V4 lift signs, the central Store/chronology signature,
common marked arithmetic--renewal transport, and regulator-uniform
estimates remain unresolved. No generation count, physical coupling,
number-field GRH result, or Einstein equation is inferred.

\paragraph{Next noncircular attack.}
Use the canonical polarized Tate module or the native graded theta tower
as the arithmetic entrance, and determine the smallest mixed interface
that connects its actual source to the already acquired theta orbit.
A complete local comparison must retain the central character and the
source-weight/marked-insertion data in addition to $P_p$. The finite
normalizer route has an explicit obstruction and should not be reopened
without changing the realization category openly.

\paragraph{Version integrity.}
This V5 continuation is appended after V4. Every V1--V4 mathematical body
is preserved byte for byte, including its historical open-problem list.
Future continuations are to append new proofs and any explicit corrections
before the bibliography and advance the version filename.
% ====================== END IMMUTABLE V5 MATHEMATICAL BODY ======================
% ===================== BEGIN IMMUTABLE V6 MATHEMATICAL BODY =====================
\clearpage
\section{V6 continuation: from the theta source to differential arithmetic}
\label{sec:v6-context}

\subsection{The upstream question and the new comparison category}
V5 separates the actual Tate-module Frobenius, the finite Heisenberg frame,
and the level-raising section map. It leaves an independently defined
source coupling between section geometry and weight-one arithmetic. The
present continuation goes upstream of a proposed unitary: which
\emph{geometric map} actually connects the squared theta section to
cohomology? There is a concrete answer at quadratic differential order.
For the same genus-two Jacobian, the second jet at the origin gives
\begin{equation}
 \boxed{
 V_0:=\ker\bigl(H^0(J,L^2)\xrightarrow{\mathrm{ev}_0}L_0^2\bigr)
 \xrightarrow[\sim]{\ \mathfrak j_2\ }
 L_0^2\otimes\operatorname{Sym}^2 H^0(C,\Omega_C^1).}
 \label{eq:v6-main-jet}
\end{equation}
Here and below a superscript on a line bundle or its fibre denotes a tensor
power. This is a natural algebraic isomorphism, not a unitary identification
with the Tate module. The five displaced source rays of V1--V4 span
exactly $V_0$. Their images are the squares of the conormal lines at the
five finite Weierstrass points. The original squared theta ray maps to
the square of the conormal line at infinity.

The classical ingredients are the Abel--Jacobi immersion, the theta divisor
of a genus-two Jacobian, and second-order jets; the tangent evaluation
formula is recorded in \cite{MilneAV}, Chapter III, Section~2, and the
Gauss-map interpretation is classical \cite{deJongGauss}. We prove the
particular isomorphism and its source identification below. We do not
claim that the general geometry of the Kummer conic is new. The advance
relative to the supplied work is the explicit connection of the
\emph{actual previously acquired source} to a quadratic Hodge carrier,
its native Cartier transfer at the two computed primes, and its exact
incompatibility with fixed-order section Frobenius.

There is a second, independent arithmetic advance. Using the already
computed $P_3$ and $P_5$, rather than collecting more local traces, we prove
\begin{equation}
 \boxed{\End_{\overline{\Q}}\operatorname{Jac}(C_{\mathrm{tet}})=\Z.}
 \label{eq:v6-main-End}
\end{equation}
Consequently the new Hodge line is not a hidden rank-one or rank-two
rational weight-one submotive. It is a differential marking of an
absolutely simple arithmetic object.

\subsection{Dependency and nonduplication ledger}
The following inputs are inherited. V1--V4 identify the five displacement
rays, their pure-source rank, their marked overlaps, and the actual
Heisenberg orbit. V5 supplies the two Frobenius polynomials, the two
Cartier matrices in their declared transpose convention, and the native
section Frobenius. The earlier arithmetic--spacetime V9 notes supply the
rational curve, its marked Weierstrass points, and the theta translations
\cite{AD9}. None of the Peano loading, torsion quadratic-space bridge,
integral $K_4$ lattice, pure-Gram classification, point counts, finite
normalizer enumeration, or abstract Gram/Krylov comparison is rederived.
The old symmetric-square obstruction to a full colour group is also an
inherited boundary, not a proposed new result \cite{AD9}.

The genuinely new calculations are the conormal/second-jet identification,
Cartier cyclicity of the actual jet source, the Frobenius jet-order law,
and the two-prime geometric endomorphism certificate. The general
specialization method for endomorphism rings is standard; see
\cite{LombardoEnd}. We give the complete finite certificate for this
particular curve without assuming a database classification. Its algebraic
monodromy is then identified using the separately cited, established
abelian-surface classification.

\subsection{Fixed arithmetic model and fibre lines}
Throughout,
\[
 C:\ y^2=f(x)=x^5-x^2-x,\qquad
 J=\operatorname{Jac}(C),\qquad
 \kappa_\infty=\mathcal O_C(\infty).
\]
Let $i:C\to J$ send $P$ to $[P-\infty]$, put
$\Theta=i(C)$ and $L=\mathcal O_J(\Theta)$, and let $\theta$ denote a
nonzero section with divisor $\Theta$. Its square is the actual source
section $s_\infty=\theta^2$ up to nonzero scalar. Write
\[
 F=H^0(C,\Omega_C^1),\quad \lambda=L|_0,\quad
 \omega_0=dx/y,\quad\omega_1=x\,dx/y.
\]
The identification $T_0^*J\simeq F$ is the Abel--Jacobi pullback of
translation-invariant differentials. The line $\lambda$ is retained:
removing it would silently introduce a theta-fibre trivialization.

The integral model is taken over
$R=\Z[1/(2\cdot283)]$. The discriminant and good reduction are inherited
from \cite{AD9}. The smooth proper curve, its Jacobian, its marked
infinity section, and the theta line extend over $R$. All fibrewise
statements used below apply in characteristic zero and at the good odd
primes. For a finite Weierstrass abscissa $r$, put
$\delta_r=[W_r-\infty]\in J[2]$ and
$\Theta_r=t_{\delta_r}(\Theta)$. The squares of the corresponding
translated theta sections are rays $[s_r]\subset H^0(J,L^2)$.
The requisite line-bundle identification is unique only up to scalar,
which is exactly the ray freedom already retained in V1--V4.

\section{The actual theta-source span is a quadratic Hodge carrier}
\label{sec:v6-jet}

\subsection{Conormal lines determined by the curve, not by a fitted metric}
\begin{proposition}[Explicit arithmetic Gauss lines]
\label{prop:v6-Gauss}
The divisors $\Theta_r$ are smooth at the origin. Under $T_0^*J\simeq F$,
their conormal lines are
\begin{equation}
 \boxed{
 \ell_\infty=\Span(\omega_0),\qquad
 \ell_r=\Span(\omega_1-r\omega_0)\quad(f(r)=0).}
 \label{eq:v6-Gauss-lines}
\end{equation}
These six lines are distinct. They are functorial under isomorphisms of
the marked curve, and the labelled collection is Galois equivariant.
\end{proposition}
\begin{proof}
The Abel--Jacobi map is a closed immersion and identifies invariant
one-forms with $F$; these are the classical statements cited above. Its
differential at $W$ is dual to evaluation of holomorphic differentials
at $W$. Since $2\delta_r=0$, translating $i(C)$ by $\delta_r$ takes
$i(W_r)=\delta_r$ to zero. Translation does not change invariant
differentials. Hence the conormal line at zero is
$\ker(F\to\Omega_C^1|_{W_r})$.

At a simple finite branch point, choose a local parameter $t$ with
$x=r+t^2$ after a harmless extension of the residue field. Then
$y=t\sqrt{f'(r)+O(t^2)}$, so $dx/y$ has nonzero constant coefficient
and $x\,dx/y$ has constant coefficient $r$ times that coefficient.
The evaluation kernel is therefore $\Span(\omega_1-r\omega_0)$.
At infinity, $x=t^{-2}$ and
\[
 y=t^{-5}\sqrt{1-t^6-t^8},\qquad
 \omega_0=\frac{-2t^2\,dt}{\sqrt{1-t^6-t^8}},\qquad
 \omega_1=\frac{-2\,dt}{\sqrt{1-t^6-t^8}}.
\]
Thus $\omega_0$ vanishes at infinity whereas $\omega_1$ does not.
The same valuation argument works in every good characteristic other
than two. Distinct abscissae give distinct lines, and the infinity line
is not one of the finite ones. Evaluation, translation and the
Abel--Jacobi map commute with field automorphisms, proving naturality
and Galois equivariance.
\end{proof}

\begin{remark}[Lines, vectors and central signs]
The display uses the equation to give coordinates; its intrinsic output
is a collection of lines in $F$. Fixing the signed function $y$ can
rigidify particular vectors, but that is stronger marking than the curve,
branch points and squared-theta ray. The distinction is essential for
the central-inversion test in \cref{sec:v6-parity}.
\end{remark}

\subsection{The second jet is a complete map on the source span}
For the maximal ideal $\mathfrak m$ at zero, the second leading jet of a
section vanishing to order at least two is its image in
\[
 \lambda^2\otimes\mathfrak m^2/\mathfrak m^3
 \simeq\lambda^2\otimes\operatorname{Sym}^2T_0^*J.
\]
We use this leading coefficient rather than the Hessian, thereby avoiding
an unnecessary factor of two. It is independent of a local trivialization:
a change of trivialization acts on the leading term by its value at zero.

\begin{lemma}[Even sections and the vanishing subspace]
\label{lem:v6-even}
With the square symmetric linearization, inversion $[-1]$ acts as $+I$
on $H^0(J,L^2)$. Every section in $V_0=\ker\mathrm{ev}_0$ therefore
vanishes to order at least two, and $\dim V_0=3$.
\end{lemma}
\begin{proof}
In characteristic zero, use the second-order theta basis of V2:
\[
 \Theta_\sigma(z)=
 \theta\!\begin{bmatrix}\sigma/2\\0\end{bmatrix}(2z\mid2\tau).
\]
Changing the summation index $n$ to $-n-\sigma$ shows that each basis
section is even. The square of the symmetric linearization has fibre
value $+1$ at zero, so this is the asserted action, not a freely chosen
central sign. If the constant term of an even section is zero, its
linear term equals its negative and hence vanishes.

This identity also holds in the good integral model. Ample-line
cohomology and base change on the abelian scheme make $H^0(L^2)$ a
locally free rank-four $R$-module, commuting with the relevant field
extensions. The endomorphism $[-1]^*-I$ vanishes on its generic fibre,
and hence on the torsion-free module and on every special fibre.

The evaluation map is nonzero. One elementary basepoint-free argument
chooses $a$ with $a,-a\notin\Theta$ over an algebraic closure.
The product of the translated theta sections at $a$ and $-a$ is a
section of $L^2$ by the theorem of the square, and is nonzero at zero.
Thus evaluation has rank one. Abelian Riemann--Roch gives $h^0(L^2)=4$,
so its kernel has dimension three. These classical line-bundle inputs
are the ones already used for the theta levels in V5 \cite{MilneAV}.
\end{proof}

\begin{theorem}[Natural second-jet source isomorphism]
\label{thm:v6-jet-isomorphism}
The leading second-jet map is an isomorphism
\begin{equation}
 \boxed{\mathfrak j_2:V_0\xrightarrow{\sim}\lambda^2\otimes\operatorname{Sym}^2F.}
 \label{eq:v6-jet-iso}
\end{equation}
For each of the six Weierstrass points,
\begin{equation}
 \boxed{[\mathfrak j_2(s_r)]=[\lambda^2\otimes\ell_r^2].}
 \label{eq:v6-ray-jet}
\end{equation}
The five nonidentity displacement sources of V1--V4 span $V_0$.
The isomorphism is compatible with the good-reduction base changes of
the fixed model and with isomorphisms preserving the marked theta data.
\end{theorem}
\begin{proof}
Each $s_r$ is the square of a section vanishing simply on the smooth
translated theta divisor $\Theta_r$. If its local equation is
$a_r+O(\mathfrak m^2)$, with nonzero linear term $a_r$, then its square
has leading second term $a_r^2$. The line of $a_r$ is precisely the
conormal line in \cref{prop:v6-Gauss}. This proves
\eqref{eq:v6-ray-jet}, including its nonvanishing and fibre factor.

In the coordinates $(\omega_0^2,\omega_0\omega_1,\omega_1^2)$, the
square of a finite conormal line has a representative
\[
 q(r)=(r^2,-2r,1)^{\mathsf T},\qquad
 q(\infty)=(1,0,0)^{\mathsf T}.
\]
For three distinct finite abscissae $r_1,r_2,r_3$, direct expansion gives
\begin{equation}
 \det[q(r_1),q(r_2),q(r_3)]
 =2(r_2-r_1)(r_3-r_1)(r_3-r_2)\ne0.
 \label{eq:v6-Veronese-det}
\end{equation}
Multiplication of columns by their nonzero theta-fibre coefficients does
not change independence. Any three of the five finite branch abscissae
therefore give three independent images of sections in $V_0$.
Both $V_0$ and $\lambda^2\otimes\operatorname{Sym}^2F$ have dimension
three, so $\mathfrak j_2$ is an isomorphism and the five sources span $V_0$.
The nonidentity singular displacements in the earlier notes are exactly
these five translates; no new source has been substituted.

The construction by the maximal-ideal filtration is intrinsic. Evaluation,
translation, and tensoring line bundles commute with base change.
Over the integral model, evaluation is fibrewise onto and its kernel is
locally free. The second-jet map is fibrewise an isomorphism by the same
argument after a splitting-field extension of the branch points. A map
between locally free rank-three modules which is an isomorphism on all
fibres is an isomorphism. This proves the asserted good-reduction and
functorial compatibility.
\end{proof}

\begin{corollary}[The six source rays form a complete marked conic]
\label{cor:v6-conic}
Under $\mathbf P(\mathfrak j_2)$ the six source rays lie on the Veronese conic
\begin{equation}
 \boxed{[A:B:C]\in\mathbf P^2,\qquad B^2=4AC.}
 \label{eq:v6-conic}
\end{equation}
The labelled points are
$[r^2:-2r:1]$ for $f(r)=0$ and $[1:0:0]$ at infinity.
Their cross-ratios on the conic are the branch-point cross-ratios.
Over an algebraically closed field they recover the hyperelliptic curve
up to isomorphism. Over the base field they do not by themselves recover
the quadratic-twist character retained in V5.
\end{corollary}
\begin{proof}
The displayed quadratic equation is the equation of a square binary
quadratic. Its parametrization by $r$ is an isomorphism from
$\mathbf P^1$ to the conic, so the labelled cross-ratios agree.
The function field of a double cover of $\mathbf P^1$ branched at the
six points is $k(x)(\sqrt{c f(x)})$ after a projective coordinate
choice. Over an algebraically closed field $c$ is a square, proving
uniqueness. Over the base field its square class is precisely the
remaining quadratic-twist freedom. This is consistent with V5's signed
Frobenius counterexample and does not reprove its source comparison.
\end{proof}

\begin{scope}[What the quadratic bridge does and does not identify]
The rank three already proved in V1 is now identified geometrically as
$\operatorname{Sym}^2F$, not just counted. The fibre coefficients in
\eqref{eq:v6-ray-jet} are not set to one by a source-fitting convention.
Nor is the transported Heisenberg-unitary metric asserted to equal the
symmetric square of the Hodge metric. The map is an algebraic isomorphism;
its numerical metric and relative comparison with a renewal carrier
remain additional data. In particular this symmetric-square carrier is
not an independently generated $\mathrm{SL}_3$ colour object, as already
explained in \cite{AD9}.
\end{scope}

\section{A native Cartier transfer on the actual quadratic source}
\label{sec:v6-Cartier}

\subsection{The transfer is defined geometrically before it is computed}
Fix $p=3$ or $p=5$, and now work over $k=\F_p$. The Cartier operator
$\mathcal C_p:F_p\to F_p$ is $\F_p$-linear. Over a larger perfect field
it is inverse-Frobenius semilinear; no complex-linear extension is
implicit. Since the fibre line $\lambda_p^2$ is retained, the map
\begin{equation}
 \boxed{
 \mathcal K_p=\mathfrak j_{2,p}^{-1}
  (\id_{\lambda_p^2}\otimes\operatorname{Sym}^2\mathcal C_p)\mathfrak j_{2,p}
  :V_{0,p}\longrightarrow V_{0,p}}
 \label{eq:v6-K-def}
\end{equation}
is an intrinsic $\F_p$-linear operation. A fibre-line trivialization is
not required to define it. This is the \emph{Cartier-induced jet transfer};
it is not relabelled as the native section Frobenius or the rank-four
Tate-module Frobenius.

The matrices $M_3,M_5$ on $(\omega_0,\omega_1)$ are inherited from
\cref{prop:v5-Cartier}. We use column vectors throughout. For
$M=\begin{psmallmatrix}a&b\\c&d\end{psmallmatrix}$, direct multiplication
of binary quadratics gives
\begin{equation}
 \operatorname{Sym}^2M=
 \begin{pmatrix}
 a^2&ab&b^2\\
 2ac&ad+bc&2bd\\
 c^2&cd&d^2
 \end{pmatrix}
 \label{eq:v6-Sym2-matrix}
\end{equation}
in the unnormalized monomial basis used above.

\begin{theorem}[Complete Cartier-source calculation at three and five]
\label{thm:v6-Cartier-cyclic}
The jet source $\mathfrak j_2(s_\infty)$ is a nonzero fibre scalar times
$e_0=(1,0,0)^{\mathsf T}$. In these coordinates,
\begin{equation}
 \boxed{
 K_3=\begin{pmatrix}1&1&1\\1&2&0\\1&0&0\end{pmatrix}\pmod3,
 \qquad
 K_5=\begin{pmatrix}1&2&4\\0&0&0\\0&0&0\end{pmatrix}\pmod5.}
 \label{eq:v6-K-matrices}
\end{equation}
At three,
\begin{equation}
 \boxed{
 K_3^3=I,\quad (K_3-I)^3=0,\quad (K_3-I)^2\ne0,\quad
 \det[e_0,K_3e_0,K_3^2e_0]=1.}
 \label{eq:v6-K3-cyclic}
\end{equation}
Thus the actual theta jet source is cyclic on all of the three-dimensional
quadratic carrier, with minimum polynomial $(T-1)^3$.
At five,
\begin{equation}
 \boxed{K_5^2=K_5,\quad\rank K_5=1,\quad K_5e_0=e_0.}
 \label{eq:v6-K5-cyclic}
\end{equation}
The same source has cyclic dimension one and misses a two-dimensional
kernel. These are exact finite-field ranks, not positive Hilbert-space
gaps or generation counts.
\end{theorem}
\begin{proof}
The source identification follows from \eqref{eq:v6-Gauss-lines} and
\eqref{eq:v6-ray-jet}. Substitution of the two inherited matrices in
\eqref{eq:v6-Sym2-matrix} gives \eqref{eq:v6-K-matrices}.
At three the three source columns are
\[
 e_0,\qquad(1,1,1)^{\mathsf T},\qquad(0,0,1)^{\mathsf T};
\]
the next column is $e_0$. Their determinant is one. Consequently they
form a cyclic basis and the operation has minimum polynomial $T^3-1$,
which equals $(T-1)^3$ in characteristic three. Its degree is three,
so the square of $K_3-I$ cannot vanish. At five the displayed matrix
has rank one and square equal to itself; its first column is $e_0$.
All statements are preserved by a nonzero scaling of the source and
by the intrinsic isomorphism $\mathfrak j_2$.
\end{proof}

\begin{proposition}[The underlying conormal line and the dual tangent source]
\label{prop:v6-linear-Cartier}
On $F_3$ the conormal source has orbit
\begin{equation}
 \omega_0\longmapsto2\omega_0+\omega_1
 \longmapsto2\omega_1\longmapsto\omega_0,
 \label{eq:v6-Cartier-orbit}
\end{equation}
which spans $F_3$. On $F_5$ it is fixed by Cartier. The Serre-dual tangent
line $\Span((0,1)^{\mathsf T})\subset H^1(C_p,\mathcal O)$ is killed
by Hasse--Witt at five, while its Hasse--Witt orbit spans the full
tangent space at three.
\end{proposition}
\begin{proof}
Multiply $e_0=(1,0)^{\mathsf T}$ by $M_3$ and its powers, and by
$M_5$. The tangent line is the annihilator of $\ell_\infty$.
With the dual basis, the Hasse--Witt matrices are $M_p^{\mathsf T}$ by
V5's convention. The products are
$M_5^{\mathsf T}(0,1)^{\mathsf T}=0$ and
$M_3^{\mathsf T}(0,1)^{\mathsf T}=(1,0)^{\mathsf T}$.
The two tangent vectors at three are independent.
\end{proof}

\begin{proposition}[The jet transfer is not a branch-permutation lift]
\label{prop:v6-Cartier-not-frame}
At three the projective Cartier action takes the conormal line at infinity
to the conormal line with abscissa $1$. That abscissa is not a branch point
of $C_3$. Consequently $\mathcal K_3$ does not permute the six marked
squared-theta rays and is not a theta-frame normalizer action.
\end{proposition}
\begin{proof}
The first image in \eqref{eq:v6-Cartier-orbit} is
$2\omega_0+\omega_1=\omega_1-\omega_0$ modulo three. By
\eqref{eq:v6-Gauss-lines} this is the conormal label $r=1$.
But $f(1)=-1=2\pmod3$, so it is not among the six branch abscissae.
The Veronese map is injective; its square is therefore not another marked
source ray. This verifies directly why the newly constructed transfer
does not contradict V5's finite-frame obstruction.
\end{proof}

\section{Section Frobenius, jet order and the central parity boundary}
\label{sec:v6-parity}

\subsection{Why the fixed second-jet square cannot commute}
\begin{theorem}[Exact Frobenius jet-order law]
\label{thm:v6-jet-Frobenius}
Let $s$ be a nonzero section of a line bundle on a smooth integral variety
in characteristic $p$, with vanishing order $m>0$ at a point. Then
\begin{equation}
 \boxed{\operatorname{ord}(s^{p^r})=mp^r,\qquad
 \operatorname{in}_{mp^r}(s^{p^r})=
             \operatorname{in}_m(s)^{p^r}.}
 \label{eq:v6-jet-order}
\end{equation}
For the actual squared theta section $m=2$. Its leading second jet after
any positive Frobenius iterate is zero, although its leading
$2p^r$-jet is nonzero. On the subspace of Frobenius-power initial forms,
inverse Frobenius over a perfect field recovers the original leading
quadratic form, semilinearly and with its coefficient line restored.
\end{theorem}
\begin{proof}
After local trivialization, let the first nonzero homogeneous term of
$s$ be $h_m$. Smoothness identifies the associated graded local ring
with a polynomial ring on the cotangent space. In characteristic $p$,
raising to $p^r$ sends the first term to $h_m^{p^r}$ and every higher
term to strictly higher degree. The polynomial ring is reduced, so this
first term is nonzero. This proves both formulas intrinsically after
restoring the tensor power of the fibre line.
Over a perfect field, the Frobenius-power subspace consists of forms
whose exponents are divisible by $p^r$ and whose coefficients are
$p^r$th powers. Dividing the exponents and taking the unique inverse
Frobenius of the coefficients is additive and inverse-Frobenius
semilinear. It is inverse to the indicated leading-term operation;
it is not an inverse to a fixed-order differential of Frobenius.
\end{proof}

\begin{corollary}[A concrete noncommuting arithmetic comparison at three]
\label{cor:v6-noncommuting}
The fixed second jet of $j_{2,3}(s_\infty)=s_\infty^3$, where
$j_{2,3}$ is the \emph{level map} of V5, is zero. In contrast,
$\mathcal K_3s_\infty$ is nonzero and has a different source ray from
$s_\infty$. Even recovery from the leading sixth jet followed by
inverse Frobenius returns the original quadratic ray, not the
Cartier-transferred one. Thus the two transfers cannot be identified by
this canonical jet construction.
\end{corollary}
\begin{proof}
Apply \cref{thm:v6-jet-Frobenius} with $m=2$, $p=3$.
The Cartier assertions are \cref{thm:v6-Cartier-cyclic} and
\cref{prop:v6-Cartier-not-frame}. The recovered ray is
$[\omega_0^2]$, whereas its Cartier image is
$[(2\omega_0+\omega_1)^2]$. They are distinct.
\end{proof}

\begin{remark}[Different maps and different targets]
The inherited map $j_{2,3}$ sends a level-two section to level six.
The new map $\mathfrak j_{2,p}$ is the second-jet isomorphism on
$V_{0,p}$. The distinct notation retains their different domains,
codomains and roles. Neither is identified with physical time evolution.
\end{remark}

\subsection{What even theta data cannot generate linearly}
\begin{theorem}[Central-inversion obstruction for a linear weight-one source]
\label{thm:v6-parity}
Work in the category of polarized marked curves/Jacobians retaining the
Weierstrass labels and squared-theta ray, but not an additional signed
coordinate or an odd rigidification. The automorphism $[-1]$ acts as
$+I$ on $H^0(J,L^2)$ and as $-I$ on Betti, de Rham or Tate weight-one
cohomology in characteristic zero. Therefore every functorial linear
intertwiner from $H^0(J,L^2)$, or from tensor constructions generated by
it using duals, sums and subquotients, into weight-one cohomology is zero.
The same statement holds for the even quadratic carrier in
\eqref{eq:v6-jet-iso}.
\end{theorem}
\begin{proof}
Evenness is \cref{lem:v6-even}. Inversion is multiplication by $-1$
on the fundamental lattice, invariant differentials and Tate module;
by duality it is also $-I$ on weight-one cohomology. On tensor products,
duals, direct sums and subquotients of an even representation, its action
remains $+I$. An intertwiner $B$ must satisfy
$B=(+I)B=(-I)B$, hence $2B=0$ and $B=0$.
On $\lambda^2\otimes\operatorname{Sym}^2F$ the fibre-square action is
$+1$ and the two negative differential signs cancel. The same argument
applies.
\end{proof}

\begin{scope}[The obstruction does not exclude a projective Gauss line]
A line is unchanged by multiplication by $-1$. Thus
\cref{prop:v6-Gauss} does not contradict \cref{thm:v6-parity}.
Nor does the theorem rule out a chosen source vector after an odd
rigidification or a fixed signed $y$ coordinate has been included in the
input. It says that such a choice cannot be erased from the input while
still claiming a functorial \emph{linear} coupling. Odd theta
first derivatives are line-valued objects of the right parity, but their
normalization and comparison with a weight-one realization must still
be retained. This is a different issue from the four Heisenberg frame
signs and the quadratic-twist character isolated earlier.
\end{scope}

\section{A two-prime certificate for the global geometric endomorphism ring}
\label{sec:v6-End}

\subsection{The classical reduction inputs are explicit}
We use three standard facts about abelian varieties, not hypotheses of
an arithmetic--physics duality. Homomorphisms of abelian varieties act
faithfully on the prime-to-characteristic Tate module, with an injective
map even after tensoring with $\Q_\ell$; see \cite{MilneAV},
Chapter I, Theorem~10.15. Rational idempotents split an abelian variety
up to isogeny by Poincare complete reducibility. Finally, good reduction
gives an injective specialization of geometric endomorphism rings.
For the latter, all endomorphisms are defined over one finite number-field
extension; extend them over the good abelian scheme and identify their
prime-to-$p$ torsion with that of the special fibre. If the specialization
vanishes, its Tate action vanishes and faithfulness makes the generic
endomorphism zero. This also proves injection after tensoring with $\Q$.
The use of several reductions in this fashion is classical
\cite{LombardoEnd}, especially Section~4. No classification of the present
curve is imported.

Write $\End^0=\End\otimes_\Z\Q$. If a Frobenius operation has distinct
eigenvalues on a $2g$-dimensional Tate module, its centralizer has dimension
$2g$ over $\Q_\ell$. Hence an endomorphism algebra commuting with it has
dimension at most $2g$. This upper bound uses only the preceding faithful
injection, not the surjectivity part of Tate's isogeny theorem.

\subsection{All powers at three retain a simple spectrum}
Let $P=P_3$ be the inherited polynomial
$T^4-2T^3+4T^2-6T+9$, and let $\alpha$ be one of its roots.

\begin{lemma}[The quartic field and its unique quadratic subfield]
\label{lem:v6-K3-field}
The polynomial $P$ is irreducible over $\Q$. The field
$K=\Q(\alpha)$ has degree four, is not Galois over $\Q$, and has
exactly one quadratic subfield, namely $K^+=\Q(\sqrt3)$.
\end{lemma}
\begin{proof}
Dividing $P(\alpha)=0$ by $\alpha^2$ gives
\[
 (\alpha+3/\alpha)^2-2(\alpha+3/\alpha)-2=0.
\]
Thus $s=\alpha+3/\alpha=1\pm\sqrt3$ and $K$ contains $K^+$.
Over $K^+$, $\alpha$ satisfies $T^2-sT+3=0$ with discriminant
$d=-8\pm2\sqrt3$. Under both real embeddings this discriminant is
negative, so it cannot be a square in the real field $K^+$. The degree
is therefore four, proving irreducibility.

Choose $d=-8+2\sqrt3$, with conjugate $d'=-8-2\sqrt3$.
Their product is $52$. In a quadratic extension $K^+(\sqrt d)$, a
square root of $d'$ can occur only if $d'$ or $d'/d$ is a square in
$K^+$. Indeed expanding $(a+b\sqrt d)^2=d'$ forces $2ab=0$.
The first alternative fails by negativity; the second is equivalent to
$13$ being a square in $\Q(\sqrt3)$. If $(a+b\sqrt3)^2=13$ for
$a,b\in\Q$, then $ab=0$, leading to $a^2=13$ or $3b^2=13$, neither
possible. Hence the nontrivial automorphism of $K^+$ does not extend to
$K$, so $K/\Q$ is not normal. If $K$ contained two distinct quadratic
subfields, their compositum would be the biquadratic Galois field $K$.
Thus $K^+$ is its unique quadratic subfield.
\end{proof}

\begin{lemma}[No quotient of distinct Frobenius roots is a root of unity]
\label{lem:v6-root-ratios}
If $\alpha_i,\alpha_j$ are distinct roots of $P$, then
$\alpha_i/\alpha_j$ is not a root of unity. In particular,
$\pi_3^n$ has four distinct eigenvalues for every $n\ge1$.
\end{lemma}
\begin{proof}
The exact resultant is
\begin{align}
 \operatorname{Res}_T(P(T),Y^4P(T/Y))
 &=81(Y-1)^4 A(Y)^2B(Y),\\
 A(Y)&=3Y^4+2Y^3+2Y^2+2Y+3,\\
 B(Y)&=9Y^4+12Y^3+10Y^2+12Y+9.
 \label{eq:v6-resultant}
\end{align}
It follows either by the Sylvester determinant or by eliminating the two
quadratic factors over $\Q(\sqrt3)$; the accompanying script verifies
the identity over $\Z$. Its roots are exactly the pairwise quotients
with their multiplicities. The polynomial $P$ is separable, and neither
$A(1)$ nor $B(1)$ vanishes, so the four factors $Y-1$ account only for
the diagonal quotients.

For a root $\zeta$ of $A$, put $u=\zeta+\zeta^{-1}$. Division by
$\zeta^2$ gives $3u^2+2u-4=0$. This polynomial is primitive and
irreducible over $\Q$, since its discriminant $52$ is not a square,
and its monic minimal polynomial is not integral. Thus $u$ is not an
algebraic integer. For a root of $B$, the corresponding polynomial is
$9u^2+12u-8$, with nonsquare discriminant $432$, and the same conclusion
holds. If $\zeta$ were a root of unity, both $\zeta$ and $\zeta^{-1}$,
and hence $u$, would be algebraic integers. This excludes every
nondiagonal quotient. Equality $\alpha_i^n=\alpha_j^n$ would give an
excluded root of unity, proving the last statement.
\end{proof}

\begin{proposition}[The geometric endomorphism algebra at three]
\label{prop:v6-End3}
For the actual reduction,
\begin{equation}
 \boxed{\End^0_{\overline{\F}_3}(J_3)=\Q[\pi_3]\cong K.}
 \label{eq:v6-End3}
\end{equation}
In particular, $J_3$ is absolutely simple.
\end{proposition}
\begin{proof}
An endomorphism defined over $\F_{3^n}$ commutes with $\pi_3^n$.
By \cref{lem:v6-root-ratios} and Tate faithfulness, its rational
endomorphism algebra has dimension at most four. It contains the
four-dimensional field $\Q[\pi_3]$ by
\cref{lem:v6-K3-field}; therefore it equals that field.
Every geometric endomorphism is defined over a finite extension of the
finite field, so taking the union over $n$ gives
\eqref{eq:v6-End3}. A geometric product decomposition would give a
nontrivial idempotent by complete reducibility, impossible in a field.
\end{proof}

\subsection{The ordinary elliptic factor at five excludes the quadratic field}
The other inherited polynomial has the exact factorization
\begin{equation}
 P_5(T)=(T^2+5)(T^2+4T+5).
 \label{eq:v6-P5-factor}
\end{equation}

\begin{proposition}[A geometric endomorphism quotient at five]
\label{prop:v6-End5}
Up to $\F_5$-isogeny, $J_5$ is the product of elliptic curves
$E_0\times E_1$ with Frobenius polynomials $T^2+5$ and
$T^2+4T+5$, respectively. The two factors remain nonisogenous over
$\overline{\F}_5$, and
\begin{equation}
 \boxed{\End^0_{\overline{\F}_5}(E_1)=\Q(i).}
 \label{eq:v6-EndE1}
\end{equation}
Consequently the geometric endomorphism algebra of $J_5$ has a unital
algebra projection onto $\Q(i)$.
\end{proposition}
\begin{proof}
The two factors of \eqref{eq:v6-P5-factor} are coprime. Chinese
remainder idempotents in $\Q[\pi_5]$ split the abelian variety up to
isogeny. Each idempotent has rank two on a Tate module, hence its image
has dimension one. This proves the asserted elliptic decomposition
without assuming explicit elliptic equations.

The roots for $E_1$ are $\gamma=-2+i$ and $\bar\gamma=-2-i$.
Their quotient is $(3-4i)/5$, which is not an algebraic integer in
$\Q(i)$ and therefore not a root of unity. All powers have two
 distinct eigenvalues. The centralizer argument of
\cref{prop:v6-End3}, now in dimension two, gives
\eqref{eq:v6-EndE1}.

A root $\beta$ for $E_0$ satisfies $\beta^2=-5$. If over a finite
extension the two elliptic curves were isogenous, some powers of their
Frobenius eigenvalues would coincide. Squaring such an equality would
make
\[
 \gamma^2/(-5)=(-3+4i)/5
\]
a root of unity (or its conjugate for $\bar\gamma$). This is again
not an algebraic integer in $\Q(i)$. Thus they remain nonisogenous.
A nonzero homomorphism between elliptic curves is an isogeny, so their
geometric cross-Hom spaces vanish. The geometric endomorphism algebra
of their product is consequently the product of their two endomorphism
algebras, and projection onto the second factor is unital with target
$\Q(i)$. No description of the supersingular factor's full quaternion
algebra is needed.
\end{proof}

\begin{theorem}[Geometric endomorphism rigidity of the actual arithmetic curve]
\label{thm:v6-EndZ}
The Jacobian used throughout V1--V5 satisfies
\begin{equation}
 \boxed{\End_{\overline{\Q}}(J)=\Z.}
 \label{eq:v6-EndZ}
\end{equation}
In particular it is absolutely simple and has neither geometric real
multiplication by a quadratic field nor geometric complex multiplication.
\end{theorem}
\begin{proof}
Specialization at three embeds
$E=\End^0_{\overline{\Q}}(J)$ into the field $K$ of
\cref{prop:v6-End3}. A finite-dimensional unital $\Q$-subalgebra of a
field is itself a field: multiplication by a nonzero element is injective,
hence surjective, and supplies its inverse. Thus $E$ is a number field.

Specialization at five followed by the projection in
\cref{prop:v6-End5} is a unital homomorphism $E\to\Q(i)$. A unital
homomorphism out of a field is injective. Hence $[E:\Q]\le2$.
If the degree were two, the embedding in $K$ would identify $E$ with
its unique quadratic subfield $\Q(\sqrt3)$, whereas the embedding in
$\Q(i)$ would identify it with $\Q(i)$. These fields are not
isomorphic over $\Q$. Therefore $E=\Q$.
The integral endomorphism ring is a finite-rank $\Z$-module containing
$\Z$ and consisting of elements integral over $\Z$; its image in
$\Q$ is consequently contained in $\Z$. This proves equality.
A nontrivial geometric isogeny decomposition would give a nontrivial
idempotent in $E$, and extra multiplication would give a larger
subfield, proving the consequences.
\end{proof}

\begin{corollary}[The jet source does not define a rational weight-one splitting]
\label{cor:v6-Hodge-nosplit}
The rational weight-one Hodge structure of $J$ has no nonzero proper
polarizable rational sub-Hodge structure. In particular its Hodge
conormal line and the rank-two space $F$ do not arise from a proper
rational abelian subvariety or a rational algebraic projector on $H^1$.
\end{corollary}
\begin{proof}
A rational morphism of weight-one Hodge structures of complex abelian
varieties is induced, up to an integer denominator, by a homomorphism of
abelian varieties: on analytic lattices this is the condition that a
rational lattice map be complex linear. A polarizable rational sub-Hodge
structure has a rational orthogonal Hodge complement, and therefore
would yield a nontrivial rational idempotent in $\End^0(J)=\Q$.
There are no such idempotents. The assertions about abelian subvarieties
also follow from complete reducibility. The line in $F$ remains a
perfectly legitimate differential marking; a Hodge-filtration marking
is not a rational direct summand of the full weight-one object.
\end{proof}

\subsection{The actual tensor closure is symplectic, not a selected rank-two factor}
The next step uses an additional classical theorem rather than an inference
from a scalar commutant alone. For an abelian surface over a number field
with geometric endomorphism algebra $\Q$, the Mumford--Tate group is
$\operatorname{GSp}_4$ and the $\ell$-adic algebraic monodromy group is
connected and equals its $\Q_\ell$-base change. This is the abelian-surface
case of the established Mumford--Tate/endomorphism classification, stated
in \cite{LombardoEnd}, Section~2.2, case~(1); see also the genus-two
classification in \cite{FKRS}. It is used as a cited theorem. No
Sato--Tate equidistribution assertion is needed.

\begin{corollary}[Full algebraic monodromy of the actual arithmetic object]
\label{cor:v6-GSp4}
Let $V_\ell=T_\ell J\otimes\Q_\ell$ for the present rational Jacobian,
and let $G_\ell$ be the Zariski closure of its Galois image. Then
\begin{equation}
 \boxed{
 \operatorname{MT}(J)=\operatorname{GSp}_{4,\Q},\qquad
 G_\ell=\operatorname{GSp}_{4,\Q_\ell},\qquad
 G_\ell^{\mathrm{der}}=\operatorname{Sp}_{4,\Q_\ell}.}
 \label{eq:v6-GSp4}
\end{equation}
The same algebraic monodromy is retained after any finite extension of
the number field.
\end{corollary}
\begin{proof}
Theorem~\ref{thm:v6-EndZ} verifies the endomorphism hypothesis for this
specific curve. Apply the stated classical abelian-surface theorem.
For the last assertion, a finite-index subgroup of a Zariski-dense group
has Zariski closure containing the identity component of the original
closure. Indeed the original algebraic group is a finite union of cosets
of the smaller closure, so they have equal dimension and their identity
components agree. The group in \eqref{eq:v6-GSp4} is connected, proving
that the smaller closure is the whole group.
\end{proof}

\begin{corollary}[No nonabelian rank-two or rank-three object in this Tate tensor category]
\label{cor:v6-small-ranks}
Form a representation from $V_\ell$ by finite tensor products, duals,
direct sums and subquotients. If its dimension is at most three, its
restriction to the derived group $\operatorname{Sp}_4$ is trivial.
Thus it factors through the one-dimensional similitude torus and has
abelian connected algebraic monodromy. The conclusion persists after
finite base extension and after twisting by a character.
\end{corollary}
\begin{proof}
Each indicated construction carries an algebraic representation of
$G_\ell=\operatorname{GSp}_4$. In characteristic zero the split Lie
algebra $\mathfrak{sp}_4$ is simple and has dimension ten. A Lie map
$\mathfrak{sp}_4\to\mathfrak{gl}_n$ for $n\le3$ must either be
injective or zero, because its kernel is an ideal. Injectivity is
impossible since $n^2\le9<10$. Hence the Lie map is zero. The image
of the connected group $\operatorname{Sp}_4$ is consequently connected
of dimension zero, and is trivial. Quotienting
$\operatorname{GSp}_4$ by its derived group gives the similitude
torus. Finite base extension leaves the monodromy unchanged by
\cref{cor:v6-GSp4}; a character is trivial on the derived group.
\end{proof}

\begin{scope}[The Hodge carrier and the Tate tensor category are different]
The rank-three space $\operatorname{Sym}^2F$ in the jet theorem is a
Hodge-filtration construction, not a rank-three subquotient representation
of the full $\operatorname{GSp}_4$ arithmetic local system. The two
corollaries make that distinction testable rather than terminological.
They exclude extracting a full nonabelian rank-two weak/spin or rank-three
colour system from this one Tate object by the listed tensor operations.
They do not exclude additional independently constructed arithmetic
objects, enriched markings, or a different comparison category. The
monodromy calculation also does not produce a common positive Hermitian
metric, a mixed source kernel, or a physical gauge connection. Those
remain separate interfaces.
\end{scope}


\section{V6 verification, proved scope and the next source problem}
\label{sec:v6-status}

\subsection{Executable certificate and its limits}
The V6 verifier is a stand-alone exact SymPy calculation. It checks the
Veronese determinant, the binary-quadratic matrix formula, both Cartier
source orbits and their Krylov determinants, the rank-one branch at five,
the nonbranch image at three, and characteristic-$p$ polynomial jet
examples. It separately checks the degree-sixteen resultant
\eqref{eq:v6-resultant}, its coefficients and trace quadratics, the
quartic factorization over $\Q(\sqrt3)$, the norm $dd'=52$, the
Chinese-remainder elliptic idempotents at five, and the nonintegral
Gaussian ratios used to separate the geometric factors. These checks
supply finite algebraic witnesses for the displayed proofs. They do not
compute a period matrix or assume one.

The Abel--Jacobi, line-bundle, base-change, Tate-faithfulness,
specialization and complete-reducibility inputs are cited classical
geometry; the abelian-surface Mumford--Tate classification is a further
cited input to the monodromy corollary. These are not claims proved by the
script. The script is not a
proof-assistant formalization. It hashes and compares every earlier
immutable mathematical body and records the new polynomial and matrix
certificate in a JSON file. No earlier theorem or historical scope
statement is rewritten.

\subsection{Proved in this continuation}
\begin{enumerate}[label=\textup{(P6.\arabic*)},leftmargin=3.6em]
\item The actual five-displacement source span is exactly the vanishing
      theta-section space $V_0$, and its second jet gives the natural
      isomorphism $V_0\simeq\lambda^2\otimes\operatorname{Sym}^2F$.
      Its six marked rays form the explicit branch-point Veronese conic.
\item This geometric map gives a native Cartier-induced transfer on
      the actual quadratic source. It is source-cyclic of rank three
      at $p=3$ and source-cyclic of rank one at $p=5$, with exact
      matrices and recurrence certificates.
\item Section Frobenius raises the first nonzero jet order from $2$
      to $2p^r$. Fixed second jets vanish, while inverse-Frobenius
      recovery of the higher jet does not equal the Cartier transfer.
\item Even squared-theta data cannot supply a natural nonzero linear
      weight-one source without retaining odd marking. The projective
      Gauss line is allowed and has been constructed explicitly.
\item The two already acquired local Frobenius polynomials prove
      $\End_{\overline{\Q}}(J)=\Z$ by a complete reduction certificate.
      There is no hidden rational weight-one projector selecting the
      new Hodge marking.
\item With the explicitly imported abelian-surface classification, the
      actual algebraic monodromy is $\operatorname{GSp}_4$. Its Tate tensor
      category has no nonabelian object of dimension two or three;
      the quadratic Hodge source is not such a subquotient.
\end{enumerate}

\subsection{Inherited inputs}
The normalized arithmetic loading and the companion spacetime, spectral,
SM action--multiplicity and variational results remain antecedents
\cite{GT77,Arithmetic,Spacetime,Spectral,Duality,Variational}.
The actual source orbit and frame signs, local Frobenius polynomials,
Cartier convention, central twist example and theta-level maps are
exactly those of V1--V5. Classical Gauss and Kummer geometry underlie
\cref{sec:v6-jet}; the source-specific identification and transfer
calculation are the present contribution. The two-prime endomorphism
argument is a new certificate for this object within a standard method,
not a claim to have invented reduction bounds for endomorphism rings.

\subsection{Unresolved, and a narrower next comparison}
The programme now has an \emph{actual geometric source map}, but its
target is quadratic Hodge data, not a chosen rank-four Hilbert carrier.
The next admissible mixed interface must retain
\[
 \boxed{(V_0,\mathfrak j_2,\lambda^2,F,\ell_\infty;
          \text{metric, period/crystalline comparison, central descent}).}
\]
The equation gives a rational coordinate representative $dx/y$ for
$\ell_\infty$ once the signed $y$ marking is retained. Betti periods
would give its vector in $H^1(C,\C)$; a crystalline comparison would
instead give its $p$-adic realization. Neither the numerical periods,
a crystalline Frobenius matrix in that marked basis, nor their positive
Hermitian comparison is evaluated here. On a local spectral route the
source moments/weights remain to be acquired independently rather than
chosen from the Frobenius polynomial.

The exact Cartier calculations provide a small, testable starting point
for the crystalline route. They are not a characteristic-zero positive
transfer and do not preserve the finite theta-frame labels. The second-jet
metric is likewise not set equal to a symmetric-square Hodge metric by
definition. The four relative arithmetic--renewal frame signs, central
Store/chronology signature, actual mixed renewal comparison, and
Frobenius--Wilson tensor/interface comparison remain unresolved.
The new rank-three carrier is not an SM generation or colour derivation,
and no GRH, quantum-continuum, or Einstein conclusion is drawn.

\paragraph{Version integrity.}
This continuation is appended after the complete V5 body. Every V1--V5
mathematical body, including its previous status ledger, is retained byte
for byte. Subsequent work must append after this V6 body, record any
correction explicitly, and advance the filename.
% ====================== END IMMUTABLE V6 MATHEMATICAL BODY ======================
% ===================== BEGIN IMMUTABLE V7 MATHEMATICAL BODY =====================
\clearpage
\section{V7 continuation: the source metric and its differential completion}
\label{sec:v7-context}

\subsection{The question after V6}
The preceding continuation identifies the actual three-dimensional theta-source
span with a quadratic Hodge carrier, but deliberately leaves two different
questions open. First, what positive metric does the actual theta source induce
under this algebraic map? Second, can the quadratic Hodge carrier support a
characteristic-zero arithmetic transport without leaving dimension three?
Neither question is answered by the finite-field Cartier calculation.
This continuation addresses both, without fitting a source to a Frobenius
polynomial or redoing the finite theta-frame classification.

The results are constructive before they are obstructive. Four independently
marked source rays reconstruct the conformal second-jet metric by algebraic
operations on the branch points and the already acquired overlaps. The metric
is not the symmetric square of \emph{any} positive metric on holomorphic
one-forms. Next, on the explicitly declared deformation
\[
 C_t:\quad y^2=x^5-x^2-tx,\qquad C_1=C_{\mathrm{tet}},
\]
we compute the actual algebraic Gauss--Manin connection and show that the
original quadratic source generates its entire ten-dimensional symmetric-square
completion. The connection has a nontrivial unipotent local monodromy; no
positive parallel metric can make it a fixed-metric unitary local system.

The deformation is a new, independently specified family containing the
arithmetic curve, not a family canonically selected by one isolated fibre or by
the renewal law. Its parameter is not physical time. The construction supplies
an actual differential comparison object for this family; it does not supply a
crystalline Frobenius matrix, an arithmetic--renewal period identification, or a
physical Wilson connection. The theta-fibre line remains explicit throughout.

\subsection{Inherited material and the new proof obligations}
V3 supplies the exact branch-weight expressions, the two exact zero overlaps,
and outward rational enclosures for the actual source magnitudes. V4 selects
the actual Heisenberg orbit and fixes a representative with known signs.
V6 supplies the second-jet isomorphism and its labelled Veronese rays. These
are used, not rederived. The normalized chronology loading, integral $K_4$
bridge, generic Gram recognition, point counts, Cartier calculation and
endomorphism/monodromy classification are likewise inherited. No generic
Krylov theorem is claimed as a new result here: the advance is the explicit
connection and a nonzero cyclic determinant for the actual differential source.

Classical algebraic de Rham cohomology and the Gauss--Manin interpretation are
cited inputs \cite{KatzOda,KTdeRham}. The rational connection is proved by four
exact differential identities, not inferred from a numerical period matrix.
The local regular-singular normal form needed below is proved directly by a
convergent recurrence. It does not rely on a guess about topological
monodromy. All exact rational and interval calculations have a separate
reproducible certificate.

\section{An algebraic reconstruction of the conformal second-jet metric}
\label{sec:v7-metric}

\subsection{A fixed projective frame}
Retain the V4 representative $\rho_\star=\ketbra{\eta}$ with
$\norm\eta=1$, and write
\[
 \xi_0=\gamma_0\eta,\quad \xi_1=\gamma_1\eta,\quad
 \xi_2=\gamma_2\eta,\quad \xi_3=\gamma_3\eta,\quad
 \xi_4=\gamma_4\eta.
\]
Their respective finite branch labels are $0,u,v,z,\bar z$. The original
source $\eta$ has label infinity. A change of the residual Heisenberg frame
only rephases these labelled rays together with a common unitary. The
projective constructions below do not depend on that change.

Put $b_{ab}=\Tr(\rho_\star i\gamma_a\gamma_b)$ for $a<b$ and extend
antisymmetrically. The known Gram has
\[
 G_{aa}=1,\qquad G_{ab}=-i b_{ab}\quad(a\ne b).
\]
Use the three-column frame $X=(\xi_1,\xi_2,\xi_3)$ on $V_0$. Its Gram
and its overlaps with the fourth point $\xi_0$ are
\begin{equation}
 K=X^*X=
 \begin{pmatrix}
 1&0&-ib_{13}\\
 0&1&-ib_{23}\\
 ib_{13}&ib_{23}&1
 \end{pmatrix},\qquad
 g=X^*\xi_0=i(b_{01},b_{02},b_{03})^{\mathsf T}.
 \label{eq:v7-Kg}
\end{equation}
The source coordinates $c=K^{-1}g$ satisfy $Xc=\xi_0$, since the three
frame vectors span $V_0$ by \cref{thm:v6-jet-isomorphism}.

On the de-twisted target $\operatorname{Sym}^2F$ use the unnormalized monomial
basis $(\omega_0^2,\omega_0\omega_1,\omega_1^2)$, and put
\begin{equation}
 q(r)=(r^2,-2r,1)^{\mathsf T},\qquad
 Q=(q(u),q(v),q(z)),\qquad d=Q^{-1}q(0).
 \label{eq:v7-conic-frame}
\end{equation}
Every three distinct points on the nonsingular Veronese conic are linearly
independent. Consequently $Q$ and $K$ are invertible and all entries of $c$
and $d$ are nonzero: a zero coordinate would put the fourth point in the span
of only two other conic points. This assertion also has an interval check for
the actual data below.

\begin{theorem}[Period-free conformal second-jet metric]
\label{thm:v7-metric-reconstruction}
Define
\begin{equation}
 D=\operatorname{diag}(d_1/c_1,d_2/c_2,d_3/c_3),\qquad L=QD,\qquad
 \boxed{H_\Theta=L^{-*}KL^{-1}.}
 \label{eq:v7-Htheta}
\end{equation}
Then $H_\Theta$ is the unique normalized representative of the positive
Hermitian conformal class transported by the actual second-jet map. The
normalization is $q(0)^*H_\Theta q(0)=1$, or equivalently the choice that
$\xi_0$ maps to $q(0)$ after a theta-fibre trivialization. Every entry of this
representative is algebraic over $\Q$ and is determined by the already
acquired arithmetic source packet. No numerical period is needed.

The formula is invariant under independent rephasing of the source frame and
the fourth source. Under a change of the differential basis it transforms as
a Hermitian form. It fixes neither an absolute physical metric scale nor a
preferred theta-fibre trivialization.
\end{theorem}
\begin{proof}
The actual second-jet map takes $\xi_j$ to a nonzero scalar multiple of the
corresponding $q(r_j)$, by \cref{thm:v6-jet-isomorphism}. After choosing the
single target-line scalar so that $\xi_0\mapsto q(0)$, write the unknown
three column scalars as $\ell_j$. The equation $Xc=\xi_0$ gives
\[
 Q\operatorname{diag}(\ell_1,\ell_2,\ell_3)c=q(0)=Qd.
\]
Invertibility of $Q$ and nonvanishing of each $c_j$ give
$\ell_j=d_j/c_j$ uniquely. Thus the actual map in source coordinates is
exactly $L$. For a target vector $y=Lx$, the inherited source norm is
$x^*Kx=y^*L^{-*}KL^{-1}y$, proving the formula and positivity.
Since $Lc=q(0)$, its squared norm is $c^*Kc=\norm{\xi_0}^2=1$.

The branch points are algebraic. V3's squared moments and V4's selected signs
make every $b_{ab}$ real algebraic. Matrix inversion and division by the
nonzero $c_j$ therefore give algebraic entries in \eqref{eq:v7-Htheta}.
For the phase assertion, replace $X$ by $XE$ and $\xi_0$ by $e_0\xi_0$,
where $E$ is diagonal unitary and $|e_0|=1$. Then
$K\mapsto E^*KE$, $c\mapsto e_0E^*c$, and $L\mapsto e_0^{-1}LE$.
Substitution in \eqref{eq:v7-Htheta} leaves $H_\Theta$ unchanged.
An arbitrary scalar change of the actual target-fibre identification instead
rescales the transported metric. This is precisely the conformal ambiguity
removed by the stated normalization, not a new determination of a physical
unit. The usual covariant change-of-basis law follows from the same norm
identity.
\end{proof}

\begin{proposition}[Reality and an exact orthogonal endpoint pair]
\label{prop:v7-metric-real}
In the real algebraic differential basis above, $H_\Theta$ is real symmetric
and
\begin{equation}
 \boxed{(H_\Theta)_{02}=0,\qquad (H_\Theta)_{22}=1.}
 \label{eq:v7-metric-zero}
\end{equation}
Indices in this display are monomial indices $0,1,2$, not the five Clifford
indices. The equality $(H_\Theta)_{00}=1$ is \emph{not} needed below or
inferred merely from its narrow numerical enclosure.
\end{proposition}
\begin{proof}
The curve, infinity marking, differential basis and jet map are real. V4's
arithmetic real descent is antiunitary for the Heisenberg metric and acts by
ordinary conjugation on the real jet target, after retaining the real
one-dimensional theta fibre. A Hermitian form invariant under conjugation is
real symmetric; a complex scalar change of the fibre leaves that conclusion
unchanged. V3 proves $\langle\eta,\gamma_0\eta\rangle=0$ exactly. Their
jet rays are $q(\infty)=e_0$ and $q(0)=e_2$, with nonzero column scalars, so
$(H_\Theta)_{02}=0$. The final diagonal identity is the normalization in
\cref{thm:v7-metric-reconstruction}.
\end{proof}

\subsection{Certified values, not an equality inferred from rounding}
Substituting the already certified root disks and branch weights into
\eqref{eq:v7-Kg}--\eqref{eq:v7-Htheta} gives the following outward rational
intervals. Every displayed decimal endpoint is an exact rational number.
\begin{equation}
\begin{array}{c|cc}
\text{entry}&\text{lower bound}&\text{upper bound}\\ \hline
(H_\Theta)_{00}&0.999999999999&1.000000000001\\
(H_\Theta)_{01}&-0.048518016314&-0.048518016313\\
(H_\Theta)_{11}&0.541265554005&0.541265554006\\
(H_\Theta)_{12}&-0.176537203420&-0.176537203419
\end{array}
\label{eq:v7-metric-intervals}
\end{equation}
The remaining entries follow from \cref{prop:v7-metric-real}.
The actual source-frame determinant is
\begin{equation}
 \det K=1-b_{13}^2-b_{23}^2>\frac{51}{100},
 \label{eq:v7-K-floor}
\end{equation}
and interval inversion certifies $|c_j|^2>1/3$ for all three entries.
There is therefore no near-zero denominator hidden in this reconstruction.

\begin{proposition}[A finite coordinate metric certificate]
\label{prop:v7-H-bounds}
The formulas above imply the strict coordinate bounds
\begin{equation}
 \boxed{\frac{31}{100}I\prec H_\Theta\prec\frac{59}{50}I.}
 \label{eq:v7-H-floor}
\end{equation}
The enclosures in \eqref{eq:v7-metric-intervals} and
\eqref{eq:v7-K-floor} have a finite rational certificate.
\end{proposition}
\begin{proof}
The inherited root disks enclose exactly the four quartic roots. V3's positive
branch-distance products yield enclosing intervals for all required squared
moments. Take the specified positive square roots, install V4's signs, form
$K,g,Q$, and perform outward complex interval inversion in the three-by-three
formula. At every inversion step the squared modulus of the pivot has a
strictly positive rational lower bound. All operations are finite rational
operations, with square roots bounded by integer-square-root inequalities.
The accompanying verifier gives the individual pivot and coordinate bounds
and returns the intervals above; this is an evaluation of the exact algebraic
formula, not a numerical estimate of a period.

For the operator bounds, use reality and the exact zero entry in
\eqref{eq:v7-metric-zero}. The smallest diagonal-minus-row-sum bound is
larger than
$0.541265554005-0.048518016314-0.176537203420>0.31$.
Every diagonal-plus-row-sum bound is smaller than $1.18$. The elementary
Hermitian row-sum bound, obtained by estimating each off-diagonal term with
$2|x_i x_j|\le |x_i|^2+|x_j|^2$, proves
\eqref{eq:v7-H-floor}.
\end{proof}

\begin{scope}
The result determines the \emph{theta-transported conformal metric}. It does
not compute the independently defined Hodge metric, an absolute theta-fibre
norm, or a mixed renewal metric. It is valid in the algebraic differential
coordinates fixed by the curve, rather than in a coordinate system chosen to
whiten the answer. The metric bounds are finite coordinate bounds at $t=1$,
not regulator-uniform physical coercivity estimates.
\end{scope}

\section{The actual theta metric is not a symmetric-square Hodge metric}
\label{sec:v7-product-metric}

\subsection{An exact obstruction in the marked algebraic coordinates}
For clarity, the symmetric product $e_0e_1$ is represented in the tensor square
by $(e_0\otimes e_1+e_1\otimes e_0)/2$. If a positive Hermitian form on
$F$ has matrix
\[
 h=\begin{pmatrix}a&c\\\bar c&b\end{pmatrix},
\]
its tensor-induced form on the monomial basis has matrix
\begin{equation}
 h^{[2]}=
 \begin{pmatrix}
 a^2&ac&c^2\\
 a\bar c&(ab+|c|^2)/2&bc\\
 \bar c^2&b\bar c&b^2
 \end{pmatrix}.
 \label{eq:v7-product-matrix}
\end{equation}
This formula follows by expanding tensor-product inner products; it makes no
assumption about the Hodge metric.

\begin{theorem}[Nonfactorizability of the actual quadratic theta metric]
\label{thm:v7-nonproduct}
There is no positive Hermitian form $h$ on $F$ and no positive scalar $s$ with
\begin{equation}
 \boxed{H_\Theta=s h^{[2]}.}
 \label{eq:v7-forbidden-product}
\end{equation}
In particular the actual second-jet map is not conformally unitary for the
symmetric square of the Hodge metric, nor for the symmetric square of any
other positive one-form metric.
\end{theorem}
\begin{proof}
The exact entry $(H_\Theta)_{02}=0$ would force $sc^2=0$, hence $c=0$.
Equation \eqref{eq:v7-product-matrix} would then force both entries $01$ and
$12$ to vanish. Each is strictly negative by
\eqref{eq:v7-metric-intervals}. This is a contradiction. The statement is
coordinate independent because both the Veronese identification and the
symmetric-square metric transform under the same change of basis of $F$.
\end{proof}

\subsection{A phase-independent separation directly in the source Grams}
The next estimate does not require a jet trivialization or even the exact
orthogonality of an approximating source.

\begin{lemma}[Three-ray inequality for every symmetric-square source]
\label{lem:v7-product-inequality}
Let $a,b,c$ be unit vectors in a two-dimensional positive Hermitian space,
and let $\widehat G$ be the normalized Gram of $a^{\otimes2},b^{\otimes2},
c^{\otimes2}$, allowing independent ray phases. Then
\begin{equation}
 \boxed{|\widehat G_{ac}|+|\widehat G_{bc}|
      \le 1+\sqrt{|\widehat G_{ab}|}.}
 \label{eq:v7-product-inequality}
\end{equation}
If $\widehat G_{ab}=0$, the left side equals one.
\end{lemma}
\begin{proof}
The left side is $|\langle a,c\rangle|^2+|\langle b,c\rangle|^2$.
The positive operator $aa^*+bb^*$ has eigenvalues
$1\pm|\langle a,b\rangle|$. Taking its expectation on $c$ proves the
inequality, since $|\widehat G_{ab}|=|\langle a,b\rangle|^2$.
When $a,b$ are orthogonal they form an orthonormal basis of the two-dimensional
space, so the sum is exactly one. Ray phases do not change any modulus.
\end{proof}

\begin{theorem}[A strict source-Gram gap to every positive product metric]
\label{thm:v7-Gram-gap}
For the actual three rays with labels $u,v,0$, put
\begin{equation}
 \delta=|G_{u0}|+|G_{v0}|-1
       =\sqrt{\mu_{0u}}+\sqrt{\mu_{0v}}-1.
 \label{eq:v7-delta}
\end{equation}
Here $\mu$ is V3's branch-pair squared-moment packet, not a new measure.
One has $G_{uv}=0$ and the certified bounds
\begin{equation}
 0.067622858309<\delta<0.067622858310.
 \label{eq:v7-delta-interval}
\end{equation}
For every normalized three-ray Gram $\widehat G$ arising from squares of
vectors in any positive two-dimensional Hermitian space,
\begin{equation}
 \boxed{
 \max_{ij\in\{u0,v0,uv\}}
  \bigl||G_{ij}|-|\widehat G_{ij}|\bigr|
 \ge \left(\frac{\sqrt{1+8\delta}-1}{4}\right)^2
 >0.00364.}
 \label{eq:v7-gap}
\end{equation}
The bound survives arbitrary ray rephasing and a common unitary comparison.
It is a certified lower bound, not a claim that the optimum is attained.
\end{theorem}
\begin{proof}
V3 supplies the exact zero and the algebraic values in
\eqref{eq:v7-delta}. The same outward rational square-root procedure used
above gives \eqref{eq:v7-delta-interval}. Let $\varepsilon$ denote the
maximum in \eqref{eq:v7-gap}. Then
$|\widehat G_{uv}|\le\varepsilon$ and
\[
 1+\delta-2\varepsilon
 \le |\widehat G_{u0}|+|\widehat G_{v0}|
 \le 1+\sqrt\varepsilon.
\]
Thus $\delta\le2\varepsilon+\sqrt\varepsilon$.
The increasing quadratic function $2x^2+x$ on $x\ge0$ gives the displayed
bound after setting $x=\sqrt\varepsilon$. Rational interval evaluation
places its value strictly between $0.003641015451$ and $0.003641015452$.
Moduli are invariant under the stated phase and unitary changes.
\end{proof}

\begin{remark}[Why this is a different obstruction from V1]
V1 excludes ordinary pointed-$S_4$ covariance. The present statement excludes
an \emph{independently natural product metric} on the correctly identified
quadratic Hodge carrier, even without demanding any tetrahedral symmetry.
The algebraic second-jet bridge remains an isomorphism; what fails is the
replacement of its actual correlated metric by an unverified product metric.
\end{remark}

\section{An explicit Gauss--Manin connection through the arithmetic curve}
\label{sec:v7-GM}

\subsection{The declared family and its cohomological basis}
Let
\begin{equation}
 B=\operatorname{Spec}\Q[t,(t\Delta(t))^{-1}],\qquad
 \Delta(t)=256t^3+27,\qquad f_t(x)=x^5-x^2-tx.
 \label{eq:v7-base}
\end{equation}
The discriminant is $-t^2\Delta(t)$, so the smooth projective completion of
$y^2=f_t(x)$ is a genus-two family over $B$, with the marked infinity
section and the additional branch section at zero. Its fibre at $t=1$ is the
fixed arithmetic curve. Put
\[
 E=H^1_{\mathrm{dR}}(C/B),\qquad
 e_i=[x^i dx/y]\quad(0\le i\le3),\qquad
 F=\Span(e_0,e_1).
\]
The notation $e_i$ in this section denotes de Rham classes, not the Clifford
or monomial vectors of earlier sections.

\begin{lemma}[A complete second-kind basis]
\label{lem:v7-basis}
The four classes $e_0,e_1,e_2,e_3$ are a basis over $\Q(t)$. The Hodge
subspace is $F=H^0(C,\Omega^1)=\Span(e_0,e_1)$.
\end{lemma}
\begin{proof}
The four forms are regular away from infinity. With $x=z^{-2}$ their
expansions there are
\[
 \frac{x^i dx}{y}
   =-2z^{2-2i}(1-z^6-tz^8)^{-1/2}\,dz.
\]
There is no residue, so all four define compact de Rham classes of the second
kind. The first two are holomorphic. If a linear combination of the four were
an exact rational differential $dh$, then $h$ could have no finite pole:
in characteristic zero differentiation strictly raises the order of a pole.
Taking the anti-invariant part under the hyperelliptic involution, one may
write $h=yR(x)$ with $R$ a polynomial, because the affine coordinate ring is
$\Q(t)[x,y]/(y^2-f_t)$. If $R$ has degree $m\ge0$, the numerator of
\[
 d(yR)=\bigl(f_tR'+\tfrac12f_t'R\bigr)\,dx/y
\]
has degree $m+4$ and nonzero leading coefficient $m+5/2$. It cannot be a
nonzero polynomial of degree at most three. The four classes are independent.
Compact de Rham cohomology has dimension $2g=4$, proving completeness;
its holomorphic subspace has dimension two. The identification of second-kind
forms modulo exact forms with de Rham cohomology is the classical curve
comparison used here \cite{KTdeRham}.
\end{proof}

The Gauss--Manin connection differentiates cohomology classes in the parameter
\cite{KatzOda}. Use column conventions:
\[
 \nabla_{\partial_t}e_j=\sum_i e_i A_{ij}(t),\qquad
 \nabla_{\partial_t}\Bigl(\sum_i e_i v_i(t)\Bigr)
     =\sum_i e_i\bigl(v_i'+(Av)_i\bigr).
\]

\begin{theorem}[Exact rational connection of the marked family]
\label{thm:v7-GM-matrix}
In this basis the actual connection matrix is
\begin{equation}
 \boxed{
 A(t)=\frac1{2t\Delta(t)}
 \begin{pmatrix}
 -192t^3&-27t&36t^2&-48t^3\\
 48t^2&-64t^3&-9t&12t^2\\
 36t&-48t^2&64t^3&9t\\
 -81&108t&-144t^2&192t^3
 \end{pmatrix}.}
 \label{eq:v7-A}
\end{equation}
No period matrix or desired transfer eigenvalue enters this formula.
\end{theorem}
\begin{proof}
Differentiate at fixed $x$. Since $\partial_t f_t=-x$,
\[
 \partial_t(x^j dx/y)=\frac{x^{j+1}}{2y^3}\,dx.
\]
Set $R_j(x)=\sum_{k=1}^4 P_{kj}x^k$, with columns indexed by $j=0,1,2,3$
and
\begin{equation}
 P(t)=\frac1{t\Delta(t)}
 \begin{pmatrix}
 -(64t^3+27)&27t&-36t^2&48t^3\\
 48t^2&-64t^3&-9t&12t^2\\
 -36t&48t^2&-64t^3&-9t\\
 27&-36t&48t^2&-64t^3
 \end{pmatrix}.
 \label{eq:v7-primitives}
\end{equation}
Direct polynomial multiplication gives, for each of the four columns,
\begin{equation}
 x^{j+1}
 =2f_t\sum_{i=0}^3A_{ij}x^i+2f_tR_j'-f_t'R_j.
 \label{eq:v7-reduction}
\end{equation}
Division by $2y^3$ turns the final two terms into $d(R_j/y)$.
Thus the differentiated form is cohomologous to the column specified in
\eqref{eq:v7-A}. Each rational identity can be checked after multiplication
by $2t\Delta$; no infinite series or approximate reduction is used.
The exact-differential quotient is compatible with differentiation of the
smooth family, giving the claimed Gauss--Manin connection.
\end{proof}

\subsection{Pairing and the first failure of three-dimensional closure}
Define the alternating matrix
\begin{equation}
 \Omega=
 \begin{pmatrix}
 0&0&0&4/3\\
 0&0&4&0\\
 0&-4&0&0\\
 -4/3&0&0&0
 \end{pmatrix}.
 \label{eq:v7-Omega}
\end{equation}
It is the de Rham pairing in the convention
$\langle\alpha,\beta\rangle=\operatorname{Res}_\infty((\int\alpha)\beta)$
for these second-kind representatives; the overall sign is conventional.
The above expansions give $4/3$ for $(e_0,e_3)$, $4$ for $(e_1,e_2)$,
and zero for the other upper-triangular entries. In particular the matrix
is constant in $t$ and has determinant $256/9$.

\begin{proposition}[Symplecticity and a rank-one Hodge variation]
\label{prop:v7-Higgs}
One has
\begin{equation}
 \boxed{A^{\mathsf T}\Omega+\Omega A=0.}
 \label{eq:v7-symplectic}
\end{equation}
The second fundamental form $F\to E/F$ in the direction $\partial_t$ is
\begin{equation}
 \boxed{B(t)=\frac1{\Delta(t)}
 \begin{pmatrix}18&-24t\\-81/(2t)&54\end{pmatrix}.}
 \label{eq:v7-B}
\end{equation}
It has rank one everywhere on $B$, with kernel
$\Span(4t e_0+3e_1)$. The corresponding conormal abscissa $-4t/3$ is not a
branch point on any smooth fibre. In particular $F$ and
$\operatorname{Sym}^2F$ are not connection-stable.
\end{proposition}
\begin{proof}
Matrix multiplication proves \eqref{eq:v7-symplectic}; it also verifies
flatness of the stated constant pairing directly. The bottom-left block of
$A$ is \eqref{eq:v7-B}. Its second column is $-4t/3$ times its first,
which is nonzero. This proves the rank and kernel assertions.
Since $e_1-r e_0$ has abscissa $r$, the kernel line has
$r=-4t/3$. Direct substitution gives
\[
 r^4-r-t=\frac{t\Delta(t)}{81}\ne0.
\]
Finally the induced map on $\operatorname{Sym}^2F$ sends $vw$ modulo
$\operatorname{Sym}^2F$ to $v\otimes B(w)+w\otimes B(v)$. A nonzero
rank-one $B$ on a two-dimensional $F$ makes this map have rank two: in a
basis $a,b$ with $B(a)=0$ and $B(b)\ne0$, its values on $ab$ and $b^2$
are independent, whereas its value on $a^2$ is zero. Thus the initial rank-three
quadratic Hodge carrier grows to rank five at first derivative order.
\end{proof}

\section{The actual quadratic source has a minimal rank-ten completion}
\label{sec:v7-rank10}

\subsection{The retained theta line and the induced connection}
V6's familywise jet construction gives
\[
 V_0\otimes\lambda^{-2}\simeq\operatorname{Sym}^2F
       \subset W:=\operatorname{Sym}^2E.
\]
This is the de-twisted comparison. No connection on the theta-fibre line
$\lambda^2$ is silently postulated. Over a local trivialization its original
source is a nonvanishing scalar times $e_0^2$. Multiplying a source by a
nonzero scalar function does not alter its differential cyclic span: the
successive derivative columns change by an invertible triangular matrix.
Adding a scalar connection on the fibre line has the same property.
Thus the following rank statement is intrinsic to the projective source even
though an absolute fibre-line connection remains additional data.

Order the monomial basis of $W$ as
\begin{equation}
 (00,01,02,03,11,12,13,22,23,33),\qquad ij=e_i e_j.
 \label{eq:v7-W-basis}
\end{equation}
Let $M(t)$ be the infinitesimal symmetric-square matrix of $A(t)$, defined
\emph{without} a new choice by
\begin{equation}
 M(t)(e_i e_j)=\sum_k A_{ki}(t)e_ke_j+
                         \sum_k A_{kj}(t)e_ie_k.
 \label{eq:v7-M}
\end{equation}
A section with coefficient column $v$ has derivative $v'+Mv$.

\begin{theorem}[An exact cyclic certificate for the original theta jet]
\label{thm:v7-rank10}
Put $v_0=(1,0,\ldots,0)^{\mathsf T}$ and recursively
\begin{equation}
 v_{n+1}(t)=v_n'(t)+M(t)v_n(t).
 \label{eq:v7-cyclic-recursion}
\end{equation}
Then $v_0(1),\ldots,v_9(1)$ are linearly independent over $\Q$.
Consequently the smallest Gauss--Manin-stable differential submodule over
$\Q(t)$ containing the actual quadratic source line is all of $W$, of rank
ten. The same conclusion holds on a germ around $t=1$ and for any nonzero
local scalar representative of that source line.
\end{theorem}
\begin{proof}
We give a finite modular certificate for the rational determinant. All entries
of the ten columns obtained from \eqref{eq:v7-cyclic-recursion} at $t=1$
have denominators divisible only by $2$ and $283$. They therefore reduce
modulo $101$. In the basis \eqref{eq:v7-W-basis}, their matrix is
\begin{equation}
\left(\begin{array}{rrrrrrrrrr}
1&50&95&98&98&34&55&33&90&37\\
0&38&25&62&38&34&36&95&50&35\\
0&79&38&52&12&81&60&61&36&49\\
0&100&54&65&48&66&41&52&44&67\\
0&0&15&84&45&45&24&3&54&58\\
0&0&73&87&8&15&71&57&77&45\\
0&0&63&4&32&28&78&46&89&67\\
0&0&40&18&41&38&98&74&70&75\\
0&0&22&12&69&73&27&94&87&65\\
0&0&51&46&11&83&100&93&51&97
\end{array}\right)\pmod{101}.
\label{eq:v7-cyclic-matrix}
\end{equation}
Its determinant is $8\pmod{101}$. For a fully explicit recurrence checking
the displayed entries, set $t=1+u$, $d(t)=t\Delta(t)$, write
$d(1+u)=\sum_{j=0}^4d_ju^j$ and
$d(1+u)M(1+u)=\sum_{j\ge0}N_ju^j$. Then $d_0=283$ and
\[
 M_n=\frac{N_n-\sum_{j=1}^{\min(n,4)}d_jM_{n-j}}{283}
\]
are the Taylor coefficients of $M$. If $v_k=\sum_r c_{k,r}u^r$, the
coefficient rule is
\begin{equation}
 c_{k+1,r}=(r+1)c_{k,r+1}+\sum_{j=0}^rM_j c_{k,r-j},
 \quad c_{0,0}=v_0,\quad c_{0,r}=0\ (r>0).
 \label{eq:v7-Taylor-recursion}
\end{equation}
Computing only $k+r\le9$ gives \eqref{eq:v7-cyclic-matrix}. This is a
finite computation in $\F_{101}$ derived from the stated rational connection;
the companion verifier independently performs it first over $\Q$ and then
reduces. A rational number with a nonzero reduction cannot be zero, so the
rational determinant is nonzero.

Every connection-stable submodule containing $e_0^2$ contains all its
successive derivatives. Their determinant is a rational function nonzero at
$t=1$, hence nonzero over $\Q(t)$. The ambient module has rank ten, proving
minimality. Nonvanishing persists locally near the smooth point $1$.
For a nonzero scalar function $a$, the leading term of $\nabla^n(as)$ is
$a\nabla^n s$; the remaining terms involve lower derivatives. The resulting
triangular transformation is invertible, proving the source-line statement.
\end{proof}

\begin{proposition}[The full quadratic head reaches rank ten at depth four]
\label{prop:v7-joint-ranks}
Starting with $S=(e_0^2,e_0e_1,e_1^2)$, the ranks at $t=1$ of the complete
banks $(S,\nabla S,\ldots,\nabla^kS)$ for $k=0,1,2,3,4$ are respectively
\begin{equation}
 \boxed{3,\quad5,\quad7,\quad9,\quad10.}
 \label{eq:v7-joint-ranks}
\end{equation}
These are differential-span ranks, not finite-state transfer powers or
additional physical source occurrences.
\end{proposition}
\begin{proof}
Apply \eqref{eq:v7-Taylor-recursion} to the three source columns. An exact
rational certificate is given by the following independent-column sets, with
columns counted from zero in the concatenated bank:
\[
\begin{array}{c|l}
k&\text{independent columns}\\\hline
0&0,1,2\\
1&0,1,2,3,4\\
2&0,1,2,3,4,6,7\\
3&0,1,2,3,4,6,7,9,10\\
4&0,1,2,3,4,6,7,9,10,12
\end{array}
\]
Independent row sets are, in order,
$(0,1,4)$, $(0,1,2,4,5)$, $(0,1,2,3,4,5,6)$,
$(0,1,2,3,4,5,6,7,8)$, and all ten rows. The resulting square minors are
\[
\begin{gathered}
 1,\quad -\frac{648}{80089},\quad
 -\frac{13286025}{102627966736},\\[0.4em]
 -\frac{19613162255625}{1316562444545690311712},\quad
 -\frac{9651146816936671875}{13496597711516133295961903104}.
\end{gathered}
\]
For the upper rank, let $J_k$ be the full bank and let $I_k,R_k$ be the stated
column and row sets. Rational substitution into the same recurrence gives
\[
 J_k=J_k[:,I_k]\,
          (J_k[R_k,I_k])^{-1}J_k[R_k,:].
\]
This identity specifies and verifies every dependent-column coefficient; the
certificate records the complete coefficient matrices. Thus the minors certify
the lower ranks and these exact spanning identities certify the upper ranks.
The first increment also follows intrinsically from
\cref{prop:v7-Higgs}.
\end{proof}

\subsection{The completed object is a symplectic adjoint source}
\begin{proposition}[An intrinsic even source in the symplectic action algebra]
\label{prop:v7-adjoint}
The map
\begin{equation}
 \boxed{
 \Phi(vw)(x)=\Omega(w,x)v+\Omega(v,x)w
 }
 \label{eq:v7-adjoint}
\end{equation}
identifies $\operatorname{Sym}^2E$ with $\mathfrak{sp}(E,\Omega)$ and
intertwines their induced connections. The source $e_0^2$ becomes the
rank-one square-zero endomorphism $2e_0e_0^{\mathsf T}\Omega$.
Its differential span is therefore the entire rank-ten symplectic adjoint
bundle near $t=1$.
\end{proposition}
\begin{proof}
In the fixed basis, \eqref{eq:v7-adjoint} is
$(vw^{\mathsf T}+wv^{\mathsf T})\Omega$. Symmetric matrices have dimension
ten, and multiplication by the invertible $\Omega$ is injective. For a
symmetric matrix $B$, the matrix $B\Omega$ satisfies
$(B\Omega)^{\mathsf T}\Omega+\Omega(B\Omega)=0$; hence the map has image
in $\mathfrak{sp}_4$, also of dimension ten, and is an isomorphism.
The constancy and connection compatibility of $\Omega$ give
\[
 [A,\Phi(vw)]=\Phi((Av)w+v(Aw)).
\]
Together with differentiation of coefficients this is connection
intertwining. For a square, the endomorphism has rank one and its square
vanishes because $\Omega(v,v)=0$. The final assertion follows from
\cref{thm:v7-rank10}.
\end{proof}

\begin{scope}[Precisely what has been completed]
The even theta jet now has a genuine characteristic-zero differential
realization in a rank-ten arithmetic adjoint object. This does not contradict
V6's prohibition on obtaining a nonabelian rank-two or rank-three Tate
subquotient. It does not identify this one-parameter family's geometric
monodromy group with the arithmetic Galois monodromy at its rational fibre.
The rank-ten result is proved by the explicit source derivatives, without any
such identification. There is still no claimed source isometry with a renewal
carrier, and no distinguished physical flow is attached to $\partial_t$.
\end{scope}

\section{Local monodromy and the limits of positive-metric transport}
\label{sec:v7-monodromy}

\subsection{A convergent local normal form}
At the deleted point $t=0$, the rational matrix has the form
\begin{equation}
 A(t)=\frac Nt+B(t),\qquad
 \boxed{N=-\frac32 E_{30},\quad N^2=0,\quad\rank N=1,}
 \label{eq:v7-residue}
\end{equation}
with $B$ holomorphic near zero. Here $E_{30}$ is the four-by-four matrix
with its only nonzero entry in row $3$, column $0$.

\begin{lemma}[Holomorphic gauge for the displayed nilpotent residue]
\label{lem:v7-local-gauge}
There is an invertible holomorphic $H(t)$ near zero with $H(0)=I$ and
\begin{equation}
 H^{-1}AH+H^{-1}H'=N/t.
 \label{eq:v7-gauge}
\end{equation}
\end{lemma}
\begin{proof}
Write $B(t)=\sum_{j\ge0}B_jt^j$ and $H(t)=I+\sum_{n\ge1}H_nt^n$.
The gauge equation is equivalent to
\begin{equation}
 (nI+\operatorname{ad}_N)H_n
   =-\sum_{j=0}^{n-1}B_jH_{n-1-j}.
 \label{eq:v7-gauge-recurrence}
\end{equation}
Since $N^2=0$, one has $\operatorname{ad}_N^3=0$ and
\[
 (nI+\operatorname{ad}_N)^{-1}
 =n^{-1}I-n^{-2}\operatorname{ad}_N+n^{-3}\operatorname{ad}_N^2.
\]
This proves unique formal solvability. In any submultiplicative matrix norm,
its norm is at most $C/n$ for $n\ge1$, with
$C=1+2\norm N+4(\norm N)^2$.
Choose $r>0$ and $M$ with $\norm{B_j}\le Mr^{-j}$. The scalar recurrence
\[
 nh_n=C\sum_{j=0}^{n-1}Mr^{-j}h_{n-1-j},\qquad h_0\ge\norm I,
\]
majorizes all matrix coefficients. Its generating function solves
$h'=CM(1-t/r)^{-1}h$, so it is analytic for $|t|<r$. Hence the matrix
series converges. Continuity of its determinant and $H(0)=I$ give local
invertibility, and the convergent identity proves \eqref{eq:v7-gauge}.
\end{proof}

\begin{theorem}[Actual nonsemisimple local transport]
\label{thm:v7-monodromy}
The horizontal equation $v'+A(t)v=0$ has local monodromy conjugate to
\begin{equation}
 \boxed{T_0=\exp(-2\pi iN)=I-2\pi iN\ne I.}
 \label{eq:v7-T0}
\end{equation}
It has one nontrivial Jordan block of size two. On the completed quadratic
carrier the nilpotent logarithm $N^{[2]}$ has ranks
\begin{equation}
 \boxed{\rank N^{[2]}=4,\qquad
        \rank (N^{[2]})^2=1,\qquad (N^{[2]})^3=0.}
 \label{eq:v7-N2-ranks}
\end{equation}
The quadratic local monodromy has Jordan block sizes
$3,2,2,1,1,1$. Neither local system admits a positive parallel Hermitian
metric around that loop, even after a scalar line twist.
\end{theorem}
\begin{proof}
By \cref{lem:v7-local-gauge}, horizontal solutions are
$H(t)t^{-N}c$. Analytic continuation once around zero changes $\log t$ by
$2\pi i$, proving \eqref{eq:v7-T0} up to the base-point conjugation. Since
$N\ne0$ and $N^2=0$, this is a nonidentity unipotent with a size-two block.

Choose $a,b,c,d$ with $Na=b$ and $Nb=Nc=Nd=0$, absorbing the nonzero
coefficient into $b$. On the symmetric square,
\[
 a^2\mapsto2ab\mapsto2b^2\mapsto0,\quad
 ac\mapsto bc\mapsto0,\quad ad\mapsto bd\mapsto0,
\]
while $c^2,cd,d^2$ are killed. This gives precisely the three nontrivial
chains and three trivial chains, proving the ranks and block sizes.
Exponentiation of a nilpotent preserves its Jordan block sizes because
$\exp(zN)-I=N(zI+\cdots)$ with invertible second factor when $z\ne0$.

If a positive Hermitian form $G$ were preserved by such a monodromy $T$, then
$G^{1/2}TG^{-1/2}$ would be unitary, hence diagonalizable over $\C$.
The nontrivial Jordan block contradicts this. Multiplication of $T$ by any
nonzero scalar leaves its nontrivial block sizes unchanged, so a scalar
line twist does not remove the obstruction. The loop can be joined to the
actual fibre $t=1$ inside the connected smooth complex base, which only
conjugates monodromy and preserves the conclusion.
\end{proof}

\subsection{What can still be done along a path}
\begin{proposition}[Positive transport locally, but not around the arithmetic loop]
\label{prop:v7-path-metric}
On a simply connected path domain, let $P(t)$ be the fundamental horizontal
solution normalized at a base point and let $G_*>0$ be an independently
specified positive metric there. Then
\begin{equation}
 \boxed{G(t)=P(t)^{-*}G_*P(t)^{-1}}
 \label{eq:v7-parallel-metric}
\end{equation}
is the unique parallel positive metric with that initial value. It descends
around a loop exactly when that loop's monodromy satisfies
$T^*G_*T=G_*$. Consequently no such metric on $E$ or $W$ descends around
$t=0$ for this family. Choosing a pathwise metric is not a proof of a
fixed-metric unitary arithmetic--Wilson equivalence.
\end{proposition}
\begin{proof}
By construction the pairing of horizontal vectors $P(t)x,P(t)y$ is
$x^*G_*y$, independently of $t$, which proves parallelism and positivity.
Every parallel metric has this value on all horizontal vectors and therefore
has the displayed matrix, proving uniqueness. After a loop, $P$ changes to
$PT$; single-valuedness of the matrix is equivalent to $T^*G_*T=G_*$.
\Cref{thm:v7-monodromy} excludes this condition for every $G_*>0$.
\end{proof}

\begin{remark}[The independent Hodge metric has not been replaced]
The positive Hodge metric is not asserted to be parallel for Gauss--Manin.
The result therefore does not prohibit its existence. Nor does it prohibit a
nonflat unitary connection, a varying metric with extra differential data, or
a noncompact/indefinite comparison category. What is excluded is the automatic
identification of the actual flat arithmetic deformation transport with a
compact unitary Wilson local system. At the base point, the actual theta
metric is already non-product by \cref{thm:v7-nonproduct}; away from that
point, a rank-ten positive extension and its relation to any Hodge or renewal
metric remain independent comparison data.
\end{remark}

\section{V7 verification, proved scope and the revised comparison}
\label{sec:v7-status}

\subsection{Finite verification and its exact role}
The standalone script
\path{Arithmetic_Loading_Duality_Research_Notes_V7_checks.py} reuses V3's
rational root and interval primitives, without a new floating-point root or
period estimate. It certifies the nonzero frame coordinates, reconstructs the
three-by-three metric with outward intervals, and verifies the strict product
separation. The exact reality and zero identities are proved geometrically
above; an interval containing zero is not used as an equality proof.

Independently, the script verifies the four polynomial differential identities,
the discriminant, constant alternating form, rank-one Hodge variation,
symmetric-square connection, and the ten-column determinant certificate. It
checks both lower minors and full rational upper spanning identities for the
joint derivative banks. It verifies the residue, induced nilpotent ranks and
the symmetric-square/adjoint intertwining on every basis element. The
rational Taylor recurrence and the displayed modulo-$101$ matrix make the
principal cyclicity proof reproducible without a large opaque symbolic
expression. Neither the script nor the notes claim a numerical period
matrix, a crystalline Frobenius evaluation or a proof-assistant
formalization. The general de Rham interpretation is a separately cited
classical input.

\subsection{Proved in this continuation}
\begin{enumerate}[label=\textup{(P7.\arabic*)},leftmargin=3.6em]
\item The actual conformal second-jet metric is recovered algebraically from
      four marked source rays, their overlaps and the branch coordinates.
      It has explicit rational interval bounds and is invariant under the
      residual source-frame phases.
\item The metric is not a symmetric square of any positive one-form metric.
      A three-ray, phase-independent Gram witness gives a strictly positive
      separation greater than $0.00364$ from all such product-source Grams.
\item The declared deformation through $C_{\mathrm{tet}}$ has the explicit
      rational Gauss--Manin matrix \eqref{eq:v7-A}. Its Hodge variation has
      rank one with an explicitly identified nonbranch kernel line.
\item The original quadratic source generates all ten dimensions of
      $\operatorname{Sym}^2H^1_{\mathrm{dR}}$ under this connection. The
      completed source is naturally a rank-one nilpotent in the symplectic
      adjoint bundle whose derivatives span that bundle.
\item Actual local transport at $t=0$ has nontrivial unipotent Jordan blocks.
      It cannot preserve any positive parallel Hermitian metric, including
      after a scalar line twist. Pathwise positive transport is possible
      but cannot be silently promoted to a single-valued global metric.
\end{enumerate}

\subsection{Inherited inputs and unresolved data}
The companions and both supplied monographs remain antecedents
\cite{GT77,Arithmetic,Spacetime,Spectral,Duality,Variational}. The source
orbit, arithmetic frame and real descent, the finite-prime Frobenius and
Cartier operators, and the second-jet isomorphism are exactly the results of
V1--V6. No earlier failure is renamed a positive physical result.

The most economical next mixed comparison now retains
\[
 \boxed{
 (C_t,\lambda^2,\mathfrak j_2,[H_\Theta],
       E,\nabla,\operatorname{Sym}^2E,[e_0^2];
       \text{central descent and independent renewal maps}).}
\]
The algebraic family and connection are explicit. They still require an
independent Betti or crystalline comparison to become the corresponding
marked arithmetic realization, and the actual Hodge metric has not been
numerically evaluated. A positive extension of the source metric from three
to ten dimensions has not been selected from arithmetic occurrence data.
The deformation is not uniquely forced by the isolated fibre and does not
supply a physical clock. Its holonomy is complex symplectic rather than
fixed-metric unitary.

The original four relative arithmetic--renewal frame signs, the central
Store/chronology signature scalar, the mixed physical source and transfer
panels, and the finite-prime Frobenius--Wilson tensor/interface comparison
remain unresolved. The new rank ten is a differential completion count, not a
particle multiplicity. No Standard-Model numerical prediction, Einstein limit,
GRH estimate, or quantum-field construction is inferred.

\paragraph{Version integrity.}
The complete V1--V6 mathematical bodies, including every historical status
ledger, are retained byte for byte. The present work is appended after V6;
future corrections must likewise be stated in a new continuation rather than
silently changing an earlier body.
% ====================== END IMMUTABLE V7 MATHEMATICAL BODY ======================
% ===================== BEGIN IMMUTABLE V8 MATHEMATICAL BODY =====================
\clearpage
\section{V8 continuation: marked periods, arithmetic monodromy, and a positive completion}
\label{sec:v8-context}

V7 left the independently marked Betti comparison and the actual Hodge metric
unevaluated. Those are genuine upstream data: choosing a positive metric by
parallel transport would neither identify the arithmetic metric nor resolve
its monodromy obstruction. This continuation supplies a marked comparison
for the actual curve, without changing its theta source or fitting a renewal
operator.

There are three concrete advances. The scalar equation generated by the
actual differential is an explicit hypergeometric equation, whose monodromy
is an arithmetic symplectic group. A smooth auxiliary presentation with a
beta-integral reference fibre then gives all sixteen entries of a marked
period matrix with a convergent rational recurrence and rigorous error
bounds. This acquires the normalized genus-two period matrix, the actual
Hodge metric, and all four theta-source frame signs in an independently
specified arithmetic marking. Finally, these data give an explicit positive
rank-ten completion under a stated minimum-change rule. That rule is a
mathematical completion convention, not a physical occurrence theorem.

\paragraph{Dependency and novelty ledger.}
The V7 Gauss--Manin matrix, V6 second-jet map, V7 conformal theta metric,
V3--V4 overlap magnitudes and phase orbit, and V2 theta acquisition formula
are inherited. Their proofs are not repeated. The new calculations are the
scalar cyclic relation, the nonsingular reference-fibre period transport,
the integral cycle/characteristic marking, the certified actual periods and
theta signs, and the resulting Hodge-based completion. Classical de
Rham--Betti comparison and Riemann bilinear relations are inputs
\cite{KatzOda,KTdeRham,MilneAV}; Euler's beta integral is used in its elementary
convergent form \cite{DLMFBeta}. The Levelt description and the
Singh--Venkataramana arithmeticity theorem are explicitly cited inputs
\cite{BeukersHeckman,SinghVenkataramana}, not new general theorems.

No finite-prime crystalline matrix, renewal source, physical clock, or
central Store/chronology signature is supplied by this calculation. Betti
conjugation below is also distinguished from V4's antiunitary real descent on
theta sections. All statements are relative to the explicitly marked
families and paths declared before the target is evaluated.

\section{The actual scalar equation and its arithmetic monodromy}
\label{sec:v8-hypergeometric}

\subsection{Scalarization of the inherited connection}
Retain $C_t:y^2=x^5-x^2-tx$ and the V7 basis $e_j=[x^jdx/y]$. Set
$\Delta=256t^3+27$, $\Theta=t\nabla_{\partial_t}$ on sections, and
$\theta_t=t\,d/dt$ on scalar period functions. A flat cycle $\delta_t$
has period $p(t)=\int_{\delta_t}e_0$.

\begin{theorem}[An exact hypergeometric Picard--Fuchs realization]
\label{thm:v8-PF}
Every period of $e_0$ satisfies
\begin{equation}
 \boxed{\left[
 \theta_t^2(\theta_t-1)(\theta_t-2)
 +\frac{256t^3}{27}\!\!\prod_{k\in\{3,9,15,21\}}
 (\theta_t+\tfrac{k}8)\right]p=0.}
 \label{eq:v8-PFt}
\end{equation}
The four sections $e_0,\Theta e_0,\Theta^2e_0,\Theta^3e_0$ form a basis
over $\Q(t)$. Consequently this is a realization of the full rank-four
period system, not an equation for a proper visible quotient.

With $z=-256t^3/27$ and $\theta=z\,d/dz$, the descended scalar equation is
\begin{equation}
 \boxed{\left[
 \theta^2(\theta-\tfrac13)(\theta-\tfrac23)
 -z(\theta+\tfrac18)(\theta+\tfrac38)
 (\theta+\tfrac58)(\theta+\tfrac78)
 \right]p=0.}
 \label{eq:v8-PFz}
\end{equation}
Its hypergeometric parameters are
$\alpha=(1/8,3/8,5/8,7/8)$ and $\beta=(1,1,1/3,2/3)$, where the first
unit lower parameter is the implicit factorial parameter.
\end{theorem}
\begin{proof}
In column coordinates put $v_0=(1,0,0,0)^{\mathsf T}$ and
$v_{n+1}=t(v_n'+A(t)v_n)$, with the already proved V7 matrix $A$.
Rational multiplication gives
\begin{equation}
 \det(v_0,v_1,v_2,v_3)
   =\frac{995085t^3}{64\Delta^3}\ne0.
 \label{eq:v8-scalar-det}
\end{equation}
For an independently reproducible scalar certificate, the relation
$v_4=c_0v_0+c_1v_1+c_2v_2+c_3v_3$ has coefficients
\[
 \begin{split}
 c_0&=-\frac{8505t^3}{16\Delta},\\
 c_1&=-\frac{2376t^3}{\Delta},\\
 c_2&=-\frac{18(172t^3+3)}{\Delta},\\
 c_3&=-\frac{3(8t-3)(64t^2+24t+9)}{\Delta}.
 \end{split}
\]
Substitution of the recurrence and multiplication by a common power of
$2t\Delta$ verify each coordinate as a polynomial identity. Expanding
\eqref{eq:v8-PFt} gives precisely this relation. Differentiation of a period
over a flat cycle pairs the cycle with the covariant derivative, so the same
identity holds for $p$. Conversely, the nonzero determinant says that a
cycle functional is determined by these four scalar jets. Thus the scalar
system loses no de Rham direction. Finally $\theta_t=3\theta$; division
by $3^4$ and substitution of $z$ give \eqref{eq:v8-PFz}. This is the
standard generalized hypergeometric operator \cite[\S16.8]{DLMFHyper}.
\end{proof}

\begin{corollary}[A geometrically normalized scalar solution]
\label{cor:v8-real-period}
For real $t>0$, let $-a(t)$ be the negative quartic root, so
$a+a^4=t$, and let $\delta_-(t)$ run from $-a$ to zero on the positive
real square-root sheet and back on the other sheet. Then
\begin{equation}
 \boxed{\int_{\delta_-(t)}\frac{dx}{y}
 =2\pi\,{}_4F_3\!\left(
 \begin{matrix}1/8,3/8,5/8,7/8\\1,1/3,2/3\end{matrix};
 -256t^3/27\right).}
 \label{eq:v8-real-hypergeom}
\end{equation}
The right side means continuation from zero along the negative real $z$-axis.
Its normalization and branch are fixed by this cycle, not by a desired
transfer value.
\end{corollary}
\begin{proof}
Substitute $x=-au$ and use $t=a+a^4$. The period becomes
\[
 2\int_0^1\frac{du}
 {\sqrt{u(1-u)}\sqrt{1+a^3(1+u+u^2+u^3)}}.
\]
It tends to $2\pi$ as $t\downarrow0$. Write $a=t b(t^3)$, where the
implicit equation $b+t^3b^4=1$ gives $b$ analytic near zero. The integral
is therefore analytic in $t^3$ there, by uniform domination of its
binomial expansion. The recurrence of \eqref{eq:v8-PFz} determines a
unique solution analytic in $z$ with value one at zero: its coefficients
are the displayed hypergeometric coefficients, since for every positive
integer $n$ the coefficient $n^2(n-1/3)(n-2/3)$ is nonzero.
This proves the equality near zero. Continuation along $t>0$ encounters
no singular fibre, and uniqueness of continuation proves the identity.
\end{proof}

\subsection{An actual arithmetic topological monodromy group}
Define
\begin{equation}
 A_L=\begin{pmatrix}0&0&0&-1\\1&0&0&0\\0&1&0&0\\0&0&1&0\end{pmatrix},
 \qquad
 B_L=\begin{pmatrix}0&0&0&-1\\1&0&0&1\\0&1&0&0\\0&0&1&1\end{pmatrix}.
 \label{eq:v8-Levelt}
\end{equation}
These are the companion matrices of
$f(X)=X^4+1$ and $g(X)=X^4-X^3-X+1=(X-1)^2(X^2+X+1)$.

\begin{theorem}[Arithmetic monodromy of the scalar comparison object]
\label{thm:v8-arithmetic-monodromy}
The monodromy group of \eqref{eq:v8-PFz} is, up to simultaneous complex
conjugacy, the group $\Gamma_L=\langle A_L,B_L\rangle$. It preserves
\begin{equation}
 \Omega_L=\begin{pmatrix}
 0&1&0&1\\-1&0&1&0\\0&-1&0&1\\-1&0&-1&0
 \end{pmatrix},\qquad \det\Omega_L=4,
 \label{eq:v8-Levelt-form}
\end{equation}
and has finite index in $\operatorname{Sp}_{\Omega_L}(\Z)$.
The actual $t$-family's monodromy is conjugate to a subgroup of index at
most three in this group, and hence is arithmetic and Zariski dense in
$\operatorname{Sp}_4$.
\end{theorem}
\begin{proof}
No upper parameter differs from a lower parameter by an integer. The
Levelt theorem for the hypergeometric equation therefore gives the
companion-matrix description \cite{BeukersHeckman,SinghVenkataramana}.
Only the generated group is asserted here; identifying particular marked
loops requires the usual simultaneous base-point and loop conventions.
Direct multiplication verifies the invariant form and
\[
 A_L^{-1}B_L=I+(e_0+e_2)e_3^{\mathsf T},\quad
 A_L^4=-I,\quad \rank(B_L^3-I)=1.
\]
For arithmeticity, $f,g$ are coprime, monic, reciprocal integral
polynomials with constant term one. They are a primitive pair, because
$g$ has a nonzero coefficient of $X$, so cannot be a polynomial in $X^k$
for $k\ge2$. The difference $f-g=X^3+X$ is monic. These are exactly the
hypotheses of the Singh--Venkataramana theorem
\cite[Theorem~1.1]{SinghVenkataramana}. It proves the finite-index
conclusion; we do not claim to have computed that index. In fact, up to
interchanging the two polynomials, this is the worked example of
\cite[Section~2.1]{SinghVenkataramana}. The new result here is the
identification of our actual period system with that known arithmetic group,
not a new arithmeticity theorem.
The map $t\mapsto-256t^3/27$ is an unramified degree-three cover of the
corresponding punctured bases. The image of a subgroup of index three
under any group homomorphism has index at most three in the full image.
Apply this to monodromy and use \cref{thm:v8-PF}.
\end{proof}

\begin{scope}[Two different meanings of arithmetic]
This arithmeticity is finite index in an integral \emph{topological}
symplectic monodromy group. It is not a finite-prime Frobenius calculation,
not the $\ell$-adic monodromy theorem of V6, and not compact unitary
holonomy. Moreover the displayed Levelt lattice has polarization elementary
divisors $(1,2)$; its modulo-two reduction is not automatically the
Jacobian's labelled two-torsion. The genuine unimodular Betti lattice is
constructed next. We do not identify these lattices without a comparison.
\end{scope}

\section{A nonsingular arithmetic reference fibre and a complete marked period matrix}
\label{sec:v8-periods}

\subsection{A smooth route to the fixed curve}
Use the auxiliary family
\begin{equation}
 D_q:\quad Y^2=X^5-qX^2-X,\qquad D_1=C_{\mathrm{tet}}.
 \label{eq:v8-Dq}
\end{equation}
Its discriminant is $-(256+27q^4)$. It is smooth on the complex disk
$|q|<(256/27)^{1/4}$, which contains the closed disk $|q|\le3/2$ and
in particular the path $[0,1]$. This is a comparison path, not physical
time. For $t>0$, the change
\[
 x=t^{1/4}X,\quad y=t^{5/8}Y,\quad q=t^{-3/4}
\]
relates it to V7's family and sends $x^jdx/y$ to
$t^{(2j-3)/8}X^j dX/Y$. No new isomorphism class is substituted at the
actual fibre.

Write $\epsilon_j=[X^j dX/Y]$, $0\le j\le3$, with columns as in V7.

\begin{proposition}[Regular reference-gauge connection]
\label{prop:v8-regular-connection}
On $D_q$ the connection is
\begin{equation}
 \boxed{B(q)=\frac1{2(256+27q^4)}
 \begin{pmatrix}
 -27q^3&36q^2&-48q&64\\
 -64&-9q^3&12q^2&-16q\\
 -48q&64&9q^3&-12q^2\\
 108q^2&-144q&192&27q^3
 \end{pmatrix}.}
 \label{eq:v8-B}
\end{equation}
The constant de Rham form is the V7 matrix $\Omega$ and
$B^{\mathsf T}\Omega+\Omega B=0$.
\end{proposition}
\begin{proof}
Differentiation gives $\partial_q(X^j dX/Y)=X^{j+2}dX/(2Y^3)$.
Put $d=256+27q^4$ and let $R_j(X)=\sum_{k=1}^4R_{kj}X^k$, where
\[
 (R_{kj})=\frac1d
 \begin{pmatrix}
 27q^3&-36q^2&48q&-64\\
 -64&-9q^3&12q^2&-16q\\
 48q&-64&-9q^3&12q^2\\
 -36q^2&48q&-64&-9q^3
 \end{pmatrix}.
\]
For $f=X^5-qX^2-X$, direct polynomial multiplication gives
\[
 X^{j+2}=2f\sum_iB_{ij}X^i+2fR_j'-f'R_j.
\]
Dividing by $2Y^3$ identifies the final terms as $d(R_j/Y)$ and proves
the connection. The pairing and its constancy follow either from the same
second-kind residue calculation as V7 or by multiplication with $\Omega$.
This uses the inherited de Rham method; the new point is the regular
reference gauge with explicit initial periods.
\end{proof}

\subsection{Four integral cycles fixed before continuation}
On $D_0:Y^2=X^5-X$ put $\zeta_8=e^{\pi i/4}$ and
$R(X,Y)=(iX,\zeta_8Y)$. Let $\delta_0$ run from $0$ to $1$ with
$Y=i\sqrt{X(1-X^4)}$ and return on the opposite sheet. Define
$\delta_k=R_*^k\delta_0$ for $0\le k\le3$.

\begin{lemma}[An independently marked unimodular cycle basis]
\label{lem:v8-star-cycles}
The four cycles form a $\Z$-basis of $H_1(D_0,\Z)$, with intersection
matrix
\begin{equation}
 J=\begin{pmatrix}
 0&1&1&1\\-1&0&1&1\\-1&-1&0&1\\-1&-1&-1&0
 \end{pmatrix},\qquad \det J=1.
 \label{eq:v8-J}
\end{equation}
They transport uniquely as flat integral homology classes on the disk above.
At $q=1$ their nonzero branch endpoints are respectively $v,z,u,\bar z$.
\end{lemma}
\begin{proof}
The four projected radial arcs have disjoint interiors and only the branch
point zero in common. In a local coordinate $w$ there, $X=w^2$ and
$Y=iw+O(w^9)$, the oriented lifts of the four cycles have tangent lines
$\zeta_8^k\R$, with positive direction $\zeta_8^k$. For $k<\ell\le3$
the oriented local intersection has sign
$\sin((\ell-k)\pi/4)>0$. Small perturbations make these intersections
transverse, without creating intersections elsewhere. This gives
\eqref{eq:v8-J}. The determinant is one. Nondegeneracy first proves
rational independence. The intersection lattice of a compact genus-two
surface is unimodular, so a full-rank sublattice with determinant one has
index one, proving the integral assertion.

The smooth family over the disk has a locally constant homology lattice;
simple connectivity fixes continuation. On the real interval the roots
starting at $1,-1$ remain the positive and negative real roots; the other
two remain respectively in the upper and lower half-planes. A change of
this ordering would require a collision or a real-axis crossing by a
conjugate pair, contradicting the nonzero discriminant.
\end{proof}

\begin{theorem}[Beta-normalized complete de Rham--Betti comparison]
\label{thm:v8-period-matrix}
Let $\Pi(q)_{kj}=\int_{\delta_k(q)}\epsilon_j$. Define
\begin{equation}
 \boxed{\Pi(0)_{kj}=-\frac{i}{2}
 B\!\left(\frac{2j+1}{8},\frac12\right)
 \zeta_8^{k(2j+1)}.}
 \label{eq:v8-P0}
\end{equation}
Let $T'=TB(q)$, $T(0)=I$. Then
\begin{equation}
 \boxed{\Pi(q)=\Pi(0)T(q).}
 \label{eq:v8-period-transport}
\end{equation}
The initial beta values and the rational differential equation determine all
sixteen marked periods at $q=1$. This is an actual Betti comparison on the
integral lattice of \cref{lem:v8-star-cycles}, not a fundamental matrix with
unspecified initial constants.
\end{theorem}
\begin{proof}
On $\delta_0$, the two sheets give
\[
 \int_{\delta_0}\epsilon_j
 =-2i\int_0^1\frac{X^{j-1/2}}{\sqrt{1-X^4}}\,dX
 =-\frac i2 B\!\left(\frac{2j+1}{8},\frac12\right).
\]
All four integrals converge; the poles of $\epsilon_2,\epsilon_3$ at
infinity do not meet these cycles. Since
$R^*\epsilon_j=\zeta_8^{2j+1}\epsilon_j$, the remaining rows are
\eqref{eq:v8-P0}. This matrix is invertible: its nonzero column factors
multiply a Vandermonde matrix with distinct nodes
$\zeta_8,\zeta_8^3,\zeta_8^5,\zeta_8^7$.
Differentiating periods along the flat cycles gives $\Pi'=\Pi B$.
The initial values and uniqueness of the holomorphic matrix initial-value
problem prove \eqref{eq:v8-period-transport}. Invertibility persists because
a fundamental matrix has invertible determinant. The classical de
Rham--Betti theorem identifies this integration pairing with the full
comparison isomorphism \cite{KTdeRham,MilneAV}.
\end{proof}

\subsection{A convergent rational certificate, including initial constants}
Write $(256+27q^4)B(q)=\sum_{r=0}^3P_rq^r$ and
$T(q)=\sum_{n\ge0}T_nq^n$. Then $T_0=I$ and
\begin{equation}
 \boxed{256(n+1)T_{n+1}
 =\sum_{r=0}^{\min(3,n)}T_{n-r}P_r
 -27(n-3)T_{n-3},}
 \label{eq:v8-Trec}
\end{equation}
where the last term is zero for $n<4$. Every entry of every $T_n$ is
rational.

\begin{theorem}[An explicit global truncation bound at the actual fibre]
\label{thm:v8-period-tail}
For the induced maximum-row-sum matrix norm,
\begin{equation}
 \boxed{\left\|T(1)-\sum_{n=0}^{N}T_n\right\|_\infty
 \le2187(2/3)^{N+1}.}
 \label{eq:v8-Ttail}
\end{equation}
Each initial beta constant also has a positive rational-series enclosure.
For $a\in\{1/8,3/8,5/8,7/8\}$, put
$c_n=(1/2)_n/n!$, $d_n=(1-a)_n/n!$. Then
\begin{align}
 B(a,\tfrac12)&=2^{-a}\sum_{n\ge0}\frac{c_n2^{-n}}{a+n}
  +2^{-1/2}\sum_{n\ge0}\frac{d_n2^{-n}}{n+1/2},\\
 0\le\text{tail after }N
 &\le 2^{-a}\frac{2^{-N}}{a+N+1}
       +2^{-1/2}\frac{2^{-N}}{N+3/2}.
 \label{eq:v8-beta-tail}
\end{align}
Thus the period computation requires rational arithmetic, square roots,
and explicit tails, not an uncertified gamma evaluation.
\end{theorem}
\begin{proof}
Coefficient comparison in $ (256+27q^4)T'=T\sum P_rq^r$ gives
\eqref{eq:v8-Trec}. On $|q|\le3/2$, the denominator modulus is at least
$256-27(3/2)^4=1909/16$. The largest sum of the absolute numerator
coefficients of $B$ evaluated at $3/2$ is $5937/16$; hence
$\|B(q)\|_\infty\le5937/1909<4$. Along every radius, the integral
series for $T'=TB$ gives $\|T(q)\|_\infty\le e^{4|q|}\le e^6<729$.
Cauchy's coefficient estimate gives
$\|T_n\|_\infty\le729(2/3)^n$. Summing the tail proves
\eqref{eq:v8-Ttail}.

Split Euler's beta integral at $1/2$ and substitute $u\mapsto1-u$ in
its second half. Expand respectively $(1-u)^{-1/2}$ and
$(1-u)^{a-1}$ on $[0,1/2]$, where both positive binomial series converge
uniformly away from the integrable factor at zero. Termwise integration,
or monotone convergence, gives the two sums. Both coefficient sequences
lie in $(0,1]$. Bounding all remaining denominators from below and summing
a geometric series proves the tail. Finally $2^{-a}$ is obtained from
three nested square roots of two and integer powers.
\end{proof}

\begin{remark}[What is and is not an exact numerical claim]
The certificate uses $N=220$ for \eqref{eq:v8-Trec}, whose tail is less
than $2.66\times10^{-36}$, and $N=140$ for each beta sum. The finite sums
are evaluated with outward rational intervals. All decimal displays below
are rounded summaries of enclosing rational intervals. A computed interval
containing zero is never used as a proof of an exact zero or symmetry.
\end{remark}

\section{The acquired period matrix, Hodge metric, and actual theta-frame signs}
\label{sec:v8-acquisition}

\subsection{A symplectic marking and the actual Siegel matrix}
Define the new cycle order $(a_1,a_2,b_1,b_2)$ by the columns of
\begin{equation}
 S=\begin{pmatrix}1&1&0&1\\0&-1&1&-1\\0&1&0&0\\0&0&0&1\end{pmatrix}.
 \label{eq:v8-S}
\end{equation}
Thus $a_1=\delta_0$, $a_2=\delta_0-\delta_1+\delta_2$,
$b_1=\delta_1$, and $b_2=\delta_0-\delta_1+\delta_3$. One has
$S^{\mathsf T}JS=\begin{psmallmatrix}0&I_2\\-I_2&0\end{psmallmatrix}$.
Let $\Pi_F$ be the first two columns of $\Pi(1)$ and put
\begin{equation}
 S^{\mathsf T}\Pi_F=\begin{pmatrix}\mathcal A\\\mathcal B\end{pmatrix},
 \qquad \tau=\mathcal B\mathcal A^{-1}.
 \label{eq:v8-tau}
\end{equation}

\begin{theorem}[Certified actual periods and Hodge pairing]
\label{thm:v8-actual-periods}
In this independently declared marking the actual period matrix is
\begin{equation}
 \boxed{\tau\simeq
 \begin{pmatrix}
 0.700823987+0.957931596i&-0.411683081-0.502230254i\\
 -0.411683081-0.502230254i&0.843099771+0.987614458i
 \end{pmatrix}.}
 \label{eq:v8-tau-numeric}
\end{equation}
The Hodge inner product on the actual algebraic differentials is
\begin{equation}
 \boxed{h_F=\frac i2\Pi_F^*J^{-1}\Pi_F
 =\mathcal A^*(\operatorname{Im}\tau)\mathcal A.}
 \label{eq:v8-Hodge}
\end{equation}
Numerically,
\[h_F\simeq\begin{pmatrix}
 17.969468821990&-1.937425189159\\
 -1.937425189159&17.645797424338
 \end{pmatrix}.\]
Here $h_F$ is real symmetric, $15I\prec h_F\prec20I$, and
$\operatorname{Im}\tau\succ(9/20)I$. In particular this computes an
independent Hodge metric, not V7's theta metric declared under another name.
\end{theorem}
\begin{proof}
The matrix $S$ is integral unimodular and has the stated symplectic pairing,
by direct multiplication. The Riemann bilinear relations for this integral
basis imply that $\mathcal A$ is invertible, $\tau$ is symmetric, and
$\operatorname{Im}\tau$ is positive definite. In the star basis the
bilinear integral is $-p_\alpha^{\mathsf T}J^{-1}p_\beta$. Consequently
\[
 (h_F)_{ij}=\frac i2\int_{C_{\mathrm{tet}}}
            \omega_j\wedge\overline{\omega_i}
          =\frac i2\overline{p_i}^{\mathsf T}J^{-1}p_j.
\]
Changing to the symplectic basis and substituting
$\mathcal B=\tau\mathcal A$ gives the second equality of
\eqref{eq:v8-Hodge}. These standard comparison and positivity facts use
\cite{MilneAV}. Reality follows from the real curve and real algebraic
forms: its antiholomorphic real involution conjugates the forms and reverses
the surface orientation, leaving these Hermitian pairings real.

For the numerical assertions, use \eqref{eq:v8-P0},
\eqref{eq:v8-Trec}--\eqref{eq:v8-beta-tail}, and outward interval matrix
inversion. The resulting enclosures include
\[
 \begin{split}
 17.969468821990235&<(h_F)_{00}<17.969468821990236,\\
 -1.937425189159362&<(h_F)_{01}<-1.937425189159361,\\
 17.645797424337921&<(h_F)_{11}<17.645797424337922.
 \end{split}
\]
The accompanying rational certificate lists all sixteen periods, all six
independent real entries of $\tau$, and every inversion pivot. The smallest
diagonal-minus-off-diagonal row bound for $h_F$ exceeds $15$; the largest
row sum is below $20$. The corresponding lower row bound for
$\operatorname{Im}\tau$ exceeds $9/20$. This proves the strict numerical
bounds. Exact symmetry and reality come from the preceding geometric
argument, not from overlapping numerical intervals.
\end{proof}

For orientation, the first two columns of the complete star-period matrix are
\begin{equation}
 \Pi_F\simeq\begin{pmatrix}
 -4.385176035420i&-2.009174802474i\\
 3.108299915617-3.534918869719i&1.361274712945+1.526782236944i\\
 4.724528327053&-1.583921110483\\
 3.108299915617+3.534918869719i&1.361274712945-1.526782236944i
 \end{pmatrix}.
 \label{eq:v8-hol-periods}
\end{equation}
The positive entry $4.724528327053\ldots$ is the real cycle in
\cref{cor:v8-real-period} at $t=1$.

\subsection{The integral marking fixes the theta characteristic}
Use binary characteristic notation $(a\mid b)$ for the half-period
$(\tau a+b)/2$. At the actual fibre the four star cycles correspond to
paths from zero to $v,z,u,\bar z$. The class of each such path is half its
closed lifted cycle. In standard cycle coordinates, if a star vector has
coordinates $S^{-1}v=(b,a)$, its half-period is $(a\mid b)$.

\begin{proposition}[An independently derived characteristic dictionary]
\label{prop:v8-characteristics}
In the marking \eqref{eq:v8-S}, the five classes $[W_r-W_\infty]$ for
$r=0,u,v,z,\bar z$ are respectively
\begin{equation}
 \begin{array}{c|ccccc}
 r&0&u&v&z&\bar z\\\hline
 (a\mid b)&(11\mid11)&(01\mid00)&(11\mid01)&(01\mid11)&(00\mid01).
 \end{array}
 \label{eq:v8-torsion-dictionary}
\end{equation}
The actual odd characteristic associated with $\mathcal O(W_\infty)$ is
\begin{equation}
 \boxed{\kappa=(\alpha\mid\beta)=(10\mid10).}
 \label{eq:v8-kappa}
\end{equation}
\end{proposition}
\begin{proof}
Write $v_\infty=(1,1,1,1)^{\mathsf T}$ in the star lattice modulo two.
The relation among the five nonzero Weierstrass differences based at zero
gives $[W_\infty-W_0]=v_\infty$. One may also check this relation from
$\operatorname{div}(Y)=W_0+\sum W_r-5W_\infty$: the sum of the five
nonbase differences is $6(W_\infty-W_0)=0$ in two-torsion. Adding each
star coordinate to $v_\infty$ and applying $S^{-1}$ gives
\eqref{eq:v8-torsion-dictionary}.

For a characteristic $(\alpha\mid\beta)$ the quadratic refinement is
$q_\kappa(a,b)=a\cdot b+a\cdot\beta+b\cdot\alpha$ modulo two.
The five displayed nonzero classes are the singular classes of the odd
characteristic at infinity, by the inherited Weierstrass/theta dictionary.
Requiring $q_\kappa=0$ on those five classes is a linear system in its four
bits. Its unique solution is $(10\mid10)$, of odd parity. Explicitly the
$u$ and $\bar z$ columns give $\beta_2=\alpha_2=0$; the $v$ column gives
$\beta_1=1$, and the $z$ column gives $\alpha_1=1$. The temporal column
then agrees automatically. Nothing is fitted against a theta expectation or
renewal frame.
\end{proof}

\subsection{Validated acquisition of the four previously undetermined signs}
In the standard second-order theta basis indexed by
$\sigma=(00,01,10,11)$, V2's already proved formula now has entirely
specified inputs:
\begin{equation}
 t_\sigma=\sum_{n\in\Z^2}
 e^{2\pi i(n+\sigma/2)^{\mathsf T}\tau(n+\sigma/2)},\qquad
 r_\sigma=(-1)^{\sigma_1}t_{\sigma+(1,0)},\qquad
 \rho=rr^*/(r^*r).
 \label{eq:v8-actual-r}
\end{equation}
For each half-period $h=(a\mid b)$ fix its Hermitian displacement by
\begin{equation}
 (P_h)_{\sigma,\sigma+a}
 =i^{\,a\cdot b\bmod2}(-1)^{\sigma\cdot b},
 \label{eq:v8-Pauli-convention}
\end{equation}
with every other entry zero. The exponent of $i$ is explicitly reduced
modulo two; replacing it by the unreduced dot product would change certain
lift signs. Let $\Gamma_0,\ldots,\Gamma_4$ be these displacements in
$(0,u,v,z,\bar z)$ order, and define
$b_{ij}=\Tr(\rho i\Gamma_i\Gamma_j)$ for $i<j$.

\begin{theorem}[Actual marked theta signs, acquired without a renewal target]
\label{thm:v8-theta-signs}
The matrices in \eqref{eq:v8-Pauli-convention} anticommute pairwise on the
five labels. In this independently marked arithmetic frame, the actual
source satisfies
\begin{equation}
 \boxed{\operatorname{sgn}(b_{01},b_{02},b_{03},b_{04})=(-,-,+,+).}
 \label{eq:v8-four-signs}
\end{equation}
Its complete sign pattern, in lexicographic pair order, is
\begin{equation}
 \boxed{(-,-,+,+,0,-,+,+,-,+).}
 \label{eq:v8-all-signs}
\end{equation}
The zero is the exact inherited identity $b_{12}=0$. The coefficient vector
in \eqref{eq:v8-actual-r} has the certified approximation
\begin{equation}
 r\simeq\begin{pmatrix}
 0.184630258660+0.383024497202i\\
 0.177521657346+0.402300202291i\\
 -0.999755090484+0.013316136795i\\
 -0.079173495922-0.411209398344i
 \end{pmatrix}.
 \label{eq:v8-r-numeric}
\end{equation}
Its squared norm satisfies
\[1.549204661650477<r^*r<1.549204661650478.\]
The nine nonzero bivectors have the certified values summarized by
\[
 \begin{array}{c|r@{\qquad}c|r}
 01&-0.550651987671264&02&-0.516970870638065\\
 03& 0.463423945963848&04& 0.463423945963848\\
 13&-0.540205734684735&14& 0.540205734684735\\
 23& 0.444411510854950&24&-0.444411510854950\\
 34& 0.146124421185146&&
 \end{array}
\]
where each display is rounded from the rational enclosures in the certificate.
\end{theorem}
\begin{proof}
The pairwise symplectic pairings of the five binary labels equal one, so the
Hermitian displacements anticommute. Their convention is fixed before
\eqref{eq:v8-actual-r} is evaluated. Apply V2's coordinate formula with
\cref{thm:v8-actual-periods,prop:v8-characteristics}; it is the actual
squared theta section, not an arbitrary vector with desired overlaps.

For a fully validated evaluation, truncate the two theta indices to
$[-5,5]^2$. The proved bound $\operatorname{Im}\tau\succeq(9/20)I$
gives $c=2\pi(9/20)>27/10$. In the notation of V2's tail theorem,
$B(c)<4$, $2+1/(c(11/2))<3$, and $c(11/2)^2>81$. Since $e>2$, each
coordinate tail is at most
\begin{equation}
 2B(c)T_5(c)<24\,2^{-81}.
 \label{eq:v8-small-theta-tail}
\end{equation}
The finite exponentials themselves are evaluated by rational complex
intervals. Machin's identity
$\pi=16\arctan(1/5)-4\arctan(1/239)$ with alternating-series tails gives
a rational interval for $\pi$. For every occurring exponent $z$, the
certificate checks $|\Re z|+|\Im z|\le2048$. Expand
$\exp(z/4096)$ through degree $40$. Its complex absolute remainder is
at most $2\,2^{-41}/41!$, since $|z/4096|\le1/2$ and $e^{1/2}<2$.
Twelve interval squarings recover $\exp(z)$. This accounts for truncation,
period uncertainty, and arithmetic roundoff, with no unvalidated
floating-point exponential.

Assemble the four coefficients, add \eqref{eq:v8-small-theta-tail}, and
form the exact rational functions $r^*(i\Gamma_i\Gamma_j)r/(r^*r)$
using outward intervals. The denominator lower bound exceeds $1.54$.
Every asserted nonzero sign is separated from zero by more than $0.1$;
the four signs in \eqref{eq:v8-four-signs} are separated by more than
$0.46$. The certificate lists the enclosing intervals. For $b_{12}=0$
and the temporal first expectation, we use the exact V3 identities;
intervals containing zero only cross-check them. The acquired squares
independently overlap every algebraic Thomae enclosure from V3, but that
numerical consistency is not substituted for V3's exact proof.
\end{proof}

\begin{scope}[Which ambiguity has actually been closed]
This closes the four-sign acquisition in a fully specified arithmetic
symplectic and Hermitian-lift marking. Relative to V4's canonical abstract
representative, its sign reorientation is $(+,-,-,+,+)$ on the five
Clifford labels. It does not acquire the four \emph{relative}
signs against an unspecified independent renewal frame. The actual source
has not been symmetrized, mixed, whitened, or redefined. The theta-section
phase is fixed by the analytic coefficient formula; its overall ray phase
remains immaterial. In particular no central Store/chronology expectation
has been computed here.
\end{scope}

\section{A Hodge-based positive rank-ten completion with an explicit selection rule}
\label{sec:v8-positive-completion}

\subsection{The previously missing positive comparison data now occur explicitly}
At the actual fibre put $E=H^1_{\mathrm{dR}}(C_{\mathrm{tet}})\otimes\C$.
The full acquired period matrix determines its Betti real structure:
\begin{equation}
 \boxed{\mathfrak c_B(v)=C_B\bar v,\qquad
 C_B=\Pi(1)^{-1}\overline{\Pi(1)}.}
 \label{eq:v8-Betti-real}
\end{equation}
Indeed integration of $\mathfrak c_B(v)$ is the conjugate of integration
of $v$. Thus $C_B\overline{C_B}=I$. This real structure interchanges
$F=H^{1,0}$ with $\bar F=H^{0,1}$; it is not V4's source-space arithmetic
conjugation operator. All entries of $C_B$ have certified enclosures in the
same finite certificate.

The actual polarized Hodge metric $h_E$ makes $E=F\oplus\bar F$ an
orthogonal direct sum, with metrics $h_F$ and its conjugate. Hence
$W=\operatorname{Sym}^2E$ has the orthogonal Hodge decomposition
\begin{equation}
 W=W^{2,0}\oplus W^{1,1}\oplus W^{0,2},\qquad
 \dim W^{2,0},\dim W^{1,1},\dim W^{0,2}=3,4,3.
 \label{eq:v8-Hodge-split}
\end{equation}
Write $h_0=\operatorname{Sym}^2h_E$ for the induced product metric.
Throughout this subsection the theta-fibre line is either retained with a
fixed positive norm or trivialized with V7's normalization
$q(0)^*H_\Theta q(0)=1$. Its absolute physical normalization is not newly
derived.

\begin{theorem}[A unique minimum-change positive completion]
\label{thm:v8-extension}
On $W^{2,0}=\operatorname{Sym}^2F$ prescribe the actual metric $H_\Theta$
of V7, and on $W^{0,2}$ prescribe its Betti conjugate. Among Hermitian
metrics $g$ satisfying these two compression constraints, the functional
\begin{equation}
 \mathcal E(g)=\|L_g-I\|_{\mathrm{HS}(h_0)}^2,
 \qquad g(v,w)=h_0(v,L_gw),
 \label{eq:v8-minchange}
\end{equation}
has a unique minimizer. It is positive and is explicitly
\begin{equation}
 \boxed{g_\star=H_\Theta\ \oplus\ h_0|_{W^{1,1}}\ \oplus\
                  \overline{H_\Theta}.}
 \label{eq:v8-gstar}
\end{equation}
It is invariant under Betti conjugation and functorial under marked
isomorphisms preserving the polarization, source metric, and fibre norm.
All its coordinates are determined by the acquired periods and V7's
algebraic metric. No single-valued positive parallel extension of it
around the V7 unipotent loop exists.
\end{theorem}
\begin{proof}
Use the $h_0$-orthogonal decomposition \eqref{eq:v8-Hodge-split}. The
compression constraints fix the first and third diagonal blocks of the
self-adjoint operator $L_g$. The squared Hilbert--Schmidt norm is the sum
of the squared norms of all its blocks after subtracting the identity.
It is uniquely minimized by setting every unprescribed off-diagonal block
to zero and the middle diagonal block to the identity. The two prescribed
blocks are positive because they represent positive metrics relative to
positive $h_0$ restrictions. Thus the minimizing operator is positive,
so the unconstrained Hermitian minimizer is already a genuine metric.
This proves \eqref{eq:v8-gstar}, uniqueness, and reality. The norm and
constraints are invariant under the stated marked isomorphisms, giving
functoriality.

For an explicit coordinate recipe let
$U=(e_0,e_1,\mathfrak c_Be_0,\mathfrak c_Be_1)$. It is invertible by
Hodge decomposition. Apply $\operatorname{Sym}^2U$ to move between the
algebraic monomial basis and the Hodge basis. In the unnormalized symmetric
product convention $vw=(v\otimes w+w\otimes v)/2$, the middle block in
the Hodge basis is $\tfrac12h_F\otimes\bar h_F$. The other two blocks
are $H_\Theta$ and its conjugate. Congruence by
$(\operatorname{Sym}^2U)^{-1}$ gives the algebraic-coordinate matrix of
$g_\star$. Every required entry is now specified by \eqref{eq:v8-Betti-real}
and the acquired metric.

If the same construction is made smoothly along the actual family, any
claim of parallelism around V7's unipotent loop would give a positive
parallel metric on its rank-ten local system, contradicted by
\cref{thm:v7-monodromy}. Even pointwise, its restriction to $W^{2,0}$ is
V7's nonproduct theta metric, so it cannot equal the product metric $h_0$
there. No flat unitary comparison is asserted.
\end{proof}

\begin{remark}[The rule is additional, but no unknown matrix is hidden in it]
The minimal block-projection argument is elementary and is not advertised as
new general metric theory. Its use here is constructive because the actual
polarization, Betti conjugation, and Hodge metric have now been acquired,
whereas V7 left them unresolved. A different completion principle could
choose different middle or cross-Hodge blocks. The result says that the
stated minimum-change principle selects one unique, executable arithmetic
comparison metric; it does not say that the theta source alone forces this
principle, or that a physical renewal instrument realizes its metric.
\end{remark}

\section{V8 verification, completed gates, and remaining comparisons}
\label{sec:v8-status}

\subsection{What the finite certificate proves}
The accompanying standalone verifier checks the scalar cyclic determinant
and operator relation, the four regular-gauge exact differential identities,
the discriminant, pairings and cycle change of basis, the Levelt matrices
and hypotheses used in the arithmeticity specialization, and the complete
homology-to-characteristic dictionary. It evaluates the rational transport
recurrence and beta sums with the analytic tails proved above. Outward
rational interval inversion produces the actual period matrix, Hodge metric,
and Betti-conjugation matrix. Finally it evaluates all four actual theta
constants, including exponential and Gaussian tails, and returns strict
sign intervals for the nine nonzero bivectors. The many exponential norm
checks are implementation checks, not separate geometric theorems.

It verifies byte-for-byte preservation of every earlier mathematical body.
The Betti comparison, Riemann bilinear identities, theta-characteristic
interpretation, Levelt theorem, and Singh--Venkataramana theorem remain
separately declared classical inputs. The verifier does not claim to
formalize those theorems or to verify a physical occurrence row.

\subsection{Proved or explicitly constructed in V8}
\begin{enumerate}[label=\textup{(P8.\arabic*)},leftmargin=3.6em]
\item The actual scalar period system is a rank-four hypergeometric pullback.
      Its companion group is arithmetic, by a verified specialization of
      the cited arithmeticity theorem. This strengthens the global
      topological transport description without conflating it with finite-prime
      or compact-unitary monodromy.
\item An independently declared smooth comparison path starts at a fibre
      with beta-integral periods and an explicit unimodular cycle basis.
      A convergent rational recurrence and explicit tails acquire the full
      marked de Rham--Betti matrix at the actual curve.
\item The actual normalized genus-two period matrix, polarized Hodge metric,
      and cohomological Betti real structure are now numerically certified.
      They were explicitly unresolved in V7.
\item The integral characteristic marking and four actual theta-source signs
      have been acquired without a renewal target. The previous algebraic
      magnitudes are inherited, and the pure source is unchanged.
\item Once the acquired polarization and source metric are retained, a
      declared minimum-change rule gives one unique positive, real rank-ten
      completion. It is an explicit comparison construction, not a new
      physical source-selection theorem.
\end{enumerate}

\subsection{Inherited inputs and unresolved data}
The normalized chronology loading and the companion geometric/SM/continuum
results remain antecedents \cite{GT77,Arithmetic,Spacetime,Spectral,Duality,Variational}.
The actual theta source, its signed orbit, the finite-prime Frobenius data,
Cartier transfer, endomorphism certificate, second-jet bridge, nonproduct
metric, rank-ten differential completion, and local nonunitarity obstruction
remain the unchanged conclusions of V1--V7.

The phrase ``marked Betti comparison remains to be acquired'' in V7 is now
superseded for the specific comparison path and marking above. It is not
superseded for a different undeclared marking. In particular the independent
renewal packet has not been supplied in this Betti/theta frame. Its mixed
source maps, four relative signs, central Store/chronology signature scalar,
positive transport and physical clock therefore remain to be identified.
An arithmetic period is not an operational source occurrence.

The $p$-adic/crystalline comparison and a Frobenius matrix in this marked
source basis are still not computed. The complex source and integral lattice
do not by themselves choose a $p$-adic comparison map or recover the central
quadratic twist from discarded rational descent data. The exact finite-prime
polynomials of V5 and the full arithmetic realization of V6 remain the
appropriate checks for that next step. No mass gap, Standard-Model numerical
prediction, Einstein limit, GRH estimate, or Frobenius--Wilson equivalence is
inferred here.

\paragraph{Version integrity.}
Every V1--V7 mathematical body, including each historical status ledger, is
preserved byte for byte. This continuation adds new results after V7; it
does not retroactively erase earlier uncertainties. The next version must
append after this body and state any correction explicitly.
% ====================== END IMMUTABLE V8 MATHEMATICAL BODY ======================
% ===================== BEGIN IMMUTABLE V9 MATHEMATICAL BODY =====================
\clearpage
\section{V9 continuation: two native compact holonomies and their comparison boundary}
\label{sec:v9-context}

The V7 obstruction concerns a \emph{flat} Gauss--Manin connection and a
positive parallel metric. V8 acquired the actual periods and two different
positive source metrics. Neither result answers the following upstream
question: what are the holonomies of the \emph{nonflat metric connections}
that these already constructed arithmetic bundles themselves determine?
This continuation answers that question, rather than choosing a unitary
matrix to imitate one Frobenius conjugacy class.

Fix a sufficiently small simply connected complex disk $U$ about $t=1$ in
\cref{eq:v7-base}. The family is still
$C_t:y^2=x^5-x^2-tx$. Continue the V8 integral marking from its actual fibre.
There are two previously constructed holomorphic bundles on $U$:
\[
 F=H^{1,0}(C_t),\quad \rank F=2,
 \qquad
 V=\ker\bigl(H^0(J_t,L_t^2)\xrightarrow{\mathrm{ev}_0}L_t^2|_0\bigr),
 \quad\rank V=3.
\]
Use the actual polarized Hodge metric on $F$ and the Heisenberg-unitary
metric on the level-two theta bundle, normalized to make the marked
second-order theta basis orthonormal. Restrict the latter to $V$. These
metrics define Chern connections $D_F,D_V$ without a renewal target.
The rank-three bundle is the \emph{actual vanishing theta-source span} of
V6, not an auxiliary three-dimensional representation appended here.

We prove, on every sufficiently small such disk,
\begin{equation}
 \boxed{\operatorname{Hol}^0(D_F)=U(2),\qquad
        \operatorname{Hol}^0(D_V)=U(3),\qquad
        \operatorname{Hol}^0(D_F\oplus D_V)=U(2)\times U(3).}
 \label{eq:v9-headline}
\end{equation}
Here restricted holonomy uses all contractible piecewise smooth loops in
the parameter disk. It is not arithmetic Galois monodromy, not the finite
theta-frame group, and not the representation of the fundamental group of
a flat system. The compact derived group is $SU(2)\times SU(3)$, but the
actual two determinant curvatures are independent. In particular this
construction does \emph{not} select $S(U(3)\times U(2))$ or establish a
physical Standard-Model realization.

\paragraph{Dependency and novelty ledger.}
We retain the V7 rational connection and quadratic-source identification,
the V8 marked period matrix and theta basis, and all previous source signs
and arithmetic inputs. We do not repeat their acquisition. The new
calculations are an exact Hodge-kernel rotation, a third-order theta-jet
certificate, the resulting two full compact holonomies, a determinant
curvature independence certificate, and a projective-curvature obstruction
to identifying the two rank-three connections. General Chern and Hodge
curvature theory is classical \cite{GreenGriffiths,SimpsonHiggs}; the
curvature-to-holonomy inclusion is classical \cite{AmbroseSinger}. The
specific holonomy computations are proved below. No general nonabelian
Hodge existence theorem, conjectural automorphy, or prime equidistribution
is used.

\paragraph{Metric and scope convention.}
A different smooth scalar normalization of the theta metric can change its
central connection. Our full-product conclusion refers to the stated
orthonormal theta-frame normalization; its projective and derived-group
conclusions do not depend on scalar metric changes. The family parameter
is a complex deformation parameter, not a physical clock. The Chern
connections need not make the distinguished theta section parallel.
All V1--V8 mathematical bodies remain unchanged.

\section{The two Chern connections are explicit functions of acquired data}
\label{sec:v9-Chern}

\subsection{Curvature convention and the Hodge connection}
For a holomorphic frame with positive Hermitian Gram $h$, set
\begin{equation}
 D=d+h^{-1}\partial h,\qquad
 F_D=R_D\,dt\wedge d\bar t,\qquad
 R_D=-\partial_{\bar t}(h^{-1}\partial_t h).
 \label{eq:v9-curv-convention}
\end{equation}
Then $R_D$ is $h$-self-adjoint. On real tangent vectors
$(\partial_x,\partial_y)$, where $t=x+iy$, curvature is $-2iR_D$.
This fixes all curvature signs and removes a possible interchange of
$dt\wedge d\bar t$ with $d\bar t\wedge dt$.

Let $\Pi(t)$ be the continuation of the V8 marked period matrix in the V7
family, so $\Pi'=\Pi A$. Put
\[
 u(t)=\Pi(t)_{[:,0:2]},\qquad Q=\frac i2J^{-1},\qquad h=u^*Qu.
\]
The constant Hermitian form $Q$ has signature $(2,2)$. The column space of
$u$ is positive and its $Q$-orthogonal complement is negative, by the
polarized Hodge decomposition already used in V8.

\begin{proposition}[An executable Hodge connection and curvature]
\label{prop:v9-Hodge-curvature}
Define $P=u h^{-1}u^*Q$ and $b=(I-P)u'$. Then
\begin{align}
 D_F&=d+h^{-1}u^*Qu'\,dt,\\
 R_F&=h^{-1}\bigl((u')^*Qu\,h^{-1}u^*Qu'-(u')^*Qu'\bigr)
     =-h^{-1}b^*Qb. \label{eq:v9-Hodge-R}
\end{align}
Consequently $R_F\succeq_h0$. On the actual family it has rank one
throughout $U$, and its kernel is
\begin{equation}
 \boxed{\ker R_F=\C k(t),\qquad k(t)=4t e_0+3e_1.}
 \label{eq:v9-kernel}
\end{equation}
Writing $\lambda=\Tr R_F$, one has $\lambda>0$ and
$\spec R_F=\{0,\lambda\}$. Every matrix in these formulas is determined by
V8's periods and V7's rational connection.
\end{proposition}
\begin{proof}
Since $u$ is holomorphic, $\partial_t h=u^*Qu'$,
$\partial_{\bar t}h=(u')^*Qu$, and
$\partial_{\bar t}\partial_t h=(u')^*Qu'$. Differentiating
$h^{-1}\partial_t h$ proves the first expression for $R_F$.
The matrix $P$ is the $Q$-orthogonal projection onto $\operatorname{Ran}u$.
Expanding $b^*Qb$ proves the second expression. Negativity of $Q$ on the
complement gives positivity. Under the de Rham--Betti isomorphism, $b$
represents the second fundamental form $F\to E/F$. Its rank and kernel
are exactly those of the bottom-left block in \cref{eq:v7-B}. That block
has rank one and kernel $(4t,3)^{\mathsf T}$. This proves every rank and
kernel assertion without a numerical eigenvalue test.
\end{proof}

\begin{remark}[What the nonflat unitary connection retains and omits]
In the smooth Hodge splitting $E=F\oplus\bar F$, the flat Gauss--Manin
connection decomposes as
\[
 \nabla=D_E+\psi+\psi^{\dagger},\qquad D_E=D_F\oplus D_{\bar F},
\]
where $\psi:F\to\bar F\otimes\Omega_U^{1,0}$ is the second fundamental
form and $\psi^{\dagger}$ is its Hodge adjoint of type $(0,1)$.
Compatibility with the flat polarized form, of signs $(+,-)$ on the two
summands, gives the plus sign between these off-diagonal terms. The diagonal
part of $\nabla^2=0$ gives $F_{D_E}+[\psi,\psi^\dagger]=0$, and its
$F$ block is \cref{eq:v9-Hodge-R}. This is the usual Hodge-bundle/Higgs
calculus \cite{GreenGriffiths,SimpsonHiggs}, not a physical Higgs-field
identification. Omitting $\psi$ changes the connection: $D_F$ is not a
rank-two flat subquotient of $E$. Thus \cref{eq:v9-headline} does not
contradict V6's small-rank arithmetic tensor obstruction or V7's unipotent
monodromy obstruction.
\end{remark}

\subsection{The actual theta-source bundle is a moving hyperplane}
In the marked holomorphic second-order theta basis, write
\[
 \mathcal T=H^0(J_t,L_t^2)\cong U\times\C^4,
 \qquad
 a(t)=(t_{00}(\tau(t)),t_{01}(\tau(t)),t_{10}(\tau(t)),t_{11}(\tau(t)))^{\mathsf T}.
\]
The notation $t_{\sigma}$ is V2's second-order theta constant and does not
mean differentiation in the deformation parameter. The evaluation row is
$a(t)^{\mathsf T}$, not $a(t)^*$. Thus
\begin{equation}
 \boxed{V_t=\{v\in\C^4:a(t)^{\mathsf T}v=0\}.}
 \label{eq:v9-hyperplane}
\end{equation}
The vector $a$ is nonzero near $t=1$, since V8 proved $a^*a>1.54$ there.
Heisenberg invariance forces the Hermitian matrix of $\mathcal T$ to be a
positive scalar times $I_4$, by irreducibility. We use the normalization
$I_4$ in this holomorphic frame. The ambient Chern connection is therefore
flat on $U$. Arbitrary scalar normalizations are discussed separately.

\begin{proposition}[An explicit source-bundle connection and its curvature]
\label{prop:v9-theta-curvature}
On a chart $a_0\ne0$, let $f=(a_1/a_0,a_2/a_0,a_3/a_0)$ and use the
holomorphic frame $w_j=e_{j+1}-f_j e_0$ of $V$. Its metric and connection
are
\begin{equation}
 \boxed{g=I_3+f^*f,\qquad
 D_V=d+g^{-1}f^*f'\,dt,\qquad
 R_V=-\frac1{1+ff^*}\,g^{-1}(f')^*f'.}
 \label{eq:v9-theta-connection}
\end{equation}
In particular $R_V$ has rank at most one. Its only possible nonzero
eigenvalue is $-\kappa$, where
\begin{equation}
 \boxed{\kappa=
 \frac{\|a'\|^2}{\|a\|^2}
 -\frac{|a^*a'|^2}{\|a\|^4}\ge0.}
 \label{eq:v9-kappa}
\end{equation}
The coefficient $\kappa$ is unchanged under a nowhere-zero holomorphic
rescaling of the evaluation row.
\end{proposition}
\begin{proof}
The displayed frame has Gram $g=I+f^*f$. Hence $g_t=f^*f'$ and
$g_{t\bar t}=(f')^*f'$. Formula \cref{eq:v9-curv-convention} gives
\[
 R_V=-g^{-1}(f')^*\bigl(1-fg^{-1}f^*\bigr)f'
     =-(1+ff^*)^{-1}g^{-1}(f')^*f',
\]
where the last identity is the scalar Schur identity, or follows by
multiplying $g=I+f^*f$. This proves the rank assertion and negativity.
Its negative trace is the squared norm of the second fundamental map
$V\to\mathcal T/V$. The orthogonal unit normal is $\bar a/\|a\|$;
differentiating $a^{\mathsf T}v=0$ identifies that squared norm with
\cref{eq:v9-kappa}. Equivalently, substitute $a=a_0(1,f^{\mathsf T})^{\mathsf T}$
in the formula. A holomorphic scalar changes $\log\|a\|^2$ by the
real part of a holomorphic function, so its mixed derivative, and hence
$\kappa$, is unchanged.
\end{proof}

\begin{scope}[The same source, but a different connection from the Hodge square]
V6's second-jet isomorphism $\mathfrak j_2:V\to\lambda^2\otimes
\operatorname{Sym}^2F$ is holomorphic on this local smooth family. Its
fibrewise bijectivity was proved in V6; leading jets commute with holomorphic
base change, which gives the bundle statement. Transporting the restricted
theta metric gives V7's conformal source metric. It does not give the
product Hodge metric. Likewise $D_V$ is the Chern connection of the actual
restricted theta metric, not the connection induced by $D_F$ through the
symmetric-square operation. A quantitative distinction is proved below.
\end{scope}

\section{Full rank-two compact holonomy from an exact kernel rotation}
\label{sec:v9-U2}

We first record the elementary Lie-algebra fact needed twice below.
The only general holonomy input is that curvature, transported back to a
fixed fibre along any path, belongs to the restricted holonomy Lie algebra
\cite{AmbroseSinger}. It follows already by differentiating the parallel
transport around small rectangles and conjugating by the path.

\begin{lemma}[An irreducible unitary Lie algebra containing rank one]
\label{lem:v9-rankone-Lie}
Let $\mathfrak h\subseteq\mathfrak u(n)$ act irreducibly on $\C^n$.
If $iP\in\mathfrak h$ for a rank-one orthogonal projection $P$, then
$\mathfrak h=\mathfrak u(n)$.
\end{lemma}
\begin{proof}
Let $G$ be the connected Lie subgroup with Lie algebra $\mathfrak h$.
If $P$ projects onto $\C v$, then $iP_{gv}\in\mathfrak h$ for every
$g\in G$. Their orbit is connected and spans $\C^n$, by irreducibility.
Initially $\mathfrak u(\C v)\oplus0\subseteq\mathfrak h$.
Suppose $\mathfrak u(W)\oplus0\subseteq\mathfrak h$ and $W$ is proper.
Connectedness and spanning of the orbit give an orbit vector with nonzero
components both in $W$ and $W^\perp$: otherwise the continuous projection
norm could not leave one while reaching a vector outside $W$.
Acting by $U(W)$ and choosing a unit vector in the new complementary line,
write it $ae_1+be_{r+1}$ with $a,b>0$.
The commutator of $iE_{11}$ with $iP_{ae_1+be_{r+1}}$, and its further
commutator with $iE_{11}$, give the two independent skew-Hermitian
cross matrices between these lines. Their commutator gives
$i(E_{11}-E_{r+1,r+1})$. Thus $iE_{r+1,r+1}$ also belongs to
$\mathfrak h$. Commutation with $\mathfrak u(W)$ supplies every cross
matrix from $W$ to the new line. Consequently
$\mathfrak u(W+\C e_{r+1})\oplus0\subseteq\mathfrak h$.
Induction proves the assertion. This is a finite-dimensional Lie argument,
not a classification of all compact holonomy groups.
\end{proof}

\begin{theorem}[The actual Hodge Chern connection has full $U(2)$ holonomy]
\label{thm:v9-U2}
For every sufficiently small disk $U$ about $1$,
$\operatorname{Hol}^0(D_F)=U(2)$.
\end{theorem}
\begin{proof}
Use the inherited matrix $A$ and the kernel vector $k=(4t,3,0,0)^{\mathsf T}$.
A direct rational identity gives
\begin{equation}
 \boxed{k'+Ak=(5/2,0,0,0)^{\mathsf T},\qquad
 \det\begin{pmatrix}4t&5/2\\3&0\end{pmatrix}=-15/2.}
 \label{eq:v9-krotation}
\end{equation}
Because $k$ is in the kernel of the second fundamental form, $\nabla k$
is already in $F$. Its Hodge projection is therefore unchanged:
$D_{F,\partial_t}k=(5/2)e_0$. The nonzero determinant shows that the
curvature-kernel line is not parallel.

If the restricted holonomy were reducible, its invariant line and its
orthogonal complement would give parallel line subbundles on $U$.
Curvature would be diagonal in this splitting. Since $R_F$ has rank one
and nonzero trace everywhere, its kernel would be one of these two
parallel lines; continuity fixes the choice on the connected disk.
This contradicts \cref{eq:v9-krotation}. Hence the holonomy representation
is irreducible. Curvature at any point is a nonzero multiple of $iP$ for
a rank-one projection, as an element of the holonomy Lie algebra.
\Cref{lem:v9-rankone-Lie} yields $\mathfrak u(2)$. A connected Lie subgroup
of $U(2)$ with this full Lie algebra is $U(2)$ itself.
\end{proof}

\begin{corollary}[A natural compact factor, not a flat arithmetic factor]
The projective holonomy of $F$ is $PU(2)$ and its derived holonomy is
$SU(2)$. These are generated by small contractible parameter loops and do
not require a conjecture about the global monodromy group. They do not
identify a weak-isospin bundle, a physical gauge coupling, or a
rank-two subquotient of the Tate representation.
\end{corollary}

\section{Third-order theta acquisition and full rank-three compact holonomy}
\label{sec:v9-U3}

\subsection{Only new jets are acquired}
All derivatives below are in $C_t$ at $t=1$, with V8's cycles as initial
marking. V8's auxiliary parameter $q$ is not used as a derivative coordinate.
Starting with the inherited rational matrix $A(t)$, define
\begin{equation}
 M_0=I_4,\qquad M_{j+1}=M_j'+AM_j.
 \label{eq:v9-Mrec}
\end{equation}
Then $\Pi^{(j)}(1)=\Pi(1)M_j(1)$, so the first four Taylor coefficients
of $\Pi$ require no new integration. Form the cycle blocks
$\mathcal A(s),\mathcal B(s)$ of $S^{\mathsf T}\Pi(1+s)$ and write
$\mathcal A(s)^{-1}=\sum I_ns^n$. The finite recurrences are
\begin{equation}
 I_0=\mathcal A_0^{-1},\qquad
 I_n=-I_0\sum_{j=1}^n\mathcal A_j I_{n-j},\qquad
 \tau_n=\sum_{j=0}^n\mathcal B_jI_{n-j}.
 \label{eq:v9-taujets}
\end{equation}
Thus $\tau(1+s)=\sum_{n=0}^3\tau_ns^n+O(s^4)$ has certified entries.

For $v\in\Z^2+\sigma/2$, put $E_j(v)=2\pi i v^{\mathsf T}\tau_jv$.
Expanding the theta summands gives
\begin{equation}
 \begin{aligned}
 a_{\sigma,0}&=\sum_v e^{E_0(v)},\\
 a_{\sigma,1}&=\sum_v e^{E_0(v)}E_1(v),\\
 a_{\sigma,2}&=\sum_v e^{E_0(v)}\bigl(E_2(v)+E_1(v)^2/2\bigr),\\
 a_{\sigma,3}&=\sum_v e^{E_0(v)}
       \bigl(E_3(v)+E_1(v)E_2(v)+E_1(v)^3/6\bigr).
 \end{aligned}
 \label{eq:v9-theta-jets}
\end{equation}
These are Taylor coefficients: $a^{(j)}(1)=j!a_{[:,j]}$.

\begin{lemma}[A common rational tail for the four theta-jet columns]
\label{lem:v9-jet-tail}
Suppose $\operatorname{Im}\tau_0\succeq(9/20)I$ and
$\|\tau_j\|_{\op}<1$ for $1\le j\le3$. Restrict every sum in
\cref{eq:v9-theta-jets} to $n_1,n_2\in[-N,N]\cap\Z$ before the
half-characteristic shift. Each complex Taylor coefficient then has error
at most
\begin{equation}
 \boxed{\varepsilon_N=
 302880\,2^{-\left\lfloor\frac32(N+\frac12)^2\right\rfloor}.}
 \label{eq:v9-theta-tail}
\end{equation}
In particular $\varepsilon_8=302880\,2^{-108}<10^{-27}$.
\end{lemma}
\begin{proof}
Set $r^2=\|v\|^2$ and $c=2\pi(9/20)>27/10$.
Then $|e^{E_0}|\le e^{-cr^2}$ and $|E_j|<7r^2$.
Every polynomial multiplying $e^{E_0}$ in \cref{eq:v9-theta-jets} is
bounded by $1+7r^2+49r^4+343r^6$. Since $c/2>1$ and
$\sup_{s\ge0}s^ke^{-s}\le k^k$,
\[
 (1+7r^2+49r^4+343r^6)e^{-cr^2/2}\le
 1+7+49\cdot4+343\cdot27=9465.
\]
The remaining Gaussian is bounded by $2^{-3r^2/2}$, since
$c/2>1.35>(3/2)\log2$. A full shifted one-dimensional sum is less than
four, using the weaker majorant $2^{-r^2}$. On either omitted tail the
first absolute value is at least $N+1/2$ and subsequent values increase
by one. A geometric bound gives a total one-dimensional tail at most
$4\,2^{-\lfloor(3/2)(N+1/2)^2\rfloor}$. The union of the two
coordinate tails is therefore at most $32$ times that power. Multiplication
by $9465$ proves the result. Uniform Gaussian domination on a smaller
parameter disk justifies differentiating the theta series.
\end{proof}

\begin{theorem}[Certified full osculation of the actual theta evaluation row]
\label{thm:v9-osculation}
Let $Z=(a_0,a_1,a_2,a_3)$ be the four-by-four Taylor-column matrix in
\cref{eq:v9-theta-jets}, with row order $00,01,10,11$. Then
\begin{equation}
 \boxed{
 \begin{aligned}
 0.000014769003284&<\operatorname{Re}\det Z<0.000014769003285,\\
 -0.000002346320446&<\operatorname{Im}\det Z<-0.000002346320445,\\
 2.23\,10^{-10}&<|\det Z|^2<2.24\,10^{-10}.
 \end{aligned}}
 \label{eq:v9-osculation-bounds}
\end{equation}
Consequently $a(1),a'(1),a''(1),a'''(1)$ are independent. There is no
nonzero constant $v\in\C^4$ with $a(t)^{\mathsf T}v=0$ on any
neighborhood of $1$.
\end{theorem}
\begin{proof}
Use the V8 rational interval matrix $\Pi(1)$ as a certified input, not a
new fitted value. Apply \cref{eq:v9-Mrec,eq:v9-taujets} with outward
rational arithmetic through order three. The coefficient matrices are
symmetric, and their maximum absolute row-sum upper bounds for orders
one through three are below one. Their zeroth imaginary matrix retains
the $9/20$ margin. Thus \cref{lem:v9-jet-tail} applies.
For each of the $4\cdot17^2=1156$ retained exponential values, divide the
exponent by $2^{14}$, whose complex one-norm is below $1/2$, take its Taylor
polynomial through degree $44$, and enclose the complex remainder in radius
$2/(2^{45}45!)$. Fourteen interval squarings recover the original
exponential. A rational Machin enclosure supplies $\pi$, as in V8.
Evaluate \cref{eq:v9-theta-jets} and add the tail
$302880\,2^{-108}$ to each real and imaginary component. The twenty-four
terms of the determinant formula give precisely the rational inequalities
\cref{eq:v9-osculation-bounds}; the complete interval inputs and outputs are
in the accompanying certificate and the executable formulas are specified
above.

This is a strict nonvanishing certificate. It does not infer an exact
identity from an interval containing zero. Multiplying the columns by
$0!,1!,2!,3!$ preserves nonvanishing, proving the derivative claim.
A constant vector in every kernel would annihilate these four derivative
rows, and hence would be zero.
\end{proof}

\begin{theorem}[The actual theta-source Chern connection has full $U(3)$ holonomy]
\label{thm:v9-U3}
For every sufficiently small disk $U$ about $1$,
$\operatorname{Hol}^0(D_V)=U(3)$.
\end{theorem}
\begin{proof}
The preceding osculation implies that $a,a'$ are independent at $1$.
Thus $\kappa>0$ there and, after shrinking $U$, everywhere on it.
Suppose holonomy is reducible. A nontrivial orthogonal invariant splitting
of a fibre gives parallel subbundles $V_1\oplus V_2=V$ on the simply
connected disk. Curvature is block diagonal in that splitting. Its rank
is exactly one, so its image lies in exactly one block, continuously on
$U$. The other block $V_0$ has zero curvature and is contained in the
kernel of the second fundamental form $\beta:V\to\mathcal T/V$.

The flat restriction to $V_0$ has local nonzero parallel sections. For any
such section $v$, its ambient flat derivative is the sum of its tangential
Chern derivative and its normal second fundamental part. Both vanish:
$dv=D_Vv+\beta v=0$. The second fundamental form has type $(1,0)$,
because $V$ is a holomorphic subbundle; there is no omitted normal
$(0,1)$ term. Thus $v$ is a nonzero constant vector in all the hyperplanes
\cref{eq:v9-hyperplane}, contradicting \cref{thm:v9-osculation}.
The holonomy representation is irreducible. Its Lie algebra contains the
nonzero rank-one skew-Hermitian curvature at $1$, so
\cref{lem:v9-rankone-Lie} gives $\mathfrak u(3)$ and hence $U(3)$.
\end{proof}

\begin{remark}[Why the finite theta frame does not bound this holonomy]
The ambient marked theta bundle is flat in the chosen local frame, but the
hyperplane $V_t$ moves inside it. Orthogonal projection of ambient
differentiation to that moving hyperplane is nonflat. Its holonomy need not
normalize the fixed finite Heisenberg frame and is not required to fix one
pure source ray. Therefore this theorem neither evades nor contradicts
V1--V5's fixed-ray and finite-normalizer obstructions. It constructs a
different, geometrically specified transport on the same source bundle.
\end{remark}

\section{Independent determinant curvatures and the full joint compact group}
\label{sec:v9-joint}

\subsection{The central comparison is acquired, not assumed}
The two factor holonomies do not by themselves prove their product for the
\emph{same} loops. A diagonal or determinant-linked joint subgroup must
be excluded separately. This requires a mixed comparison of curvature
functions on the same parameter disk.

\begin{proposition}[Certified curvature scales and their first real derivatives]
\label{prop:v9-curvature-numbers}
At $t=1$, let $\lambda=\Tr R_F$ and $\kappa=-\Tr R_V$. A subscript $x$
denotes differentiation along $t=1+x$ with real $x$. Then
\begin{equation}
 \begin{array}{c|rr}
 &\text{strict lower bound}&\text{strict upper bound}\\ \hline
 \lambda&0.007223718617020870&0.007223718617020871\\
 \lambda_x&-0.028255459688082396&-0.028255459688082395\\
 \kappa&0.006952530120285544&0.006952530120285545\\
 \kappa_x&-0.027264360587641939&-0.027264360587641938
 \end{array}
 \label{eq:v9-curv-table}
\end{equation}
In particular,
\begin{align}
 -0.000000503134614213
 &<\lambda\kappa_x-\kappa\lambda_x
 <-0.000000503134614212, \label{eq:v9-central-W}\\
 0.000271188496735326
 &<\lambda-\kappa
 <0.000271188496735327. \label{eq:v9-det-difference}
\end{align}
\end{proposition}
\begin{proof}
For the Hodge calculation put $u_j=\Pi^{(j)}(1)_{[:,0:2]}$ and
\[
 h=u_0^*Qu_0,\quad H=u_0^*Qu_1,\quad K=u_1^*Qu_1,\quad
 L=H^*h^{-1}H-K.
\]
Then $R_F=h^{-1}L$. The real derivatives needed are the finite products
\[
 h_x=H+H^*,\qquad H_x=K+u_0^*Qu_2,\qquad
 K_x=u_2^*Qu_1+u_1^*Qu_2,
\]
\[
 (h^{-1})_x=-h^{-1}h_xh^{-1},\qquad
 L_x=H_x^*h^{-1}H+H^*(h^{-1})_xH+H^*h^{-1}H_x-K_x.
\]
Thus $\lambda_x=\Tr((h^{-1})_xL+h^{-1}L_x)$. Only the first two
period derivatives are used. Exact reality follows from Hermitian
curvature, not a small imaginary numerical residual.

For the theta calculation use the jets already certified above and set
\[
 N=a^*a,\quad b=a^*a',\quad q=(a')^*a',\quad
 N_x=2\operatorname{Re}b,\quad b_x=q+a^*a'',\quad
 q_x=2\operatorname{Re}((a')^*a'').
\]
Differentiating $\kappa=q/N-|b|^2/N^2$ gives
\begin{equation}
 \kappa_x=\frac{q_x}{N}-\frac{qN_x}{N^2}
 -\frac{2\operatorname{Re}(\bar b b_x)}{N^2}
 +\frac{2|b|^2N_x}{N^3}.
 \label{eq:v9-kappax}
\end{equation}
The denominator is strictly positive by V8. Outward interval substitution
in these formulas proves \cref{eq:v9-curv-table}; direct interval
multiplication then proves \cref{eq:v9-central-W,eq:v9-det-difference}.
All input periods are the unchanged V8 acquisition and all newly evaluated
theta derivatives carry \cref{eq:v9-theta-tail}. No finite difference
approximation to a derivative is used.
\end{proof}

\begin{theorem}[Full product holonomy on the actual direct-sum bundle]
\label{thm:v9-joint-group}
For the stated metrics and every sufficiently small disk about $1$,
\begin{equation}
 \boxed{\operatorname{Hol}^0(D_F\oplus D_V)=U(2)\times U(3).}
 \label{eq:v9-product-group}
\end{equation}
Its derived group is $SU(2)\times SU(3)$ and its projective image is
$PU(2)\times PU(3)$. The two central directions are genuinely independent
in this normalization.
\end{theorem}
\begin{proof}
Let $\mathfrak h\subseteq\mathfrak u(2)\oplus\mathfrak u(3)$ be the
restricted joint holonomy Lie algebra. Its projections are surjective by
\cref{thm:v9-U2,thm:v9-U3}. Its derived algebra is therefore a subdirect
product of $\mathfrak{su}(2)$ and $\mathfrak{su}(3)$. A subdirect
product of two nonisomorphic simple Lie algebras is their product:
the kernels of the projections are ideals, and a proper subdirect product
with both kernels zero would identify the two simple quotients. Their
dimensions $3$ and $8$ exclude that possibility. Hence
\[
 [\mathfrak h,\mathfrak h]=\mathfrak{su}(2)\oplus\mathfrak{su}(3).
\]

The quotient by this derived algebra is detected by the two traces.
Curvature transported from $t$ to $1$ belongs to $\mathfrak h$;
unitary transport does not change either trace. Thus its central image
contains all vectors proportional to $(\lambda(t),-\kappa(t))$.
The same finite-dimensional vector subspace contains their real derivative
at $1$. The determinant of these two vectors is nonzero by
\cref{eq:v9-central-W}. Its central image is therefore all of $\R^2$.
Together with the full derived algebra this proves
$\mathfrak h=\mathfrak u(2)\oplus\mathfrak u(3)$.
The corresponding connected subgroup is the full product.
\end{proof}

\begin{corollary}[The natural determinant line is not parallel]
\label{cor:v9-no-det-reduction}
In the displayed normalization the connection on
$\det F\otimes\det V$ has curvature
$(\lambda-\kappa)dt\wedge d\bar t$, which is nonzero at the actual
fibre. Therefore no gauge change on $F\oplus V$ reduces this connection
to $S(U(2)\times U(3))$, even on a sufficiently small parameter disk.
More generally, there is no nontrivial parallel determinant-character
relation $(\det F)^m\otimes(\det V)^n$ with fixed integers $(m,n)\ne(0,0)$.
\end{corollary}
\begin{proof}
Curvature of a determinant connection is the trace of curvature. Equation
\cref{eq:v9-det-difference} excludes a parallel total determinant
trivialization. A determinant character as stated would have curvature
$m\lambda-n\kappa$. Its value and real derivative at $1$ cannot both
vanish unless $m=n=0$, by \cref{eq:v9-central-W}. Gauge conjugation does
not change determinant curvature.
\end{proof}

\begin{remark}[Scalar choices and the unresolved determinant incidence]
Multiplying the theta metric by $e^\varphi$ changes $R_V$ by
$-\partial_t\partial_{\bar t}\varphi\,I_3$. It does not change its
projective curvature or projective holonomy, but can change determinant
curvature. A separately prescribed line connection can therefore alter
central relations. Such an alteration is extra connection data, not a gauge
change of the fixed system. The determinant incidence of the companion
Standard-Model paper is a separate physical construction \cite{Duality};
it has not been identified with this arithmetic line. Matching the compact
derived group is a concrete structural result, not a derivation of
hypercharges, matter incidence, field equations, or physical gauge dynamics.
The fixed Jacobian and family are still additional arithmetic inputs to the
basic chronology loading; they have not been selected uniquely by chronology.
\end{remark}

\section{The quadratic source isomorphism is not a connection duality}
\label{sec:v9-curvature-obstruction}

There are now two explicitly constructed rank-three metric connections:
$D_V$ on the actual theta-source span, and
$D_2=\operatorname{Sym}^2D_F$ on the quadratic Hodge bundle. The V6
holomorphic second-jet map compares their underlying holomorphic bundles
after the theta-fibre line is retained. The following theorem quantifies
why it cannot simply be declared an equivalence of their metric transports.

\begin{theorem}[Projective-curvature separation, including all scalar line twists]
\label{thm:v9-projective-gap}
Write $\widehat R=R-(\Tr R/3)I$ on a rank-three fibre. At every point
where $\lambda>0$ one has
\begin{equation}
 \spec\widehat R_2=\{-\lambda,0,\lambda\},\qquad
 \spec\widehat R_V=\{-2\kappa/3,\kappa/3,\kappa/3\}.
 \label{eq:v9-projective-spectra}
\end{equation}
For every isometry $U$ between their respective positive Hermitian fibres,
\begin{equation}
 \boxed{
 \|\widehat R_2-U\widehat R_VU^{-1}\|_{\HS}^2
 \ge 2\lambda^2-2\lambda\kappa+\frac23\kappa^2
 =\frac{\lambda^2}{2}+\frac23(\kappa-3\lambda/2)^2.}
 \label{eq:v9-projective-gap}
\end{equation}
The lower bound is the exact optimum over such isometries. At $t=1$ its
square root lies in
\begin{equation}
 \boxed{(0.006011912798837785,\;0.006011912798837786).}
 \label{eq:v9-projective-gap-numeric}
\end{equation}
In particular no unitary connection isomorphism, even after scalar line
twists on either side, identifies $D_V$ with $D_2$ on a neighborhood of
$1$. No change of source-ray phases removes this obstruction.
\end{theorem}
\begin{proof}
Curvature is functorial under the symmetric-square representation. Since
$R_F$ has eigenvalues $0,\lambda$, $R_2$ has eigenvalues
$0,\lambda,2\lambda$. Proposition~\ref{prop:v9-theta-curvature} gives
$R_V$ eigenvalues $-\kappa,0,0$. Removing scalar traces proves
\cref{eq:v9-projective-spectra}. Scalar line twists change only those traces.

For Hermitian matrices with ordered spectra $r_i,s_i$, the squared
Hilbert--Schmidt norm after unitary comparison is
$\sum r_i^2+\sum s_i^2-2\sum_{ij}r_i s_j|U_{ij}|^2$.
The matrix $(|U_{ij}|^2)$ is doubly stochastic. The rearrangement inequality,
or its convex extension from permutation matrices, bounds its final sum
by $\sum r_i s_i$. Equality is attained by matching ordered eigenbases.
Substitution gives
$(\lambda-2\kappa/3)^2+(\kappa/3)^2+(\lambda-\kappa/3)^2$, which
expands to \cref{eq:v9-projective-gap}. The numerical interval follows
from \cref{eq:v9-curv-table}. A unitary connection equivalence would
intertwine curvature, contradicting its strictly positive lower bound.
Ray rephasings are included among unitary changes of frame and cannot
change the spectra.
\end{proof}

\begin{remark}[A constructive distinction rather than another static no-go]
The conclusion is stronger than V7's pointwise nonproduct-metric theorem:
it compares two actual connections, includes arbitrary scalar fibre-line
curvatures, and gives a strict positive bound in the fixed parameter
coordinate. It does not remove the second-jet algebraic isomorphism.
Instead, the same arithmetic object now supports two inequivalent compact
geometric transports, with full joint holonomy. One must specify which
transport is to be compared with an independent renewal system.
\end{remark}

\section{V9 audit, stability, and the next arithmetic--physical comparison}
\label{sec:v9-status}

\subsection{Reproducible finite evidence and its exact logical role}
The accompanying verifier takes the unchanged V8 certified period intervals,
checks their SHA256, and propagates them through explicit rational
recurrences. It checks \cref{eq:v9-krotation}, the factorized polarized
second fundamental form, the first period derivative recurrences, the
theta-jet Gaussian bounds, the determinant nonvanishing, both curvature
scales and their real derivatives, central independence, the projective
spectral bound, and the byte-for-byte preservation of V1--V8.
It also checks all $1156$ range conditions for the validated complex
exponentials. These are numerical evaluation safeguards, not $1156$
different geometric theorems.

The Lie-generation arguments, the invariant-subspace arguments, and the
connection comparison proofs are in the text. General de Rham--Betti and
theta geometry remain inherited inputs. Curvature inclusion in restricted
holonomy is the cited classical theorem. The V8 certificate can be
independently regenerated with its included verifier; V9 does not silently
replace it with floating-point central values. This is computer-assisted
verification with explicit inequalities, not a proof-assistant formalization.

\begin{proposition}[The two decisive finite tests have nonzero tolerance]
\label{prop:v9-robust}
The osculation and central independence certificates are stable under small
acquisition errors. More explicitly, if each of
$\lambda,\kappa,\lambda_x,\kappa_x$ changes by at most $10^{-6}$, their
central determinant remains nonzero. If every complex entry of the theta
Taylor-column matrix changes by at most $10^{-9}$ in absolute value, its
determinant remains nonzero.
\end{proposition}
\begin{proof}
The sum of the absolute values of the four quantities in
\cref{eq:v9-curv-table} is less than $0.07$. The change of
$\lambda\kappa_x-\kappa\lambda_x$ is at most
$0.07\varepsilon+2\varepsilon^2$ when each error is at most
$\varepsilon$. At $\varepsilon=10^{-6}$ this is less than
$7.1\,10^{-8}$, whereas the actual determinant has magnitude above
$5\,10^{-7}$.

The certified absolute values of all entries of the Taylor-column matrix
are below two. If they change by at most $\delta$, multilinearity and the
24-term determinant expansion bound the determinant change by
$24((2+\delta)^4-2^4)$. At $\delta=10^{-9}$ this is less than
$8\,10^{-7}$. The actual determinant has magnitude greater than
$\sqrt{2.23\,10^{-10}}>1.49\,10^{-5}$, so it cannot vanish.
These tolerances refer to numerical acquisition within the fixed geometric
bundle construction. Arbitrary changes of the family or the connections
still require checking the same structural hypotheses.
\end{proof}

\subsection{Proved in this continuation}
\begin{enumerate}[label=\textup{(P9.\arabic*)},leftmargin=3.6em]
\item The already acquired Hodge bundle carries an explicitly computable
      Chern connection with full local $U(2)$ holonomy. The exact reason
      is rotation of its rank-one curvature kernel, not a count of two
      basis vectors.
\item The actual vanishing theta-source bundle carries a different Chern
      connection with full local $U(3)$ holonomy. Third-order theta
      osculation is strictly certified at the actual arithmetic fibre.
\item On the same loop space their restricted joint holonomy is the full
      product $U(2)\times U(3)$. Its compact derived group has factors
      $SU(2)$ and $SU(3)$, while the two central determinant directions
      are independent in the stated theta normalization.
\item The natural total determinant curvature is strictly nonzero. No
      gauge choice of this fixed system yields a determinant-one
      reduction or a nontrivial parallel determinant-character relation.
\item The theta connection and the symmetric-square Hodge connection are
      quantitatively inequivalent even projectively. Scalar line twists
      do not remove their different curvature multiplicities.
\item All new numerical steps have explicit rational error bounds, and
      the two decisive nondegeneracy tests have stated positive tolerance.
\end{enumerate}

\subsection{Inherited inputs and unresolved comparisons}
The source spectrum, arithmetic phase class and signs, marked period matrix,
second-jet isomorphism, finite-prime Frobenius polynomials, Cartier maps,
endomorphism ring, flat symplectic monodromy, and rank-ten differential
completion remain the unchanged V1--V8 results. The chronology loading and
companion spacetime, spectral, SM, and continuum results remain antecedents
\cite{GT77,Arithmetic,Spacetime,Spectral,Duality,Variational}.

The new compact connections give genuine unitary parallel transport and
Wilson-type traces around parameter loops. They do not identify those loops
with spacetime loops, their holonomy with physical gauge transport, or the
second fundamental form with a Standard-Model Higgs field. The two bundle
ranks and a nondegenerate parameter curve are important ingredients of the
compact-group calculation; the theorem is not a universal arithmetic
selection principle for particle physics.

The next source-level interface should now compare the independently given
renewal rank-two and rank-three transport with $D_F$ and $D_V$, retaining
actual source sections, mixed Gram data, curvature, and determinant-line
incidence. It must not use the desired $S(U(3)\times U(2))$ answer to select
a scalar metric or an added line connection. The central Store/chronology
signature and the physical clock remain unacquired. Finite-prime Frobenius
still acts on its arithmetic realization; no marked crystalline matrix or
Frobenius--Wilson intertwiner has been computed here. The appearance of
compact unitary holonomy for a \emph{nonflat} connection is not such an
intertwiner.

\paragraph{Version integrity.}
All eight earlier mathematical bodies and historical status ledgers are
preserved byte for byte. Future work must append after this body. Any
correction to an inherited result must be identified explicitly rather than
silently changing the preserved text.
% ====================== END IMMUTABLE V9 MATHEMATICAL BODY ======================
% APPEND V10 AND LATER ITERATIONS HERE.

% ===================== BEGIN IMMUTABLE V10 MATHEMATICAL BODY =====================
\clearpage
\section{V10: the native cotangent determinant line, before a gauge reduction}
\label{sec:v10-context}

V9 constructed the actual Chern connections on the rank-two Hodge bundle
$F$ and the rank-three vanishing theta-source bundle $V$, with restricted
joint holonomy $U(2)\times U(3)$. It also proved that their determinant
curvatures are independent. Merely projecting that connection to a desired
Lie algebra would change the problem. We instead go upstream to the
\emph{algebraic determinant line already carried by the marked curve}.

The new calculation retains the cotangent line at the distinguished
Weierstrass point. If $\mathscr P$ denotes that line, then the actual
second-jet source map of V6 gives
\begin{equation}
 \boxed{\det F\simeq\mathscr P^4,\qquad
 \det V\simeq\mathscr P^{-6},\qquad
 \mathscr D:=\det F\otimes\det V\simeq\mathscr P^{-2}.}
 \label{eq:v10-main-lines}
\end{equation}
These are natural holomorphic, and algebraic, line isomorphisms. They are
\emph{not} identities of the previously chosen Hermitian metrics.
There are two useful consequences. A third independently marked
Weierstrass point gives an explicit determinant tensor on the original
pair. Without selecting a residual root, the existing cotangent line gives
\begin{equation}
 \boxed{W:=F\otimes\mathscr P,\qquad C:=V,\qquad
 \beta:\det C\otimes\det W\xrightarrow{\sim}\mathcal O.}
 \label{eq:v10-main-beta}
\end{equation}
The colour letter $C$ in this display names the rank-three bundle only;
we write $\mathcal C$ for the algebraic curve in the general-family
arguments. No physical colour interpretation is included.

The canonical tensor $\beta$ is obtained from conormal, evaluation and
leading-jet maps, not by requiring a preferred gauge group. With the V9
metrics on $F$ and $V$ fixed, it determines a unique metric on the
\emph{already existing} line $\mathscr P$ for which $\beta$ has unit norm.
The corresponding Chern connection has restricted holonomy
\begin{equation}
 \boxed{\operatorname{Hol}^{0}
       (D_{F\otimes\mathscr P}^{\rm bal}\oplus D_V)
       =S(U(2)\times U(3)).}
 \label{eq:v10-main-holonomy}
\end{equation}
This is a determinant-compatible \emph{completion}, with its metric rule
stated explicitly. The actual Hodge quotient metric on $\mathscr P$ does
not pass that compatibility condition; a strict numerical certificate is
proved below. Thus \cref{eq:v10-main-holonomy} is not a gauge transformation
of the V9 connection and is not asserted for every natural metric on the
new line.

\subsection{Dependency and novelty ledger}
The V1--V9 mathematical bodies are left unchanged. We use V6's second-jet
isomorphism, V8's marked period comparison, and V9's Chern formulas and
full projective joint holonomy. We do not rederive the chronology loading,
$K_4$ lattice, theta-state signs, point counts, Cartier matrices, period
series, theta-jet osculation or factor holonomies. The fully alternating
five-slot incidence and its unitary stabilizer are antecedents of the
companion duality manuscript \cite{Duality}; the new question is whether
an independently defined arithmetic tensor actually supplies that line.

The conormal exact sequence is standard algebraic geometry
\cite{StacksCartier}. The general determinant-of-theta-bundles problem has
classical precedents \cite{KouvidakisTheta}. The exact orders of the
cotangent class on stacks of Weierstrass-pointed curves are imported from
Edidin--Hu \cite{EdidinHu}, not presented as a new Chow-ring calculation.
Our contribution is the explicit specialization of these line structures
to the source map in this programme, the native determinant tensor and
its root-cover descent, the complete metric-compatibility test, and the
holonomy theorem for the resulting specified arithmetic connection.

\begin{scope}[Three assertions that must remain distinct]
The existence of a holomorphic determinant trivialization, its parallelism
for an existing unitary connection, and the physical occurrence of the
corresponding five-slot operation are different statements. V10 proves the
first after the indicated line twist (or extra marking), tests the second,
and constructs a unique determinant-compatible metric completion. It does
not prove the third. The base is still an arithmetic parameter disk, not
reconstructed spacetime; the original curve and family are not selected by
chronology alone.
\end{scope}

\section{The exact cotangent weight of the arithmetic source}
\label{sec:v10-lines}

Let $f:\mathcal C\to S$ be a smooth proper family of genus-two curves in
characteristic zero, with a Weierstrass section $w$. Let $J\to S$ be its
Jacobian, $i(P)=[P-w]$, $\Theta=i(\mathcal C)$, and
$L=\mathcal O_J(\Theta)$. We retain this divisor-defined line, rather
than replacing it by a line rigidified to be trivial at zero. Put
\[
 F=f_*\omega_{\mathcal C/S},\quad
 \mathscr P=w^*\omega_{\mathcal C/S},\quad
 \Lambda_\Theta=0^*L,\quad
 \mathcal T=\pi_*L^2,\quad
 V=\ker(\mathcal T\to\Lambda_\Theta^2).
\]
Here $\pi:J\to S$. All powers of lines are tensor powers. The letter
$\Lambda_\Theta$ is V6's retained theta-fibre line $\lambda$, renamed to
avoid confusion with V9's curvature eigenvalue $\lambda(t)$.

\begin{lemma}[The two successive Weierstrass evaluation jets]
\label{lem:v10-eval}
Let $\ell=\ker(F\to\mathscr P)$ be the kernel of evaluation at $w$.
There are canonical exact sequences and isomorphisms
\begin{equation}
 0\longrightarrow\ell\longrightarrow F\longrightarrow\mathscr P
 \longrightarrow0,\qquad
 \ell\xrightarrow{\;j^2_w\;}\mathscr P^3,\qquad
 \det F\xrightarrow{\sim}\mathscr P^4.
 \label{eq:v10-eval}
\end{equation}
The middle map is the leading coefficient of a differential vanishing at
$w$, not a chosen coordinate normalization.
\end{lemma}
\begin{proof}
On a genus-two fibre, the canonical system has no base point. Evaluation
is therefore onto a one-dimensional space, and its kernel is a line.
At a Weierstrass point every differential in that kernel vanishes to
order two. One can see this either from the canonical divisor $2w$ or
from a local hyperelliptic coordinate: the two holomorphic differentials
have leading terms $du$ and $u^2du$. The order cannot exceed two for a
nonzero holomorphic differential, since its divisor has degree two.
The leading coefficient is consequently a nonzero element of
\[
 (\mathfrak m_w^2/\mathfrak m_w^3)\otimes
 \omega_{\mathcal C/S}|_w\simeq\mathscr P^3.
\]
It is independent of the local parameter and is nonzero on every fibre.
Cohomology and base change, or local constant-rank evaluation, make the
kernel a line subbundle and the leading-jet map an isomorphism of lines.
Taking the determinant of the exact sequence proves the last identity.
Every map used here commutes with an isomorphism of pointed families.
\end{proof}

\begin{lemma}[The theta-fibre line is not an extra free scalar line]
\label{lem:v10-adjunction}
There is a canonical isomorphism
\begin{equation}
 \boxed{\Lambda_\Theta\simeq\ell^{-1}
 \simeq\mathscr P^{-3}
 \simeq\mathscr P\otimes(\det F)^{-1}.}
 \label{eq:v10-adjunction}
\end{equation}
\end{lemma}
\begin{proof}
The Abel--Jacobi image is a smooth relative divisor through zero. Its
conormal sequence, restricted to $w$, is
\[
 0\longrightarrow\Lambda_\Theta^{-1}
 \longrightarrow 0^*\Omega^1_{J/S}=F
 \longrightarrow\mathscr P\longrightarrow0.
\]
The last map is evaluation of invariant differentials along the
Abel--Jacobi tangent, hence the same map as in \cref{eq:v10-eval}.
The conormal line is its kernel. This proves
$\Lambda_\Theta^{-1}\simeq\ell$, and the preceding lemma gives the
remaining assertions. Equivalently, relative adjunction gives
$\Lambda_\Theta\otimes\det F\simeq\mathscr P$.
The canonical section of $\mathcal O_J(\Theta)$ realizes the conormal
map; choosing an unrelated scalar multiple of an analytic theta function
without changing the line identification would not be the same marking.
\end{proof}

\begin{theorem}[Native determinant identity for the actual theta source]
\label{thm:v10-line}
For the actual vanishing theta-source bundle,
\begin{align}
 \det V&\xrightarrow{\sim}\mathscr P^{-6},&
 \det\mathcal T&\xrightarrow{\sim}\mathscr P^{-12},\notag\\
 \mathscr D:=\det V\otimes\det F
   &\xrightarrow{\sim}\mathscr P^{-2}.&&
 \label{eq:v10-det}
\end{align}
These isomorphisms are functorial in the pointed curve and defined over its
field of definition. They therefore apply to the rational curve of this
programme before an analytic period matrix or a Hermitian metric is chosen.
\end{theorem}
\begin{proof}
Use the inherited, natural V6 isomorphism
\[
 V\xrightarrow{\mathfrak j_2}
 \Lambda_\Theta^2\otimes\operatorname{Sym}^2F.
\]
For a rank-two bundle, the determinant of its symmetric square is the
third power of its determinant. In a local $2\times2$ transition matrix
this is the polynomial identity
$\det(\operatorname{Sym}^2g)=(\det g)^3$; it fixes the line map in the
ordinary monomial convention used in V6. Hence
\[
 \det V\simeq\Lambda_\Theta^6\otimes(\det F)^3
 \simeq\mathscr P^{-18}\otimes\mathscr P^{12}
 =\mathscr P^{-6}.
\]
The evaluation exact sequence
$0\to V\to\mathcal T\to\Lambda_\Theta^2\to0$
gives $\det\mathcal T\simeq\mathscr P^{-6-6}$.
Multiplying $\det V$ by \cref{eq:v10-eval} gives the final identity.
The construction uses only maps already proved natural, so it descends
with the pointed algebraic curve.
\end{proof}

\begin{remark}[This is not an identity of the V9 determinant curvatures]
A holomorphic line isomorphism becomes an identity of Chern connections
only when its metrics are transported as well. For example,
$\det F\simeq\mathscr P^4$ does not assert that V9's curvature trace
is four times the curvature of the Hodge \emph{quotient} metric on
$\mathscr P$. The strict values in \cref{sec:v10-metrics} explicitly
exclude that inference. The same caveat applies to adjunction and the
second-jet line factors.
\end{remark}

\begin{proposition}[An automorphism witnesses the original determinant obstruction]
\label{prop:v10-isotropy}
On the two-pointed reference curve
$Y^2=X^5-X$ with $w=\infty$ and $w_0=0$, let
\[
 R(X,Y)=(iX,\zeta_8Y),\qquad\zeta_8=e^{\pi i/4}.
\]
The eigencharacters of $R^*$ on $F$ are $(\zeta_8,\zeta_8^3)$;
its character on $\mathscr P$ is $\zeta_8^3$, and on $\mathscr D$ it
is $\zeta_8^2=i$. Consequently the original determinant has no
trivialization natural under all two-pointed isomorphisms.
\end{proposition}
\begin{proof}
Pulling back $dX/Y$ and $X\,dX/Y$ gives the stated characters. The
rational uniformizer $u=X^2/Y$ satisfies $R^*u=\zeta_8^3u$, so its
cotangent character is $\zeta_8^3$. The determinant identity gives
$(\zeta_8^3)^{-2}=\zeta_8^2$. A natural nonzero determinant
trivialization would be fixed by this automorphism, contrary to its
nontrivial character. This is a statement on the marked moduli problem,
not a claim that a one-dimensional fibre or a small disk is abstractly
nontrivializable.
\end{proof}

\section{Two rational Weierstrass marks, the order-four line, and the root cover}
\label{sec:v10-cover}

The original arithmetic data include both $w=\infty$ and $w_0=0$.
It would overstate the global obstruction to forget the second marking.
For reference, the exact integral Chow calculation of Edidin--Hu says
\cite[Theorems 1.3--1.5]{EdidinHu}, in genus two and characteristic zero,
\[
 A^*(\mathcal H_{2,1}^{w})=\Z[p]/(40p),\quad
 A^*(\mathcal H_{2,2}^{w})=\Z[p]/(8p),\quad
 A^*(\mathcal H_{2,n}^{w})=\Z[p]/(2p)\quad(3\le n\le6),
\]
where $p=c_1(\mathscr P)$. These are smooth stacks, so their degree-one
Chow groups identify the Picard groups. Combining this imported theorem
with \cref{thm:v10-line} gives the following \emph{source-line}
consequence.

\begin{corollary}[Exact order of the determinant class]
\label{cor:v10-order}
The line $\mathscr D$ has exact orders $20$, $4$, and $1$ on the
one-, two-, and at-least-three-Weierstrass-marked genus-two stacks,
respectively. In particular the original two-marked arithmetic source has
an order-four determinant obstruction; adding an independently specified
residual Weierstrass point removes its algebraic line obstruction.
\end{corollary}
\begin{proof}
The class is $-2p$ by \cref{thm:v10-line}. Its additive order in a
cyclic group of order $N$ is $N/\gcd(N,2)$, yielding the displayed
orders. The Chow/Picard-group computation is the cited input, not a
finite-matrix verification performed here.
\end{proof}

\begin{theorem}[An explicit root-cover determinant tensor]
\label{thm:v10-root-section}
For a third marked Weierstrass point $w_r$, let $g_r$ be the unique
meromorphic function with
\[
 \operatorname{div}(g_r)=2w-2w_0,\qquad g_r(w_r)=1.
\]
Its leading quadratic term at $w$ is a nowhere-zero section
\begin{equation}
 \boxed{\sigma_r=j_w^2(g_r)\in\mathscr P^2
   \simeq\mathscr D^*.}
 \label{eq:v10-sigma}
\end{equation}
This is a determinant tensor on the \emph{original} pair $F,V$, defined
on the root cover without choosing a phase or a metric. The four possible
residual roots give four covariant tensors, not a preferred rational one.
Their product descends to a canonical nonzero section of
$\mathscr P^8\simeq\mathscr D^{*4}$.
\end{theorem}
\begin{proof}
The hyperelliptic quotient identifies $2w$ and $2w_0$, so a meromorphic
function with the stated divisor exists and is unique up to scalar.
Its nonzero finite value at the distinct third branch point fixes that
scalar. Since its zero at $w$ has order exactly two, its leading term
is a nonzero element of $\mathfrak m_w^2/\mathfrak m_w^3=\mathscr P^2$.
This is parameter independent and functorial in the three marked points.
The line identification in \cref{thm:v10-line} then gives a nowhere-zero
functional on $\mathscr D$. The four-fold product is invariant under
permutations of the four residual marks and therefore descends. This is
ordinary finite descent of a symmetric tensor product, not averaging the
four linear functionals, whose sum need not be nonzero.
\end{proof}

For the actual family $y^2=x(x^4-x-t)$, put $u=x^2/y$ at infinity and
$e=du|_\infty$. Then
\begin{equation}
 \boxed{g_r=r/x,\qquad\sigma_r=r e^2,\qquad
 \prod_{r^4-r-t=0}\sigma_r=-t\,e^8.}
 \label{eq:v10-root-product}
\end{equation}
Indeed $u^2=x^{-1}+O(x^{-4})$, and the product of the four roots is
$-t$. Under a monic coordinate change preserving the two marked points,
$x'=a^2x$, $y'=a^5y$, one has $r'=a^2r$ and $e'=a^{-1}e$;
therefore both $r e^2$ and $(\prod r)e^8$ are unchanged. An arbitrary
monic-model isomorphism preserving $0$ and infinity has this form: writing
$x'=c x$, $y'=d y$ gives $d^2=c^5$, and $a=d/c^2$ has
$a^2=c$, $a^5=d$.

\begin{scope}[Relation to the earlier quartic frame cover]
The older arithmetic--spacetime notes already introduced a quartic root
cover to select a residual Store frame \cite{AD9}. That frame-cover
calculation is not repeated. The new statement is that the same extra
\emph{arithmetic marking} gives the explicit determinant tensor
\cref{eq:v10-sigma}. It still does not identify the arithmetic root with
an independently given physical Store label, and it does not make that
tensor parallel for V9's fixed connection. A fixed complex marking, as
used in V8, is not a rational descent of a preferred root.
\end{scope}

\begin{proposition}[A native finite-flat line metric, and its residual character]
\label{prop:v10-flat-line}
The descended section $s_8=\prod_r\sigma_r$ gives a unique Hermitian
metric on $\mathscr P$ for which $\|s_8\|=1$. Its Chern connection
is flat on the smooth two-marked base. On the actual $t$-family,
\begin{equation}
 \boxed{\|e\|^2_{\rm ff}=|t|^{-1/4},\qquad
 A_{\mathscr P}^{\rm ff}=-\frac18\frac{dt}{t}.}
 \label{eq:v10-flat-metric}
\end{equation}
A positive loop about $t=0$ has cotangent holonomy $e^{2\pi i/8}$ and
induced $\mathscr D$ holonomy $e^{-2\pi i/4}$. Thus the flat determinant
character on this family has order exactly four. This connection is not
the V9 determinant connection.
\end{proposition}
\begin{proof}
In a holomorphic frame with $s_8=c e^8$, the norm condition is
$|c|^2\|e\|^{16}=1$, giving $\|e\|^2=|c|^{-1/4}$.
Its logarithm is locally the real part of a holomorphic function, so its
mixed derivative vanishes. For $c=-t$, the Chern coefficient is
$\partial\log\|e\|^2=-dt/(8t)$. A horizontal coefficient acquires
the exponential of the negative integral of this form. Raising to the
power $-2$ gives the determinant holonomy. Its value has order four,
whereas the fourth power is trivial by the global section $s_8$.
\end{proof}

\begin{remark}[The one-mark discriminant has a different weight]
If only infinity is marked, the usual quintic discriminant yields a
section of $\mathscr P^{40}$. On this equation it is
$-t^2(256t^3+27)e^{40}$. The two-mark section $-t e^8$ is stronger
because the second point is retained. Neither nonzero high tensor power
provides a chosen nonzero section of $\mathscr P^2$ on the original
moduli problem. Holomorphic torsion, a flat finite-character metric,
and the nonflat V9 metric are distinct data.
\end{remark}

\section{A root-free arithmetic determinant incidence on the cotangent twist}
\label{sec:v10-incidence}

\begin{theorem}[The native determinant-compensating bundle]
\label{thm:v10-beta}
Set $W=F\otimes\mathscr P$ and retain $C=V$ with its actual theta
source. There is a canonical nowhere-zero algebraic functional
\begin{equation}
 \boxed{\beta:\det C\otimes\det W\xrightarrow{\sim}\mathcal O_S.}
 \label{eq:v10-beta}
\end{equation}
It is defined with only the original marked infinity and descends over the
rational curve; no residual root or square root of a modular form is
selected. It gives a nonzero five-slot incidence
\begin{equation}
 T_\beta\in\operatorname{Hom}
 (C\otimes W\otimes W,\Lambda^2C^*),
 \quad
 T_\beta(c,w_1,w_2)(c_2,c_3)
 =\beta(c\wedge c_2\wedge c_3,\,w_1\wedge w_2).
 \label{eq:v10-incidence}
\end{equation}
\end{theorem}
\begin{proof}
The determinant of $W$ is
$\det F\otimes\mathscr P^2\simeq\mathscr P^6$.
Pair it with $\det V\simeq\mathscr P^{-6}$ using the exact line maps
of \cref{thm:v10-line} and ordinary evaluation of a line and its dual.
That composition is $\beta$. Each constituent is natural, so the
functional is defined over the original field, not merely after choosing
an analytic trivialization. It is everywhere nonzero because each line
map is an isomorphism. Substitution into the displayed alternating formula
then gives a nonzero incidence. This formula is the standard determinant
incidence of the companion manuscript \cite{Duality}; what is newly
supplied here is its arithmetic source functional.
\end{proof}

\begin{proposition}[What is, and is not, minimal in this twist]
\label{prop:v10-twist}
For twists $F\otimes\mathscr P^a$ and $V\otimes\mathscr P^b$,
their determinant has cotangent weight
\begin{equation}
 \boxed{-2+2a+3b.}\label{eq:v10-weight}
\end{equation}
Cancelling this weight using the jet and conormal maps alone gives
$a=1+3k$, $b=-2k$, $k\in\Z$. Keeping the actual rank-three source
bundle untwisted selects $a=1$, $b=0$ in this zero-weight family.
On the original two-marked moduli stack, abstract line-class triviality is
weaker: it requires $2a+3b\equiv2\pmod8$ by the imported Picard calculation.
\end{proposition}
\begin{proof}
A rank-$r$ bundle twisted by a line $M$ changes its determinant by $M^r$.
Use \cref{thm:v10-line}. The integer solutions of $2a+3b=2$ have $b$
even; writing $b=-2k$ gives the stated family. The congruence follows from
the exact order eight of $\mathscr P$. These are statements within the
specified class of line twists, not a universal minimization over all
possible bundles or new arithmetic inputs.
\end{proof}

\begin{proposition}[A local coefficient formula retaining the actual theta scalar]
\label{prop:v10-coefficient}
At the actual family use $\omega_0=dx/y$, $\omega_1=x\,dx/y$, and
$\eta=\omega_1|_\infty$ as a local frame of $\mathscr P$. Choose a
local frame $v_1,v_2,v_3$ of $V$ and a local frame $l$ of
$\Lambda_\Theta$. Let $J$ be the matrix of $\mathfrak j_2$ in the
bases $(v_i)$ and $l^2(\omega_0^2,\omega_0\omega_1,\omega_1^2)$.
Write the canonical theta section's first term as $c\,l\omega_0$;
then $c\ne0$, and the local coefficient of \cref{eq:v10-beta} is
\begin{equation}
 \boxed{B_\beta=
 \beta(v_1\wedge v_2\wedge v_3,
       (\omega_0\otimes\eta)\wedge(\omega_1\otimes\eta))
       =\frac{16\det J}{c^6}.}
 \label{eq:v10-beta-coefficient}
\end{equation}
Equivalently $c^2$ is the $\omega_0^2$ coefficient of
$\mathfrak j_2(\theta^2)$, so no square-root phase needs to be chosen.
The formula uses the canonical divisor section and the corresponding
line identification; an arbitrary rescaling of a source ray alone is not
an allowed replacement for that section.
\end{proposition}
\begin{proof}
At infinity, with $u=x^2/y$, the leading terms are
$\omega_0=-2u^2du+\cdots$ and $\omega_1=-2du+\cdots$.
Consequently $j^2_w(\omega_0)=\eta^3/4$ and
$\omega_0\wedge\omega_1\mapsto\eta^4/4$ under
\cref{eq:v10-eval}. The conormal map sends $l^{-1}$ to $c\omega_0$,
so $l\mapsto(4/c)\eta^{-3}$. Taking the determinant of $J$ gives
\[
 v_1\wedge v_2\wedge v_3
 \longmapsto(\det J)(4/c)^6(1/4)^3\eta^{-6}.
\]
The determinant of the $W$ frame maps to $\eta^6/4$. Multiplication
leaves $16\det J/c^6$. Squaring the first theta term proves the
alternative acquisition of $c^2$.
\end{proof}

\begin{scope}[An arithmetic operation, not historical physical occurrence]
The tensor $T_\beta$ is a genuine algebraic bundle morphism. It is not
identified with a renewal history word, a Yukawa matrix, or an accepted
physical operation. Its coordinate-free source comes from the arithmetic
jet and conormal maps. The companion manuscript's historical-ancestry and
common-carrier requirements are still required to identify it with a
physical incidence. This separation prevents an algebraic construction
from silently supplying a missing operational event.
\end{scope}

\section{The actual quotient metric fails the determinant compatibility test}
\label{sec:v10-metrics}

Return to V9's small complex disk about $t=1$. Fix its Hodge metric $h$
on $F$ and its restricted Heisenberg-unitary metric $g$ on $V$ exactly as
before. For every Hermitian bundle we use V9's convention
\[
 F_D=R_D\,dt\wedge d\bar t,\qquad
 R_D=-\partial_{\bar t}(h^{-1}\partial_t h).
\]
Thus $\operatorname{Tr}R_F=\lambda>0$ and
$\operatorname{Tr}R_V=-\kappa<0$, with both functions and their
strict bounds inherited from V9.

The evaluation quotient $F\twoheadrightarrow\mathscr P$ gives a
geometrically independent candidate metric on $\mathscr P$. In the frame
$\eta=\omega_1|_\infty$ its squared norm is
\begin{equation}
 \boxed{h_{\mathscr P}^{Q}=\frac{\det h}{h_{00}}.}
 \label{eq:v10-quotient-metric}
\end{equation}
Indeed the least norm lift of the coset of $\omega_1$ is its orthogonal
projection away from $\Span(\omega_0)$; the scalar Schur complement is
the displayed ratio. This is not a new choice of a source metric.

\begin{proposition}[An exact curvature test requiring no new periods]
\label{prop:v10-quotient-curvature}
Put $a=h_{00}$, $b=(h_t)_{00}$, $c=(h_{t\bar t})_{00}$ and let
$r_\ell$ be the curvature coefficient of the restricted Hodge metric on
$\ell=\Span(\omega_0)$. Then
\begin{equation}
 \boxed{r_\ell=-\frac ca+\frac{|b|^2}{a^2},\qquad
 r_Q:=R_{\mathscr P}^{Q}=\lambda-r_\ell.}
 \label{eq:v10-quotient-curvature}
\end{equation}
The total determinant curvature of the pair
$(F\otimes\mathscr P,h\otimes h_{\mathscr P}^{Q})$ and $(V,g)$ is
\begin{equation}
 \boxed{\delta_Q=\lambda-\kappa+2r_Q.}
 \label{eq:v10-deltaQ}
\end{equation}
At the actual fibre the following strict rational enclosures hold:
\begin{align}
 0.010875118695194633777&<r_Q<0.010875118695194633778,
       \label{eq:v10-rQ-number}\\
 0.022021425887124593831&<\delta_Q<0.022021425887124593832.
       \label{eq:v10-deltaQ-number}
\end{align}
In particular the canonical algebraic tensor $\beta$ is not parallel for
this original quotient-metric extension.
\end{proposition}
\begin{proof}
The curvature of a line of squared norm $a$ is
$-\partial_t\partial_{\bar t}\log a$. This gives $r_\ell$.
Taking the logarithm of \cref{eq:v10-quotient-metric} gives
$r_Q=\lambda-r_\ell$. Tensoring a rank-two connection by a line of
curvature $r_Q$ adds $2r_Q$ to its trace, proving \cref{eq:v10-deltaQ}.

For the acquisition, no numerical differentiation is used. V9's certified
matrices $h,A_F,R_F$ at $1$ give the exact identities
\[
 h_t=hA_F,\qquad
 h_{t\bar t}=h_t^*h^{-1}h_t-hR_F=h_t^*A_F-hR_F.
\]
Substitute their rational complex interval enclosures in
\cref{eq:v10-quotient-curvature}. All divisions have the inherited
strict denominator $h_{00}>17$. The standalone verifier rounds every
rational operation outward to a grid of size $10^{-70}$ and proves the
two displayed enclosures. It also obtains
\[
 17.436908870047675674811<h_{\mathscr P}^{Q}(1)
 <17.436908870047675674812.
\]
Reality follows from the Hermitian metric identities, not merely from
an imaginary enclosure containing zero. A nonzero determinant curvature
precludes a nonzero parallel determinant tensor.
\end{proof}

\begin{corollary}[The native finite-flat metric also fails the same test]
\label{cor:v10-flat-fails}
For the line metric of \cref{prop:v10-flat-line}, the determinant
curvature of $F\otimes\mathscr P$ together with $V$ is
$\lambda-\kappa\ne0$. Thus neither the Hodge quotient metric nor the
metric selected by the algebraic section $s_8$ makes $\beta$ parallel.
\end{corollary}
\begin{proof}
The finite-flat metric has line curvature zero. Apply the trace formula
and V9's strict positive bound for $\lambda-\kappa$.
\end{proof}

\section{The unique determinant-compatible metric and its full compact holonomy}
\label{sec:v10-balanced}

\begin{theorem}[A canonical metric once determinant compatibility is required]
\label{thm:v10-balanced-metric}
Fix the actual metrics $h$ on $F$ and $g$ on $V$, together with the
arithmetic functional $\beta$ of \cref{thm:v10-beta}. There is exactly
one Hermitian metric $h_{\mathscr P}^{\rm bal}$ on $\mathscr P$ for
which the induced norm of $\beta$ is one. It is defined without a new
choice of a line bundle, a branch root, or a scalar curvature function.
For any preliminary line metric $h_0$, if
$q_0=\|\beta\|^2_{h,g,h_0}$, then
\begin{equation}
 \boxed{h_{\mathscr P}^{\rm bal}=q_0^{1/2}h_0.}
 \label{eq:v10-balanced-metric}
\end{equation}
This formula is independent of $h_0$. The induced Chern connections make
$\beta$ parallel. In local holomorphic frames where its coefficient is
$B_\beta$, the connection on $\mathscr P$ is
\begin{equation}
 \boxed{A_{\mathscr P}^{\rm bal}
 =\frac12\bigl(d\log B_\beta-\Tr A_F-\Tr A_V\bigr),\qquad
 R_{\mathscr P}^{\rm bal}=\frac{\kappa-\lambda}{2}.}
 \label{eq:v10-balanced-connection}
\end{equation}
Here $d\log B_\beta$ is the holomorphic differential of the nonzero
coefficient, so the formula does not require choosing a logarithm globally.
\end{theorem}
\begin{proof}
Multiplying a line metric by $e^\varphi$ multiplies the determinant
metric of $F\otimes\mathscr P$ by $e^{2\varphi}$. The squared norm
of the dual determinant functional $\beta$ is therefore multiplied by
$e^{-2\varphi}$. Precisely one positive scalar at each point makes its
squared norm equal to one, giving \cref{eq:v10-balanced-metric}.
If the preliminary metric changes to $e^\varphi h_0$, then
$q_0$ changes to $e^{-2\varphi}q_0$, and the resulting metric is
unchanged. Smoothness and positivity are immediate.

A nowhere-zero holomorphic section of constant norm is parallel for the
Chern connection on a Hermitian line. Applied to $\beta$, this makes
its determinant functional parallel. In local frames its covariant
constancy reads
$dB_\beta-B_\beta(\Tr A_F+\Tr A_V+2A_{\mathscr P})=0$,
which proves the connection formula. Taking curvature gives
$\lambda-\kappa+2R_{\mathscr P}=0$.
The formula glues because it is the Chern connection of the globally
constructed metric. Equivalently, use the canonical isomorphism
$\mathscr P^2\simeq\mathscr D^{-1}$ and take the unique metric and
connection square root on the \emph{specified} line $\mathscr P$.
No abstract choice among square-root line bundles is involved.
\end{proof}

\begin{scope}[The compatibility rule is not a conclusion about the old metric]
The requirement that the independently derived tensor $\beta$ be a unit
parallel incidence is an explicit geometric normalization. It is enough to
determine the additional line metric, but it is not implied by Heisenberg
invariance or by the original Hodge quotient. Thus the construction is
stronger than choosing a scalar adjustment without a source, and weaker
than a proof that an unspecified physical process selects that adjustment.
All subsequent holonomy claims refer to this declared completion.
\end{scope}

\begin{proposition}[The required modification has a strict, acquired size]
\label{prop:v10-correction-size}
At $t=1$,
\begin{align}
 -0.000135594248367663139
 &<R_{\mathscr P}^{\rm bal}
 <-0.000135594248367663138,\label{eq:v10-bal-number}\\
 0.011010712943562296915
 &<\partial_t\partial_{\bar t}
       \log(h_{\mathscr P}^{\rm bal}/h_{\mathscr P}^{Q})
 <0.011010712943562296916.\label{eq:v10-correction-number}
\end{align}
In particular the determinant-compatible metric is not the Hodge quotient
metric, a constant rescaling of it, or a unitary gauge change of it.
The rank-two curvature of the balanced bundle has eigenvalues
\begin{equation}
 \boxed{\frac{\kappa-\lambda}{2},\qquad
        \frac{\kappa+\lambda}{2}.}
 \label{eq:v10-balanced-spectrum}
\end{equation}
They have opposite signs at the actual fibre; the balanced bundle is not
being asserted to have the original Griffiths-semipositive Hodge metric.
\end{proposition}
\begin{proof}
The first enclosure is the rational subtraction of the inherited
$\lambda,\kappa$ bounds. For any two line metrics, their curvature
difference is minus the mixed derivative of the logarithm of their ratio.
Hence the middle expression in \cref{eq:v10-correction-number} is
$r_Q-(\kappa-\lambda)/2=\delta_Q/2$, already acquired above.
A constant rescaling or unitary gauge change has zero such curvature
difference. Scalar tensoring shifts each eigenvalue of $R_F$, which is
$(0,\lambda)$, by $(\kappa-\lambda)/2$. V9 certified
$\lambda>\kappa>0$, proving the final assertion.
\end{proof}

\begin{theorem}[Full determinant-one compact holonomy on the arithmetic disk]
\label{thm:v10-SU}
For V9's fixed theta and Hodge metrics and the unique line metric of
\cref{thm:v10-balanced-metric}, let
\[
 D_W=D_F\otimes1+1\otimes D_{\mathscr P}^{\rm bal},\qquad C=V.
\]
On every sufficiently small complex disk about $t=1$,
\begin{equation}
 \boxed{\operatorname{Hol}^{0}(D_W\oplus D_V)
       =S(U(2)\times U(3)).}
 \label{eq:v10-SU}
\end{equation}
The determinant incidence $T_\beta$ is parallel and nonzero. Its group
has the global compact presentation
\begin{equation}
 \boxed{S(U(2)\times U(3))\simeq
 \frac{SU(2)\times SU(3)\times U(1)}{\mu_6}.}
 \label{eq:v10-Z6}
\end{equation}
\end{theorem}
\begin{proof}
The parallel unit functional $\beta$ puts the holonomy inside the stated
determinant-one subgroup. A scalar line twist leaves the projective
connections unchanged. V9's proved full joint holonomy therefore implies
that the projective image of the new holonomy is
$PU(2)\times PU(3)$. Let $\mathfrak h$ be its compact Lie algebra.
Its derived algebra maps onto
$\mathfrak{su}(2)\oplus\mathfrak{su}(3)$. The kernel of that map is
central, and a compact semisimple Lie algebra has zero-dimensional centre.
Thus
\[
 [\mathfrak h,\mathfrak h]=\mathfrak{su}(2)\oplus\mathfrak{su}(3)
\]
as the standard traceless blocks. The total determinant relation leaves
at most one central direction. Curvature of the new pair has traces
\[
 (\Tr R_W,\Tr R_V)=(\kappa,-\kappa),
\]
with $\kappa(1)>0$. A real tangent-plane curvature endomorphism differs
from these Hermitian coefficients only by the common nonzero imaginary
factor. The classical curvature-in-holonomy theorem \cite{AmbroseSinger}
therefore supplies a nonzero central direction. Hence
$\mathfrak h=\mathfrak s(\mathfrak u(2)\oplus\mathfrak u(3))$.
The connected subgroup with this Lie algebra, already contained in the
connected determinant-one group, is the whole group.

For the last presentation use
\[
 (A,B,z)\longmapsto(z^3A,z^{-2}B).
\]
Its determinants multiply to one. It is onto by choosing a sixth root of
the determinant of the rank-two block; its kernel has $z^6=1$,
$A=z^{-3}I_2$ and $B=z^2I_3$, and is exactly $\mu_6$.
Finally \cref{eq:v10-incidence} is a natural tensor operation on
$\beta$, so a connection preserving $\beta$ preserves $T_\beta$.
This global group presentation is the familiar one used in the companion
manuscript; the new content is its realization as the holonomy of this
specified arithmetic completion.
\end{proof}

\begin{proposition}[Nonzero normalized incidence, and scalar-calibration scope]
\label{prop:v10-incidence-norm}
Give exterior powers the induced metric with orthonormal ordered wedge
bases. Then $\|\beta\|=1$ implies
\begin{equation}
 \boxed{T_\beta T_\beta^*=2I_{\Lambda^2C^*},\qquad
        \|T_\beta\|_{\HS}^2=6.}
 \label{eq:v10-incidence-norm}
\end{equation}
If the fixed metric on $V$ is instead multiplied by $e^\varphi$, while
$h$ is kept, its determinant-compatible cotangent metric is multiplied by
$e^{-3\varphi/2}$. The projective connections are unchanged, but the
surviving central curvature changes from $\kappa$ to
$\kappa+3\partial_t\partial_{\bar t}\varphi$. Therefore the full
$U(1)$ holonomy claim includes the fixed scalar normalization of V9; it
is not a scalar-calibration-independent group-selection theorem.
\end{proposition}
\begin{proof}
Choose oriented orthonormal bases in which $\beta$ is the unit volume
functional. For each of the three output wedge basis elements, the only
nonzero inputs are the complementary colour vector and the two orders of
the weak pair. Their coefficients have moduli one, opposite weak signs,
and disjoint supports. Hence the row Gram is $2I_3$.
A scalar change of $g$ multiplies its determinant by $e^{3\varphi}$,
so the squared norm of the dual functional is multiplied by
$e^{-3\varphi}$. Formula \cref{eq:v10-balanced-metric} gives the
cotangent change. Its curvature increases by
$\tfrac32\varphi_{t\bar t}$, whereas the rank-three trace decreases
by $3\varphi_{t\bar t}$. Their total remains zero and the displayed
central coefficient follows.
\end{proof}

\section{Descent, parity, and the remaining mixed arithmetic--renewal comparison}
\label{sec:v10-scope}

\begin{proposition}[The completion keeps the earlier weight-one obstruction intact]
\label{prop:v10-parity}
Hyperelliptic inversion acts as $-I$ on both $F$ and $\mathscr P$, and
therefore as $+I$ on $W=F\otimes\mathscr P$. It acts as $+I$ on
$V$ with the square theta linearization. The determinant incidence is
inversion-even. Accordingly this completion does not restore the central
quadratic-twist character that V5 proved invisible in the projective
squared-theta source, and it is not a small-rank flat Tate subquotient.
\end{proposition}
\begin{proof}
In a hyperelliptic local parameter at a ramification point inversion
sends $u$ to $-u$, so it acts by $-1$ on the cotangent. It acts by
$-1$ on holomorphic one-forms as well. Their tensor product is even.
The action on $V$ is the restriction of the even square-theta action
proved in V6. The tensor formula for $\beta$ and $T_\beta$ then
preserves inversion. No flat connection or new rational Frobenius
realization is introduced by these holomorphic tensor products and Chern
metrics; the V5 and V6 obstructions consequently have not been bypassed.
\end{proof}

There is now an explicit arithmetic source of the determinant-incidence
\emph{type} used in the companion paper. Its identification with the
physical incidence is still a separate map on a common represented
carrier. For clarity, the remaining comparison would require independently
specified maps from the arithmetic $W,V$ to the physical weak and colour
multiplicity bundles, an identification of their source sections and
metrics, and an intertwining of their actual connections and determinant
tensors. A mere isomorphism of the compact groups does not provide any of
those maps. In particular $T_\beta$ is not silently substituted for a
measured historical up-incidence word.

The new completion also does not contradict V7--V9's separation of the
rank-three source connection from the symmetric-square Hodge connection.
We keep the original theta metric and connection on $V$ unchanged; a
scalar twist of $F$ does not change that earlier traceless-curvature
obstruction. The original rank-ten flat differential completion and its
unipotent monodromy remain as before. The present holonomy is nonflat and
unitary, and is defined on the deformation base.

The useful next comparison is therefore no longer an unspecified scalar
repair. It is the \emph{actual arithmetic functional} $\beta$, its
unique determinant-compatible line metric, and its represented five-slot
morphism, tested against independently reconstructed renewal data.
If the physical metric on the candidate cotangent factor is instead the
Hodge quotient or the finite-flat metric, \cref{prop:v10-quotient-curvature,cor:v10-flat-fails}
return an explicit nonzero compatibility residual. No physical clock,
central Store signature, accepted field law, marked crystalline Frobenius,
or Frobenius--Wilson intertwiner is obtained from the determinant tensor
alone.

\section{V10 proof ledger and verification boundary}
\label{sec:v10-status}

\subsection{Proved in this iteration}
The new results are the native cotangent weight of the actual source and
its determinant, the explicit four-root determinant multisection and its
descended eighth-power cotangent section, the root-free cotangent-twisted
determinant tensor, the acquired failure of the original Hodge quotient
metric, the unique determinant-compatible line metric and connection,
and the full determinant-one compact holonomy for this stated completion.
The five-slot incidence is a nonzero parallel algebraic morphism with an
explicit normalization. Its historical physical occurrence is not claimed.

\subsection{Inherited and external inputs}
V6 supplies the second-jet isomorphism and its source identification.
V8 supplies the marked de Rham--Betti comparison. V9 supplies the actual
Chern matrices, certified rational enclosures, and full projective joint
holonomy. Their proofs are not repeated or strengthened by assumption.
Conormal and determinant functoriality are classical; the exact global
Picard orders use the stated Edidin--Hu theorem. The companion determinant
incidence/stabilizer calculation is cited rather than relabelled as a new
general representation theorem.

\subsection{What the computer checks}
\path{Arithmetic_Loading_Duality_Research_Notes_V10_checks.py} uses only
Python's exact rational arithmetic and SymPy, with no network calls and no
numerical integrator. It verifies the symmetric-square determinant,
line-weight arithmetic, local infinity coefficients, reference
automorphism characters, explicit incidence matrix, determinant and
curvature relations, the actual polynomial discriminant and root product,
and the finite-flat line-connection identities. It then reads the
SHA256-pinned V9 certificate, computes the new quotient and compatible
curvature enclosures with outward-rounded rational operations, and checks
all nine preserved mathematical bodies. The numerical input is inherited:
a hash comparison does not re-prove its V8--V9 period acquisition.

The script does not prove the cited Chow-ring theorem, general geometric
base-change assertions, or the noncomputational holonomy argument.
Those are respectively imported with a citation or proved in the text.
No proof-assistant formalization is asserted. No new values of the full
period matrix, theta ray or crystalline Frobenius are reported.

\subsection{Unresolved inputs and next attack}
The essential new distinction is between \emph{algebraically supplied}
determinant incidence and \emph{physically selected} compatible transport.
The arithmetic side now gives a concrete tensor and a unique compatibility
rule, but it does not show that renewal selects that rule. Its geometric
line twist changes the rank-two bundle and its scalar connection; it is
not an unrecorded gauge adjustment of V9. The fixed marked curve, family
and theta normalization remain declared arithmetic data, not outputs of
canonical chronology alone.

The next useful upstream attack is to identify the physical provenance and
common-source realization of this determinant tensor, or to prove that the
actual renewal metric selects (or fails) its compatibility rule. Repeating
the compact group calculation, choosing a determinant connection without
its tensor, or comparing only Lie-algebra dimensions would not advance that
question. Finite-prime Frobenius requires its own common realization and
source comparison, as the previous versions already establish.

\paragraph{Version integrity.}
All nine earlier mathematical bodies and historical proof-status ledgers
are preserved byte for byte. V10 is an appended result, not a revision of
V9's nonzero determinant-curvature theorem. Future work must append after
this block; any correction to an inherited assertion must be identified
explicitly.
% ====================== END IMMUTABLE V10 MATHEMATICAL BODY ======================
% APPEND V11 AND LATER ITERATIONS HERE.

% ===================== BEGIN IMMUTABLE V11 MATHEMATICAL BODY =====================
\clearpage
\section{V11: actual crystalline Frobenius before a physical transport comparison}
\label{sec:v11-context}

V10 supplies a concrete arithmetic determinant tensor and a specified
compatible compact connection. It does not supply a finite-prime operator
on that complex Hermitian packet. Repeating its compact group computation,
or declaring an independently unobserved renewal incidence to equal the
arithmetic tensor, would not advance the comparison. We instead acquire
an earlier missing arithmetic object: the \emph{actual crystalline
Frobenius in the marked algebraic differential basis} already used by
V6--V8. The acquired matrices then determine the exact prime-cyclic span
of the actual quadratic theta source.

Throughout this continuation the rational model is fixed, including its
central quadratic-descent choice:
\begin{equation}
 \mathcal C:\quad Y^2=f(X)=X^5-X^2-X,
 \qquad p\in\{3,5\}.
 \label{eq:v11-curve}
\end{equation}
No new arithmetic curve, deformation, target transfer, spectral weights,
or physical clock is selected. The two primes and their characteristic
polynomials were acquired in V5; we do not count their points again.

There are three principal conclusions. First, Frobenius is computed to
forty $p$-adic digits in one integral basis, by a fully specified rational
reduction and an analytic omitted-tail bound. Second, the original
quadratic source has prime-cyclic dimension nine at three and eight at
five, whereas \emph{the two already marked rational branch sources}
at infinity and zero generate all ten dimensions. Two sources and word
depth four are both minimal. Third, the acquired initial matrices and
V7's actual Gauss--Manin connection determine a unique convergent
horizontal Frobenius germ near the actual fibre. This closes the local
marked differential--Frobenius comparison, not the arithmetic--Wilson
comparison.

\subsection{Dependencies and what is genuinely new}
We inherit the actual source/second-jet identification from V6, the
rational connection and pairing from V7, and the marked complex comparison
from V8. V9--V10's nonflat compact connections are not substituted for
crystalline Frobenius. The standard Monsky--Washnitzer/crystalline
comparison and the basic Frobenius lift are cited from Kedlaya
\cite{KedlayaMW}; the integral-lattice description is the genus-two
specialization of the principal-parts construction of van den Bogaart
\cite{BogaartBasis}. We give the specialized exact reductions, precision
proof, actual matrices, source-rank certificates, and the resulting
minimal source completion. The finite Krylov recognition machinery in
the supplied monographs \cite{GT77,Arithmetic} is not re-proved under a
new name.

\begin{scope}[The coefficient field and the source line remain visible]
All Frobenius matrices below act on $\Q_p$-cohomology. Congruences and
valuation margins are not positive Hilbert-space estimates. The quadratic
source is compared after the explicit de-twist
$V\otimes\Lambda_\Theta^{-2}\simeq\operatorname{Sym}^2F$ of V6--V7.
No Frobenius on the retained theta-fibre line is inferred. Nonzero changes
of either marked source's scalar representative do not change any span
statement. No complex-linear operator on the four-dimensional theta
representation is asserted.
\end{scope}

\section{The integral marked cohomology lattice at three and five}
\label{sec:v11-lattice}
Put $e_i=[X^i\dd X/Y]$, $0\le i\le3$, and write
\begin{equation}
 \boxed{b=(b_0,b_1,b_2,b_3)=(e_0,e_1,e_2,3e_3),\qquad
 D=\operatorname{diag}(1,1,1,3).}
 \label{eq:v11-basis}
\end{equation}
Matrices use column coordinates: $\Phi_p(b_j)=\sum_i b_i(F_p)_{ij}$.
Here $\Phi_p$ is crystalline pullback by absolute $p$-power Frobenius,
with its usual action on $H^1$; it has the Weil polynomial of the
$p$-power Jacobian endomorphism used in V5. Since the residue field is
$\F_p$, its action on the scalar field $\Q_p$ is the identity.

\begin{proposition}[A common integral differential basis]
\label{prop:v11-integral}
The classes in \cref{eq:v11-basis} form a $\Z_p$-basis for
$H^1_{\mathrm{cris}}(\mathcal C_{\F_p}/\Z_p)$ under the smooth-lift
de Rham comparison, for both $p=3$ and $p=5$. In this basis the
alternating cup form is
\begin{equation}
 \boxed{\Omega_b=
 \begin{pmatrix}
 0&0&0&4\\0&0&4&0\\0&-4&0&0\\-4&0&0&0
 \end{pmatrix},\qquad \det\Omega_b=256.}
 \label{eq:v11-omega}
\end{equation}
Consequently $F_p\in M_4(\Z_p)$ and
\begin{equation}
 \boxed{F_p^{\mathsf T}\Omega_bF_p=p\Omega_b.}
 \label{eq:v11-similitude}
\end{equation}
\end{proposition}
\begin{proof}
The discriminant of $f$ is $-283$, a unit at the two primes, so the
standard odd-degree projective model is smooth there. Its anti-invariant
meromorphic differentials of pole order at most four at infinity have
basis $e_0,e_1,e_2,e_3$. There are no anti-invariant functions of pole
order at most three. The principal-parts integral de Rham description
\cite[Sections 3--5]{BogaartBasis} therefore identifies the cohomology
lattice with the kernel of the principal-part integration obstruction.
We compute that kernel here rather than silently regarding the raw
meromorphic basis as integral.

Set $z=Y/X^3$ and $u=X^{-1}$. Then
\[
 z^2=u-u^4-u^5,\qquad u=z^2+O(z^8),\qquad
 e_1=-2(1-4u^3-5u^4)^{-1}\dd z.
\]
Thus the only possible negative terms are
\[
 e_2=-2z^{-2}\dd z+O(z^4)\dd z,\qquad
 e_3=-2z^{-4}\dd z+O(z^2)\dd z.
\]
The primitive of the first principal part is integral. The second
has denominator three. At $p=3$ the kernel is exactly the span of
$e_0,e_1,e_2,3e_3$; at $p=5$ the raw basis is integral and multiplication
of its last vector by three is a unit change. This proves the basis
claim using the stated integral comparison theorem. Transporting V7's
pairing gives $\Omega_b=D^{\mathsf T}\Omega D$, exactly
\cref{eq:v11-omega}. Its determinant is a unit at both primes.
Crystalline Frobenius preserves the integral lattice and acts on $H^2$
by $p$. Cup-product functoriality gives \cref{eq:v11-similitude}.
\end{proof}

\begin{remark}[No source weights are inferred from a polynomial]
V5 proved
\[
 P_3(T)=T^4-2T^3+4T^2-6T+9,\qquad
 P_5(T)=T^4+4T^3+10T^2+20T+25.
\]
These polynomials determine a semisimple conjugacy class but not its
matrix relative to the actual differential sources. The calculation
below does not use either polynomial to choose a matrix entry. Their
agreement with the acquired matrices is an independent consistency check.
\end{remark}

\section{An exact Frobenius reduction with a proved infinite-tail bound}
\label{sec:v11-acquisition}

Let $\operatorname{red}$ denote reduction of an anti-invariant meromorphic
differential to its four cohomology coordinates in $e$. Work in the weak
completion with $Y$ inverted, and put
\begin{equation}
 \delta_p(X)=f(X^p)-f(X)^p\in p\Z[X],\qquad
 m_{p,j}=pj+(p-1)/2.
 \label{eq:v11-defect}
\end{equation}
The standard lift has $X\mapsto X^p$ and
$Y\mapsto Y^p(1+\delta_p/f^p)^{1/2}$. Its action on the basis is
\begin{equation}
 \Phi_p(e_i)=
 \sum_{j\ge0}p\binom{-1/2}{j}
 \operatorname{red}\left(
  \frac{X^{p(i+1)-1}\delta_p(X)^j\dd X}
       {Y^{p(2j+1)}}\right),\qquad0\le i\le3.
 \label{eq:v11-series}
\end{equation}
The expression denotes the cohomology coordinates of the sum. The
cohomological identification with proper crystalline $H^1$ is the
anti-invariant comparison in \cite{KedlayaMW,BogaartBasis}; it is not
replaced by a matching characteristic polynomial.

\subsection{Completely specified rational recurrences}
The inverse of $f'$ modulo $f$ is
\begin{equation}
 \boxed{h(X)=
 \frac{374X^4-27X^3+36X^2-422X-283}{283}.}
 \label{eq:v11-bezout}
\end{equation}
For a polynomial $a$ of degree at most four, set
\[
 B=(ha)\bmod f,\qquad A=(a-Bf')/f.
\]
Then $a=Af+Bf'$, $\deg B\le4$, and $\deg A\le3$.
For $m\ge1$ the exact-differential identity is
\begin{equation}
 \frac{a\dd X}{Y^{2m+1}}
 \equiv
 \frac{(A+\frac2{2m-1}B')\dd X}{Y^{2m-1}}
 \quad\pmod{\dd\bigl(\Q(X,Y)\bigr)}.
 \label{eq:v11-pole}
\end{equation}
Indeed the difference is
$-\frac2{2m-1}\dd(B/Y^{2m-1})$.
At $m=0$ the only nonbasis monomial obeys
\begin{equation}
 \operatorname{red}(X^4\dd X/Y)
 =\tfrac15 e_0+\tfrac25 e_1.
 \label{eq:v11-degree4}
\end{equation}
For any remaining polynomial degree $n\ge4$, use
\begin{equation}
 \begin{split}
 2\dd(X^{n-4}Y)=
 \bigl((2n-3)X^n-(2n-6)X^{n-3}-(2n-7)X^{n-4}\bigr)
 \frac{\dd X}{Y}.
 \end{split}
 \label{eq:v11-horizontal-red}
\end{equation}

These formulas give an exact finite evaluator: expand each numerator in
powers of the monic $f$, with digits of degree at most four; lower the
$Y$ denominator by \cref{eq:v11-pole}; and use
\cref{eq:v11-degree4,eq:v11-horizontal-red} for nonnegative powers of
$f$ remaining at denominator $Y$. Equivalently let $R_m$ be the $4$-by-$5$
map for the denominator $Y^{2m+1}$. Its first four columns at $m=0$
are $I_4$ and its last column is $(1,2,0,0)^{\mathsf T}/5$.
Equation \cref{eq:v11-pole} recursively gives every $R_m$ by a
$5$-by-$5$ rational multiplication. Thus no choice of reduction
complement or fitted Frobenius frame remains.

\begin{lemma}[Specialized denominator bound]
\label{lem:v11-denominator}
Let $a\in\Z_p[X]$ have degree at most four. Reduction of
$a\dd X/Y^{2m+1}$ for $m\ge1$ loses at most
$\lfloor\log_p(2m+1)\rfloor$ digits in the $e$ coordinates.
At $m=0$ at most $v_p(5)$ digits are lost. Changing the output to the
$b$ basis loses at most $v_p(3)$ further digits.
\end{lemma}
\begin{proof}
We spell out the pole argument behind the standard reduction estimate
\cite[Lemma 2]{KedlayaMW}, separating the $m=0$ degree reduction.
After an unramified extension the five roots of $f$ lift to distinct
integral roots. At each ramification point $Y$ is a parameter,
$\dd X=2Y\dd Y/f'(X)$, and the negative terms of the given differential
have integral coefficients in even powers $Y^{-2m},\ldots,Y^{-2}$.
Their primitives have odd denominators of absolute value at most
$2m-1$. Multiplication by
$p^{\lfloor\log_p(2m+1)\rfloor}$ makes all these principal parts
integral.

The finite reduction primitive is a sum
$\sum_{k=1}^m B_k(X)/Y^{2k-1}$ with $\deg B_k\le4$.
Start with its highest pole. Integrality of the five values of the
scaled $B_m$ at the five lifted roots implies coefficient integrality,
because their Vandermonde determinant is a unit. Subtract that term and
repeat downward. All scaled $B_k$ are integral. Subtracting the resulting
integral exact differential leaves an integral polynomial numerator at
denominator $Y$. For $m\ge1$ the explicit last step in
\cref{eq:v11-pole} has degree at most three, so it is already in the
$e$ basis. This proves the asserted bound and descends from the
unramified extension. For $m=0$, \cref{eq:v11-degree4} displays the
only possible additional denominator. Finally $b_3=3e_3$, so the output
coordinate change divides only the last coordinate by three.
\end{proof}

\begin{theorem}[Certified forty-digit marked acquisition]
\label{thm:v11-tail}
The $j$th term of \cref{eq:v11-series}, expressed in the $b$ basis,
has every coordinate of valuation at least
\begin{equation}
 \boxed{j-\lfloor\log_p(2j+1)\rfloor-2\qquad(j\ge2).}
 \label{eq:v11-tail-floor}
\end{equation}
In particular the first $48$ terms, $0\le j<48$, determine the
actual matrix $F_p$ modulo $p^{40}$. The omitted tail lies in
$p^{42}M_4(\Z_3)$ at three and in $p^{44}M_4(\Z_5)$ at five.
\end{theorem}
\begin{proof}
The binomial coefficient is
$(-1)^j\binom{2j}{j}/4^j\in\Z_p$. Every numerator coefficient,
including the initial factor $p$, is therefore divisible by $p^{j+1}$.
For the present polynomial $\deg\delta_p=5p-3$. Even in column $i=3$,
its degree before reduction is at most
\[
 4p-1+(5p-3)j<5(m_{p,j}+1)\qquad(j\ge2,\ p=3,5).
\]
Consequently the monic $f$-adic expansion of every omitted numerator
uses only denominators $Y^{2m+1}$ with $0\le m\le m_{p,j}$.
Monic division introduces no denominator. Apply
\cref{lem:v11-denominator}, allowing the sum of the pole bound, the
possible degree-four division, and the output division by three.
Since $2m_{p,j}+1=p(2j+1)$ and $v_p(15)=1$ here, this even gives
$j-\lfloor\log_p(2j+1)\rfloor-1$; the displayed estimate retains
one extra digit of reserve.

The bound tends to infinity. The same integration argument bounds the
exact primitives with pole and degree growing only linearly in $j$;
their coefficients grow in valuation linearly up to a logarithm.
Thus the primitives converge in a weak completion, so reduction of the
convergent Frobenius lift is legitimate, as in the standard
Monsky--Washnitzer construction. For $j\ge48$ the displayed lower
bound is nondecreasing: increasing $j$ by one increases
$\lfloor\log_p(2j+1)\rfloor$ by at most one. Its values at $48$
are $42$ and $44$. The nonarchimedean triangle inequality gives the
stated tail bounds and hence the forty-digit assertion.
\end{proof}

\begin{remark}[A finite sum is not the exact Frobenius matrix]
The rational partial matrix is only a congruence representative of $F_p$.
It may have different exact rational eigenvalues. All rank conclusions
below combine a strict $p$-adic nonvanishing witness with an independently
proved exact upper bound. In particular, the small nonzero determinant
of a lifted ten-column one-source bank is \emph{not} evidence that its
true rank is ten.
\end{remark}

\section{The actual marked matrices and their first Hodge comparison}
\label{sec:v11-matrices}

\begin{theorem}[Frobenius in the acquired arithmetic source basis]
\label{thm:v11-matrices}
The crystalline operators in \cref{eq:v11-basis} satisfy
\begin{equation}
 \boxed{F_3\equiv
 \begin{pmatrix}
 612&525&435&603\\393&717&420&558\\
 177&705&0&68\\381&471&646&131
 \end{pmatrix}\pmod{3^6},}
 \label{eq:v11-F3}
\end{equation}
\begin{equation}
 \boxed{F_5\equiv
 \begin{pmatrix}
 12875&14820&8436&4740\\1740&7265&7702&11619\\
 925&675&12115&7692\\2625&13695&13670&14616
 \end{pmatrix}\pmod{5^6}.}
 \label{eq:v11-F5}
\end{equation}
The accompanying rational certificate gives all sixteen entries modulo
$p^{40}$ at each prime. Their characteristic polynomials are the
inherited $P_3,P_5$, and their exact pairing relation is
\cref{eq:v11-similitude}.
\end{theorem}
\begin{proof}
Use \cref{eq:v11-series} with $48$ terms, the inverse
\cref{eq:v11-bezout}, and the reduction recurrences. Apply $D^{-1}$
on the output and $D$ on the input to obtain the $b$ matrix. All
operations in this finite calculation take place in $\Q$, and
\cref{thm:v11-tail} identifies its output with the actual Frobenius
to more than the stated precision. The executable calculation, its
full forty-digit entries, and the determinant witnesses used below are
included in the verification bundle. Reduction of those entries gives
\cref{eq:v11-F3,eq:v11-F5}. The cohomological comparison and V5's
point counts identify the exact characteristic polynomial. Independently,
the program checks both characteristic polynomials and the pairing
similitude modulo $p^{40}$; neither check was used to fit the entries.
\end{proof}

\begin{proposition}[Actual Frobenius is not the rank-two Cartier operator]
\label{prop:v11-hodge}
Let $F^{\mathrm{Hdg}}=\Span_{\Q_p}(b_0,b_1)$ be the holomorphic
differential plane of the fixed characteristic-zero lift. Then
\begin{equation}
 \boxed{F^{\mathrm{Hdg}}\cap\Phi_p F^{\mathrm{Hdg}}=0
 \quad(p=3,5).}
 \label{eq:v11-hodge-transverse}
\end{equation}
More precisely the bottom-left $2$-by-$2$ block of $F_p$ has determinant
of valuation $2$ at three, with leading unit $1\pmod3$, and valuation
$3$ at five, with leading unit $3\pmod5$.
Also $b_0$ is a cyclic vector for $\Phi_p$ on the full rank-four
cohomology. The determinant of
$(b_0,\Phi_pb_0,\Phi_p^2b_0,\Phi_p^3b_0)$ has valuation $4$ at
three and $5$ at five, with leading unit one at both primes.
\end{proposition}
\begin{proof}
All four determinants are computed from the forty-digit matrix by integral
polynomials in its entries. Their valuations are strictly below forty,
so they are the exact valuations of the actual determinants. The
nonzero bottom block makes the map
$F^{\mathrm{Hdg}}\to E_p/F^{\mathrm{Hdg}}$ induced by $\Phi_p$
an isomorphism, proving transversality. The four-column determinant
proves cyclicity. The valuations and leading digits are listed in the
certificate and verified by the rational evaluator.
\end{proof}

\begin{scope}[The earlier Cartier calculation is not being changed]
V6's Cartier map acts on differentials over $\F_p$, and its induced
quadratic map was defined through the second jet. Here $\Phi_p$ acts on
characteristic-zero crystalline cohomology with its retained Hodge
filtration. Those are different objects. The transverse-plane result is
not a contradiction to the Cartier ranks in V6; it forbids identifying
the two transfers without a separately specified filtered comparison.
The oriented $b_0$ used in the last assertion retains the rational curve
model; no inversion-even theta datum has been turned functorially into
an odd weight-one vector.
\end{scope}

\section{Two existing arithmetic branch sources minimally complete prime transport}
\label{sec:v11-cyclic}

Let
\[
 E_p=H^1_{\mathrm{cris}}(\mathcal C_{\F_p})\otimes\Q_p,
 \qquad\mathcal W_p=\operatorname{Sym}^2E_p,
 \qquad S_p=\operatorname{Sym}^2\Phi_p.
\]
Use the unscaled monomial basis in the order
\begin{equation}
 (00,01,02,03,11,12,13,22,23,33),\qquad ij=b_ib_j.
 \label{eq:v11-symbasis}
\end{equation}
The actual squared theta section at infinity and its displacement
associated with the already rational point $W_0$ have de-twisted
second jets on the rays
\begin{equation}
 \boxed{s_\infty=b_0^2,
 \qquad s_0=b_1^2.}
 \label{eq:v11-sources}
\end{equation}
This follows directly from V6's conormal formula
$\ell_\infty=\Span(\omega_0)$ and
$\ell_r=\Span(\omega_1-r\omega_0)$ at $r=0$. Both sources precede
the current computation. In particular the second is not an eigenvector
chosen after seeing a missing Frobenius direction.

\begin{lemma}[Exact scalar minimal polynomials of the quadratic transfer]
\label{lem:v11-minpoly}
Define
\begin{align}
 q_3(T)&=(T^4+2T^3+6T^2+18T+81)
          (T^4+4T^3+10T^2+36T+81),\notag\\
 q_5(T)&=(T^2-6T+25)(T^4+30T^2+625).
 \label{eq:v11-q}
\end{align}
Then
\begin{align}
 \det(T-S_3)&=(T-3)^2q_3(T),&
 \mu_3(T)&=(T-3)q_3(T),\notag\\
 \det(T-S_5)&=(T-5)^2(T+5)^2q_5(T),&
 \mu_5(T)&=(T-5)(T+5)q_5(T).
 \label{eq:v11-polys}
\end{align}
The displayed $\mu_p$ are the exact minimal polynomials. In particular
a single vector can generate at most nine dimensions at three and at
most eight dimensions at five.
\end{lemma}
\begin{proof}
Both $P_p$ are squarefree in characteristic zero. Any operator with
that characteristic polynomial is diagonalizable over an algebraic
closure. On its symmetric square the ten eigenvalues are the ten
products $\alpha_i\alpha_j$, $i\le j$. Computing those products by
the rational companion matrix of $P_p$ gives the two factorizations
in \cref{eq:v11-polys}; this is a finite exact polynomial calculation
recorded in the verifier. Their displayed squarefree parts are
squarefree, again by exact polynomial gcd. The symmetric square remains
diagonalizable, so these squarefree parts are exactly its minimal
polynomials. A vector's cyclic span is a quotient of
$\Q_p[T]/(\mu_p)$, giving the upper bounds.
\end{proof}

\begin{theorem}[Actual one-source and minimal two-source prime completion]
\label{thm:v11-two-sources}
For the sources \cref{eq:v11-sources},
\begin{equation}
 \boxed{
 \dim\Q_3[S_3]s_\infty=9,\qquad
 \dim\Q_5[S_5]s_\infty=8.}
 \label{eq:v11-one-ranks}
\end{equation}
At both primes the ten columns
\begin{equation}
 \boxed{
 \mathscr B_p=
 (s_\infty,S_ps_\infty,\ldots,S_p^4s_\infty,
  s_0,S_ps_0,\ldots,S_p^4s_0)
 }
 \label{eq:v11-bank}
\end{equation}
form a basis of $\mathcal W_p$. Their determinants satisfy
\begin{equation}
 \boxed{
 v_3(\det\mathscr B_3)=21,\quad
 3^{-21}\det\mathscr B_3\equiv2\pmod3;
 \qquad
 v_5(\det\mathscr B_5)=27,\quad
 5^{-27}\det\mathscr B_5\equiv4\pmod5.}
 \label{eq:v11-two-det}
\end{equation}
Thus two source rays are necessary and sufficient for the full prime
carrier. Among banks retaining those two rays and all iterates up to a
common depth, depth four is minimal; the ranks through depths
$0,1,2,3,4$ are exactly
\begin{equation}
 \boxed{2,4,6,8,10.}
 \label{eq:v11-depth}
\end{equation}
\end{theorem}
\begin{proof}
We specify strict finite witnesses for every nontrivial lower bound.
Rows and columns are numbered from zero in \cref{eq:v11-symbasis}.
At $p=3$ take the first nine iterates of $s_\infty$ and the row order
\[
 (0,4,7,3,8,5,9,2,1).
\]
This nine-by-nine determinant has valuation $31$ and leading unit
$1\pmod3$. At $p=5$ take the first eight iterates and rows
\[
 (0,4,2,3,5,1,6,8).
\]
Its determinant has valuation $31$ and leading unit $4\pmod5$.
These valuations are strictly below the acquired precision forty.
Because all matrix multiplication and determinants are integral
polynomials in the entries of $F_p$, the computed residues certify
exact nonvanishing. Combined with \cref{lem:v11-minpoly}, they give
\cref{eq:v11-one-ranks}. This argument does not regard the rational
congruence representative as the exact operator.

Apply the same integral-polynomial computation to the full bank
\cref{eq:v11-bank}. Its valuations and leading digits are
\cref{eq:v11-two-det}, again strictly below forty. Hence the ten
columns are independent. Every subset of these columns is independent,
which proves \cref{eq:v11-depth}. A single source is impossible by
the exact minimal-polynomial bounds. A two-source bank of common depth
$d$ has at most $2(d+1)$ columns, so it cannot span ten dimensions
when $d<4$. Both minimalities are therefore exact.
\end{proof}

\begin{remark}[What the valuations measure]
The valuation $21$ or $27$ is a conditioning cost of these integral
source coordinates, not a small complex determinant or a positive
physical gap. The acquired precision exceeds that cost, so the rank
claims are certified. The resulting basis is generally not a
unimodular lattice basis, and no uniform-in-prime conditioning estimate
is inferred from these two local certificates.
\end{remark}

\section{The missing directions are identified spectrally, not filled by hand}
\label{sec:v11-obstruction}

The exact upper bounds have an intrinsic symplectic explanation. Retain
V7's tensor identification
\[
 \mathfrak a(vw)(x)=\Omega_b(w,x)v+\Omega_b(v,x)w,
 \qquad\mathfrak a:\operatorname{Sym}^2E_p\xrightarrow{\sim}
 \mathfrak{sp}(E_p,\Omega_b).
\]

\begin{proposition}[Normalized quadratic Frobenius is adjoint transport]
\label{prop:v11-adjoint}
One has
\begin{equation}
 \boxed{
 \mathfrak a\circ(p^{-1}S_p)
 =\operatorname{Ad}(F_p)\circ\mathfrak a.}
 \label{eq:v11-adjoint}
\end{equation}
Its fixed space has dimension two at both primes. At five its
$-1$ eigenspace also has dimension two. At three the one-source
quotient missing from the full carrier is precisely a one-dimensional
quotient of the fixed space. At five the missing quotient has one
fixed and one $-1$ component.
\end{proposition}
\begin{proof}
For a symmetric coordinate matrix $B$ the adjoint identification sends
it to $B\Omega_b$. Equation \cref{eq:v11-similitude} gives
\[
 F_p(B\Omega_b)F_p^{-1}
  =p^{-1}(F_pBF_p^{\mathsf T})\Omega_b,
\]
which is \cref{eq:v11-adjoint}. The multiplicities of $p$ and $-p$
in \cref{eq:v11-polys} give the eigenspace dimensions, because the
operator is semisimple. Every other eigenvalue there is simple.
The maximal one-source ranks in \cref{thm:v11-two-sources} show
that $s_\infty$ has a nonzero component at every distinct eigenvalue.
Its orbit can occupy just one line in each two-dimensional repeated
eigenspace and all of every simple eigenspace. The missing quotient
is therefore exactly as stated. No orthogonal complement or canonical
missing vector is asserted without a chosen metric.
\end{proof}

\begin{theorem}[The second rational branch resolves each repeated component]
\label{thm:v11-projectors}
The spectral projectors onto the repeated eigenspaces can be computed
without choosing eigenvectors. At three,
\begin{equation}
 E_{3,+}=\frac{q_3(S_3)}{151632}.
 \label{eq:v11-proj3}
\end{equation}
At five put $\mu_5$ as in \cref{eq:v11-polys}; then
\begin{equation}
 E_{5,+}=\frac{(\mu_5/(T-5))(S_5)}{400000},\qquad
 E_{5,-}=\frac{(\mu_5/(T+5))(S_5)}{-1600000}.
 \label{eq:v11-proj5}
\end{equation}
For each of these rank-two projectors $E$, the two vectors
$Es_\infty$ and $Es_0$ are independent. Thus the additional source
resolves the missing eigendirections rather than enlarging an
unobserved spectator space.
\end{theorem}
\begin{proof}
If $\mu=(T-a)q$ is squarefree, $q(S)/q(a)$ acts as one at $a$ and
zero at every other eigenvalue. This proves the projector formulas;
the displayed denominators are exact polynomial evaluations.
Independence also follows from the complete two-source span, but there
are direct lower-dimensional witnesses. Before dividing by the projector
denominator, use its numerator on the two-column source bank
$(s_\infty,s_0)$. In the column order $(s_0,s_\infty)$ the following
minors are nonzero:
\begin{center}
\begin{tabular}{ccccc}
\toprule
Prime & Eigenspace & Rows & Determinant valuation & Leading unit\\
\midrule
$3$ & $+3$ & $(2,3)$ & $13$ & $2\pmod3$\\
$5$ & $+5$ & $(0,2)$ & $13$ & $1\pmod5$\\
$5$ & $-5$ & $(0,1)$ & $12$ & $3\pmod5$\\
\bottomrule
\end{tabular}
\end{center}
All are integral-polynomial witnesses of precision forty and are
verified in the certificate. Division by a nonzero scalar does not
change their independence.
\end{proof}

\begin{scope}[Differential closure and prime closure answer different questions]
V7 proved that one quadratic source generates ten dimensions under
successive Gauss--Manin derivatives. Here one prime operator necessarily
leaves the explicit one- or two-dimensional quotient above. The two
statements are compatible: differentiation moves the eigenspaces as
well as the source, whereas powers of a fixed semisimple operator cannot
separate two directions carrying the same eigenvalue. The current result
closes the prime source defect with the already existing second rational
branch source, not by changing the meaning of cyclicity.
\end{scope}

\section{A convergent marked Frobenius germ compatible with the actual connection}
\label{sec:v11-horizontal}

We now use the acquired matrix rather than an unspecified initial
Frobenius in the differential equation. On V7's family
\[
 \mathcal C_t:Y^2=X^5-X^2-tX,
\]
let $A_b(t)=D^{-1}A(t)D$, with the exact $A(t)$ of
\cref{eq:v7-A}. For reference it is
\begin{equation}
 A_b(t)=\frac1{2t(256t^3+27)}
 \begin{pmatrix}
 -192t^3&-27t&36t^2&-144t^3\\
 48t^2&-64t^3&-9t&36t^2\\
 36t&-48t^2&64t^3&27t\\
 -27&36t&-48t^2&192t^3
 \end{pmatrix}.
 \label{eq:v11-Ab}
\end{equation}
Take the base Frobenius lift $\sigma(t)=t^p$, which fixes $t=1$.
A Frobenius structure is a horizontal map from $\sigma^*E_p$ to
$E_p$, not an endomorphism of every parameter fibre with the base
ignored \cite{KedlayaFrobBook,KedlayaTuitman}.

\begin{theorem}[The marked local differential--Frobenius comparison]
\label{thm:v11-horizontal}
There is a unique convergent germ $\mathcal F_p(t)$ at $t=1$ satisfying
\begin{equation}
 \boxed{
 \mathcal F_p'=p t^{p-1}\mathcal F_p A_b(t^p)
                 -A_b(t)\mathcal F_p,
 \qquad \mathcal F_p(1)=F_p.}
 \label{eq:v11-horizontal}
\end{equation}
It is the germ of the actual cohomological Frobenius induced by
\[
 X\mapsto X^p,\qquad
 Y\mapsto Y^p\left(1+
 \frac{f_{t^p}(X^p)-f_t(X)^p}{f_t(X)^p}\right)^{1/2}.
\]
Writing $u=t-1$ and $\mathcal F_p(1+u)=\sum_{n\ge0}F_{p,n}u^n$,
one has
\begin{equation}
 \boxed{n! F_{p,n}\in M_4(\Z_p).}
 \label{eq:v11-factorial}
\end{equation}
In particular it converges for $|u|_p<p^{-1/(p-1)}$, including the
closed disk $u\in p\Z_p$. The Taylor tail beyond degree $d$ on that
disk has all entries of valuation at least
\begin{equation}
 \inf_{n>d}\{n-v_p(n!)\}.
 \label{eq:v11-u-tail}
\end{equation}
It also satisfies the exact identity
\begin{equation}
 \boxed{\mathcal F_p(t)^{\mathsf T}\Omega_b\mathcal F_p(t)
       =p\Omega_b.}
 \label{eq:v11-germ-similitude}
\end{equation}
\end{theorem}
\begin{proof}
Since $283$ and $2$ are units, the coefficients of the Taylor series
of $A_b(1+u)$ belong to $\Z_p$. The same holds for
$p(1+u)^{p-1}A_b((1+u)^p)$. Write these two series as
$\sum A_nu^n$ and $\sum B_nu^n$. Equation
\cref{eq:v11-horizontal} is the recursive relation
\begin{equation}
 (n+1)F_{p,n+1}
 =\sum_{j=0}^{n}\bigl(F_{p,n-j}B_j-A_jF_{p,n-j}\bigr).
 \label{eq:v11-horizontal-rec}
\end{equation}
It determines a unique formal series from the actual $F_p$.
Inductively, multiplying this relation by $n!$ proves
$(n+1)!F_{p,n+1}$ integral, because
$n!/(n-j)!$ is an integer. Legendre's bound
$v_p(n!)\le n/(p-1)$ proves convergence and
\cref{eq:v11-u-tail}.

The displayed lift has numerator defect divisible by $p$, so it defines
the standard relative weak-completion Frobenius. The cohomological
Frobenius structure is horizontal for the Gauss--Manin connection
\cite{KedlayaFrobBook,KedlayaTuitman}; in our column convention,
$\nabla(b\mathcal F)=b(A_b\mathcal F+\mathcal F')\dd t$,
while pullback of the connection on the input gives
$b\mathcal F A_b(t^p)p t^{p-1}\dd t$.
This yields exactly \cref{eq:v11-horizontal}, with specialization
$F_p$ acquired above. Uniqueness identifies the convergent germ with
the actual one; no new initial period or matrix is chosen.

Finally $A_b^{\mathsf T}\Omega_b+\Omega_bA_b=0$ by V7 and the
fixed basis change. Differentiating
$\mathcal F^{\mathsf T}\Omega_b\mathcal F$ shows that its
quadratic pairing obeys the pulled-back homogeneous pairing equation.
The constant $p\Omega_b$ obeys the same equation and has the same
initial value by \cref{eq:v11-similitude}. Uniqueness proves
\cref{eq:v11-germ-similitude}. Invertibility follows as well.
\end{proof}

\begin{proposition}[An explicitly acquired first horizontal coefficient]
\label{prop:v11-firstjet}
The first coefficient is fixed by the marked matrix, not a new
measurement:
\begin{equation}
 \boxed{F_{p,1}=pF_pA_b(1)-A_b(1)F_p.}
 \label{eq:v11-firstjet}
\end{equation}
In particular
\begin{equation}
 F_{3,1}\equiv
 \begin{pmatrix}27&27&27&0\\66&66&39&9\\24&75&54&14\\45&45&45&54\end{pmatrix}
 \pmod{3^4},
 \label{eq:v11-der3}
\end{equation}
\begin{equation}
 F_{5,1}\equiv
 \begin{pmatrix}535&510&561&270\\105&595&470&213\\350&145&345&267\\250&140&315&50\end{pmatrix}
 \pmod{5^4}.
 \label{eq:v11-der5}
\end{equation}
\end{proposition}
\begin{proof}
Use $n=0$ in \cref{eq:v11-horizontal-rec} and the acquired matrix.
All denominators in $A_b(1)$ are units at the two primes, so its
forty-digit congruence immediately determines these four-digit displays.
\end{proof}

\begin{remark}[Tensor lift and its remaining line datum]
Applying $\operatorname{Sym}^2$ to the horizontal structure gives a
rank-ten horizontal Frobenius on the same quadratic differential system
that carried V7's source. The initial matrix is exactly $S_p$, and its
two-source completion is \cref{thm:v11-two-sources}. Restoring
$\Lambda_\Theta^2$ would require its actual Frobenius linearization.
The present construction does not substitute the Chern connection or
V10's determinant-compatible cotangent metric for that missing line map.
\end{remark}

\section{V11 proof ledger, reproducibility, and the sharper next comparison}
\label{sec:v11-status}

\subsection{Proved in this continuation}
The new arithmetic acquisition is the marked crystalline matrix at the
two already selected primes, including its integral lattice and a
convergent exact reduction with a forty-digit certificate. On that
actual realization the Hodge plane is transverse to its Frobenius image.
The original quadratic theta ray has maximal one-source ranks nine and
eight; the second already rational branch ray minimally completes the
full ten-dimensional prime carrier at common depth four. The repeated
spectral directions are identified explicitly by exact polynomial
projectors and strict source-separation minors. Finally the actual
Gauss--Manin connection and these acquired matrices determine a
convergent marked horizontal Frobenius germ with a factorial bound.

\subsection{Inherited and external inputs}
We use V5's point counts and polynomials rather than redoing them, V6's
source jets, V7's connection and cup form, and V8's independently marked
complex comparison. The standard proper crystalline/anti-invariant weak
cohomology identification, integral principal-parts description, and
horizontality of cohomological Frobenius are the cited classical inputs.
The rational recurrences, source determinants, and local Taylor bound
are specialized and proved in the present continuation. The geometric
endomorphism and compact-holonomy conclusions of V6 and V9--V10 are not
premises of the new source-rank calculation.

\subsection{What the executable certificate establishes}
\path{Arithmetic_Loading_Duality_Research_Notes_V11_checks.py} uses only
exact integer/rational arithmetic and SymPy. It evaluates the forty-eight
terms of the actual Frobenius lift separately at the two primes; no
characteristic polynomial, theta Gram, renewal operator, or source
weights are used to fit the entries. It checks the specialized Bezout
and exact-differential identities, pairing and characteristic-polynomial
congruences, exact quadratic characteristic factorizations, source
minors, repeated-eigenspace minors, and horizontal first jets.

The full residues, term valuations, row indices, precision, and leading
nonzero digits are stored in
\path{Arithmetic_Loading_Duality_Research_Notes_V11_certificate.json}.
All decisive determinants have valuation strictly below forty. Upper
rank bounds follow from the exact minimal polynomials, not from testing
whether a finite residue vanishes. The omitted-tail proof is
\cref{thm:v11-tail}; printing more matching numerical digits is not
used as a substitute. The verifier compares all ten preserved bodies
with the supplied V10 source. It does not machine-prove the cited
cohomological theorems or the general source-jet comparison. No
proof-assistant formalization is asserted.

\subsection{Unresolved: a common metric realization is not a coefficient-field convention}
There are now two acquired arithmetic realizations in compatible
\emph{algebraic markings}: V8's complex period pairing and this
continuation's local crystalline Frobenius. Their scalar fields and
natural comparison structures remain different. A $p$-adic residue
cannot be inserted as a complex unitary matrix, and an embedding of a
splitting field does not transport the acquired theta/Hodge metric by
itself. The original four-dimensional theta normalizer obstruction,
quadratic-twist descent distinction, and nonflat Chern/flat differential
separation remain in force.

In particular the two-source quadratic system remains insensitive to a
central quadratic twist after symmetric squaring. This does not recover
the signed weight-one arithmetic data from the projective theta packet.
The rational model in \cref{eq:v11-curve} fixed those data before
acquisition. Its parity was not removed by a computational convention.

The next actual comparison should retain the two independently occurring
rational branch sources, their common arithmetic realization, and the
physical mixed source and transport maps. Their existence is not supplied
by equality of group types or by the new prime-cyclic ranks. The
arithmetic determinant functional of V10 still has to be identified
with the reconstructed physical incidence, and the theta-fibre
Frobenius requires its own line-level comparison. The central Store
signature, physical clock, accepted field law, and Frobenius--Wilson
intertwiner remain unacquired. No Einstein limit, Standard-Model
coupling, GRH estimate, or analytic continuation assertion is inferred.

\paragraph{Version integrity.}
All V1--V10 mathematical bodies and historical proof-status ledgers are
preserved byte for byte. V11 updates the previously unresolved marked
crystalline acquisition and local prime source completion; it does not
silently revise any earlier finite-field, metric, or holonomy statement.
Future additions must follow this block.
% ====================== END IMMUTABLE V11 MATHEMATICAL BODY ======================
\clearpage
% =============================================================================
\section{V12: local reconstruction of genus-two moduli from the marked arithmetic theta source}
\label{sec:v12-source-moduli}
% =============================================================================

\subsection{Purpose of this continuation and dependency audit}
V11 acquired the marked crystalline Frobenius on the de-twisted quadratic
cohomology and isolated one remaining line-level datum: restoring the actual
theta-section carrier requires a Frobenius linearization of the retained
theta-fibre line.  Before introducing such a scalar datum, the present
continuation goes one layer upstream and asks how much of the underlying
arithmetic geometry is already determined by the \emph{projective} pure theta
source itself.

The answer is stronger than was previously used.  In the independently fixed
symplectic/theta marking of V8, the four second-order theta constants form a
local coordinate system on the full three-complex-dimensional Siegel space at
the actual arithmetic period matrix.  Since V2 identifies the squared theta
source with a fixed signed permutation of those four constants, the marked
rank-one source projector locally determines the marked period matrix.  Thus
the static source does not merely remember a selected one-parameter family: it
locally remembers the whole genus-two deformation germ.  What it does \emph{not}
determine is a tangent direction through that point, nor the absolute scalar
linearization of the theta-fibre line.

We inherit V8's certified period matrix and theta marking, V2's exact relation
between second-order theta constants and the squared theta section, and V10's
canonical line identity
\[
 \Lambda_\Theta\simeq\mathscr P^{-3}.
\]
No period integral, branch-weight computation, crystalline matrix, or finite
Krylov theorem is repeated.  The only numerical work below is the derivative
of the already acquired theta-null packet, with a new independent Gaussian
tail certificate.

\begin{scope}[What is and is not reconstructed]
The statements below are local in the independently marked Siegel chart.  A
marked source ray determines a nearby marked principally polarized abelian
surface; on the Jacobian locus the corresponding curve-level statement uses
the classical Torelli theorem.  Forgetting the theta structure reintroduces
the finite Heisenberg ambiguity classified in V2--V4.  A static source point
does not select a preferred path, clock, Gauss--Manin direction, renewal
transport, or Wilson connection.  The absolute theta-line scalar also remains
separate from the projective source.
\end{scope}

\subsection{The affine theta-null chart and its exact differential}
Let \(\mathfrak H_2\) be Siegel upper half-space and retain V8's convention
for the four second-order theta constants
\begin{equation}
 t_\sigma(\tau)
 =\sum_{n\in\mathbb Z^2}
 \exp\!\left(
 2\pi i\,(n+\sigma/2)^{\mathsf T}\tau(n+\sigma/2)
 \right),
 \qquad \sigma\in\mathbb F_2^2.
 \label{eq:v12-theta-null}
\end{equation}
This is the zero-value of the second-order theta function used in V2 and V8;
absolute convergence and termwise differentiation on compact Siegel sets are
standard \cite{DLMFTheta} and also follow directly from the Gaussian estimate
below.

Write the symmetric period coordinate as
\[
 q=(\tau_{11},\tau_{12},\tau_{22}).
\]
At the actual period matrix \(\tau_*\) of V8, \(t_{00}\ne0\), so define
\begin{equation}
 \boxed{
 \mathcal U(\tau)
 =\left(
 \frac{t_{01}}{t_{00}},
 \frac{t_{10}}{t_{00}},
 \frac{t_{11}}{t_{00}}
 \right)\in\mathbb C^3.}
 \label{eq:v12-affine-map}
\end{equation}

\begin{lemma}[Exact theta-null differential]
\label{lem:v12-theta-diff}
Put \(v=n+\sigma/2\) and
\(E_{\sigma,n}=\exp(2\pi i v^{\mathsf T}\tau v)\).  Then
\begin{align}
 \partial_{\tau_{11}}t_\sigma
 &=2\pi i\sum_n v_1^2 E_{\sigma,n},\notag\\
 \partial_{\tau_{12}}t_\sigma
 &=4\pi i\sum_n v_1v_2 E_{\sigma,n},\notag\\
 \partial_{\tau_{22}}t_\sigma
 &=2\pi i\sum_n v_2^2 E_{\sigma,n}.
 \label{eq:v12-theta-derivatives}
\end{align}
Consequently, for \(\sigma\ne00\),
\begin{equation}
 \boxed{
 \partial_q\!\left(\frac{t_\sigma}{t_{00}}\right)
 =\frac{t_{00}\,\partial_qt_\sigma
       -t_\sigma\,\partial_qt_{00}}{t_{00}^2}.}
 \label{eq:v12-ratio-diff}
\end{equation}
\end{lemma}
\begin{proof}
On a compact subset on which
\(\operatorname{Im}\tau\succeq\mu I\) with \(\mu>0\), the summand in
\cref{eq:v12-theta-null} has modulus at most
\(e^{-2\pi\mu\|v\|^2}\).  Multiplication by any fixed polynomial in \(v\)
still gives an absolutely and locally uniformly convergent series.  Termwise
differentiation is therefore valid and gives
\cref{eq:v12-theta-derivatives}; the quotient rule gives
\cref{eq:v12-ratio-diff}.
\end{proof}

\subsection{A certified nondegenerate theta-null Jacobian}
Let
\[
 J_\Theta=D_q\mathcal U(\tau_*).
\]
The following theorem is the finite quantitative input needed for the inverse
problem.

\begin{theorem}[The marked pure theta source is locally moduli-complete]
\label{thm:v12-local-inverse}
At the actual period matrix \(\tau_*\) acquired in V8, the Jacobian of
\cref{eq:v12-affine-map} is invertible.  In the coordinate order
\((\tau_{11},\tau_{12},\tau_{22})\), it has the certified enclosure whose
midpoint is
\begin{equation}
 J_\Theta\approx
 \begin{pmatrix}
 -0.00924665-0.17718746i&
 -0.00830117-0.17936331i&
 -0.64864148+0.09226261i\\
 -0.62327113+0.26078322i&
  0.06587879-0.12555628i&
  0.07646983-0.12510563i\\
 -0.65311689+0.24794197i&
 -1.35594690+0.61712727i&
 -0.64174809+0.24639827i
 \end{pmatrix}.
 \label{eq:v12-J}
\end{equation}
Moreover
\begin{align}
 -0.317165077970180297
 &<\Re\det J_\Theta
 <-0.317165077970180296,\notag\\
 0.665403123897688791
 &<\Im\det J_\Theta
 <0.665403123897688792,
 \label{eq:v12-detJ}
\end{align}
so
\begin{equation}
 \boxed{|\det J_\Theta|>0.73.}
 \label{eq:v12-det-margin}
\end{equation}
Its smallest singular value obeys the purely certified bound
\begin{equation}
 \boxed{\sigma_{\min}(J_\Theta)>0.35,
 \qquad \|J_\Theta^{-1}\|_2<2.86.}
 \label{eq:v12-inverse-margin}
\end{equation}
Hence \(\mathcal U\) is a biholomorphism between neighborhoods of
\(\tau_*\) and \(\mathcal U(\tau_*)\).
\end{theorem}
\begin{proof}
V8 proves
\(\operatorname{Im}\tau_*\succ0.45I\).  Truncate each series in
\cref{eq:v12-theta-null,eq:v12-theta-derivatives} to
\(-6\le n_1,n_2\le6\).  Outside this square, set
\(m=\max(|n_1|,|n_2|)\ge7\).  There are at most \(8m+4\) points in the
corresponding shell, while
\[
 \|n+\sigma/2\|^2\ge(m-1/2)^2.
\]
Since \(0.9\pi>3\log2\), every undifferentiated summand is bounded by
\(2^{-3(m-1/2)^2}\), and every first derivative summand is bounded by
\[
 4\pi(m+1/2)^2\,2^{-3(m-1/2)^2}.
\]
For \(m\ge7\), the ratio of successive shell bounds is less than
\(2^{-41}\).  Using \(\pi<22/7\), twice the first shell is less than
\(10^{-33}\).  We retain the weaker common tail radius \(10^{-32}\).

The certified V8 interval for \(\tau_*\), rational interval evaluation of
all retained exponential terms, and \cref{eq:v12-ratio-diff} then give
\cref{eq:v12-J,eq:v12-detJ}.  In particular the determinant interval excludes
zero.  The same certificate gives
\[
 \|J_\Theta\|_F^2<4.18.
\]
If \(s_1\ge s_2\ge s_3\) are the singular values, then
\[
 |\det J_\Theta|=s_1s_2s_3,
 \qquad
 s_1s_2\le\frac{s_1^2+s_2^2}{2}
 \le\frac{\|J_\Theta\|_F^2}{2}.
\]
Therefore
\[
 s_3\ge\frac{2|\det J_\Theta|}{\|J_\Theta\|_F^2}>0.35,
\]
which proves \cref{eq:v12-inverse-margin}.  The holomorphic inverse function
theorem now gives the local biholomorphism.
\end{proof}

\begin{remark}[Agreement with the global theta picture is not used in the proof]
For genus two with the corresponding theta marking, the second-order theta-null
map is classically known to give global projective moduli coordinates.  The
present theorem does not import that stronger statement: the local conclusion
is proved directly at the actual arithmetic fibre by the nonzero Jacobian
\cref{eq:v12-detJ}.
\end{remark}

\subsection{From the pure source projector back to the marked period matrix}
V2 proved, for the characteristic fixed independently in V8,
\begin{equation}
 r_\sigma=(-1)^{\sigma_1}
 t_{\sigma+(1,0)}.
 \label{eq:v12-r-t}
\end{equation}
Thus \(r=P_\kappa t\) for one fixed signed permutation matrix
\(P_\kappa\).  The normalized pure source is
\[
 \rho_\Theta=\frac{rr^*}{r^*r}.
\]

\begin{theorem}[Local reconstruction from the actual marked pure source]
\label{thm:v12-source-reconstruction}
Fix the V8 theta marking.  There are neighborhoods \(\mathscr U\) of the
actual marked genus-two period point and \(\mathscr R\) of the actual pure
source projector such that
\begin{equation}
 \boxed{
 \tau\longmapsto\rho_\Theta(\tau)
 }
 \label{eq:v12-source-map}
\end{equation}
is injective on \(\mathscr U\), with a real-analytic inverse on its image.
Equivalently, the marked rank-one theta source locally determines the complete
marked principally polarized abelian surface.  On the Jacobian locus it
therefore locally determines the marked curve by the classical Torelli theorem.
\end{theorem}
\begin{proof}
A rank-one Hermitian projector determines its complex line uniquely; hence
\(\rho_\Theta\) determines \([r]\in\mathbb P^3\).  The fixed matrix
\(P_\kappa\) in \cref{eq:v12-r-t} determines \([t]\).  The actual
\(t_{00}\) is nonzero: the certificate gives
\[
 |t_{00}|^2>0.99.
\]
Thus \([t]\) determines the affine coordinate \(\mathcal U(\tau)\) near the
actual point.  Apply \cref{thm:v12-local-inverse}.  The final curve-level
statement is the classical Torelli implication from the principally polarized
Jacobian to the curve; no new curve reconstruction theorem is claimed here.
\end{proof}

\begin{corollary}[Infinitesimal deformation is readable from source motion]
\label{cor:v12-tangent}
Let \(s\mapsto\rho_\Theta(s)\) be a differentiable marked source path lying in
the local theta-source image, and let \(u(s)=\mathcal U(\tau(s))\) be its
affine projective coordinate.  At the actual fibre,
\begin{equation}
 \boxed{
 \dot q(0)=J_\Theta^{-1}\dot u(0),
 \qquad
 \|\dot q(0)\|_2<2.86\,\|\dot u(0)\|_2.}
 \label{eq:v12-tangent-reconstruction}
\end{equation}
Thus the tangent direction of any independently occurring source motion is
recoverable.  A static source point alone, however, supplies no preferred
\(\dot u\) and therefore no preferred one-parameter transport.
\end{corollary}
\begin{proof}
Differentiate \(u=\mathcal U(q)\) and use
\cref{eq:v12-inverse-margin}.  The last sentence is a logical distinction:
the inverse function theorem reconstructs a tangent from a supplied tangent;
it does not manufacture a tangent vector from a point.
\end{proof}

\section{The projective source metric and a quantitative observability margin}
\label{sec:v12-source-metric}
The projective source itself supplies a positive metric on the marked moduli
germ, independently of the Hodge and theta-section metrics compared in
V7--V10.  In the affine coordinate \(u\in\mathbb C^3\), the Fubini--Study
matrix is
\begin{equation}
 G_{\rm FS}(u)
 =\frac{(1+\|u\|^2)I-\bar u\,u^{\mathsf T}}
        {(1+\|u\|^2)^2}.
 \label{eq:v12-fs}
\end{equation}
Define the pullback source metric
\begin{equation}
 \boxed{g_{\rm src}=J_\Theta^*G_{\rm FS}(\mathcal U(\tau_*))J_\Theta.}
 \label{eq:v12-source-metric}
\end{equation}

\begin{theorem}[A positive moduli metric reconstructed from source variation]
\label{thm:v12-source-metric}
At the actual arithmetic fibre, \(g_{\rm src}\) has certified midpoint
\begin{equation}
 g_{\rm src}\approx
 \begin{pmatrix}
 0.50026223&0.50969507-0.00776473i&
 0.13230742-0.06551860i\\
 0.50969507+0.00776473i&1.32836825&
 0.54227318-0.05880551i\\
 0.13230742+0.06551860i&
 0.54227318+0.05880551i&0.48064349
 \end{pmatrix}.
 \label{eq:v12-gsrc-matrix}
\end{equation}
Its leading principal minors satisfy
\begin{align}
 (g_{\rm src})_{11}&>0.500262233366886981,\notag\\
 \det(g_{\rm src})_{1:2,1:2}&>0.404683109114289521,\notag\\
 \det g_{\rm src}&>0.094211873998027970.
 \label{eq:v12-principal-minors}
\end{align}
Hence it is positive definite.  More quantitatively,
\begin{equation}
 \boxed{
 0.0706\,I\prec g_{\rm src}\prec2.31\,I
 }
 \label{eq:v12-source-floor}
\end{equation}
in the declared Euclidean period coordinates.

The source metric is not a scalar multiple of the standard Siegel metric at
this point.  Indeed
\begin{equation}
 \boxed{
 -0.065518601724154544
 <\Im(g_{\rm src})_{13}
 <-0.065518601724154543,
 }
 \label{eq:v12-imaginary-witness}
\end{equation}
whereas the Siegel metric
\(\operatorname{Tr}(Y^{-1}\delta\tau\,Y^{-1}\delta\bar\tau)\),
\(Y=\operatorname{Im}\tau_*\), has a real coefficient matrix in the real
symmetric basis
\((E_{11},E_{12}+E_{21},E_{22})\).
\end{theorem}
\begin{proof}
The rank-one projector map to \(\mathbb P^3\) carries the standard
Fubini--Study metric, whose affine matrix is \cref{eq:v12-fs}.  Pulling it back
through the exact differential \cref{eq:v12-J} gives
\cref{eq:v12-source-metric}.  The same rational interval evaluation used in
\cref{thm:v12-local-inverse} gives
\cref{eq:v12-gsrc-matrix,eq:v12-principal-minors}; Sylvester's criterion proves
positivity.

Let \(0<\lambda_1\le\lambda_2\le\lambda_3\) be its eigenvalues.  The
certificate gives
\(\operatorname{Tr}g_{\rm src}<2.31\).  Since
\[
 \lambda_2\lambda_3
 \le\left(\frac{\lambda_2+\lambda_3}{2}\right)^2
 \le\left(\frac{\operatorname{Tr}g_{\rm src}}2\right)^2,
\]
we obtain
\[
 \lambda_1
 =\frac{\det g_{\rm src}}{\lambda_2\lambda_3}
 \ge\frac{4\det g_{\rm src}}{(\operatorname{Tr}g_{\rm src})^2}
 >0.0706.
\]
The upper bound follows from \(\lambda_3\le\operatorname{Tr}g_{\rm src}<2.31\).
Finally \cref{eq:v12-imaginary-witness} excludes proportionality to the real
Siegel coefficient matrix.
\end{proof}

\begin{remark}[What the new metric does]
The inequality \cref{eq:v12-source-floor} is a finite local observability
margin: every nonzero infinitesimal period deformation changes the marked pure
source by a quantitatively nonzero Fubini--Study amount.  It does not identify
this projective-source metric with the Hodge metric on \(F\), the algebraic
metric on the quadratic source of V7, the Chern metrics of V9, or a physical
renewal source metric.  Those comparisons remain separate tests.
\end{remark}

\section{The theta-fibre Frobenius ambiguity is one-dimensional}
\label{sec:v12-line-lift}
The local inverse removes a possible ambiguity in the interpretation of the
missing V11 theta-line datum.  The marked projective source already determines
the local moduli point.  What remains at line level is not another matrix of
moduli parameters.

\begin{lemma}[Lifts of a projective source transport form a line torsor]
\label{lem:v12-line-torsor}
Let \(X\) be a connected local source chart and let \(\mathscr S\to X\) be
the tautological source line.  Let \(f:X\to X\) be a fixed projective source
transport.  If a nowhere-singular linear lift
\[
 \widetilde f:f^*\mathscr S\longrightarrow\mathscr S
\]
exists, then every other such lift is uniquely
\(g\widetilde f\) for one \(g\in\mathcal O_X^\times\).  At a fixed fibre the
ambiguity is one nonzero scalar.
\end{lemma}
\begin{proof}
The ratio of two nonzero maps between the same pair of line bundles is a
nowhere-vanishing scalar function.  Conversely multiplication by any such
function preserves the underlying projective map.  Fibrewise this is exactly
the free transitive action of the multiplicative group on the nonzero maps
between two one-dimensional vector spaces.
\end{proof}

\begin{theorem}[Projective completeness reduces the unresolved theta Frobenius to one scalar]
\label{thm:v12-line-scalar}
Near the actual arithmetic fibre, once an arithmetic projective transport of
the marked pure theta source has been independently specified, its induced
transport of the marked period matrix is fixed by
\cref{thm:v12-source-reconstruction}.  Any compatible linear transport of the
actual theta source line then differs from any other by one nonzero scalar at
each fixed fibre.

Moreover V10's canonical line identity
\[
 \Lambda_\Theta\simeq\mathscr P^{-3}
\]
shows that a chosen line transport on the native cotangent line \(\mathscr P\)
with scalar \(c\) induces scalars
\begin{equation}
 \boxed{
 c_{\Lambda_\Theta}=c^{-3},
 \qquad
 c_{\Lambda_\Theta^2}=c^{-6}.}
 \label{eq:v12-line-weight}
\end{equation}
Thus the missing theta-fibre Frobenius in V11 is equivalent, at line level, to
one multiplicative cotangent calibration; it is not an unobserved
four-by-four or ten-by-ten transport matrix.
\end{theorem}
\begin{proof}
The first statement combines the local inverse
\cref{thm:v12-source-reconstruction} with
\cref{lem:v12-line-torsor}.  The second follows functorially by taking the
indicated tensor powers of V10's canonical line isomorphism.  No claim is made
that crystalline Frobenius itself supplies the required same-line
linearization without this datum; V5 already showed that native section
Frobenius changes theta degree.
\end{proof}

\begin{theorem}[Projective source data cannot determine the missing scalar]
\label{cth:v12-projective-no-scalar}
The complete marked projective source, all of its Fubini--Study moments, and
its locally reconstructed period matrix do not determine an absolute
linearization of the theta-fibre line.
\end{theorem}
\begin{proof}
Fix any nonzero line map covering the same projective transport.  Multiplying
it by a constant \(a\in\mathbb C^\times\) (or by \(a\in\mathbb Q_p^\times\)
in a fixed local arithmetic fibre) leaves the projective source action and all
projective source moments unchanged, while changing the line-level map unless
\(a=1\).  Therefore no invariant of the projective source alone can recover
that scalar.  An additional line normalization, metric, divisor linearization,
or physical calibration is necessary.
\end{proof}

\begin{remark}[Why this is progress rather than another obstruction]
V11 left the theta-fibre Frobenius as an unspecified line-level comparison.
The present continuation proves that after the actual marked source has been
retained, this missing datum is exactly one multiplicative scalar, and V10
identifies the native arithmetic line on which it may be measured.  The large
matrix part of the local geometry is already reconstructible from the source;
only the scalar lift remains independent.
\end{remark}

\section{V12 proof ledger and the next comparison}
\label{sec:v12-status}

\subsection{Proved in this continuation}
The marked second-order theta-null map has a certified nonzero Jacobian at the
actual arithmetic period matrix, with a quantitative inverse margin.  The
marked pure rank-one source therefore locally reconstructs the full
three-complex-dimensional period germ.  A projective Fubini--Study metric on
the source pulls back to a positive moduli metric with an explicit lower
margin, and this source metric is not the Siegel metric.  Finally the remaining
theta-fibre lift is classified as one multiplicative line scalar, equivalently
one cotangent-line scalar through V10's native line identity.

\subsection{Inherited and external inputs}
We use V2's exact signed-permutation relation between the squared theta source
and second-order theta constants, V8's certified period matrix and marking,
and V10's line identity.  The theta-series notation and standard convergence
background are cited in \cite{DLMFTheta}; the local Jacobian and all displayed
bounds are independently certified here.  The curve-level interpretation of
the reconstructed principally polarized Jacobian uses the classical Torelli
theorem and is not reproved.  No crystalline matrix from V11 is used to fit the
complex Jacobian calculation.

\subsection{Executable certificate}
\path{Arithmetic_Loading_Duality_Research_Notes_V12_checks.py} evaluates the
four theta constants and their three first derivatives from the certified V8
period intervals using rational interval arithmetic.  The common truncation
error is bounded above by \(10^{-32}\), larger than the analytic shell bound
proved in \cref{thm:v12-local-inverse}.  It certifies the determinant,
singular-value margin, source metric, principal minors, non-Siegel witness,
and byte-for-byte preservation of the V11 mathematical body.  The finite
certificate does not machine-prove Torelli or any Frobenius/Wilson theorem.

\subsection{Sharper unresolved interface}
The chosen one-parameter family of V7 is no longer needed to identify the
\emph{point} in genus-two moduli: the marked pure source already does that
locally.  What remains independent is a source \emph{motion} or transport law,
together with one theta-line scalar when a linear rather than projective
transport is required.  Thus the next noncircular arithmetic--renewal
comparison should acquire, on the physical side, the marked source velocity or
finite projective transport and its one line normalization, then compare these
with the arithmetic source metric and the already acquired Frobenius/Gauss--Manin
objects.  Equality of compact group types is still insufficient.

The central Store signature, physical clock, accepted field law, actual mixed
renewal source Gram, and Frobenius--Wilson intertwiner remain unacquired.  No
Standard-Model numerical coupling, Einstein limit, GRH estimate, or quantum
continuum assertion is inferred.

\paragraph{Version integrity.}
All V1--V11 mathematical bodies and their historical proof-status ledgers are
preserved byte for byte.  V12 adds only the local moduli inverse, projective
source metric, and scalar-lift classification described above.  Future
continuations must follow this block.
% ====================== END IMMUTABLE V12 MATHEMATICAL BODY ======================
% ===================== BEGIN IMMUTABLE V13 MATHEMATICAL BODY =====================
\section{V13: upstream curve selection and native cotangent Frobenius}
\label{sec:v13-context}

\subsection{Purpose and route correction}
V12 closed a local inverse problem: once the marked pure theta source of the
chosen arithmetic curve is present, its projective source ray locally
reconstructs the complete genus-two period point, and only one line scalar is
left when a linear theta-fibre transport is requested.  Two logically earlier
questions remained.

First, the arithmetic curve
\[
 C_{\rm tet}:\quad y^2=x(x^4-x-1)
\]
was an explicit arithmetic realization of the pointed tetrahedral packet, but
it had not been proved to be selected by the structural requirements that
motivated it.  If those requirements admit a family of inequivalent curves,
then a duality theorem may use $C_{\rm tet}$ as an independently loaded
arithmetic source, but it may not claim that tetrahedral symmetry alone has
reconstructed the curve.

Second, V11 acquired the crystalline Frobenius on the rank-four cohomology and
its rank-ten quadratic lift, while V12 reduced the missing theta-fibre
linearization to one scalar.  Absolute Frobenius does not naturally preserve a
theta line at fixed tensor degree, so choosing this scalar directly from the
projective source would be circular.  The two-pointed curve nevertheless
already carries the descended eighth-power cotangent tensor of V10.  We show
that requiring Frobenius to preserve this existing tensor fixes the scalar
canonically on the $p$-adic disk of the actual fibre.

The result is therefore a route correction in both directions: the isolated
curve is less canonical than the earlier structural language might suggest,
whereas the residual line scalar is more rigid once the already occurring
cotangent tensor is retained.

\section{The pointed tetrahedral packet does not select a unique arithmetic curve}
\label{sec:v13-family}

For a parameter $t$ in a characteristic-zero field, put
\begin{equation}
 C_t:\quad y^2=x(x^4-x-t),
 \qquad g_t(x)=x^4-x-t.
 \label{eq:v13-family}
\end{equation}
We retain the two marked Weierstrass points
\[
 W_\infty=\infty,\qquad W_0=(0,0).
\]
The previously used curve is $C_1=C_{\rm tet}$.

\begin{lemma}[An infinite arithmetic $S_4$ family]
\label{lem:v13-S4-family}
Let $t\in\mathbb Z$ satisfy
\[
 t\equiv1\pmod{14}.
\]
Then $C_t$ is a smooth genus-two curve, and
\begin{equation}
 \boxed{\operatorname{Gal}(g_t/\mathbb Q)\cong S_4.}
 \label{eq:v13-S4}
\end{equation}
Moreover the marked point $W_0$ determines a nonzero rational two-torsion
class
\[
 \tau_0=[W_0-W_\infty]\in\operatorname{Jac}(C_t)[2](\mathbb Q),
\]
and the rational odd theta characteristic
$\kappa_\infty=\mathcal O_{C_t}(W_\infty)$ satisfies
$q_{\kappa_\infty}(\tau_0)=0$.
\end{lemma}

\begin{proof}
The quartic discriminant is
\begin{equation}
 \operatorname{disc}(g_t)=-(256t^3+27).
 \label{eq:v13-quartic-discriminant}
\end{equation}
For $t\equiv1\pmod{14}$ it is nonzero, and $t\ne0$; hence
$xg_t(x)$ is square-free.  Its degree is five, so the smooth projective
hyperelliptic model has genus two.

Modulo two, every such $g_t$ is
\[
 x^4+x+1.
\]
It has no root in $\mathbb F_2$.  The only monic irreducible quadratic over
$\mathbb F_2$ is $x^2+x+1$, whose square is $x^4+x^2+1$; therefore
$x^4+x+1$ is irreducible.  Since the discriminant is odd, Frobenius at two
has cycle type $(4)$ in the action on the four roots.

Modulo seven, $t\equiv1$ gives
\[
 g_t(x)=(x-3)(x^3+3x^2+2x-2).
\]
The cubic has no root in $\mathbb F_7$: its values at
$0,1,\ldots,6$ are respectively
\[
 5,4,1,2,6,5,5.
\]
The discriminant is nonzero modulo seven because
$256+27\equiv3\pmod7$.  Frobenius at seven therefore has cycle type $(1)(3)$.
The four-cycle already makes the Galois action transitive.  A transitive
subgroup of $S_4$ containing a three-cycle is either $A_4$ or $S_4$ once a
four-cycle is also present, and the four-cycle is odd.  Thus the group is
$S_4$.

Finally
\[
 \operatorname{div}(x)=2W_0-2W_\infty,
\]
so $\tau_0$ is rational two-torsion and is nonzero because $W_0\ne W_\infty$.
The odd-degree genus-two model has
$K_{C_t}\sim2W_\infty$, hence
$\kappa_\infty=\mathcal O(W_\infty)$ is a theta characteristic with one
section and is odd.  Moreover
\[
 \kappa_\infty\otimes\tau_0\simeq\mathcal O(W_0),
\]
which also has one section.  The usual theta quadratic refinement therefore
gives
\[
 q_{\kappa_\infty}(\tau_0)
 =h^0(\mathcal O(W_0))+h^0(\mathcal O(W_\infty))
 \equiv0\pmod2.
\]
\end{proof}

\begin{theorem}[Two-pointed isomorphism classification of the family]
\label{thm:v13-family-isom}
Let $k$ be an algebraically closed field of characteristic zero and let
$s,t\in k^\times$ be such that $C_s$ and $C_t$ are smooth.  An isomorphism
\[
 (C_t,W_0,W_\infty)\simeq(C_s,W_0,W_\infty)
\]
exists if and only if
\begin{equation}
 \boxed{s=a t\quad\text{for some }a\in\mu_3.}
 \label{eq:v13-isom-classification}
\end{equation}
When this holds, one such isomorphism is
\begin{equation}
 (x,y)\longmapsto(ax,ay).
 \label{eq:v13-explicit-isom}
\end{equation}
Consequently the scalar
\begin{equation}
 \boxed{\mathfrak m(C_t,W_0,W_\infty):=-t^3}
 \label{eq:v13-moduli-scalar}
\end{equation}
is a complete invariant of this two-pointed one-parameter family over
$\overline{\mathbb Q}$.  Over $\mathbb Q$, distinct rational parameters give
nonisomorphic two-pointed curves.
\end{theorem}

\begin{proof}
For a genus-two curve the hyperelliptic involution is intrinsic, so every
isomorphism descends to a projective transformation of the hyperelliptic
quotient.  Preserving $W_0$ and $W_\infty$ forces the descended map to fix
$0$ and $\infty$, hence it has the form $x\mapsto ax$ with $a\in k^\times$.
The monic polynomial whose roots are the images of the four residual branch
points is
\[
 a^4 g_t(X/a)=X^4-a^3X-a^4t.
\]
It must equal $g_s(X)=X^4-X-s$.  Therefore $a^3=1$ and
$s=a^4t=at$.

Conversely assume $a^3=1$ and $s=at$.  Substituting
$X=ax$, $Y=ay$ gives
\[
 X(X^4-X-s)
 =a x\bigl(a^4x^4-a x-a t\bigr)
 =a^2x(x^4-x-t)
 =Y^2,
\]
so \cref{eq:v13-explicit-isom} is an isomorphism preserving both marked
points.  Equality of $-t^3$ and $-s^3$ is equivalent to $s/t\in\mu_3$,
proving completeness.  Over $\mathbb Q$ the only rational cube root of unity
is one.
\end{proof}

\begin{corollary}[Tetrahedral symmetry is not a curve-selection theorem]
\label{cor:v13-nonselection}
There are infinitely many pairwise nonisomorphic two-pointed curves over
$\mathbb Q$ carrying all of the following data simultaneously:
\begin{enumerate}[label=(\roman*)]
\item a rational odd theta characteristic at infinity;
\item a rational nonzero singular two-torsion class $[W_0-W_\infty]$;
\item residual Galois group $S_4$ on the other four finite Weierstrass points.
\end{enumerate}
In particular these structural requirements do not select $C_{\rm tet}=C_1$.
\end{corollary}

\begin{proof}
Take the infinite progression $t=1+14n$ with $n\ge0$.  The data are supplied
by \cref{lem:v13-S4-family}, while pairwise nonisomorphism over $\mathbb Q$
follows from \cref{thm:v13-family-isom}.
\end{proof}

\begin{scope}[Correction to the role of the explicit curve]
The earlier arithmetic tetrahedron remains a valid explicit arithmetic source,
and every theorem proved for its actual theta section remains unchanged.  The
new point is logical: $C_{\rm tet}$ is an independently loaded arithmetic
realization, not an output of the bare $S_4$-tetrahedral and pointed-torsion
conditions.  A future paper should therefore distinguish
``canonical arithmetic loading from chronology'' from
``selection of the particular genus-two arithmetic source.''
\end{scope}

\section{The marked pure source gives a local selector once it actually occurs}
\label{sec:v13-source-selector}

V12 proved that the marked pure theta source is locally moduli-complete at
$C_1$.  Restrict that local inverse to the family \cref{eq:v13-family}.

\begin{corollary}[A source-readable local curve coordinate]
\label{cor:v13-source-modulus}
Let $\rho_\Theta(t)$ be the independently marked pure theta-source projector
for $C_t$ in the V12 marking, for $t$ near one.  There exists a neighbourhood
of $\rho_\Theta(1)$ on the source image and a holomorphic scalar
$\mathfrak m_\Theta$ such that
\begin{equation}
 \boxed{
 \mathfrak m_\Theta(\rho_\Theta(t))=-t^3.
 }
 \label{eq:v13-source-modulus}
\end{equation}
Its derivative along the actual family is
\begin{equation}
 \boxed{
 \frac{d}{dt}\mathfrak m_\Theta(\rho_\Theta(t))\bigg|_{t=1}=-3.
 }
 \label{eq:v13-source-modulus-derivative}
\end{equation}
Thus the marked pure source distinguishes the nearby two-pointed curves and
recovers the family parameter modulo the global $\mu_3$ ambiguity of
\cref{thm:v13-family-isom}.
\end{corollary}

\begin{proof}
By \cref{thm:v12-source-reconstruction}, the marked projective pure source
locally reconstructs the full marked period point and hence the corresponding
local genus-two moduli point.  Restrict this local inverse to the holomorphic
curve $t\mapsto C_t$ and compose with the algebraic function
$\mathfrak m=-t^3$ from \cref{thm:v13-family-isom}.  This gives
\cref{eq:v13-source-modulus}.  Differentiation gives
\cref{eq:v13-source-modulus-derivative}.
\end{proof}

\begin{remark}[Occurrence versus recognition]
The corollary is a recognition theorem, not a generation theorem.  Once the
actual marked source occurs, it identifies the nearby curve.  It does not show
that pointed chronology, the abstract $K_4$ lattice, or the bare arithmetic
loading produces the value $\mathfrak m=-1$ corresponding to $C_1$.
This is the upstream occurrence/selection datum that remains after the local
inverse has been closed.
\end{remark}

\section{The native cotangent tensor fixes the residual Frobenius line scalar}
\label{sec:v13-line-frobenius}

We now return to the one-dimensional ambiguity isolated in V12.  The
construction uses no target Wilson system.  It uses only the two-pointed
cotangent line and the descended eighth-power tensor already constructed in
V10.

Let $p\in\{3,5\}$, work over $\mathbb Q_p$, and use the open $p$-adic disk
\[
 \mathbb D_p=\{t\in\mathbb Q_p:|t-1|_p<1\}.
\]
The base Frobenius lift is
\begin{equation}
 \sigma_p(t)=t^p.
 \label{eq:v13-base-frobenius}
\end{equation}
Let $\mathscr P$ be the cotangent line at $W_\infty$.  In V10's native frame
$e=du|_\infty$, the descended tensor and finite-flat connection are
\begin{equation}
 s_8=-t e^8,
 \qquad
 \nabla^{\rm ff}e=-\frac18\frac{dt}{t}\,e.
 \label{eq:v13-native-line}
\end{equation}

\begin{lemma}[The native connection is characterized by the descended tensor]
\label{lem:v13-native-connection}
On a simply connected open set where $s_8$ is nonzero, there is a unique line
connection on $\mathscr P$ for which $s_8$ is parallel.  In the frame $e$ it
is the connection in \cref{eq:v13-native-line}.
\end{lemma}

\begin{proof}
Write a line connection as $\nabla e=Ae$.  Then
\[
 0=\nabla^{\otimes8}(-t e^8)
 =-dt\,e^8-8t A e^8,
\]
so necessarily $A=-dt/(8t)$.  This coefficient indeed makes the displayed
expression vanish.
\end{proof}

\begin{theorem}[Unique normalized native Frobenius on the cotangent line]
\label{thm:v13-native-frobenius}
There exists a unique analytic isomorphism
\begin{equation}
 \varphi_p^{\rm ff}:\sigma_p^*\mathscr P\longrightarrow\mathscr P
 \label{eq:v13-line-frob-map}
\end{equation}
that satisfies the three conditions
\begin{enumerate}[label=(\roman*)]
\item $\varphi_p^{\rm ff}$ is horizontal for $\nabla^{\rm ff}$;
\item $(\varphi_p^{\rm ff})^{\otimes8}$ carries $\sigma_p^*s_8$ to $s_8$;
\item its scalar at $t=1$ is one.
\end{enumerate}
It is given in the native frame by
\begin{equation}
 \boxed{
 \varphi_p^{\rm ff}(\sigma_p^*e)
 =c_p(t)e,
 \qquad
 c_p(t)=t^{(1-p)/8}
 =\exp\!\left(\frac{1-p}{8}\log t\right).
 }
 \label{eq:v13-cp}
\end{equation}
Without condition (iii), the compatible maps form the torsor
$\mu_8\,c_p(t)$ after extending coefficients to contain the required eighth
roots of unity.
\end{theorem}

\begin{proof}
Because $p$ is odd, $8$ is a unit in $\mathbb Z_p$.  On $\mathbb D_p$ the
$p$-adic logarithm converges and belongs to the maximal ideal; the exponential
therefore defines the analytic function in \cref{eq:v13-cp}.  It satisfies
\begin{equation}
 c_p(t)^8=t^{1-p}.
 \label{eq:v13-eighth-root}
\end{equation}
Now
\[
 \sigma_p^*s_8=-t^p(\sigma_p^*e)^8,
\]
so preservation of $s_8$ is equivalent to
$-t^pc_p^8=-t$, precisely \cref{eq:v13-eighth-root}.

For horizontality, the source connection coefficient is
\[
 \sigma_p^*\!\left(-\frac18\frac{dt}{t}\right)
 =-\frac p8\frac{dt}{t}.
\]
A line map with scalar $c_p$ is horizontal exactly when
\[
 d\log c_p-\frac18\frac{dt}{t}
 =-\frac p8\frac{dt}{t},
\]
or
\[
 d\log c_p=\frac{1-p}{8}\frac{dt}{t},
\]
which holds for \cref{eq:v13-cp}.

If $c$ is another section-compatible analytic scalar, then
$(c/c_p)^8=1$.  Its image lies in the finite group of eighth roots of unity,
so on the connected disk it is constant.  Condition (iii) makes that constant
one, proving uniqueness.
\end{proof}

\begin{corollary}[The theta-fibre scalar is fixed by the native line rule]
\label{cor:v13-theta-line-frobenius}
Use the algebraic line identity inherited from V10,
\[
 \Lambda_\Theta\simeq\mathscr P^{-3}.
\]
The normalized native Frobenius of \cref{thm:v13-native-frobenius} induces
\begin{equation}
 \boxed{
 c_{\Lambda_\Theta,p}(t)=c_p(t)^{-3}=t^{3(p-1)/8},
 \qquad
 c_{\Lambda_\Theta^2,p}(t)=c_p(t)^{-6}=t^{3(p-1)/4}.
 }
 \label{eq:v13-theta-line-scalars}
\end{equation}
In particular both scalars equal one at the actual fibre $t=1$.  Explicitly,
\begin{align*}
 p=3:&\quad c_3=t^{-1/4},\quad
 c_{\Lambda_\Theta,3}=t^{3/4},\quad
 c_{\Lambda_\Theta^2,3}=t^{3/2},\\
 p=5:&\quad c_5=t^{-1/2},\quad
 c_{\Lambda_\Theta,5}=t^{3/2},\quad
 c_{\Lambda_\Theta^2,5}=t^3.
\end{align*}
\end{corollary}

\begin{proof}
Tensor powers multiply the line exponent.  Substitution of
$c_p=t^{(1-p)/8}$ gives \cref{eq:v13-theta-line-scalars}; the two displayed
prime cases are immediate.
\end{proof}

\begin{corollary}[Line-completed rank-ten horizontal Frobenius]
\label{cor:v13-rank10-completion}
Let $\mathcal F_p(t)$ be V11's acquired horizontal Frobenius germ on
$E=H^1_{\rm dR}$ and let
\[
 \mathcal S_p(t)=\operatorname{Sym}^2\mathcal F_p(t)
\]
be its rank-ten lift.  On
\[
 \widehat{\mathcal W}
 =\Lambda_\Theta^2\otimes\operatorname{Sym}^2E
\]
the native line rule determines the horizontal Frobenius
\begin{equation}
 \boxed{
 \widehat{\mathcal F}_{p,\Theta}(t)
 =t^{3(p-1)/4}\,\mathcal S_p(t).
 }
 \label{eq:v13-completed-frobenius}
\end{equation}
At $t=1$ its matrix is exactly the rank-ten matrix $S_p$ acquired in V11;
there is no remaining scalar ambiguity in this normalization.
\end{corollary}

\begin{proof}
The finite-flat cotangent connection has coefficient
$A_{\mathscr P}=-dt/(8t)$.  Hence the induced connection on
$\Lambda_\Theta^2\simeq\mathscr P^{-6}$ has coefficient
\[
 A_{\Lambda_\Theta^2}=\frac34\frac{dt}{t}.
\]
Its normalized Frobenius scalar $c_{\Lambda_\Theta^2,p}$ satisfies
\[
 d\log c_{\Lambda_\Theta^2,p}
 =\sigma_p^*A_{\Lambda_\Theta^2}-A_{\Lambda_\Theta^2}
 =\frac34(p-1)\frac{dt}{t},
\]
which is exactly the derivative of $t^{3(p-1)/4}$.  Tensoring this horizontal
line map with V11's horizontal $\mathcal S_p$ gives
\cref{eq:v13-completed-frobenius}.  Evaluation at $t=1$ removes the scalar.
\end{proof}

\begin{theorem}[Why this does not contradict the V12 scalar no-go]
\label{thm:v13-no-contradiction}
The projective theta source and its locally reconstructed period point still do
not determine an absolute theta-line Frobenius scalar.  The scalar in
\cref{eq:v13-theta-line-scalars} becomes canonical only after one retains the
additional two-pointed line datum $s_8=-t e^8$ and requires its normalized
horizontal preservation.  Forgetting $s_8$ restores exactly the multiplicative
line torsor of V12.
\end{theorem}

\begin{proof}
Multiplying a lift of the projective source transport by an arbitrary
nonvanishing scalar leaves every projective source observable unchanged, as
proved in V12.  By contrast, preservation of $s_8$ forces the eighth-power
equation \cref{eq:v13-eighth-root}; on the normalized disk its solutions are
the finite $\mu_8$ torsor of \cref{thm:v13-native-frobenius}, and the value at
the actual fibre fixes the remaining constant.  Thus the new rigidity comes
from additional line-level arithmetic data rather than from information that
V12 had declared absent.
\end{proof}

\begin{scope}[Native finite-flat Frobenius versus theta-section Frobenius]
The map \cref{eq:v13-line-frob-map} is the normalized Frobenius structure on
the native two-pointed cotangent line and its tensor powers.  Absolute
Frobenius on sections of a theta line naturally raises the theta tensor degree,
as emphasized in V5.  The present theorem therefore supplies the missing
\emph{line normalization} for V11's de-twisted quadratic comparison; it does
not claim a same-level Frobenius endomorphism of the four-dimensional complex
theta-section space, nor a Wilson holonomy.
\end{scope}

\section{V13 proof ledger and the sharper remaining interface}
\label{sec:v13-status}

\subsection{Proved in this continuation}
The tetrahedral pointed-torsion conditions used to motivate $C_{\rm tet}$ are
satisfied by infinitely many pairwise nonisomorphic two-pointed curves over
$\mathbb Q$.  Within the explicit family, the scalar $-t^3$ is the complete
two-pointed isomorphism invariant over $\overline{\mathbb Q}$, and V12's
source inverse makes that scalar locally readable from an actually occurring
marked pure theta source.  Separately, V10's descended eighth-power cotangent
tensor selects a unique normalized horizontal $p$-adic Frobenius on the
cotangent line at the two acquired primes.  Through the algebraic identity
$\Lambda_\Theta\simeq\mathscr P^{-3}$ this fixes the line scalar left open in
V11--V12 and completes the rank-ten horizontal Frobenius after the declared
finite-flat line normalization.

\subsection{Inherited statements not reproved}
The actual theta-source reconstruction and local inverse are V12 results; the
cotangent/theta line identities and the descended section $s_8=-t e^8$ are V10
results; the rank-four crystalline and rank-ten symmetric-square horizontal
Frobenius germs are V11 results.  The present continuation uses these objects
without repeating their period, theta, crystalline, or source-rank
computations.

\subsection{What remains open}
The strongest upstream selection problem is now explicit.  Bare chronology
and the tetrahedral $S_4$ packet do not select the value
$\mathfrak m=-1$ of the curve $C_1$; an occurrence or selection principle for
the marked pure source is still needed if the final paper is to claim that the
specific arithmetic curve is derived rather than independently loaded.

On the transport side, the native line scalar is no longer an unknown once the
finite-flat two-pointed normalization is declared.  The remaining cross-domain
problem is the independently occurring renewal source/transport itself: its
marked projective source, mixed metric, central Store signature, and physical
transfer must be compared with the arithmetic object.  Matching the compact
group $S(U(2)\times U(3))$ or the rank-ten carrier is not sufficient.
The physical clock, accepted law, Frobenius--Wilson intertwiner, and cofinal
coercive comparison remain unproved.

\paragraph{Next high-value attack.}
The next iteration should first test whether the renewal/chronology packet
contains an independently defined scalar or source occurrence that can be
compared with the curve selector $\mathfrak m_\Theta=-t^3$.  If not, the
honest paper-level statement should treat $C_{\rm tet}$ as an independently
chosen arithmetic realization and move directly to the finite mixed-source
comparison.  On the arithmetic side, further manipulation of the theta-line
scalar would now be circular: the normalized native scalar has already been
fixed by the strongest available intrinsic two-pointed line datum.

\paragraph{Version integrity.}
All V1--V12 mathematical bodies and their historical proof-status ledgers are
preserved byte for byte.  V13 adds only the curve-selection firewall, the
source-readable family invariant, and the native cotangent/theta-line
Frobenius normalization proved above.  Future continuations must follow this
block.
% ====================== END IMMUTABLE V13 MATHEMATICAL BODY ======================
% APPEND V14 AND LATER ITERATIONS HERE.











% ====================== BEGIN IMMUTABLE V14 MATHEMATICAL BODY ======================
\section{V14: the minimal two-triplet renewal completion of the arithmetic moduli source}
\label{sec:v14-context}

\subsection{Why the programme changes level here}
V13 left two logically different questions.  The first is global selection of
one arithmetic curve inside the family $C_t$; that question was reduced to an
additional scalar occurrence row and should not be reopened by fitting the
curve to the renewal target.  The second is the actual mixed-source question:
which renewal source can carry the full local information of the marked
arithmetic theta source already acquired in V8--V12?

The second question has a finite answer before any Wilson or continuum
comparison.  The marked pure theta source is locally moduli-complete by
\Cref{thm:v12-local-inverse}; hence its infinitesimal source has complex rank
three and real rank six.  By contrast, the resolved primitive one-cell renewal
score bank in the supplied Grand-Tensor monograph has real source rank three.
The same monograph also contains a strictly larger typed half-edge source with
two orthogonal copies of the tetrahedral standard module.  This continuation
proves that those two copies are not merely sufficient by dimension: they are
the \emph{minimal $S_4$-equivariant real carrier on which the arithmetic
complex structure can live}.  It then reduces an exact local arithmetic--
renewal duality to one finite mixed Gram and one complex-structure residual.

Nothing in this section identifies the two sources by declaration.  The
arithmetic marginal and the renewal marginal are already independently known;
the missing cross Gram remains an actual mixed occurrence measurement.

\subsection{The arithmetic local source has six real directions}
Let $\mathfrak H_2$ be the genus-two Siegel upper half space in the marked
period chart of V8--V12.  At the acquired arithmetic point $\tau_*$ put
\[
 \mathcal T_{\rm ar}:=T_{\tau_*}\mathfrak H_2\cong\C^3.
\]
The marked pure-source map to $\mathbf P^3$ has differential $J_\Theta$ and
pullback Fubini--Study metric $g_{\rm src}$ from
\Cref{eq:v12-source-metric}.  Regard $\mathcal T_{\rm ar}$ as a six-
dimensional real space and denote by $G_{\rm ar}^{\R}$ the realification of
$g_{\rm src}$:
\begin{equation}
 \boxed{
 G_{\rm ar}^{\R}
 =\begin{pmatrix}
    \Re g_{\rm src}&-\Im g_{\rm src}\\
    \Im g_{\rm src}& \Re g_{\rm src}
   \end{pmatrix}.}
 \label{eq:v14-realification}
\end{equation}
Multiplication by $i$ becomes the orthogonal complex structure
\begin{equation}
 \boxed{
 J_{\rm ar}=\begin{pmatrix}0&-I_3\\ I_3&0\end{pmatrix},
 \qquad J_{\rm ar}^2=-I_6.}
 \label{eq:v14-Jar}
\end{equation}

\begin{lemma}[Realification retains the arithmetic source floor]
\label{lem:v14-realification-spectrum}
If $0<\lambda_1\le\lambda_2\le\lambda_3$ are the eigenvalues of
$g_{\rm src}$, then the eigenvalues of $G_{\rm ar}^{\R}$ are
\[
 \lambda_1,\lambda_1,\lambda_2,\lambda_2,\lambda_3,\lambda_3.
\]
In particular V12 gives
\begin{equation}
 \boxed{G_{\rm ar}^{\R}\succ0.0706\,I_6.}
 \label{eq:v14-real-floor}
\end{equation}
\end{lemma}

\begin{proof}
Choose a unitary $U$ diagonalizing the Hermitian matrix $g_{\rm src}$.
Realification sends $U$ to an orthogonal transformation and sends each scalar
Hermitian block $[\lambda_j]$ to $\lambda_j I_2$.  The spectrum therefore
repeats each Hermitian eigenvalue twice.  The lower bound is
\Cref{eq:v12-source-floor}.
\end{proof}

\begin{theorem}[A one-triplet renewal source cannot encode the local arithmetic source]
\label{thm:v14-one-triplet-no-go}
Let $A:\R^6\to\mathcal H$ be any source synthesis with
$A^*A=G_{\rm ar}^{\R}$, and let $B:\R^3\to\mathcal H$ be any renewal source
of rank at most three on the same real Hilbert carrier.  Write $P_B$ for the
orthogonal projection onto $\operatorname{Ran} B$.  Then the arithmetic innovation after
shorting the renewal source,
\[
 \mathbb V_{{\rm ar}\mid B}
 :=A^*(I-P_B)A,
\]
satisfies
\begin{equation}
 \boxed{
 \rank\mathbb V_{{\rm ar}\mid B}\ge3,}
 \label{eq:v14-rank-obstruction}
\end{equation}
\begin{equation}
 \boxed{
 \Tr\mathbb V_{{\rm ar}\mid B}
 \ge 2\lambda_1+\lambda_2
 >3(0.0706)=0.2118.}
 \label{eq:v14-energy-obstruction}
\end{equation}
Consequently no differentiable local source duality from the complete marked
arithmetic theta source to one primitive real renewal triplet can be locally
injective or source-unitary.
\end{theorem}

\begin{proof}
Since $G_{\rm ar}^{\R}\succ0$, the map $A$ has rank six.  The component
$P_BA$ has rank at most three.  From
$A=P_BA+(I-P_B)A$ and the rank subadditivity inequality,
\[
 6=\rank A\le\rank(P_BA)+\rank((I-P_B)A),
\]
so $\rank((I-P_B)A)\ge3$.  Since
$\mathbb V_{{\rm ar}\mid B}=((I-P_B)A)^*((I-P_B)A)$, this proves
\eqref{eq:v14-rank-obstruction}.

For the trace bound, $\Tr\mathbb V_{{\rm ar}\mid B}$ is the squared
Hilbert--Schmidt error of approximating $A$ by the rank-at-most-three map
$P_BA$.  The Eckart--Young--Mirsky variational principle therefore bounds it
below by the sum of the three smallest squared singular values of $A$.  Those
squared singular values are precisely the eigenvalues of
$G_{\rm ar}^{\R}$, which by
\Cref{lem:v14-realification-spectrum} occur as
$\lambda_1,\lambda_1,\lambda_2,\lambda_2,\lambda_3,\lambda_3$.  The three
smallest sum to $2\lambda_1+\lambda_2$, proving
\eqref{eq:v14-energy-obstruction}.  The final inequality uses
$\lambda_1,\lambda_2>0.0706$.

A locally injective source identification would have injective differential,
which is impossible through a rank-three target.  A source-unitary
identification would force the displayed innovation to vanish, contradicting
the strict bound.
\end{proof}

\begin{scope}[What this obstruction applies to]
The theorem applies to one real rank-three source, including the primitive
one-cell score bank and one tetrahedral standard triplet.  It does not apply
to a renewal carrier containing several independent triplets, connected
pair-source directions, or other actual source innovations.  It therefore
selects the next source grade rather than ruling out an arithmetic--renewal
comparison altogether.
\end{scope}

% -----------------------------------------------------------------------------
\section{The two-standard renewal plane is representation-theoretically minimal}
\label{sec:v14-two-standard}

The supplied Grand-Tensor monograph proves, for the sourcewise-uniform
matching-axis half-edge packet, the existence of normalized source maps
\[
 E_0,A_0:W_0\longrightarrow\mathcal H_{\rm ha},
 \qquad W_0=\one_4^\perp\subset\R^4,
\]
which are $S_4$-equivariant and satisfy
\begin{equation}
 \boxed{
 E_0^*E_0=A_0^*A_0=I_{W_0},
 \qquad E_0^*A_0=0.}
 \label{eq:v14-inherited-EA}
\end{equation}
The two ranges are the entry and conditional matching-axis copies of the real
standard representation $[31]$ \cite{GT77}.  We use this theorem as an
inherited input and do not reprove its finite renewal calculation.

Put
\[
 \mathcal R_2:=\operatorname{Ran} E_0\oplus\operatorname{Ran} A_0.
\]

\begin{theorem}[Canonical typed complex structure on the half-edge plane]
\label{thm:v14-renewal-complex-structure}
Define $J_{\rm ren}:\mathcal R_2\to\mathcal R_2$ by
\begin{equation}
 \boxed{
 J_{\rm ren}(E_0x)=A_0x,
 \qquad
 J_{\rm ren}(A_0x)=-E_0x.}
 \label{eq:v14-Jren}
\end{equation}
Then
\[
 J_{\rm ren}^2=-I,
 \qquad J_{\rm ren}^*J_{\rm ren}=I,
 \qquad [J_{\rm ren},U_g]=0\quad(g\in S_4).
\]
Hence $\mathcal R_2$, with multiplication by $i$ represented by
$J_{\rm ren}$, is a complex three-dimensional $S_4$ module isomorphic to the
complex standard module $W_0\otimes_\R\C$.

After complexification the two eigenspace syntheses
\begin{equation}
 \boxed{
 \Gamma_+=\frac{E_0-iA_0}{\sqrt2},
 \qquad
 \Gamma_- =\frac{E_0+iA_0}{\sqrt2}}
 \label{eq:v14-chiral-sources}
\end{equation}
satisfy
\[
 \Gamma_+^*\Gamma_+=\Gamma_-^*\Gamma_-=I,
 \qquad
 \Gamma_+^*\Gamma_-=0,
\]
and have $J_{\rm ren}$ eigenvalues $+i$ and $-i$, respectively.
\end{theorem}

\begin{proof}
The orthogonal decomposition and the identities
\eqref{eq:v14-inherited-EA} make \eqref{eq:v14-Jren} well defined.  Applying
it twice gives $-I$, and the same identities show that it preserves the real
inner product.  Equivariance of $E_0$ and $A_0$ gives commutation with every
$S_4$ action.  The complex-module statement follows by using $J_{\rm ren}$ as
multiplication by $i$.  Direct expansion of
\eqref{eq:v14-chiral-sources} with
$E_0^*A_0=A_0^*E_0=0$ gives the three Gram identities; the eigenspace formulas
follow from \eqref{eq:v14-Jren}.
\end{proof}

\begin{theorem}[Two real standard copies are the minimum for an equivariant complex source]
\label{thm:v14-minimal-two-copies}
Let $W_0$ be the real irreducible standard representation of $S_4$.  On
$W_0^{\oplus m}$ an $S_4$-equivariant real complex structure exists if and
only if $m$ is even.  In particular one standard triplet cannot carry the
complex structure of the arithmetic local moduli source, whereas two copies
are sufficient and minimal.
\end{theorem}

\begin{proof}
The real standard representation is absolutely irreducible, so
$\End_{S_4}(W_0)=\R I$.  Therefore
\[
 \End_{S_4}(W_0^{\oplus m})\cong M_m(\R)\otimes I_{W_0}.
\]
An equivariant complex structure is thus a real matrix $J_m\in M_m(\R)$ with
$J_m^2=-I_m$.  If $m$ is odd, taking determinants gives
$(\det J_m)^2=(-1)^m=-1$, impossible over $\R$.  If $m=2r$ is even, the block
diagonal matrix with $r$ copies of
$\begin{psmallmatrix}0&-1\\1&0\end{psmallmatrix}$ has square $-I_m$.
The minimal case is $m=2$, realized concretely by
\Cref{thm:v14-renewal-complex-structure}.
\end{proof}

\subsection{The arithmetic branch-deformation carrier has the same complex type}
Fix a locally labelled four-tuple of nonzero residual branch points
$r=(r_1,r_2,r_3,r_4)$, with the two distinguished Weierstrass points fixed at
$0$ and $\infty$.  Infinitesimal logarithmic branch motion is
\[
 \dot\ell=(\dot r_1/r_1,\ldots,\dot r_4/r_4)\in\C^4.
\]
The residual coordinate rescaling $x\mapsto cx$ adds the same scalar to every
component.  Define
\begin{equation}
 \boxed{
 \mathcal D_r:=\C^4/\C\one_4.}
 \label{eq:v14-branch-tangent}
\end{equation}
Permutation of the four labels transports these fibres in the evident way.
This is a statement on the labelled branch-deformation bundle; it is not a
claim that the Galois group acts complex-linearly on one fixed Hermitian
fibre.

\begin{proposition}[The labelled arithmetic moduli tangent is the complex standard bundle]
\label{prop:v14-branch-standard}
For a smooth two-pointed genus-two branch configuration, the logarithmic map
from labelled branch motion modulo the residual scaling identifies the local
moduli tangent with $\mathcal D_r$.  Fibrewise,
\[
 \boxed{\mathcal D_r\cong W_0\otimes_\R\C,}
 \label{eq:v14-ar-standard}
\]
with the label permutations acting by the complexified standard
representation.  Via the local theta-source inverse of V12, the same complex
rank-three tangent is recovered from the marked pure theta source.
\end{proposition}

\begin{proof}
With $0$ and $\infty$ fixed, a nearby smooth hyperelliptic curve is specified
by the four distinct nonzero residual branch points.  The only remaining
M\"obius transformations preserving both marked points are scalings
$x\mapsto cx$.  Passing to logarithmic differentials therefore quotients
$\C^4$ by the common line $\C\one_4$, giving
\eqref{eq:v14-branch-tangent}.  The quotient of the permutation module by its
constant line is the standard module, proving
\eqref{eq:v14-ar-standard}.  V12 proves that the marked theta source is a
local holomorphic coordinate system for the same moduli germ, so its tangent
is identified with this carrier.  The qualification about different labelled
sheets is essential: permutation covariance transports tangent fibres rather
than manufacturing a fixed-fibre unitary Galois action.
\end{proof}

\begin{corollary}[Minimal module-level arithmetic--renewal match]
\label{cor:v14-module-match}
After a declared identification of the four residual arithmetic labels with
the four renewal opportunity labels, the arithmetic labelled-deformation
bundle and the typed half-edge plane have the same minimal complex
$S_4$-module type.  Equivalently, there is an $S_4$-equivariant complex-linear
bundle isomorphism
\[
 \boxed{
 \mathcal D\longrightarrow(\mathcal R_2,J_{\rm ren})}
 \]
which is unique up to multiplication by a nowhere-zero complex scalar on each
connected labelled chart.
\end{corollary}

\begin{proof}
Both sides are copies of the irreducible complex standard representation.
Schur's lemma makes the intertwiner space one dimensional over $\C$.  Any
nonzero intertwiner is an isomorphism.
\end{proof}

\begin{scope}[Module match is not source duality]
\Cref{cor:v14-module-match} is stronger than a dimension coincidence but
weaker than a physical source identification.  The arithmetic projective
source metric of V12 is not replaced by the renewal metric
\eqref{eq:v14-inherited-EA}; a common physical carrier and its mixed Gram are
still required.  Likewise, the label identification must occur as an actual
joint record or an independently justified descent map rather than being
chosen after inspecting the target.
\end{scope}

% -----------------------------------------------------------------------------
\section{A terminating mixed-Gram criterion for the first genuine local duality}
\label{sec:v14-mixed-test}

The preceding section identifies the smallest renewal carrier on which an
exact local arithmetic source duality could possibly live.  We now isolate the
single finite object that decides whether it actually does.

Let
\[
 A:\R^6\to\mathcal H_{\rm joint},
 \qquad
 B:\R^6\to\mathcal H_{\rm joint}
\]
be source syntheses for the realified arithmetic tangent and the normalized
renewal entry--axis plane on one common physical carrier.  Assume
\[
 A^*A=G_{\rm ar}^{\R},
 \qquad B^*B=I_6,
\]
and define the complete mixed Gram
\begin{equation}
 \boxed{C:=B^*A.}
 \label{eq:v14-cross-gram}
\end{equation}
Whiten the arithmetic source by
\[
 \widehat A=A(G_{\rm ar}^{\R})^{-1/2},
 \qquad
 \boxed{T:=B^*\widehat A=C(G_{\rm ar}^{\R})^{-1/2}.}
 \label{eq:v14-normalized-cross}
\]

\begin{theorem}[Exact mixed-source recognition]
\label{thm:v14-mixed-recognition}
Suppose the block Gram of $A$ and $B$ is positive, as it must be when both
sources occur on one Hilbert carrier.  Then
\begin{equation}
 \boxed{
 \widehat A^*(I-P_B)\widehat A=I-T^*T\succeq0,}
 \label{eq:v14-short-ar}
\end{equation}
\begin{equation}
 \boxed{
 B^*(I-P_{\widehat A})B=I-TT^*\succeq0.}
 \label{eq:v14-short-ren}
\end{equation}
The following are equivalent:
\begin{enumerate}[label=\textup{(\roman*)},leftmargin=2.8em]
\item the two normalized sources have the same range;
\item both source innovations in
      \eqref{eq:v14-short-ar}--\eqref{eq:v14-short-ren} vanish;
\item $T$ is unitary.
\end{enumerate}
If these conditions hold, the coefficient transport from arithmetic to
renewal is exactly $T$, since $\widehat A=BT$.
\end{theorem}

\begin{proof}
Because $B$ and $\widehat A$ are isometries,
$P_B=BB^*$ and $P_{\widehat A}=\widehat A\widehat A^*$.  Therefore
\[
 \widehat A^*(I-P_B)\widehat A
 =I-\widehat A^*BB^*\widehat A=I-T^*T,
\]
and similarly for \eqref{eq:v14-short-ren}.  Positivity follows because both
are Gram matrices.  Equality of the ranges makes $B^*\widehat A$ a unitary
change of orthonormal source frames.  Conversely, if $T$ is unitary then
\[
 \|\widehat Ax-BTx\|^2
 =\|x\|^2+\|Tx\|^2-2\Re\langle Tx,B^*\widehat Ax\rangle=0,
\]
so $\widehat A=BT$ and the ranges agree.  Vanishing of both shorts is plainly
equivalent to mutual range inclusion and hence equality.
\end{proof}

The arithmetic complex structure in the real period coordinates is
$J_{\rm ar}$ from \eqref{eq:v14-Jar}.  Use the ordered renewal coefficient
basis $(E_0,A_0)$ to identify its complex structure with
\[
 J_{\rm ren}^{\rm coef}
 =\begin{pmatrix}0&-I_3\\I_3&0\end{pmatrix}.
\]

\begin{theorem}[Terminating local arithmetic--renewal source test]
\label{thm:v14-terminating-test}
Define
\begin{equation}
 \boxed{
 \Delta_{\rm mix}:=
 \|I-T^*T\|_{\HS}^2
 +\|I-TT^*\|_{\HS}^2
 +\|TJ_{\rm ar}-J_{\rm ren}^{\rm coef}T\|_{\HS}^2.}
 \label{eq:v14-mixed-defect}
\end{equation}
Then
\[
 \boxed{\Delta_{\rm mix}=0}
\]
if and only if the two acquired six-real-dimensional sources are identical up
to a complex-linear source-fixing unitary on their common carrier.
Every nonzero term in \eqref{eq:v14-mixed-defect} returns a finite witness:
respectively arithmetic innovation, renewal innovation, or failure of the
complex structures to intertwine.
\end{theorem}

\begin{proof}
By \Cref{thm:v14-mixed-recognition}, the first two terms vanish exactly when
$T$ is unitary and the source ranges agree.  On that branch the third term
vanishes exactly when the coefficient unitary is complex linear.  Conversely,
an exact complex-linear source unitary makes all three residuals zero.  Each
term is a squared Hilbert--Schmidt norm of an explicitly reconstructed finite
matrix, so failure supplies the corresponding finite witness.
\end{proof}

\begin{corollary}[With a genuinely common $S_4$ action the mixed test reduces to a multiplicity plane]
\label{cor:v14-S4-reduction}
Assume, in addition, that one common physical $S_4$ action has been acquired
on the joint carrier and that both six-real-dimensional source presentations
transform as two copies of the same real standard module.  In standard
isotypic coordinates every equivariant normalized cross transport has the
form
\[
 \boxed{T=M\otimes I_{W_0},\qquad M\in M_2(\R).}
 \label{eq:v14-MtensorI}
\]
If it also intertwines the two complex structures, then
\[
 M=\begin{pmatrix}a&-b\\b&a\end{pmatrix}.
\]
If the source innovations vanish as well, then $a^2+b^2=1$; hence the entire
six-dimensional comparison reduces to one phase $a+ib\in U(1)$.
\end{corollary}

\begin{proof}
The commutant of two copies of the absolutely irreducible real standard
module is $M_2(\R)\otimes I_{W_0}$, giving
\eqref{eq:v14-MtensorI}.  Commutation with
$\begin{psmallmatrix}0&-1\\1&0\end{psmallmatrix}$ is equivalent to the
displayed two-parameter form.  Unitarity gives
$M^{\mathsf T}M=(a^2+b^2)I_2=I_2$.
\end{proof}

\begin{scope}[Why the symmetry reduction is conditional]
The arithmetic Galois action on the labelled branch data is naturally a
semilinear descent action, while the renewal $S_4$ action is a real linear
relabelling action on one physical carrier.  A common linear $S_4$ action is
therefore an additional comparison datum, not a consequence of equal abstract
group names.  Without that common action the complete $6\times6$ mixed Gram
in \eqref{eq:v14-cross-gram} must be acquired.
\end{scope}

% -----------------------------------------------------------------------------
\section{What V14 proves about the research programme}
\label{sec:v14-status}

\subsection{Proved in this continuation}
\begin{enumerate}[leftmargin=2.2em]
\item The full marked arithmetic theta-source germ has six real tangent
      directions with inherited floor $>0.0706$.
\item Every coupling to a single rank-three renewal source leaves at least
      three arithmetic innovation directions and source energy greater than
      $0.2118$ in the V12 normalization.
\item Two copies of the real tetrahedral standard module are the minimum
      $S_4$ carrier supporting an equivariant complex structure.
\item The already existing renewal entry--axis half-edge plane supplies
      exactly that minimum and has an explicit typed complex structure.
\item The labelled arithmetic branch-deformation bundle is the complex
      standard bundle, so the two sides have the same minimal complex module
      type without identifying their metrics or occurrence records.
\item Exact local source duality is decided by one finite normalized cross
      Gram $T$ and the residual \eqref{eq:v14-mixed-defect}.
\item If a common linear $S_4$ action is independently established, the mixed
      transport reduces to a $2\times2$ multiplicity matrix and, on the exact
      complex-unitary branch, to one $U(1)$ phase.
\end{enumerate}

\subsection{Inherited results not reproved}
V12 supplies local moduli completeness and the source metric
$g_{\rm src}$.  V13 supplies the curve-selection firewall.  The primitive
rank-three renewal tangent and the matching-axis two-standard source maps
$E_0,A_0$ with \eqref{eq:v14-inherited-EA} are imported from the supplied
Grand-Tensor monograph \cite{GT77}.  V14 uses these as fixed inputs and does
not repeat their finite score calculations.

\subsection{The next noncircular experiment}
The programme no longer needs another marginal arithmetic invariant before
attempting the first local source comparison.  The decisive missing object is
now the \emph{mixed} same-history source panel producing
\[
 C=B^*A.
\]
Separate arithmetic and renewal marginals cannot determine this matrix: two
joint Hilbert realizations can have identical marginal Grams while their
source ranges are orthogonal in one realization and identical in another.
Therefore fitting $C$ from the desired unitary is inadmissible.

The next attack should be one of the following, in this order:
\begin{enumerate}[label=\textup{(N\arabic*)},leftmargin=3em]
\item identify an actually occurring common record in the Renewal loading
      whose differentiated arithmetic and entry--axis readers give the mixed
      block $C$ directly;
\item failing that, prove a word-native compiler that reconstructs this mixed
      block from already occurring common-history words;
\item if neither exists, return the missing cross Gram as the exact source
      obstruction and move to the Frobenius--Wilson comparison only at the
      weaker representation level.
\end{enumerate}

The central Store/chronology signature may reduce the final phase or complex-
orientation ambiguity once a common mixed carrier exists, but it cannot repair
the rank-three obstruction of \Cref{thm:v14-one-triplet-no-go}.

\paragraph{Version integrity.}
All V1--V13 mathematical bodies and their historical proof-status ledgers are
preserved byte for byte.  V14 adds only the six-real-dimensional source
obstruction, the minimal two-standard completion theorem, and the terminating
mixed-Gram criterion above.  Future continuations must follow this block.
% ====================== END IMMUTABLE V14 MATHEMATICAL BODY ======================

% ====================== BEGIN IMMUTABLE V15 MATHEMATICAL BODY ======================
\section{V15: the complete mixed-source extension cone and the word-native boundary}
\label{sec:v15-context}

\subsection{Why the next step is an identifiability problem}
V14 reduced the first viable local arithmetic--renewal source comparison to a
mixed Gram
\[
 C=B^*A
\]
between the six-real-dimensional arithmetic moduli source and the normalized
renewal entry--axis plane.  Before acquiring new arithmetic invariants, one
must decide whether $C$ is already a deterministic function of the existing
history table.  The supplied Grand-Tensor monograph gives the exact criterion:
inserted mixed Grams are table-native when both source maps are fixed
word-native expressions in one represented flat history algebra; an external
projector, router, or source shadow is not reconstructed from the old word
moments alone \cite{GT77}.

The present arithmetic source is not currently supplied in that form.  Its
construction uses an independently selected genus-two source, its marked theta
section, and the period/moduli reconstruction of V8--V13.  In particular,
\Cref{cor:v13-nonselection} proves that the universal tetrahedral and pointed
arithmetic entrance does not select the curve.  The supplied Grand-Tensor
monograph correspondingly lists the cross Grams of independently loaded
physical sectors as an open occurrence problem.  This is a dependency audit,
not a theorem that no future word-native realization can exist.

We therefore first classify \emph{all} joint source realizations consistent
with the already acquired marginals.  This yields a complete finite
identifiability theorem.  On the branch where a common $S_4$ action and the
typed complex structures are independently established, the entire ambiguity
collapses to one point of the closed unit disk.  Exact source duality is the
unit-circle boundary of that disk.

% -----------------------------------------------------------------------------
\section{The complete positivity cone of joint source extensions}
\label{sec:v15-contraction-cone}

Let $E_A,E_B$ be finite-dimensional real or complex Hilbert coefficient
spaces, and let
\[
 G_A\succ0,\qquad G_B\succ0
\]
be two independently acquired source Grams.  For a proposed cross block
$C:E_A\to E_B$, put
\begin{equation}
 \mathcal G(C)
 :=\begin{pmatrix}G_B&C\\ C^*&G_A\end{pmatrix},
 \qquad
 T:=G_B^{-1/2}CG_A^{-1/2}.
 \label{eq:v15-block-gram}
\end{equation}

\begin{theorem}[Complete contraction parametrization of mixed source Grams]
\label{thm:v15-contraction-param}
The following are equivalent:
\begin{enumerate}[label=\textup{(\roman*)},leftmargin=2.8em]
\item $\mathcal G(C)\succeq0$;
\item $I-T^*T\succeq0$;
\item $I-TT^*\succeq0$;
\item $\norm{T}_{\op}\le1$.
\end{enumerate}
Moreover every contraction $T:E_A\to E_B$ occurs as the normalized cross
Gram of an explicit common-carrier realization with marginals $G_A,G_B$.
One may take
\[
 \mathcal H=E_B\oplus E_A,
\]
\begin{equation}
 Bx=(G_B^{1/2}x,0),
 \qquad
 Ay=\left(TG_A^{1/2}y,
 (I-T^*T)^{1/2}G_A^{1/2}y\right).
 \label{eq:v15-universal-realization}
\end{equation}
Then
\[
 B^*B=G_B,\qquad A^*A=G_A,\qquad B^*A=C.
\]
Thus positivity imposes no hidden compatibility condition beyond contraction
of the whitened cross block.
\end{theorem}

\begin{proof}
Congruence by
$\operatorname{diag}(G_B^{-1/2},G_A^{-1/2})$ sends
$\mathcal G(C)$ to
\[
 \widehat{\mathcal G}(T)
 =\begin{pmatrix}I&T\\T^*&I\end{pmatrix}.
\]
The upper-left Schur complement is $I-T^*T$ and the lower-right Schur
complement is $I-TT^*$.  Since the diagonal blocks are positive definite,
either Schur condition is equivalent to positivity of the block matrix.
Those conditions are equivalent to all singular values of $T$ being at most
one, hence to $\norm T_{\op}\le1$.

For \eqref{eq:v15-universal-realization}, direct multiplication gives
\[
 B^*A=G_B^{1/2}TG_A^{1/2}=C
\]
and
\[
 A^*A
 =G_A^{1/2}\bigl(T^*T+I-T^*T\bigr)G_A^{1/2}
 =G_A.
\]
The formula for $B^*B$ is immediate.
\end{proof}

\begin{corollary}[Exact mutual source innovations]
\label{cor:v15-general-shorts}
For every admissible cross block,
\begin{align}
 G_A-C^*G_B^{-1}C
 &=G_A^{1/2}(I-T^*T)G_A^{1/2},
 \label{eq:v15-short-A}\\
 G_B-CG_A^{-1}C^*
 &=G_B^{1/2}(I-TT^*)G_B^{1/2}.
 \label{eq:v15-short-B}
\end{align}
If $\dim E_A=\dim E_B$, both residuals vanish exactly when $T$ is unitary.
In that case the two source ranges coincide.  Thus the V14 unitary condition
is the boundary case of the complete positive extension cone.
\end{corollary}

\begin{proof}
The two displayed identities are the two Schur complements after substituting
$C=G_B^{1/2}TG_A^{1/2}$.  In equal dimension,
$I-T^*T=I-TT^*=0$ is equivalent to unitarity.  The common-range statement is
\Cref{thm:v14-mixed-recognition}.
\end{proof}

\begin{remark}[Why separate marginal positivity cannot imply duality]
The theorem gives an explicit common Hilbert realization for every
contraction $T$.  Therefore the already known marginal Grams, however rigid
they may be individually, do not force the two source ranges to coincide.
The missing mixed block is genuine relative-orientation information.
\end{remark}

% -----------------------------------------------------------------------------
\section{With common tetrahedral covariance the entire ambiguity is one disk}
\label{sec:v15-disk}

We now impose the strongest source symmetries that V14 identified as the
natural minimal comparison, but only conditionally on their actual joint
occurrence.  Let
\[
 E=\R^2\otimes W_0,
 \qquad \dim_\R E=6,
\]
where $W_0$ is the real standard representation of $S_4$.  Put
\[
 J=J_2\otimes I_{W_0},
 \qquad
 J_2=\begin{pmatrix}0&-1\\1&0\end{pmatrix}.
\]
Assume the two normalized coefficient spaces have been independently aligned
so that the same linear $S_4$ action and the same typed complex structure $J$
are represented on both.

\begin{theorem}[The symmetric complex mixed-source cone is the closed unit disk]
\label{thm:v15-disk}
Let $T:E\to E$ be a normalized cross transport satisfying
\[
 T\rho(g)=\rho(g)T\quad(g\in S_4),
 \qquad TJ=JT.
\]
Then there are unique real numbers $a,b$ such that
\begin{equation}
 \boxed{
 T=T_z:=(aI_2+bJ_2)\otimes I_{W_0},
 \qquad z:=a+ib.}
 \label{eq:v15-Tz}
\end{equation}
The joint block Gram
\[
 \begin{pmatrix}I_6&T_z\\T_z^{\mathsf T}&I_6\end{pmatrix}
\]
is positive if and only if
\begin{equation}
 \boxed{|z|\le1.}
 \label{eq:v15-disk-condition}
\end{equation}
On this branch,
\begin{equation}
 \boxed{
 I-T_z^{\mathsf T}T_z
 =I-T_zT_z^{\mathsf T}
 =(1-|z|^2)I_6,}
 \label{eq:v15-disk-short}
\end{equation}
and the twelve-dimensional whitened joint Gram has spectrum
\begin{equation}
 \boxed{
 \spec\widehat{\mathcal G}(T_z)
 =\{(1+|z|)^{[6]},(1-|z|)^{[6]}\}.}
 \label{eq:v15-block-spectrum}
\end{equation}
Hence exact local source duality occurs precisely on the unit circle
$|z|=1$; every interior point has a full-rank six-dimensional innovation on
both sides.
\end{theorem}

\begin{proof}
By absolute irreducibility of $W_0$,
\[
 \End_{S_4}(E)=M_2(\R)\otimes I_{W_0}.
\]
Thus $T=M\otimes I_{W_0}$ for a real $2\times2$ matrix $M$.  The condition
$TJ=JT$ says that $M$ commutes with $J_2$, whose real commutant is
$\{aI_2+bJ_2:a,b\in\R\}$.  This proves \eqref{eq:v15-Tz} and uniqueness.

Since $J_2^{\mathsf T}=-J_2$ and $J_2^2=-I_2$,
\[
 (aI_2+bJ_2)^{\mathsf T}(aI_2+bJ_2)
 =(a^2+b^2)I_2.
\]
Therefore all six singular values of $T_z$ equal $|z|$.  The contraction
criterion of \Cref{thm:v15-contraction-param} gives
\eqref{eq:v15-disk-condition}, and the Schur residuals are
\eqref{eq:v15-disk-short}.  The standard singular-value diagonalization of
$\begin{psmallmatrix}I&T\\T^*&I\end{psmallmatrix}$ gives eigenvalues
$1\pm\sigma_j(T)$, yielding \eqref{eq:v15-block-spectrum}.
\end{proof}

\begin{corollary}[The coherence radius is an exact source-gap coordinate]
\label{cor:v15-coherence}
Write $\rho:=|z|$.  For normalized source isometries $A,B$ realizing the
symmetric branch,
\begin{equation}
 \boxed{
 \|P_A-P_B\|_{\HS}^2=12(1-\rho^2),}
 \label{eq:v15-projection-hs}
\end{equation}
and
\begin{equation}
 \boxed{
 \|P_A-P_B\|_{\op}=\sqrt{1-\rho^2}.}
 \label{eq:v15-projection-op}
\end{equation}
Thus $\rho$ measures the exact Grassmannian separation of the two source
ranges.  The phase $z/|z|$ is meaningful only after $\rho>0$.
\end{corollary}

\begin{proof}
For two rank-six orthogonal projections with isometric synthesis maps,
\[
 \|P_A-P_B\|_{\HS}^2
 =\Tr P_A+\Tr P_B-2\Tr(P_AP_B)
 =12-2\|B^*A\|_{\HS}^2.
\]
Here $B^*A=T_z$ and $\|T_z\|_{\HS}^2=6\rho^2$, proving
\eqref{eq:v15-projection-hs}.  All six principal angles have cosine $\rho$,
so the operator norm of the projection difference is the sine of that common
angle, giving \eqref{eq:v15-projection-op}.
\end{proof}

\begin{corollary}[Near-duality has a sharp one-scalar certificate]
\label{cor:v15-near-duality}
On the exact symmetric complex branch, if $\rho\ge1-\varepsilon$ with
$0\le\varepsilon\le1$, then
\[
 0\le1-\rho^2\le2\varepsilon,
\]
\[
 \|P_A-P_B\|_{\HS}^2\le24\varepsilon,
 \qquad
 \|P_A-P_B\|_{\op}\le\sqrt{2\varepsilon}.
\]
If $\rho>0$ and $u=z/\rho\in U(1)$, then
\[
 \boxed{
 \|T_z-T_u\|_{\HS}=\sqrt6(1-\rho).}
 \label{eq:v15-polar-distance}
\]
Thus the distance to the exact source-dual boundary is quantitatively
controlled by the coherence deficit $1-|z|$.
\end{corollary}

\begin{proof}
The scalar inequalities follow from
$1-\rho^2=(1-\rho)(1+\rho)$ and
\Cref{cor:v15-coherence}.  Since
$T_z-T_u=(\rho-1)T_u$ and $T_u$ is orthogonal on six real dimensions,
its Hilbert--Schmidt norm is $\sqrt6(1-\rho)$.
\end{proof}

% -----------------------------------------------------------------------------
\section{Separate source data do not identify the point of the disk}
\label{sec:v15-nonidentifiability}

\begin{theorem}[A continuum of symmetry-preserving joint completions has identical marginals]
\label{thm:v15-disk-nonidentifiability}
For every $z\in\overline{\mathbb D}$ there is a common-carrier realization of
two normalized six-real-dimensional sources with
\[
 A_z^*A_z=B^*B=I_6,
 \qquad B^*A_z=T_z,
\]
which carries the same diagonal $S_4$ action and the same orthogonal complex
structure on both source types.  Explicitly, on
$\mathcal H=E\oplus E$ set
\begin{equation}
 Bx=(x,0),
 \qquad
 A_zy=\left(T_zy,\sqrt{1-|z|^2}\,y\right).
 \label{eq:v15-z-realization}
\end{equation}
Let $S_4$ and $J$ act diagonally on the two copies of $E$.  Then all source
marginals, all separate symmetry data, and all separate identity-transfer
Krylov panels are independent of $z$.

Consequently no function of the two marginal source packets, their separate
$S_4$ types, their separate complex structures, or arbitrarily deep separate
identity-transfer moments can determine the mixed-source coherence $z$.
An actual mixed occurrence row or a represented common compiler is necessary.
\end{theorem}

\begin{proof}
The Gram identities follow immediately from
$T_z^{\mathsf T}T_z=|z|^2I_6$.  Because $T_z$ commutes with the standard
$S_4$ action and with $J$, both source maps intertwine the same diagonal
actions on $\mathcal H$.  Choosing the identity as the physical transfer on
$\mathcal H$ makes every separate moment
$A_z^*I^nA_z$ and $B^*I^nB$ equal to $I_6$, independently of $z$.  Distinct
values of $z$ nevertheless give distinct cross Grams.  Hence the listed
separate data cannot identify $z$.
\end{proof}

\begin{scope}[A binary orientation cannot replace mixed coherence]
Even if an additional independent binary mark removes a sign or conjugation
ambiguity in the phase of $z$, it cannot determine the continuous radius
$|z|\in[0,1]$.  In particular the protected Store/chronology sign may become
relevant only after a nonzero mixed coupling is acquired; it cannot force the
vanishing of the positive innovation $(1-|z|^2)I_6$.
\end{scope}

% -----------------------------------------------------------------------------
\section{The exact word-native boundary}
\label{sec:v15-word-native}

The preceding nonidentifiability does not mean that every mixed panel requires
a new experiment.  The supplied Grand-Tensor monograph proves that inserted
mixed Grams are already encoded in the ordinary word table when the proposed
source is a represented word-native operation.  We record the specialization
needed here without repeating the general ancestry theorem.

\begin{proposition}[If the arithmetic source is word-native, the V14 cross Gram is table-native]
\label{prop:v15-word-native-closure}
Let $\mathcal A_{\rm hist}$ be one certified flat represented history algebra
with word synthesis $W_r$.  Suppose the renewal source and the proposed
arithmetic source have the forms
\[
 B=W_rB_{\rm ren},
 \qquad
 A=\mathscr L W_rB_{\rm ar},
\]
where $\mathscr L$ is a fixed word-native operation in the sense imported from
\cite{GT77}.  Then
\begin{equation}
 \boxed{
 C=B^*A
 =B_{\rm ren}^*W_r^*\mathscr L W_rB_{\rm ar}}
 \label{eq:v15-word-native-C}
\end{equation}
is determined by a finite deeper ordinary word Gram and, after flatness, by
the stabilized word Gram and multiplication table.  No independently executed
mixed arithmetic--renewal panel is required on this branch.
\end{proposition}

\begin{proof}
Equation \eqref{eq:v15-word-native-C} is the mixed block in the inherited
word-native ancestry compiler of \cite{GT77}.  Expanding
$\mathscr L(X)=\sum_\alpha c_\alpha A_\alpha XB_\alpha$ expresses every entry
of $W_r^*\mathscr LW_r$ as an ordinary inserted word moment.  A sufficiently
deep ordinary word Gram evaluates those moments, and certified flatness
reduces them using the reconstructed multiplication table.
\end{proof}

\begin{corollary}[On the common symmetric branch the compiled panel reduces to one complex number]
\label{cor:v15-word-native-z}
Under the additional hypotheses of \Cref{thm:v15-disk}, let $T$ be the
whitened matrix obtained from the compiled cross Gram.  Then
\begin{equation}
 \boxed{
 z
 =\frac16\Tr T
 -\frac{i}{6}\Tr(JT).}
 \label{eq:v15-z-traces}
\end{equation}
Equivalently, in the complex standard presentation a single normalized
prototype overlap
\[
 \ip{\Gamma_{\rm ren}w}{\Gamma_{\rm ar}w}
\]
for any unit coefficient vector $w$ equals $z$.  Thus one complex mixed read,
or two real quadratures, is sufficient once the common $S_4$ and complex
intertwining laws are themselves independently established.
\end{corollary}

\begin{proof}
For $T=aI+bJ$ on six real dimensions,
$\Tr T=6a$ and
$\Tr(JT)=-6b$, giving \eqref{eq:v15-z-traces}.  In the complex presentation
$T$ acts as the scalar $a+ib=z$ on the irreducible complex standard module,
so every normalized coefficient vector gives the same overlap by Schur's
lemma.
\end{proof}

\begin{theorem}[No curve-blind chronology compiler can produce the selected theta source]
\label{thm:v15-curve-blind-no-go}
Consider the family $C_t$ of V13 and suppose a proposed natural compiler from
the universal chronology-derived arithmetic packet to the marked pure theta
source uses only upstream data that are unchanged throughout the family and
contains no additional curve/source-selection coordinate.  Then such a
compiler cannot reproduce the actual marked theta sources for all nearby
$t$.

More precisely, any output determined functorially by the same upstream
packet would have the same source-readable invariant $\mathfrak m_\Theta$,
whereas V13 proves
\[
 \boxed{\mathfrak m_\Theta(\rho_\Theta(t))=-t^3}
\]
with nonzero derivative $-3$ at $t=1$.  Hence every compiler producing the
actual theta source must depend on additional data which distinguish the
curve, locally equivalent to the source coordinate $\mathfrak m_\Theta$.
\end{theorem}

\begin{proof}
A deterministic natural compiler with identical declared inputs has identical
output up to the declared natural isomorphism, so every invariant of that
output is constant along the family as long as the input packet is unchanged.
But \Cref{cor:v13-source-modulus} constructs the source invariant
$\mathfrak m_\Theta$ and proves that its restriction to the actual family is
$-t^3$, with derivative $-3$ at the arithmetic point.  It is therefore not
locally constant.  The contradiction proves that an additional
curve/source-selection coordinate is necessary.
\end{proof}

\begin{scope}[What the no-go does and does not say]
The theorem does not exclude a word-native compiler after the curve/source
selection datum has itself been physically or arithmetically loaded.  It says
that the canonical chronology-derived arithmetic packet, together with only
the structural tetrahedral data shared by the V13 family, cannot by itself
supply the actual theta source.  Even after the source is selected, the
relative mixed coherence with the renewal plane remains an independent issue
by \Cref{thm:v15-disk-nonidentifiability} unless one common represented
compiler is supplied.
\end{scope}

% -----------------------------------------------------------------------------
\section{A robust symmetry projection for an acquired mixed panel}
\label{sec:v15-projection}

If a full mixed panel is eventually measured, one should not silently impose
its desired symmetry.  The symmetry-compatible part and its failure can be
separated orthogonally.

Let $\rho(g)$ denote the common six-dimensional coefficient action and assume
$[\rho(g),J]=0$.  For any real $6\times6$ matrix $X$ define
\begin{align}
 \Pi_{S_4}(X)
 &:=\frac1{24}\sum_{g\in S_4}\rho(g)^{\mathsf T}X\rho(g),
 \label{eq:v15-PiG}\\
 \Pi_J(X)
 &:=\frac12\left(X+J^{\mathsf T}XJ\right),
 \label{eq:v15-PiJ}\\
 \Pi_{\rm tc}(X)&:=\Pi_J\Pi_{S_4}(X).
 \label{eq:v15-Pitc}
\end{align}

\begin{theorem}[Orthogonal projection onto the tetrahedral complex intertwiner line]
\label{thm:v15-symmetry-projection}
The maps $\Pi_{S_4}$ and $\Pi_J$ are commuting Hilbert--Schmidt orthogonal
projections.  Their common image is
\[
 \boxed{\operatorname{span}_{\R}\{I,J\}.}
\]
Hence
\begin{equation}
 \boxed{
 \Pi_{\rm tc}(X)=a(X)I+b(X)J,}
 \label{eq:v15-projected-X}
\end{equation}
where
\begin{equation}
 \boxed{
 a(X)=\frac16\Tr X,
 \qquad
 b(X)=-\frac16\Tr(JX).}
 \label{eq:v15-ab-projection}
\end{equation}
Moreover
\begin{equation}
 \boxed{
 \|X\|_{\HS}^2
 =6\bigl(a(X)^2+b(X)^2\bigr)
 +\|X-\Pi_{\rm tc}(X)\|_{\HS}^2.}
 \label{eq:v15-pythagoras}
\end{equation}
If $X$ is a contraction then so is $\Pi_{\rm tc}(X)$, and therefore
$a(X)^2+b(X)^2\le1$.
\end{theorem}

\begin{proof}
Finite group averaging is the Hilbert--Schmidt orthogonal projection onto the
commutant of $\rho(S_4)$.  Since $J$ is orthogonal and $J^2=-I$, the map
$X\mapsto J^{\mathsf T}XJ$ is a self-adjoint involutive isometry; averaging
with the identity gives the orthogonal projection onto matrices intertwining
$J$.  The two operations commute because $J$ commutes with the group action.
The common commutant was computed in \Cref{thm:v15-disk} and is
$\operatorname{span}_{\R}\{I,J\}$.  Since
$\|I\|_{\HS}^2=\|J\|_{\HS}^2=6$ and
$\langle I,J\rangle_{\HS}=0$, orthogonal projection gives
\eqref{eq:v15-ab-projection} and the Pythagoras identity.

Both projections are averages of left/right orthogonal conjugations and hence
are contractions in operator norm.  Therefore a contraction remains a
contraction under $\Pi_{\rm tc}$; applying
\Cref{thm:v15-disk} to the projected matrix gives
$a^2+b^2\le1$.
\end{proof}

\begin{corollary}[A terminating measured-panel decomposition]
\label{cor:v15-measured-decomposition}
For an acquired whitened cross panel $T$, define
\[
 z_{\rm tc}:=a(T)+ib(T),
 \qquad
 \Delta_{\rm tc}:=\|T-\Pi_{\rm tc}(T)\|_{\HS}^2.
\]
Then
\[
 \boxed{
 \Delta_{\rm tc}=0,\quad |z_{\rm tc}|=1}
\]
if and only if the measured panel is an exact complex-linear $S_4$-equivariant
source unitary.  If $\Delta_{\rm tc}>0$, the panel itself violates the claimed
common tetrahedral/complex symmetry.  If $\Delta_{\rm tc}=0$ but
$|z_{\rm tc}|<1$, the symmetry is exact but the two sources retain the positive
innovation
\[
 (1-|z_{\rm tc}|^2)I_6.
\]
Thus symmetry failure and source incompleteness are operationally distinct
residuals.
\end{corollary}

\begin{proof}
If $\Delta_{\rm tc}=0$, the panel is exactly of the form $T_z$ in
\Cref{thm:v15-disk}.  That theorem makes unitarity equivalent to $|z|=1$ and
identifies the Schur residual for the interior.  The converse is immediate.
\end{proof}

% -----------------------------------------------------------------------------
\section{What V15 proves about the research programme}
\label{sec:v15-status}

\subsection{Proved in this continuation}
\begin{enumerate}[leftmargin=2.2em]
\item For arbitrary positive arithmetic and renewal source marginals, the
      complete set of admissible mixed Grams is exactly the operator
      contraction ball after whitening.
\item Every contraction in that ball has an explicit common-Hilbert-space
      realization; separate marginal positivity therefore contains no hidden
      source-duality theorem.
\item On a genuinely common $S_4$-equivariant complex two-standard branch,
      the entire mixed-source cone is one closed unit disk $z\in\overline{\mathbb D}$.
      Its radius is the exact source coherence and its unit-circle boundary is
      exact source duality.
\item The projection gap between the arithmetic and renewal source ranges is
      exactly $12(1-|z|^2)$ in Hilbert--Schmidt square and
      $\sqrt{1-|z|^2}$ in operator norm.
\item There is a continuum of joint realizations with identical marginal
      Grams, separate symmetries, complex structures and separate transfer
      moments but different mixed coherence.  Hence the mixed parameter is
      not inferable from separate data.
\item If the arithmetic source is represented by a fixed word-native compiler
      on the same flat history algebra, its cross Gram is table-native by the
      existing Grand-Tensor ancestry compiler; under exact common symmetry the
      result reduces to one complex number $z$.
\item No curve-blind compiler depending only on the universal chronology
      packet and structural tetrahedral data shared across the V13 family can
      produce the actual marked theta source.  At least one curve/source
      selection coordinate is necessary upstream.
\item Any acquired mixed panel admits an exact orthogonal decomposition into
      its tetrahedral-complex scalar component and a symmetry-breaking
      residual; source incompleteness and symmetry failure are therefore
      separately testable.
\end{enumerate}

\subsection{Inherited inputs not reproved}
V13 supplies the nonselection family and the source-readable modulus
$\mathfrak m_\Theta=-t^3$.  V14 supplies the six-real-dimensional arithmetic
source, the minimal two-standard renewal plane, and the first mixed-Gram test.
The word-native ancestry compiler and the external-shadow firewall are
imported from the supplied Grand-Tensor monograph \cite{GT77}; they are not
reproved here.

\subsection{Current source-completeness verdict}
The present attached material does not supply a represented word-native map
from the renewal history algebra to the selected theta/moduli source, nor an
actual same-history arithmetic--entry--axis mixed record.  This is consistent
with the explicit physical-source open problem in the Grand-Tensor monograph.
Accordingly, the current lineage is on the \emph{external mixed-source}
branch: $C$, or equivalently $z$ after a separately certified common symmetry,
remains an acquisition datum.

This is a status statement about the supplied construction, not a theorem that
such a compiler cannot be built.  A future represented compiler immediately
moves the problem to \Cref{prop:v15-word-native-closure}.

\subsection{The next noncircular attack}
The next iteration should proceed in the following order.
\begin{enumerate}[label=\textup{(N\arabic*)},leftmargin=3em]
\item Search the actual locked Store/half-edge history grammar for a same-event
      Read or word-native operator whose deterministic pushforward is the
      selected arithmetic moduli source.  If one exists, compile $C$ from the
      stabilized word table and evaluate \Cref{cor:v15-measured-decomposition}.
\item If no such represented compiler exists, do not infer $z$ from marginal
      arithmetic or renewal data.  Declare one complex mixed overlap (under a
      certified common $S_4$ and complex structure), or the full $6\times6$
      panel otherwise, as the minimal new occurrence measurement.
\item Only on the branch $\Delta_{\rm tc}=0$ and $|z|=1$ promote the local
      source comparison to transfer/Krylov and Frobenius--Wilson tests.  For an
      interior value $|z|<1$, retain the explicit innovation
      $(1-|z|^2)I_6$ as the obstruction rather than fitting it away.
\end{enumerate}

The source problem has therefore reached a sharp stopping boundary: no further
single-sector arithmetic invariant can determine the missing relative source
coherence.  Progress now requires either an actual common-history compiler or
an actual mixed occurrence row.

\paragraph{Version integrity.}
All V1--V14 mathematical bodies and their historical proof-status ledgers are
preserved byte for byte.  V15 adds only the complete mixed-extension
classification, the symmetry-reduced disk theorem, the word-native boundary,
and the robust measured-panel decomposition above.  Future continuations must
follow this block.
% ====================== END IMMUTABLE V15 MATHEMATICAL BODY ======================
% APPEND V16 AND LATER ITERATIONS HERE.

% ====================== BEGIN IMMUTABLE V16 MATHEMATICAL BODY ======================
\section{V16: scalar mixed-source certificates and determinant-line compression}
\label{sec:v16-context}

V15 reached a sharp information boundary.  Once the arithmetic and renewal
marginal source Grams are fixed, the whitened mixed block is an arbitrary
contraction.  The present continuation asks a narrower question than full
mixed-panel reconstruction:
\begin{quote}
\emph{What is the weakest genuinely mixed datum that certifies exact equality
of the two normalized source ranges?}
\end{quote}
This distinction matters because transfer comparison cannot begin before the
source ranges have been identified, whereas reconstructing every entry of the
cross Gram may contain considerably more information than that yes/no gate
requires.

Throughout this continuation let
\[
 A:\mathbb C^n\to\mathcal H,
 \qquad
 B:\mathbb C^n\to\mathcal H
\]
be full-rank source maps on one common Hilbert carrier.  Put
\[
 G_A=A^*A\succ0,
 \qquad
 G_B=B^*B\succ0,
\]
\[
 \widehat A=A G_A^{-1/2},
 \qquad
 \widehat B=B G_B^{-1/2},
\]
so that $\widehat A$ and $\widehat B$ are isometries, and define
\[
 P_A=\widehat A\widehat A^*,
 \qquad
 P_B=\widehat B\widehat B^*,
 \qquad
 T=\widehat B^*\widehat A.
\]
By V15, $T$ is a contraction whenever the complete joint Gram is positive.
The arithmetic--renewal application has $n=3$ over $\mathbb C$, equivalently
real source dimension six.  We work first at arbitrary $n$.

\section{A single quadratic coherence decides source-range equality}
\label{sec:v16-quadratic}

\begin{theorem}[Scalar projection-overlap certificate]
\label{thm:v16-projection-overlap}
Define the normalized mixed coherence
\[
 \boxed{
 \kappa_2(A,B)
 :=\frac1n\Tr_{\mathcal H}(P_AP_B)
 =\frac1n\Tr_{\mathbb C^n}(T^*T).}
 \tag{V16.1}\label{eq:v16-kappa}
\]
Then
\[
 0\le \kappa_2(A,B)\le1,
\]
and the following are equivalent:
\begin{enumerate}[label=\textup{(\roman*)},leftmargin=2.8em]
\item $\kappa_2(A,B)=1$;
\item $T$ is unitary;
\item $P_A=P_B$;
\item $\operatorname{Ran}A=\operatorname{Ran}B$.
\end{enumerate}
Moreover
\[
 \boxed{
 \|P_A-P_B\|_{\rm HS}^2
 =2n\bigl(1-\kappa_2(A,B)\bigr).}
 \tag{V16.2}\label{eq:v16-projgap}
\]
Thus one real mixed scalar is sufficient to decide exact equality of the
normalized source ranges.
\end{theorem}

\begin{proof}
Since $\widehat A$ and $\widehat B$ are isometries,
$P_A$ and $P_B$ are rank-$n$ orthogonal projections and
\[
 \Tr(P_AP_B)
 =\Tr(\widehat A^*\widehat B\widehat B^*\widehat A)
 =\Tr(T^*T).
\]
The contraction property gives $0\preceq T^*T\preceq I_n$, so
$0\le\kappa_2\le1$.  Equality $\kappa_2=1$ forces all eigenvalues of
$T^*T$ to equal one, hence $T$ is unitary.  If $T$ is unitary, then for every
$x$,
\[
 \|P_B\widehat A x\|^2
 =\|\widehat B^*\widehat A x\|^2
 =\|Tx\|^2
 =\|x\|^2
 =\|\widehat A x\|^2.
\]
Orthogonal projection preserves the norm only on its range, so
$\operatorname{Ran}\widehat A\subseteq\operatorname{Ran}\widehat B$.  Equal finite dimensions give equality
of the ranges and hence $P_A=P_B$.  The converse implications are immediate.
Finally,
\[
 \|P_A-P_B\|_{\rm HS}^2
 =\Tr(P_A)+\Tr(P_B)-2\Tr(P_AP_B)
 =2n-2\Tr(T^*T),
\]
which is \eqref{eq:v16-projgap}.
\end{proof}

\begin{corollary}[The V15 disk radius is a projection coherence]
\label{cor:v16-disk-kappa}
On the exact common $S_4$-equivariant complex branch of V15,
$T=T_z=zI_3$ on the complex standard source.  Therefore
\[
 \boxed{\kappa_2=|z|^2.}
 \tag{V16.3}
\]
In complex Hilbert--Schmidt convention,
\[
 \|P_A-P_B\|_{\rm HS}^2=6(1-|z|^2),
\]
and realification doubles this to the V15 value
$12(1-|z|^2)$.
\end{corollary}

\begin{proof}
Substitute $T_z^*T_z=|z|^2I_3$ into
\Cref{thm:v16-projection-overlap}.  A complex matrix coefficient contributes
two real coordinates under realification, giving the last factor two.
\end{proof}

\begin{remark}[What the scalar does and does not certify]
The equality $\kappa_2=1$ certifies equality of the two normalized source
ranges.  It does not by itself identify a predeclared coefficient basis, a
complex structure, an $S_4$ action, or a transfer operator.  If those
structures are independently represented on the same carrier and both source
maps intertwine them, the unitary $T=\widehat B^*\widehat A$ automatically
intertwines the corresponding coefficient actions.  Without those occurrence
rows, range equality must not be promoted to full typed source duality.
\end{remark}

\section{Top exterior power: one alternating mixed amplitude}
\label{sec:v16-exterior}

The quadratic scalar above is basis free but forgets all phase information.
There is a complementary one-scalar certificate on the determinant line.

Choose an orthonormal basis $e_1,\ldots,e_n$ of the coefficient space and put
\[
 a_\wedge
 =\widehat A e_1\wedge\cdots\wedge\widehat A e_n,
 \qquad
 b_\wedge
 =\widehat B e_1\wedge\cdots\wedge\widehat B e_n
 \in\Lambda^n\mathcal H.
\]
Both vectors have norm one.

\begin{theorem}[Exterior determinant certificate]
\label{thm:v16-exterior-det}
The normalized top-exterior mixed amplitude is
\[
 \boxed{
 \Delta_\wedge(A,B)
 :=\langle b_\wedge,a_\wedge\rangle
 =\det T.}
 \tag{V16.4}
\]
Consequently
\[
 |\Delta_\wedge|\le1,
\]
and the following are equivalent:
\begin{enumerate}[label=\textup{(\roman*)},leftmargin=2.8em]
\item $|\Delta_\wedge|=1$;
\item $T$ is unitary;
\item $\operatorname{Ran}A=\operatorname{Ran}B$.
\end{enumerate}
Thus one complex top-exterior mixed occurrence, or only its modulus for the
range question, is sufficient to certify exact source duality.
\end{theorem}

\begin{proof}
The Gram determinant identity gives
\[
 \langle b_\wedge,a_\wedge\rangle
 =\det\bigl[\langle\widehat B e_i,\widehat A e_j\rangle\bigr]_{i,j}
 =\det(\widehat B^*\widehat A)
 =\det T.
\]
If $s_1,\ldots,s_n\in[0,1]$ are the singular values of the contraction $T$,
then
\[
 |\Delta_\wedge|=\prod_{j=1}^n s_j\le1.
\]
Equality holds exactly when every $s_j=1$, equivalently $T^*T=I_n$.
Apply \Cref{thm:v16-projection-overlap}.
\end{proof}

\begin{proposition}[Quantitative determinant bounds]
\label{prop:v16-det-bounds}
Let
\[
 d:=|\Delta_\wedge|=|\det T|.
\]
Then
\[
 \boxed{
 1-d^{2/n}
 \le
 \Tr(I-T^*T)
 \le
 n\bigl(1-d^{2/n}\bigr),}
 \tag{V16.5}\label{eq:v16-dettrace}
\]
\[
 \boxed{
 2\bigl(1-d^{2/n}\bigr)
 \le
 \|P_A-P_B\|_{\rm HS}^2
 \le
 2n\bigl(1-d^{2/n}\bigr),}
 \tag{V16.6}\label{eq:v16-deths}
\]
and
\[
 \boxed{
 \sqrt{1-d^{2/n}}
 \le
 \|P_A-P_B\|_{\rm op}
 \le
 \sqrt{1-d^2}.}
 \tag{V16.7}\label{eq:v16-detop}
\]
In particular, a top-exterior amplitude close to unit modulus forces the
entire source plane to be close, not merely its oriented volume.
\end{proposition}

\begin{proof}
Set $x_j=s_j^2\in[0,1]$ and
$g=(\prod_jx_j)^{1/n}=d^{2/n}$.  By arithmetic--geometric mean,
\[
 \frac1n\sum_jx_j\ge g,
\]
which gives the upper bound in \eqref{eq:v16-dettrace}.  If
$m=\min_jx_j$, then $m\le g$ because
$\prod_jx_j\ge m^n$.  Therefore
\[
 \sum_j(1-x_j)\ge1-m\ge1-g,
\]
proving the lower bound.  Equation \eqref{eq:v16-deths} follows from
\Cref{thm:v16-projection-overlap}.

For equal-rank subspaces the largest principal-angle sine satisfies
\[
 \|P_A-P_B\|_{\rm op}^2=1-m.
\]
The inequality $m\le g$ gives the lower bound in \eqref{eq:v16-detop}.  Since all
other $x_j\le1$,
\[
 d^2=\prod_jx_j\le m,
\]
which gives the upper bound.
\end{proof}

\begin{corollary}[Exact formulas on the tetrahedral complex branch]
\label{cor:v16-det-disk}
On the V15 symmetry-reduced branch $T_z=zI_3$,
\[
 \boxed{
 \Delta_\wedge=z^3,
 \qquad
 |\Delta_\wedge|=|z|^3,
 \qquad
 \kappa_2=|\Delta_\wedge|^{2/3}.}
 \tag{V16.8}
\]
Hence the modulus of one determinant-line mixed amplitude determines the full
V15 radial coherence.  Its phase determines the phase of $z$ only modulo the
cube-root ambiguity $\mu_3$; one independently marked lower-degree phase
anchor is needed to choose a particular cube root.
\end{corollary}

\begin{proof}
The complex determinant of $zI_3$ is $z^3$.  The remaining identities follow
from \Cref{cor:v16-disk-kappa}.  Taking a cube root has exactly the stated
$\mu_3$ ambiguity.
\end{proof}

\section{Why the determinant amplitude is tetrahedrally natural}
\label{sec:v16-s4-det}

\begin{proposition}[The top source line is the sign representation]
\label{prop:v16-sign-line}
Let $W_0=\mathbf1_4^\perp$ be the real standard representation of $S_4$ and
$W=W_0\otimes_{\mathbb R}\mathbb C$.  Then
\[
 \boxed{\Lambda^3W\cong\operatorname{sgn}.}
 \tag{V16.9}
\]
Consequently, for arithmetic and renewal source fibres each of type $W$, the
Hermitian pairing of their determinant lines is invariant under a genuinely
common diagonal $S_4$ action.
\end{proposition}

\begin{proof}
The four-dimensional permutation representation is
$\mathbf1\oplus W_0$.  Its determinant on a permutation $g$ is
$\operatorname{sgn}(g)$, while the trivial line has determinant one.  Hence
$\det(g|_{W_0})=\operatorname{sgn}(g)$.  Complexification does not change the
character, so the top exterior power is the sign line.  Pairing two sign lines
cancels the two signs.
\end{proof}

\begin{corollary}[A basis-free tetrahedral range witness]
\label{cor:v16-basisfree-det}
Assume only that the two complex rank-three source fibres have independently
identified tetrahedral label actions of standard type and that a common
same-history determinant-line pairing occurs.  Then
$|\Delta_\wedge|$ is independent of the choices of orthonormal coefficient
bases and is a tetrahedrally invariant scalar.  The condition
\[
 \boxed{|\Delta_\wedge|=1}
\]
certifies equality of the underlying normalized source ranges without first
reconstructing the full $3\times3$ complex cross Gram.
\end{corollary}

\begin{proof}
Changing the arithmetic and renewal orthonormal coefficient bases multiplies
$\Delta_\wedge$ by unit-modulus determinant factors, so its modulus is
unchanged.  Diagonal $S_4$ invariance is
\Cref{prop:v16-sign-line}; the range criterion is
\Cref{thm:v16-exterior-det}.
\end{proof}

\section{A weaker word-native entrance than the full mixed panel}
\label{sec:v16-projector-word}

V15 required a word-native realization of the full arithmetic source map in
order to compile every cross-Gram entry.  The scalar certificates above show
that the source-range gate has a weaker sufficient entrance.

\begin{proposition}[Represented source projectors suffice for the quadratic gate]
\label{prop:v16-projector-native}
Suppose the two orthogonal source projections $P_A,P_B$ occur as represented
operators on one finite historical carrier, and the declared trace used on
that carrier agrees with the Hilbert trace on their common finite support up
to a known positive scalar.  If the historical word hierarchy is flat and
$P_A,P_B$ are represented elements, or fixed word-native functional-calculus
outputs, then the single scalar
\[
 \Tr(P_AP_B)
\]
is a deterministic function of the stabilized word Gram and multiplication
table.  Therefore exact source-range equality is table-native through
\Cref{thm:v16-projection-overlap}, even when the oriented source maps themselves
are not reconstructed.
\end{proposition}

\begin{proof}
At word flatness, multiplication of represented historical elements is known.
Hence the product $P_AP_B$ is reconstructed, and its declared trace is part of
the finite represented algebraic packet.  The known trace conversion gives
$\Tr_{\mathcal H}(P_AP_B)$, and
\Cref{thm:v16-projection-overlap} decides equality of the ranges.  No
orientation of the coefficient frames is needed for this scalar question.
\end{proof}

\begin{proposition}[Represented determinant lines suffice for the exterior gate]
\label{prop:v16-det-native}
Suppose instead that the normalized top-exterior source vectors
$a_\wedge,b_\wedge$ occur on one represented same-history exterior carrier,
or are fixed word-native outputs there.  Then their one mixed matrix
coefficient
\[
 \langle b_\wedge,a_\wedge\rangle
\]
is sufficient to decide exact source-range equality by
\Cref{thm:v16-exterior-det}.  On a common tetrahedral branch its modulus is an
$S_4$-invariant scalar by \Cref{prop:v16-sign-line}.
\end{proposition}

\begin{proof}
The coefficient is exactly $\Delta_\wedge$.  Apply
\Cref{thm:v16-exterior-det} and \Cref{prop:v16-sign-line}.
\end{proof}

\begin{remark}[These are still mixed occurrences]
\label{fire:v16-mixed-still-needed}
Neither \Cref{prop:v16-projector-native} nor
\Cref{prop:v16-det-native} manufactures a cross-sector relation from separate
marginals.  The projector product or determinant-line pairing must actually be
represented on one carrier, or compiled from a previously represented
word-native relation.  V15's continuum of joint completions keeps all separate
marginal data fixed while varying both scalar witnesses.  Therefore absence of
the required common-history operation remains an exact source-completeness
obstruction.
\end{remark}

\section{Nonidentifiability of the scalar witnesses from inherited marginals}
\label{sec:v16-nonident}

\begin{theorem}[The V15 disk varies both one-scalar certificates]
\label{thm:v16-scalar-nonident}
The V15 family of common-Hilbert-space realizations with
\[
 T_z=zI_3,
 \qquad z\in\overline{\mathbb D},
\]
has identical arithmetic and renewal marginal Grams, identical separate
complex structures, identical separate $S_4$ representation types and
identical separate transfer moments, while
\[
 \boxed{
 \kappa_2(z)=|z|^2,
 \qquad
 \Delta_\wedge(z)=z^3.}
 \tag{V16.10}
\]
Hence neither scalar certificate is a function of the inherited separate
sector data.
\end{theorem}

\begin{proof}
The invariance of the separate marginals and separate structures is the
explicit realization theorem of V15.  The two formulas are
\Cref{cor:v16-disk-kappa,cor:v16-det-disk}.  Since they vary with $z$, no
function of data held fixed across that family can determine either one.
\end{proof}

\begin{corollary}[One real mixed scalar is the weakest sufficient range gate in the current lineage]
\label{cor:v16-one-scalar-gate}
For the yes/no question whether the two normalized rank-three complex source
ranges coincide, it is sufficient to acquire the single real scalar
$\kappa_2$, or equivalently to acquire only the modulus of the single complex
scalar $\Delta_\wedge$.  No theorem in the inherited material makes either
quantity table-native for the selected external arithmetic theta/moduli
source.
\end{corollary}

\begin{proof}
Sufficiency is \Cref{thm:v16-projection-overlap,thm:v16-exterior-det}.
The final statement is the V15 word-native/external-shadow audit together with
\Cref{thm:v16-scalar-nonident}: the supplied construction contains no
represented arithmetic source projector or determinant-line pairing on the
renewal history carrier.
\end{proof}

\section{Determinant firewalls: three different uses of the word ``determinant''}
\label{sec:v16-det-firewall}

The present scalar must not be conflated with two determinant constructions
already present elsewhere in the programme.

\begin{remark}[Local-source determinant versus V10 and the SM determinant incidence]
\label{fire:v16-three-dets}
There are three logically distinct objects.
\begin{enumerate}[label=\textup{(D\arabic*)},leftmargin=3em]
\item The present
\[
 \Delta_\wedge
 \in
 \operatorname{Hom}\bigl(\Lambda^3\mathcal R_{2,\mathbb C},
           \Lambda^3\mathcal T_{\rm ar}\bigr)
\]
is a mixed overlap of the two local complex rank-three source planes.
\item V10's arithmetic determinant tensor belongs to the determinant of the
rank-two/rank-three Hodge--theta bundles and is used to reduce
$U(2)\times U(3)$ to $S(U(2)\times U(3))$ after a stated metric
compatibility normalization.
\item The companion Standard-Model construction has a represented fully
alternating incidence shadow on its internal typed carrier.
\end{enumerate}
No theorem presently identifies these three determinant lines or their
physical occurrence rows.  In particular, an existing nonzero determinant
incidence in the Standard-Model packet does not supply
$|\Delta_\wedge|=1$ for the arithmetic--renewal source comparison.
\end{remark}

\section{The sharpened next gate}
\label{sec:v16-status}

\subsection{Proved in this continuation}
\begin{enumerate}[leftmargin=2.2em]
\item Equality of two normalized rank-$n$ source ranges is decided by one real
      mixed coherence scalar $\kappa_2=n^{-1}\Tr(P_AP_B)$.
\item The same equality is decided by the modulus of one complex top-exterior
      mixed amplitude $\Delta_\wedge=\det T$.
\item The determinant amplitude quantitatively controls the complete
      principal-angle gap through \eqref{eq:v16-dettrace}--\eqref{eq:v16-detop}; it is not only
      a zero/nonzero witness.
\item On the exact tetrahedral complex branch,
      $\kappa_2=|z|^2$ and $\Delta_\wedge=z^3$.  Hence one determinant-line
      mixed amplitude determines the V15 coherence radius and its phase modulo
      $\mu_3$.
\item The determinant lines of both complex standard $S_4$ sources carry the
      sign character, so the modulus of their mixed top-exterior pairing is a
      tetrahedrally invariant scalar.
\item A represented source projector is sufficient for a table-native
      quadratic source-range test; a full oriented source-map compiler is not
      required for that narrower gate.
\item The one-scalar witnesses are nevertheless genuinely mixed: the V15 disk
      varies them while preserving every inherited separate-sector datum.
\end{enumerate}

\subsection{Inherited inputs not reproved}
V14 supplies the minimal two-standard renewal plane and the complex rank-three
arithmetic tangent.  V15 supplies the complete contraction cone, the disk
reduction under common $S_4$ and complex covariance, and the external-shadow
boundary.  The inserted-word and determinant-shadow compilers in the
Grand-Tensor monograph are used only conditionally: their hypotheses require
represented word-native operations on one common historical carrier.

\subsection{Current acquisition verdict}
The full $6\times6$ mixed panel is no longer necessary merely to decide
whether the two source ranges coincide.  The weakest explicit sufficient
measurements isolated here are:
\[
 \boxed{
 \kappa_2=\frac13\Tr(P_AP_B)
 \quad\text{or}\quad
 |\Delta_\wedge|.}
\]
For the selected arithmetic theta/moduli source, neither scalar is presently
contained in the inherited renewal word table.  It must arise from an actual
same-history projector overlap, an actual determinant-line pairing, or a new
represented compiler satisfying the V15 word-native criterion.

\subsection{The next noncircular attack}
The next iteration should therefore proceed in this order.
\begin{enumerate}[label=\textup{(N\arabic*)},leftmargin=3em]
\item Search the locked history grammar for a represented arithmetic source
      \emph{projector}, which is weaker than searching for the full oriented
      arithmetic source map.  If it exists, compute $\kappa_2$ directly from
      one product trace.
\item In parallel, search for a same-history alternating three-source Read
      pairing the arithmetic and renewal determinant lines.  Its modulus alone
      decides source-range equality.
\item If neither occurrence exists, stop the source-duality proof at this
      exact one-scalar acquisition boundary.  Do not reconstruct either scalar
      from separate marginals.
\item Only if $\kappa_2=1$ (equivalently $|\Delta_\wedge|=1$) proceed to the
      symmetry/complex-structure residual and then to transfer/Krylov
      comparison.  A strict value below one is the quantitative source
      obstruction and should be retained rather than fitted away.
\end{enumerate}

The source comparison is therefore sharper than at V15: a complete cross Gram
is needed to reconstruct the relative source map, but one mixed scalar is
sufficient to certify or obstruct equality of the two source ranges.

\paragraph{Version integrity.}
All V1--V15 mathematical bodies and their historical status ledgers are
preserved byte for byte.  V16 adds only the scalar coherence theorem, the
exterior determinant criterion, quantitative principal-angle bounds, the
weaker projector-word-native gate, and the determinant firewalls above.
Future continuations must follow this block.
% ====================== END IMMUTABLE V16 MATHEMATICAL BODY ======================
% APPEND V17 AND LATER ITERATIONS HERE.



% ====================== BEGIN IMMUTABLE V17 MATHEMATICAL BODY ======================
\section{V17: source-birth rank, projector-word reconstruction, and the upstream information cost}
\label{sec:v17-context}

V16 reduced the first arithmetic--renewal range test to one mixed scalar,
provided the corresponding source projectors or determinant lines actually
occur on one common carrier.  The present continuation asks what that
requirement costs and what can be reconstructed once such projectors are
represented.  Three logically separate questions are treated:
\begin{enumerate}[label=\textup{(Q\arabic*)},leftmargin=3em]
\item Given a normalized mixed contraction $T$, what is the minimum common
      Hilbert-space dimension required to realize both sources?
\item If the two source projectors are represented, how much of the complete
      principal-angle geometry is already table native?
\item Are the needed projectors already observables of the bare arithmetic and
      renewal point records, or do they themselves require an additional
      source-processing operation?
\end{enumerate}
The answers sharpen the stopping boundary of V16.  The mixed scalar is small as
an \emph{output}, but its operational availability is not free.

Throughout the first two sections let
\[
 \widehat A,\widehat B:\mathbb C^n\longrightarrow\mathcal H
\]
be isometries and put
\[
 T=\widehat B^*\widehat A,
 \qquad
 P_A=\widehat A\widehat A^*,
 \qquad
 P_B=\widehat B\widehat B^*.
\]
Thus $T$ is a contraction by positivity of the complete mixed Gram.

\section{The minimal common carrier and exact source-birth rank}
\label{sec:v17-minimal-carrier}

\begin{theorem}[Minimal common dilation dimension]
\label{thm:v17-minimal-dilation}
Let $T\in M_n(\mathbb C)$ be a contraction.  Among all pairs of isometries
$\widehat A,\widehat B:\mathbb C^n\to\mathcal H$ satisfying
$\widehat B^*\widehat A=T$, the minimum possible dimension of the joint source
span is
\[
 \boxed{
 \dim\Span(\operatorname{Ran}\widehat A,\operatorname{Ran}\widehat B)
 =n+\rank(I-T^*T).}
 \tag{V17.1}\label{eq:v17-min-dim}
\]
Equivalently,
\[
 \boxed{
 \rank\bigl(\widehat A^*(I-P_B)\widehat A\bigr)
 =\rank(I-T^*T)}
 \tag{V17.2}\label{eq:v17-birth-rank}
\]
is the exact minimum number of additional complex source directions which must
be present beyond the renewal source range.

Every source-minimal realization of a fixed $T$ is unique up to a unitary which
fixes the two labelled source maps.
\end{theorem}

\begin{proof}
Decompose the arithmetic source into its renewal-range and orthogonal parts:
\[
 \widehat A=P_B\widehat A+(I-P_B)\widehat A
 =\widehat BT+A_\perp.
\]
Since $\widehat B$ is an isometry,
\[
 A_\perp^*A_\perp
 =\widehat A^*(I-P_B)\widehat A
 =I-T^*T.
\]
Therefore
\[
 \rank A_\perp=\rank(I-T^*T).
\]
The range of $A_\perp$ is orthogonal to $\operatorname{Ran}\widehat B$, so
\[
 \Span(\operatorname{Ran}\widehat A,\operatorname{Ran}\widehat B)
 =\operatorname{Ran}\widehat B\oplus\operatorname{Ran} A_\perp,
\]
which proves \eqref{eq:v17-min-dim} and \eqref{eq:v17-birth-rank}.

Conversely let $D=(I-T^*T)^{1/2}$ and work on
$\mathbb C^n\oplus\overline{\operatorname{Ran} D}$.  Define
\[
 \widehat B_0x=(x,0),
 \qquad
 \widehat A_0x=(Tx,Dx).
\]
Then
$\widehat B_0^*\widehat A_0=T$ and
$\widehat A_0^*\widehat A_0=T^*T+D^2=I$.
Its joint span has dimension
$n+\rank D=n+\rank(I-T^*T)$, proving attainability.

Finally, suppose two source-minimal pairs have the same $T$.  Their complete
joint source Grams are both
\[
 \begin{pmatrix}I&T\\T^*&I\end{pmatrix}.
\]
The map sending the labelled columns of one pair to the corresponding labelled
columns of the other preserves every inner product.  Quotienting the common
Gram kernel makes this map well defined and isometric; source minimality makes
it onto.  Hence it extends uniquely on the joint source span to the claimed
source-fixing unitary.
\end{proof}

\begin{corollary}[All-or-nothing birth on the tetrahedral complex disk]
\label{cor:v17-disk-birth}
On the V15--V16 common $S_4$-equivariant complex branch,
$T=T_z=zI_3$.  Therefore
\[
 \boxed{
 \dim\Span(\operatorname{Ran}\widehat A,\operatorname{Ran}\widehat B)
 =\begin{cases}
 3,&|z|=1,\\
 6,&|z|<1,
 \end{cases}}
 \tag{V17.3}\label{eq:v17-disk-dim}
\]
and
\[
 \boxed{
 \rank(I-T_z^*T_z)
 =\begin{cases}
 0,&|z|=1,\\
 3,&|z|<1.
 \end{cases}}
 \tag{V17.4}\label{eq:v17-disk-rank}
\]
Thus a strict interior value of the V15 disk is not a one-direction defect: it
requires an entire additional complex standard triplet.  There are no
source-minimal dimensions four or five on the exact symmetry branch.
\end{corollary}

\begin{proof}
Here
$I-T_z^*T_z=(1-|z|^2)I_3$.
Apply \Cref{thm:v17-minimal-dilation}.
\end{proof}

\begin{remark}[Birth rank is not particle multiplicity]
The rank in \eqref{eq:v17-birth-rank} is the minimum orthogonal innovation
required by one pair of source maps.  It is not a generation count, a new
spacetime dimension, or an additional matter multiplet.  A nonzero birth is
charged only after the corresponding mixed source occurrence is physically
loaded.
\end{remark}

\section{One-copy projector words reconstruct the complete principal-angle packet}
\label{sec:v17-projector-moments}

V16 showed that one trace $\Tr(P_AP_B)$ decides equality of the source ranges.
When both projectors are represented, three alternating projector moments
recover substantially more: the entire rank-three principal-angle geometry.

\begin{theorem}[Projector-word reconstruction of principal angles]
\label{thm:v17-projector-moments}
Let
\[
 X_A:=P_AP_BP_A\big|_{\operatorname{Ran} P_A}.
\]
Under the isometry $\widehat A:\mathbb C^n\to\operatorname{Ran} P_A$,
\[
 \boxed{X_A\simeq T^*T.}
 \tag{V17.5}\label{eq:v17-compression}
\]
Hence for every $k\ge1$,
\[
 \boxed{
 s_k:=\Tr_{\mathcal H}\bigl((P_AP_B)^k\bigr)
 =\Tr_{\mathbb C^n}\bigl((T^*T)^k\bigr).}
 \tag{V17.6}\label{eq:v17-sk}
\]
The first $n$ values $s_1,\ldots,s_n$ determine the complete multiset of
squared singular values of $T$, equivalently all principal angles between the
two source ranges.

For $n=3$, put
\[
 e_1=s_1,
 \qquad
 e_2=\frac{s_1^2-s_2}{2},
 \qquad
 e_3=\frac{s_1^3-3s_1s_2+2s_3}{6}.
\]
Then the squared singular values are exactly the roots of
\[
 \boxed{
 q(\lambda)
 =\lambda^3-e_1\lambda^2+e_2\lambda-e_3.}
 \tag{V17.7}\label{eq:v17-newton-poly}
\]
Moreover
\[
 \boxed{e_3=|\det T|^2=|\Delta_\wedge|^2.}
 \tag{V17.8}\label{eq:v17-e3}
\]
The source-birth rank is the number of roots of \eqref{eq:v17-newton-poly}
strictly below one, counted with multiplicity.
\end{theorem}

\begin{proof}
For $x\in\mathbb C^n$,
\[
 P_AP_BP_A\widehat A x
 =\widehat A\widehat A^*\widehat B\widehat B^*\widehat A x
 =\widehat A T^*Tx,
\]
which proves \eqref{eq:v17-compression}.  The nonzero cyclic trace of an
alternating product of two equal-rank projections may be taken on
$\operatorname{Ran} P_A$, giving \eqref{eq:v17-sk}.  The spectral theorem now identifies
$s_k$ as the power sums of the eigenvalues of $T^*T$.  Newton's identities
recover the elementary symmetric functions and hence the characteristic
polynomial.  For $n=3$ they give precisely
\eqref{eq:v17-newton-poly}.  Finally
$e_3=\det(T^*T)=|\det T|^2$, and
\Cref{thm:v17-minimal-dilation} identifies the multiplicity of eigenvalues
strictly below one with the source-birth rank.
\end{proof}

\begin{corollary}[Table-native source geometry once both projectors are represented]
\label{cor:v17-table-native-projectors}
Suppose $P_A$ and $P_B$ are represented historical operations on one flat
finite history algebra and the declared trace is known.  Then, for rank-three
complex sources, the three ordinary word traces
\[
 \boxed{
 \Tr(P_AP_B),
 \quad
 \Tr(P_AP_BP_AP_B),
 \quad
 \Tr(P_AP_BP_AP_BP_AP_B)}
 \tag{V17.9}\label{eq:v17-three-words}
\]
reconstruct all principal angles, $|\Delta_\wedge|$, and the exact source-birth
rank.  The first trace alone decides equality of the ranges by V16.
\end{corollary}

\begin{proof}
At word flatness the represented multiplication table reduces the three
products in \eqref{eq:v17-three-words}.  Apply
\Cref{thm:v17-projector-moments} and the V16 projection-overlap theorem.
\end{proof}

\begin{remark}[Phase still requires an oriented mixed row]
Projector words determine the singular values of $T$ and hence
$|\det T|$, but they do not determine the phase of $\det T$ or an oriented
coefficient-frame intertwiner.  That information remains in an exterior-line
or lower-degree phase occurrence.  The present theorem therefore strengthens
the range and birth audit without collapsing the V16 determinant firewall.
\end{remark}

\section{The bare matching-axis record does not represent its source projector}
\label{sec:v17-halfedge-no-projector}

V16 suggested searching for a represented renewal source projector before
asking for the full cross Gram.  The inherited matching-axis construction
supplies the source plane as an explicit subspace of the twelve-point half-edge
record.  The following theorem shows that the \emph{bare classical point
algebra} does not itself contain the corresponding orthogonal projector.

Let
\[
 \Omega=\{1,2,3,4\},
 \qquad
 \mathscr M=\{\text{three perfect matchings of }K_4\},
\]
and equip
\[
 \mathcal H_{\rm ha}=L^2(\Omega\times\mathscr M,\text{uniform})
\]
with the inherited matching-axis source maps.  Let $E_0,A_0:W_0\to\mathcal
H_{\rm ha}$ be their normalized orthogonal standard copies and put
\[
 \boxed{
 \mathcal R_2=\operatorname{Ran} E_0\oplus\operatorname{Ran} A_0,
 \qquad
 P_{\rm ren}=E_0E_0^*+A_0A_0^*.}
 \tag{V17.10}\label{eq:v17-pren}
\]
Then $\rank P_{\rm ren}=6$.

\begin{theorem}[Static half-edge multiplication cannot supply the renewal source projector]
\label{thm:v17-halfedge-no-projector}
Let
\[
 \mathcal D_{\rm ha}=L^\infty(\Omega\times\mathscr M)
\]
act on $\mathcal H_{\rm ha}$ by pointwise multiplication.  Then
\[
 \boxed{P_{\rm ren}\notin\mathcal D_{\rm ha}.}
 \tag{V17.11}\label{eq:v17-pren-notdiag}
\]
Moreover $P_{\rm ren}$ is not the conditional-expectation projection onto any
classical sub-$\sigma$-algebra of the half-edge record.

Consequently the one-time classical half-edge table alone does not instantiate
the V16 projector-word gate.  A source-processing operation, a richer history
algebra, or a direct source-overlap Read is required.
\end{theorem}

\begin{proof}
Every orthogonal projection in the finite commutative multiplication algebra
$\mathcal D_{\rm ha}$ is multiplication by an indicator $\one_S$ of a subset
$S\subseteq\Omega\times\mathscr M$.  Its range consists of functions supported
inside $S$ and its rank is $|S|$.

Assume $P_{\rm ren}=M_{\one_S}$.  For every half-edge point $(i,m)$ there is a
vector $x\in W_0$ with $x_i\ne0$.  The entry-source function
$(E_0x)(i,m)$ is a nonzero scalar multiple of $x_i$, so some vector in
$\mathcal R_2$ is nonzero at $(i,m)$.  Hence no point can lie outside $S$.
Thus $S=\Omega\times\mathscr M$, forcing
$\rank P_{\rm ren}=12$, contradicting \eqref{eq:v17-pren}.  This proves
\eqref{eq:v17-pren-notdiag}.

A conditional expectation onto a classical sub-$\sigma$-algebra is a unital
orthogonal projection; its range contains the constant function.  Both
$E_0(W_0)$ and $A_0(W_0)$ are centered standard subrepresentations and are
orthogonal to the constant line.  Therefore
$1\notin\mathcal R_2$, so $P_{\rm ren}$ cannot be such a conditional
expectation either.
\end{proof}

\begin{remark}[Mathematical reconstructibility versus represented occurrence]
\label{fire:v17-projector-occurrence}
The formula \eqref{eq:v17-pren} reconstructs $P_{\rm ren}$ mathematically once
the two source syntheses are known.  \Cref{thm:v17-halfedge-no-projector} says
something different: the projector is not already a pointwise observable or a
classical coarse-graining operation of the bare half-edge record.  A richer
Grand-Tensor history algebra may still represent it, but that is an occurrence
or compiler statement which must be checked rather than inferred from the
static source Gram.
\end{remark}

\section{A six-real-dimensional information lower bound for a universal theta-source compiler}
\label{sec:v17-moduli-cost}

V12 proved that, at the selected arithmetic fibre, the marked projective
second-order theta source is a local holomorphic coordinate system on the
three-complex-dimensional marked genus-two moduli germ.  V13 then showed that a
one-scalar family $C_t$ is only a selected one-dimensional arithmetic slice.
The next theorem turns these facts into a source-information lower bound.

\begin{theorem}[Local moduli cost of a universal arithmetic theta compiler]
\label{thm:v17-moduli-cost}
Let $\mathcal M$ be a smooth marked genus-two moduli chart through the selected
arithmetic point, so
\[
 \dim_{\mathbb R}\mathcal M=6.
\]
Let
\[
 \Theta:\mathcal M\longrightarrow\mathscr S_\Theta
\]
be the marked projective theta-source map.  Assume the V12 nondegeneracy
certificate, so $d\Theta$ has real rank six at the arithmetic point.

Let $X$ be a finite-dimensional smooth upstream writer manifold and
$F:X\to\mathscr S_\Theta$ a $C^1$ compiler whose image contains a neighborhood
of the arithmetic source in $\Theta(\mathcal M)$.  Then
\[
 \boxed{\dim_{\mathbb R}X\ge6.}
 \tag{V17.12}\label{eq:v17-six-cost}
\]
In particular, discrete tetrahedral labels together with fewer than six
independent real continuous writers cannot provide a locally universal
compiler for the marked arithmetic theta source.

A one-real-parameter writer such as the $t$ of V13 can be sufficient only
after an independent theorem restricts the arithmetic source to a
one-dimensional locus such as the chosen family $C_t$.
\end{theorem}

\begin{proof}
Because $d\Theta$ has real rank six, the inverse-function theorem identifies a
neighborhood of the arithmetic source inside $\Theta(\mathcal M)$ with an open
subset of $\mathbb R^6$.  If $\dim_{\mathbb R}X=k<6$, every point of the
$C^1$ map $F:X\to\mathbb R^6$ is critical.  Sard's theorem therefore implies
that $F(X)$ has six-dimensional Lebesgue measure zero locally and hence cannot
contain an open neighborhood.  This contradicts the assumed local
universality.  Thus $k\ge6$.

Discrete labels do not change local manifold dimension.  The final statement
follows because restricting the target to a preselected one-dimensional
submanifold changes the required local dimension; it does not make the
one-parameter family universal in the ambient marked moduli space.
\end{proof}

\begin{corollary}[The curve-selection scalar is not a universal source compiler]
\label{cor:v17-one-scalar-not-universal}
The V13 invariant
\[
 \mathfrak m_\Theta=-t^3
\]
selects the member of the chosen $C_t$ family, but it cannot by itself compile
the complete marked theta source on an open genus-two moduli neighborhood.
Thus the curve-selection problem and the local source-occurrence problem are
strictly different.
\end{corollary}

\begin{proof}
The invariant carries one real parameter on the real $C_t$ slice.  Apply
\Cref{thm:v17-moduli-cost}.
\end{proof}

\section{The revised acquisition hierarchy}
\label{sec:v17-status}

\subsection{Proved in this continuation}
\begin{enumerate}[leftmargin=2.2em]
\item For any fixed mixed contraction $T$, the source-minimal common carrier
      has dimension $n+\rank(I-T^*T)$, and this rank is exactly the orthogonal
      source-birth cost.
\item On the exact tetrahedral complex disk, source birth is all or nothing:
      the joint complex dimension is three on the unit circle and six at every
      strict interior point.
\item Once both source projectors are represented, the first three alternating
      projector-word traces reconstruct the complete rank-three principal-angle
      packet, determinant modulus, and source-birth rank.
\item The bare commutative twelve-point matching-axis record does not contain
      the renewal rank-six source projector as a multiplication observable or a
      classical conditional expectation.
\item A locally universal compiler for the marked arithmetic theta source must
      carry at least six real continuous source-selection coordinates unless an
      independent theorem first restricts the source to a lower-dimensional
      locus.
\end{enumerate}

\subsection{Inherited inputs not reproved}
The matching-axis half-edge source and its two orthogonal standard triplets are
inherited from the Grand-Tensor monograph.  V12 supplies the nonzero
three-complex-dimensional theta Jacobian.  V15 supplies the complete mixed
contraction cone, and V16 supplies the one-scalar projector and determinant
range certificates.  No result above identifies the arithmetic and renewal
projectors by assumption.

\subsection{Current acquisition verdict}
The source-range question now has a three-level audit.
\begin{enumerate}[label=\textup{(A\arabic*)},leftmargin=3em]
\item If both projectors occur as represented historical operations, the
      single trace $\Tr(P_AP_B)$ decides exact range equality, and three
      projector words recover all principal angles.
\item The static half-edge record does not by itself provide even the renewal
      projector.  One must therefore locate a richer source-localizer or
      history operation on the renewal side before claiming the projector-word
      route is physically populated.
\item On the arithmetic side, a universal theta-projector compiler cannot be
      generated from only the curve-blind chronology packet plus a few scalar
      labels: local source universality costs six real continuous coordinates,
      or else a genuine source-selection theorem must first reduce the target
      locus.
\end{enumerate}

\subsection{The next noncircular attack}
The next iteration should not compute another marginal theta invariant.  It
should proceed in the following order.
\begin{enumerate}[label=\textup{(N\arabic*)},leftmargin=3em]
\item Search the relationally completed Grand-Tensor history algebra for the
      actual noncommutative source-localizer whose range is the matching-axis
      two-standard plane.  The bare classical point record is now excluded by
      \Cref{thm:v17-halfedge-no-projector}.
\item On the arithmetic side, identify which already loaded marked moduli
      writers, if any, furnish the six-real-dimensional local theta-source
      compiler.  If only the selected $C_t$ family is physically loaded, state
      the comparison as family-relative rather than universal.
\item If both projectors become represented, evaluate the V16 scalar
      $\kappa_2$ first.  Only on the branch $\kappa_2=1$ spend the additional
      acquisition required for the complex phase, common $S_4$ action, and
      transfer/Krylov comparison.
\item If no represented source-localizers exist, retain their absence as the
      upstream source-occurrence obstruction.  Do not replace it by a
      mathematically reconstructed projector which the physical history algebra
      does not act by.
\end{enumerate}

The key change from V16 is therefore operational.  One scalar remains enough
\emph{after} the source projectors are represented, but producing those
projectors is itself a nontrivial source-completeness problem on both sides.

\paragraph{Version integrity.}
All V1--V16 mathematical bodies and their historical status ledgers are
preserved byte for byte.  V17 adds only the minimal common-dilation theorem,
the projector-word principal-angle compiler, the bare half-edge projector
obstruction, and the six-real-dimensional arithmetic source-information lower
bound above.  Future continuations must follow this block.
% ====================== END IMMUTABLE V17 MATHEMATICAL BODY ======================
% APPEND V18 AND LATER ITERATIONS HERE.



\begin{thebibliography}{99}
\setlength{\itemsep}{2pt}
\setlength{\parsep}{0pt}
\bibitem{GT77}
A.~P\'elissier,
\emph{The Renewal Grand Tensor: Relative Primitive Minimality, Operational
Spectralization, and Domain-Specific SM--Spacetime and Arithmetic Loadings},
Version V077, 19 September 2026. Supplied monograph;
\path{Pasted text(20260920-195255).txt}.

\bibitem{Arithmetic}
A.~P\'elissier,
\emph{The Arithmetic Renewal Gran Tensor: Conditional Closure Mechanisms for
the Generalized Riemann Hypothesis and Affine Prime Correlations},
12 September 2026. Supplied arithmetic monograph;
\path{Pasted text (2)(20260920-195320).txt}.

\bibitem{AD9}
Renewal Arithmetic and Geometry Programme,
\emph{Arithmetic--Spacetime--Standard-Model Duality Programme},
cumulative V9 (including V1--V8), 2026. Earlier research notes recovered from
the user's Library;
\path{Arithmetic_Spacetime_Standard_Model_Duality_Program_V9.tex}.
The precise source is \path{eq:v9-theta-source-family}.

\bibitem{Spacetime}
A.~Pelissier,
\emph{Renewal Geometry and the Emergence of Lorentzian Spacetime},
2026. Supplied companion manuscript, treated as a citable antecedent;
\path{Pasted text(20260920-200547).txt}.

\bibitem{Variational}
A.~Pelissier,
\emph{Variational Closure of Finite Einstein--Standard-Model Actions},
2026. Supplied companion manuscript;
\path{Pasted text (2)(20260920-200551).txt}.

\bibitem{Duality}
A.~Pelissier,
\emph{A Duality Between Spacetime--Gauge--Incidence Actions and Matter
Multiplicity}, 2026. Supplied companion manuscript;
\path{Pasted text (3)(20260920-200558).txt}.

\bibitem{Spectral}
A.~Pelissier,
\emph{From Predictive Dynamics to Spectral Geometry},
2026. Supplied companion manuscript;
\path{Pasted text (4)(20260920-200603).txt}.

\bibitem{DehaeneDeMoor}
J.~Dehaene and B.~De~Moor,
The Clifford group, stabilizer states, and linear and quadratic operations
over GF(2), \emph{Physical Review A} \textbf{68} (2003), 042318.
\href{https://arxiv.org/abs/quant-ph/0304125}{arXiv:quant-ph/0304125}.

\bibitem{Hasebe}
K.~Hasebe,
Perfectly spherical Bloch hyper-spheres from quantum matrix geometry,
\emph{Nuclear Physics B} \textbf{1005} (2024), 116595.
\href{https://arxiv.org/abs/2402.07149}{arXiv:2402.07149}.

\bibitem{ReichsteinShukla}
Z.~Reichstein and A.~K.~Shukla,
\emph{Essential dimension of double covers of symmetric and alternating groups},
arXiv:1906.03698 (2019), especially the presentations in Section 1.
\href{https://arxiv.org/abs/1906.03698}{arXiv:1906.03698}.

\bibitem{DLMFTheta}
B.~Deconinck,
\emph{Multidimensional Theta Functions}, Chapter 21 of the NIST Digital
Library of Mathematical Functions, Sections 21.2 and 21.6(ii),
especially Equation 21.6.8.
\href{https://dlmf.nist.gov/21.6}{DLMF Chapter 21}.

\bibitem{BernatskaThomae}
J.~Bernatska,
General derivative Thomae formula for singular half-periods,
\emph{Letters in Mathematical Physics} \textbf{110} (2020), 2983--3014.
The classical first Thomae theorem is stated in Section~2.6, Equation~(18).
\href{https://arxiv.org/abs/1904.09333}{arXiv:1904.09333};
\href{https://doi.org/10.1007/s11005-020-01315-8}{doi:10.1007/s11005-020-01315-8}.


\bibitem{MilneAV}
J.~S.~Milne,
\emph{Abelian Varieties}, course notes, version 2.00, 16 March 2008.
Chapter I, Sections~10--11; Chapter II, Section~1; Chapter III, Sections~2 and~11.
\href{https://www.jmilne.org/math/CourseNotes/AV.pdf}{Author's course notes}.

\bibitem{AchterHowe}
J.~D.~Achter and E.~W.~Howe,
Hasse--Witt and Cartier--Manin matrices: A warning and a request,
\emph{Contemporary Mathematics} \textbf{722} (2019), 1--18.
\href{https://arxiv.org/abs/1710.10726}{arXiv:1710.10726}.

\bibitem{AchingerZdanowicz}
P.~Achinger and M.~Zdanowicz,
\emph{Serre--Tate theory for Calabi--Yau varieties},
arXiv:1807.11295v2 (2019), Proposition~2.3.1.
\href{https://arxiv.org/abs/1807.11295}{arXiv:1807.11295}.


\bibitem{deJongGauss}
R.~de Jong,
Gauss map on the theta divisor and Green's functions,
In B.~Edixhoven, G.~van der Geer and B.~Moonen (eds.),
\emph{Modular Forms on Schiermonnikoog}, Cambridge University Press, 2008.
Updated preprint, arXiv:0705.0098v2 (2012).
\href{https://arxiv.org/abs/0705.0098}{arXiv:0705.0098}.

\bibitem{LombardoEnd}
D.~Lombardo,
Computing the geometric endomorphism ring of a genus 2 Jacobian,
\emph{Mathematics of Computation} \textbf{88} (2019), 889--929.
\href{https://arxiv.org/abs/1610.09674}{arXiv:1610.09674};
\href{https://doi.org/10.1090/mcom/3358}{doi:10.1090/mcom/3358}.


\bibitem{FKRS}
F.~Fit\'e, K.~S.~Kedlaya, V.~Rotger and A.~V.~Sutherland,
Sato--Tate distributions and Galois endomorphism modules in genus 2,
\emph{Compositio Mathematica} \textbf{148} (2012), 1390--1442.
\href{https://arxiv.org/abs/1110.6638}{arXiv:1110.6638};
\href{https://doi.org/10.1112/S0010437X12000279}{doi:10.1112/S0010437X12000279}.


\bibitem{KatzOda}
N.~M.~Katz and T.~Oda,
On the differentiation of de Rham cohomology classes with respect to parameters,
\emph{Journal of Mathematics of Kyoto University} \textbf{8} (1968), 199--213.
\href{https://doi.org/10.1215/kjm/1250524135}{doi:10.1215/kjm/1250524135}.

\bibitem{KTdeRham}
B.~K\"ock and J.~Tait,
On the de Rham cohomology of hyperelliptic curves,
\emph{Research in Number Theory} \textbf{4} (2018), article 19.
\href{https://arxiv.org/abs/1709.03422}{arXiv:1709.03422}.


\bibitem{DLMFBeta}
NIST Digital Library of Mathematical Functions,
Chapter 5, Section 5.12, Euler's beta integral.
\href{https://dlmf.nist.gov/5.12}{DLMF Section 5.12}.

\bibitem{DLMFHyper}
NIST Digital Library of Mathematical Functions,
Chapter 16, Section 16.8, generalized hypergeometric differential equations.
\href{https://dlmf.nist.gov/16.8}{DLMF Section 16.8}.

\bibitem{BeukersHeckman}
F.~Beukers and G.~Heckman,
Monodromy for the hypergeometric function ${}_nF_{n-1}$,
\emph{Inventiones Mathematicae} \textbf{95} (1989), 325--354.
\href{https://doi.org/10.1007/BF01393900}{doi:10.1007/BF01393900}.

\bibitem{SinghVenkataramana}
S.~Singh and T.~N.~Venkataramana,
Arithmeticity of certain symplectic hypergeometric groups,
\emph{Duke Mathematical Journal} \textbf{163} (2014), no.~3, 591--617.
Theorem~1.1 and the Levelt background in Section~1.
\href{https://arxiv.org/abs/1208.6460}{arXiv:1208.6460};
\href{https://doi.org/10.1215/00127094-2410655}{doi:10.1215/00127094-2410655}.


\bibitem{GreenGriffiths}
M.~Green and P.~Griffiths,
\emph{Positivity of vector bundles and Hodge theory}, arXiv:1803.07405.
Sections II.B and VI.B: metric subbundle and Hodge-curvature formulas.
\href{https://arxiv.org/abs/1803.07405}{arXiv:1803.07405}.

\bibitem{SimpsonHiggs}
C.~T.~Simpson,
Higgs bundles and local systems,
\emph{Publications Math\'ematiques de l'IH\'ES} \textbf{75} (1992), 5--95.
\href{https://doi.org/10.1007/BF02699491}{doi:10.1007/BF02699491}.

\bibitem{AmbroseSinger}
W.~Ambrose and I.~M.~Singer,
A theorem on holonomy,
\emph{Transactions of the American Mathematical Society} \textbf{75}
(1953), no.~3, 428--443.
\href{https://doi.org/10.1090/S0002-9947-1953-0063739-1}{doi:10.1090/S0002-9947-1953-0063739-1}.


\bibitem{StacksCartier}
The Stacks Project Authors,
\emph{The Stacks Project}, effective Cartier divisors, conormal sheaves,
and their canonical sections, Tags 0C4S and 0B4A.
\href{https://stacks.math.columbia.edu/tag/0C4S}{Tag 0C4S};
\href{https://stacks.math.columbia.edu/tag/0B4A}{Tag 0B4A}.

\bibitem{KouvidakisTheta}
A.~Kouvidakis,
\emph{Theta line bundles and the determinant of the Hodge bundle},
arXiv:alg-geom/9604017 (1996).
\href{https://arxiv.org/abs/alg-geom/9604017}{arXiv:alg-geom/9604017}.

\bibitem{EdidinHu}
D.~Edidin and Z.~Hu,
\emph{The integral Chow rings of the stacks of hyperelliptic Weierstrass
points}, arXiv:2208.00556v2, 31 July 2024,
Theorems~1.3--1.5. The genus-two specializations give cotangent-class
orders $40$, $8$, and $2$ for one, two, and at least three marks.
\href{https://arxiv.org/abs/2208.00556}{arXiv:2208.00556}.


\bibitem{KedlayaMW}
K.~S.~Kedlaya,
Counting points on hyperelliptic curves using Monsky--Washnitzer cohomology,
\emph{Journal of the Ramanujan Mathematical Society} \textbf{16}
(2001), no.~4, 323--338. Sections 3--4; the degree-zero pole reduction
is separated explicitly in the present continuation.
\href{https://arxiv.org/abs/math/0105031}{arXiv:math/0105031}.

\bibitem{BogaartBasis}
T.~van den Bogaart,
\emph{About the choice of a basis in Kedlaya's algorithm},
arXiv:0809.1243 (2008), Sections 3--5.
\href{https://arxiv.org/abs/0809.1243}{arXiv:0809.1243}.

\bibitem{KedlayaFrobBook}
K.~S.~Kedlaya,
\emph{$p$-adic Differential Equations}, Cambridge University Press,
2010, Chapter 17, Frobenius structures on differential modules.
\href{https://doi.org/10.1017/CBO9780511750922.019}{doi:10.1017/CBO9780511750922.019}.

\bibitem{KedlayaTuitman}
K.~S.~Kedlaya and J.~Tuitman,
Effective convergence bounds for Frobenius structures on connections,
\emph{Rendiconti del Seminario Matematico della Universit\`a di Padova}
\textbf{128} (2012), 7--16.
\href{https://doi.org/10.4171/RSMUP/128-2}{doi:10.4171/RSMUP/128-2};
\href{https://arxiv.org/abs/1111.0136}{arXiv:1111.0136}.

\end{thebibliography}
\end{document}
